\parte{Dualidad y principios fundamentales}{Las extensiones de funcionales convierten información parcial en mediciones globales. Desarrollamos Hahn--Banach, separación, mapeo abierto, gráfica cerrada y adjuntos, con especial atención a las hipótesis de completitud.}
\chapter{Espacios, duales y operadores: punto de partida}\label{cap:v3-01}
La dimensión finita oculta diferencias que serán decisivas: allí toda aplicación lineal es continua y toda bola cerrada es compacta. En dimensión infinita hay que distinguir esas propiedades y comprobarlas. Este capítulo fija las convenciones y reúne las herramientas previas que utilizaremos; los argumentos breves muestran cómo se conectan los resultados de los dos primeros volúmenes.

\section{Normas, completitud y ejemplos}\label{lec:v3-00}
\fuentelectura{00}{Algunos resultados preliminares}
\begin{definicion}
Una norma en un espacio vectorial $X$ sobre $\F=\R$ o $\C$ es una función $\norm\cdot:X\to[0,\infty)$ tal que $\norm x=0$ equivale a $x=0$, $\norm{\alpha x}=|\alpha|\norm x$ y $\norm{x+y}\leq\norm x+\norm y$. La distancia asociada es $d(x,y)=\norm{x-y}$. El espacio es de Banach si toda sucesión de Cauchy converge en él.
\end{definicion}
\index{Banach, espacio de}\index{norma}
La desigualdad triangular implica $|\norm x-\norm y|\leq\norm{x-y}$. Por tanto la norma es continua. También lo son las operaciones, pues
\[
 \norm{(x+y)-(x_0+y_0)}\leq\norm{x-x_0}+\norm{y-y_0},\qquad
 \norm{\alpha x-\alpha_0x_0}\leq|\alpha|\norm{x-x_0}+|\alpha-\alpha_0|\norm{x_0}.
\]
En la segunda estimación $|\alpha|$ permanece acotado cuando $\alpha\to\alpha_0$.

\begin{proposicion}\label{v3:series-banach}
Un espacio normado es de Banach si y solo si toda serie $\sum_nx_n$ con $\sum_n\norm{x_n}<\infty$ converge en el espacio.
\end{proposicion}
\begin{proof}
Las sumas parciales de una serie absolutamente convergente son de Cauchy, porque la norma de cualquier cola no supera la suma de las normas de esa cola. La completitud da el límite.

Recíprocamente, dada una sucesión de Cauchy $(y_n)$, elija $n_1<n_2<\cdots$ de modo que $\norm{y_p-y_q}<2^{-j}$ si $p,q\geq n_j$. La serie
\[
 y_{n_1}+\sum_{j\geq1}(y_{n_{j+1}}-y_{n_j})
\]
es absolutamente convergente. Sus sumas parciales son $y_{n_j}$, así que esta subsucesión converge a algún $y$. Para $\eps>0$, elija $N$ con $\norm{y_p-y_q}<\eps/2$ para $p,q\geq N$ y después $j$ con $n_j\geq N$, $\norm{y_{n_j}-y}<\eps/2$. Entonces $\norm{y_n-y}<\eps$ para $n\geq N$. Converge la sucesión completa.
\end{proof}

Los siguientes ejemplos recuperan los de la lectura, con las precisiones necesarias.
\begin{enumerate}
\item En $\F^m$ las normas $\norm\cdot_1$, $\norm\cdot_2$ y $\norm\cdot_\infty$ son equivalentes: $\norm x_\infty\leq\norm x_2\leq\norm x_1\leq m\norm x_\infty$. La convergencia de Cauchy se verifica coordenada a coordenada, y esos espacios son completos.
\item Para un espacio de medida $(\Omega,\Sigma,\mu)$, $L^p(\mu)$, $1\leq p<\infty$, consta de \emph{clases de equivalencia} de funciones medibles, con $\norm f_p=(\int|f|^p\,d\mu)^{1/p}$. La norma sería solo una seminorma si no se identificaran funciones iguales casi en todas partes. Minkowski da la desigualdad triangular. Para comprobar completitud mediante series, suponga $\sum_n\norm{f_n}_p<\infty$. Las sumas $g_N=\sum_{n\leq N}|f_n|$ satisfacen $\norm{g_N}_p\leq\sum_n\norm{f_n}_p$. Convergencia monótona da $g=\sum_n|f_n|\in L^p$, de modo que $\sum_nf_n$ converge casi en todas partes a $f$. Aplicando Fatou a cada cola se obtiene $\norm{f-\sum_{n\leq N}f_n}_p\leq\sum_{n>N}\norm{f_n}_p\to0$. El caso de medida de conteo produce $\ell^p(I)$.
\item $L^\infty(\mu)$ lleva la norma $\norm f_\infty=\inf\{C\geq0:|f|\leq C\text{ casi en todas partes}\}$. Para una serie con suma de normas finita, retire la unión numerable de los conjuntos excepcionales donde las cotas fallan. Fuera de ella la serie converge uniformemente y su cola está dominada por la cola de normas. Esto demuestra completitud. Incluye $\ell^\infty(I)$.
\item Para un espacio topológico $X$, la norma uniforme está definida en $C_b(X)$, las funciones continuas \emph{acotadas}. Una sucesión uniformemente de Cauchy tiene un límite uniforme acotado; la estimación de tres términos demuestra su continuidad. Si $X$ es compacto, toda función continua es acotada y escribimos $C(X)$. En un espacio general no puede atribuirse una norma finita a todo $C(X)$: en $X=\R$, $x\mapsto x$ es un contraejemplo.
\item Las medidas con signo o complejas de variación total finita forman un espacio de Banach con $\norm\nu=|\nu|(\Omega)$. En efecto, si $\sum_n\norm{\nu_n}<\infty$, defina $\nu(E)=\sum_n\nu_n(E)$. Para conjuntos disjuntos $E_j$, la suma doble es absolutamente convergente porque $\sum_n\sum_j|\nu_n(E_j)|\leq\sum_n\norm{\nu_n}$. Se pueden intercambiar las sumas; así $\nu$ es una medida y $\norm{\nu-\sum_{n\leq N}\nu_n}\leq\sum_{n>N}\norm{\nu_n}$. La caracterización anterior da completitud. En un espacio localmente compacto de Hausdorff, las medidas de Borel regulares de variación finita forman un subespacio cerrado: una aproximación en variación traslada las aproximaciones interior y exterior de una medida regular al límite, con un error adicional a lo sumo dos veces la distancia en variación. Por tanto también es de Banach.
\item En $C^m([a,b])$, $\norm f_{C^m}=\sum_{j=0}^m\norm{f^{(j)}}_\infty$ es completa. Si las derivadas de una sucesión de Cauchy convergen uniformemente a $g_0,\ldots,g_m$, las identidades
\[
 f_n^{(j)}(t)-f_n^{(j)}(a)=\int_a^t f_n^{(j+1)}(s)\,ds
\]
pasan al límite uniformemente. El teorema fundamental del cálculo da $g_j'=g_{j+1}$. Por tanto $g_0\in C^m$ y la convergencia ocurre en esa norma. Esta norma es adecuada para el espacio de Banach; su eventual uso como \emph{norma de álgebra} exige otra comprobación en el capítulo~\ref{cap:v3-10}.
\end{enumerate}
\begin{ejemplo}[La suavidad ilimitada no garantiza completitud]
En un intervalo con $c$ interior, $h_\eps(t)=\sqrt{(t-c)^2+\eps^2}$ converge uniformemente a $|t-c|$. Para $m\geq1$, integre $h_\eps$ sucesivamente $m$ veces desde $a$, con constantes de integración nulas. Las funciones resultantes son suaves y convergen en $C^m$ a una función cuya derivada $m$-ésima es $|t-c|$, que no es diferenciable en $c$. Así $C^\infty([a,b])$ no es completo con ninguna norma $\norm\cdot_{C^m}$ fija. Para $m=0$ basta la sucesión $h_\eps$.
\end{ejemplo}

\section{Operadores y duales}
\begin{teorema}\label{v3:acotacion}
Para una aplicación lineal $T:X\to Y$ entre espacios normados son equivalentes: continuidad; continuidad en cero; existencia de $C\geq0$ con $\norm{Tx}\leq C\norm x$. Los operadores que satisfacen estas condiciones forman $\Bcal(X,Y)$, con norma
\[
 \norm T=\sup_{\norm x\leq1}\norm{Tx}=\inf\{C\geq0:\norm{Tx}\leq C\norm x\ \forall x\}.
\]
Si $X\ne\{0\}$, el supremo puede tomarse sobre $\norm x=1$. Si $Y$ es de Banach, $\Bcal(X,Y)$ también lo es.
\end{teorema}
\begin{proof}
La continuidad implica continuidad en cero. Si $\norm{Tx}<1$ cuando $\norm x<\delta$, para $x\ne0$ aplique esa implicación a $\delta x/(2\norm x)$ y obtenga $\norm{Tx}\leq2\norm x/\delta$. Una cota global implica $\norm{Tx-Tz}\leq C\norm{x-z}$ y, por tanto, continuidad. Normalizar $x$ muestra las fórmulas del supremo y del ínfimo. Homogeneidad y desigualdad triangular pasan al supremo; si $\norm T=0$, todo $Tx$ es cero. En el dominio nulo el operador cero tiene norma cero, sin recurrir al supremo sobre una esfera vacía.

Sea $(T_n)$ de Cauchy en norma. Para cada $x$, $(T_nx)$ es de Cauchy en $Y$, de modo que $Tx=\lim_nT_nx$ está definido. La continuidad de las operaciones muestra que $T$ es lineal. Además $\sup_n\norm{T_n}=M<\infty$ implica $\norm{Tx}\leq M\norm x$. Si $\norm{T_n-T_m}<\eps$ para $m,n\geq N$, fijar $n$ y pasar $m\to\infty$ da $\norm{(T_n-T)x}\leq\eps\norm x$. Tomar supremo produce $\norm{T_n-T}\leq\eps$. Así la convergencia es en norma de operador.
\end{proof}
\index{operador!acotado}\index{dual!continuo}
El dual continuo es $X^*=\Bcal(X,\F)$; siempre es de Banach. Si $S:Y\to Z$ es acotado, $\norm{ST}\leq\norm S\norm T$ porque $\norm{STx}\leq\norm S\norm{Tx}$. En particular, $\Bcal(X)=\Bcal(X,X)$ tiene un producto continuo: si $S_n\to S$ y $T_n\to T$, entonces
\[
 \norm{S_nT_n-ST}\leq\norm{S_n}\norm{T_n-T}+\norm{S_n-S}\norm T\longrightarrow0.
\]
La acotación de $(\norm{S_n})$ procede de su convergencia.

\section{Geometría de Hilbert y proyecciones}
Un producto interno satisface positividad definida, linealidad en su primera variable y simetría conjugada $\ip xy=\overline{\ip yx}$. Cauchy--Schwarz se obtiene, si $y\ne0$, expandiendo $0\leq\norm{x-\ip xy\,y/\norm y^2}^2$. Resulta $|\ip xy|\leq\norm x\norm y$, y la expansión de $\norm{x+y}^2$ da la desigualdad triangular. Un espacio completo para esta norma es de Hilbert.

\begin{proposicion}[Paralelogramo y polarización]\label{v3:polarizacion}
Una norma proviene de un producto interno si y solo si
\[
 \norm{x+y}^2+\norm{x-y}^2=2\norm x^2+2\norm y^2.
\]
El producto interno, cuando existe, es único. En el caso complejo, lineal en la primera variable, es
\[
 \ip xy=\frac{\norm{x+y}^2-\norm{x-y}^2}{4}
 +i\frac{\norm{x+iy}^2-\norm{x-iy}^2}{4}.
\]
\end{proposicion}
\begin{proof}
La necesidad y la fórmula se verifican expandiendo los cuatro cuadrados. Para la suficiencia, ponga $q(x)=\norm x^2$ y, en el espacio real subyacente, $B(x,y)=(q(x+y)-q(x-y))/4$. Es simétrica, continua, $B(x,x)=q(x)$ y $B(0,y)=0$. Restar dos identidades del paralelogramo da
\[
 B(x+z,y)+B(x-z,y)=2B(x,y).
\]
Tomar $z=x$ muestra $B(2x,y)=2B(x,y)$. Aplicar la igualdad anterior a $x=(u+v)/2$, $z=(u-v)/2$ da aditividad. La aditividad implica homogeneidad racional; la continuidad la extiende a escalares reales. En el caso real, $B$ es el producto buscado.

En el caso complejo, $q(ix)=q(x)$ implica $B(ix,iy)=B(x,y)$ y $B(ix,y)=-B(x,iy)$. Defina $\ip xy=B(x,y)+iB(x,iy)$. Las identidades anteriores verifican $\ip{ix}y=i\ip xy$ y la simetría conjugada; la bilinealidad real verifica las demás propiedades. También $B(x,ix)=0$, pues $B(x,ix)=-B(ix,x)$, así que $\ip xx=q(x)>0$ si $x\ne0$.
\end{proof}
\begin{observacion}[Precisión del repaso original]
No es correcto afirmar que, de todos los ejemplos anteriores, solo $L^2$ admite producto interno: $\F^m$ con norma euclidiana también lo admite, y espacios de funciones sobre un solo punto son unidimensionales. El criterio exacto es el paralelogramo. Por ejemplo, en $\ell^1$ con $x=e_1$, $y=e_2$ sus dos miembros valen $8$ y $4$.
\end{observacion}

\begin{teorema}[Proyección ortogonal]\label{v3:proyeccion}
Si $M$ es un subespacio completo de un espacio con producto interno $X$, todo $x\in X$ admite una única descomposición $x=m+r$, donde $m\in M$, $r\in M^\perp$. Además
\[
 \dist(x,M)=\norm{x-m},\qquad
 \norm{x-u}^2=\norm{x-m}^2+\norm{m-u}^2\quad(u\in M).
\]
En Hilbert, esto ocurre exactamente para los subespacios cerrados.
\end{teorema}
\begin{proof}
Ponga $d=\dist(x,M)$ y elija $m_n\in M$ con $\norm{x-m_n}\to d$. Como $(m_n+m_k)/2\in M$, el paralelogramo implica
\[
 \norm{m_n-m_k}^2\leq2\norm{x-m_n}^2+2\norm{x-m_k}^2-4d^2\longrightarrow0.
\]
La completitud de $M$ da $m_n\to m\in M$, y $\norm{x-m}=d$. Para $u\in M$, $\norm{x-m-tu}^2\geq\norm{x-m}^2$ para todo $t\in\R$. Expandir ese polinomio cuadrático obliga a $\operatorname{Re}\ip{x-m}u=0$; en el caso complejo, reemplazar $u$ por $iu$ obliga también a que la parte imaginaria sea cero. Así $r=x-m\in M^\perp$.

La identidad pitagórica da la fórmula para $u$ y la unicidad: cualquier otro minimizador debe satisfacer $\norm{m-u}=0$. Si todos los vectores admiten la descomposición, $M=(M^\perp)^\perp$ y es cerrado, pues las funcionales $x\mapsto\ip xy$ son continuas. Un cerrado de Hilbert es completo, por lo que el recíproco ya está probado.
\end{proof}
\index{proyeccion ortogonal@proyección ortogonal}
Escribimos $P_Mx=m$. La unicidad verifica que $P_M$ es lineal; Pitágoras da $\norm{P_Mx}\leq\norm x$, y $P_M^2=P_M$. El término \emph{distinguido} de las lecturas equivale a ortogonalmente complementado. La condición correcta exige $m\in M$; y la identidad de distancias es $\dist(x,M)^2=\norm x^2-\norm m^2$, con el cuadrado en la distancia.

Una familia ortonormal $E$ satisface $\ip ef=\delta_{ef}$. Es una base ortonormal si su espacio generado es denso. Zorn da una familia ortonormal maximal; si su espacio generado no fuera denso, la proyección sobre su cierre produciría un vector unitario perpendicular a toda la familia, contradicción. Para un subconjunto finito $F\subseteq E$, Pitágoras da
\[
 \sum_{e\in F}|\ip xe|^2\leq\norm x^2.
\]
Solo puede haber un número numerable de coeficientes no nulos: para cada $n$, los que tienen módulo al menos $1/n$ forman un conjunto finito. Las sumas de Fourier convergen por completitud y tienen resto perpendicular a $E$, de modo que, si $E$ es total, el resto es cero.

\section{Separabilidad y dimensión algebraica}
\index{separabilidad}\index{base!de Hamel}
Una base de Hamel expresa cada vector mediante una combinación \emph{finita}; una familia total solo exige densidad del espacio generado. No debe llamarse base de Schauder a una familia total arbitraria, pues esta última noción exige además expansiones convergentes únicas en un orden dado. Un espacio normado es separable si contiene un conjunto numerable denso, o equivalentemente un subespacio denso de dimensión algebraica a lo sumo numerable. Para ver la equivalencia, tome el espacio generado por el conjunto denso; en la otra dirección use coeficientes racionales, o racionales complejos, en todas las combinaciones finitas de una familia generadora numerable.

El teorema de Baire prohíbe una base de Hamel infinita numerable en un espacio de Banach: si $x_1,x_2,\ldots$ fuera tal base, el espacio sería la unión de los cerrados $\spanop(x_1,\ldots,x_n)$, todos de interior vacío. Son cerrados porque, en dimensión finita, las coordenadas son equivalentes a cualquier norma: la función $a\mapsto\norm{\sum_ja_jx_j}$ tiene mínimo positivo sobre $\max_j|a_j|=1$. Un subespacio con interior no vacío contiene una bola en cero y, por escalamiento, todo el espacio.

Esto \emph{no} prohíbe espacios de Banach separables infinitos dimensionales. $\ell^2$ es uno: las sucesiones de soporte finito con entradas racionales complejas son numerables y densas. La frase contraria en la lectura0 confunde una familia total numerable con una base de Hamel numerable. El mismo cuidado se aplica a los monomios en $C([a,b])$: son linealmente independientes y totales por Weierstrass, pero solo generan algebraicamente los polinomios.

\section{Adjuntos, compacidad y dimensión finita}
\begin{teorema}[Representación y adjunto de Hilbert]\label{v3:riesz-adjunto}
Para cada $\varphi\in H^*$ existe un único $y\in H$ tal que $\varphi(x)=\ip xy$ y $\norm\varphi=\norm y$. Para cada $T\in\Bcal(H,K)$ existe un único $T^*\in\Bcal(K,H)$ con $\ip{Tx}y=\ip{x}{T^*y}$; satisface
\[
 (\alpha S+\beta T)^*=\bar\alpha S^*+\bar\beta T^*,\quad
 (UT)^*=T^*U^*,\quad (T^*)^*=T,\quad\norm{T^*}=\norm T.
\]
\end{teorema}
\begin{proof}
Si $\varphi\ne0$, su núcleo $M$ es cerrado y tiene codimensión uno: si $\varphi(u)=1$, entonces $x-\varphi(x)u\in M$. Elija un vector unitario $e\in M^\perp$. Todo $x$ es $m+\ip xe e$, de modo que $\varphi(x)=\ip xe\varphi(e)=\ip{x}{\overline{\varphi(e)}e}$. Cauchy--Schwarz da la cota de norma, y evaluar en $y/\norm y$ da igualdad. Para $\varphi=0$ tome $y=0$. Si dos vectores representan la misma funcional, evaluar su diferencia en sí misma demuestra unicidad.

Para $y\in K$, aplique lo anterior a $x\mapsto\ip{Tx}y$, cuya norma es a lo sumo $\norm T\norm y$. Su representante es $T^*y$. Las identidades del producto interno y la unicidad muestran que $T^*$ es lineal y acotado. La misma identidad demuestra las tres fórmulas algebraicas. Además
\[
 \norm T=\sup_{\norm x\leq1,\,\norm y\leq1}|\ip{Tx}y|
 =\sup_{\norm x\leq1,\,\norm y\leq1}|\ip{x}{T^*y}|=\norm{T^*}.
\]
La igualdad $\norm z=\sup_{\norm y\leq1}|\ip zy|$ se verifica con $y=z/\norm z$ si $z\ne0$ y directamente si $z=0$.
\end{proof}

\begin{definicion}
Un operador lineal $T:X\to Y$ es compacto si manda cada sucesión acotada a una sucesión con una subsucesión convergente en $Y$. Equivalentemente, manda conjuntos acotados a conjuntos relativamente compactos. Denotamos su espacio por $\Kcal(X,Y)$.
\end{definicion}
\index{operador!compacto}
La equivalencia usa la caracterización secuencial de compacidad métrica. Todo operador compacto es acotado: de lo contrario existirían $\norm{x_n}\leq1$ con $\norm{Tx_n}>n$, sin subsucesión convergente. Un operador de rango finito es compacto \emph{si es acotado}: su bola imagen es acotada en un espacio finito dimensional y tiene cierre compacto. La continuidad no se deduce del rango finito cuando el dominio es infinito dimensional.

La suma de dos compactos es compacta: primero extraiga una subsucesión para uno y después una subsucesión de esta para el otro. Componer con un operador acotado a cualquiera de los lados conserva compacidad, usando respectivamente la continuidad o la conservación de acotación. Por eso $\Kcal(X)$ es un ideal de $\Bcal(X)$.

\begin{proposicion}\label{v3:aproximacion-compactos}
Todo operador compacto entre espacios de Hilbert es límite en norma de operadores acotados de rango finito. No es necesario que los espacios sean separables.
\end{proposicion}
\begin{proof}
Sea $T:H\to K$ compacto y $\eps>0$. El cierre de $T(\overline B_H)$ es compacto; elija una red finita $y_1,\ldots,y_m$ de radio $\eps$ y ponga $F=\spanop(y_1,\ldots,y_m)$. Para $\norm x\leq1$, la propiedad de mejor aproximante implica
\[
 \norm{Tx-P_FTx}=\dist(Tx,F)\leq\min_j\norm{Tx-y_j}<\eps.
\]
Así $P_FT$ tiene rango finito y $\norm{T-P_FT}\leq\eps$. Elegir $\eps=1/n$ produce la sucesión requerida.
\end{proof}

El encaje $C^1([a,b])\to C([a,b])$ es compacto: la bola unitaria tiene funciones uniformemente acotadas y, por el teorema del valor medio, equicontinuas; Arzelà--Ascoli da la conclusión. Para el operador de Volterra $Vf(t)=\int_0^tf(s)\,ds$ en $L^2([0,1])$, Cauchy--Schwarz da $|Vf(t)|\leq\norm f_2$ y
\[
 |Vf(t)-Vf(s)|\leq|t-s|^{1/2}\norm f_2.
\]
La bola imagen es relativamente compacta en $C([0,1])$ por Arzelà--Ascoli, y su inclusión continua en $L^2$ conserva esa propiedad. Estos ejemplos tienen rango infinito: las primitivas de monomios distintos siguen siendo linealmente independientes.

\begin{teorema}[Teorema espectral en dimensión finita]\label{v3:espectral-finito}
Un operador normal en un espacio complejo de Hilbert de dimensión finita admite una base ortonormal de autovectores. Un operador autoadjunto real satisface la misma conclusión con autovalores reales.
\end{teorema}
\begin{proof}
Para un operador normal $A$, expandir cuadrados y usar $A^*A=AA^*$ da $\norm{(A-\lambda I)x}=\norm{(A^*-\bar\lambda I)x}$. Por tanto $Ax=\lambda x$ implica $A^*x=\bar\lambda x$. En dimensión compleja positiva, el polinomio característico tiene una raíz y existe un autovector unitario $e$. Si $x\perp e$, entonces $\ip{Ax}e=\ip{x}{A^*e}=0$ y también $A^*x\perp e$. Así $e^\perp$ reduce a $A$; la restricción sigue siendo normal. Inducción sobre la dimensión da la base ortonormal. Para un autoadjunto real, se puede complejificar el espacio: un autovalor $\lambda$ cumple $\lambda\norm z^2=\ip{Az}z\in\R$. Si $z=u+iv$ es su autovector, al menos uno de $u,v$ es un autovector real no nulo. La misma inducción termina la prueba.
\end{proof}

\section{Balance de las herramientas}
La norma controla la continuidad; la completitud permite sumar correcciones; el producto interno convierte la aproximación en perpendicularidad. En adelante preguntaremos cuándo una medición definida en un subespacio puede extenderse, cuándo la imagen de un operador es cerrada y qué impide invertir $T-\lambda I$. Las respuestas exigirán distinguir cuidadosamente las diferentes formas de compacidad y convergencia.

\chapter{Hahn--Banach y la geometría del dual}\label{cap:v3-02}
Una funcional definida en un subespacio puede medir solo una parte del espacio. El teorema de Hahn--Banach permite extender esa medición sin empeorar su cota. Desarrollaremos primero el mecanismo de una dimensión, después la extensión global y finalmente las consecuencias para distancias, cocientes y el bidual. El ejemplo de los límites de Banach mostrará por qué una propiedad adicional debe formar parte de la definición y no suponerse automáticamente.

\section{Extender una funcional conservando una cota}\label{lec:v3-01}
\fuentelectura{01}{El teorema de Hahn--Banach}
\begin{definicion}
En un espacio vectorial real $X$, una funcional $p:X\to\R$ es sublineal si $p(x+y)\leq p(x)+p(y)$ y $p(tx)=tp(x)$ para $t\geq0$. Una funcional lineal y una norma son ejemplos. No se exige que $p$ sea no negativa ni simétrica.
\end{definicion}
\index{funcional!sublineal}
Si la homogeneidad se enuncia solo para $t>0$, tomar $x=0$, $t=2$ demuestra $p(0)=0$ y recupera el caso $t=0$. La subaditividad implica siempre $-p(-x)\leq p(x)$.

\begin{teorema}[Hahn--Banach real]\label{v3:hb-real}
Sea $M\leq X$, sea $p:X\to\R$ sublineal y sea $f:M\to\R$ lineal con $f\leq p$ en $M$. Existe una funcional lineal $F:X\to\R$ que extiende a $f$ y satisface $F\leq p$ en $X$.
\end{teorema}
\index{Hahn--Banach, teorema de}
\begin{proof}
\emph{Paso de una dimensión.} Si $x_0\notin M$, la representación $m+tx_0$ en $M+\R x_0$ es única: dos representaciones darían un múltiplo no nulo de $x_0$ en $M$. Toda extensión lineal tiene entonces la forma $f_c(m+tx_0)=f(m)+tc$. Para que $f_c\leq p$, los casos $t>0$ y $t<0$ requieren, respectivamente,
\[
 c\leq p(z+x_0)-f(z)\quad(z\in M),\qquad
 c\geq f(y)-p(y-x_0)\quad(y\in M).
\]
En efecto, divida por $t$ en el primer caso y escriba $t=-s$, $s>0$, en el segundo. Para $t=0$ ya tenemos la hipótesis original. El intervalo de valores admisibles es no vacío porque
\[
 f(y)+f(z)=f(y+z)\leq p(y+z)\leq p(y-x_0)+p(z+x_0).
\]
Así cada cota inferior es menor o igual que cada cota superior. El supremo de las inferiores y el ínfimo de las superiores son reales: tomar $y=0$ y $z=0$ da cotas finitas en ambos sentidos. Elija $c$ entre ellos.

\emph{Paso global.} Considere los pares $(N,g)$ donde $M\subseteq N\leq X$, $g$ extiende a $f$ y $g\leq p$ en $N$. Ordénelos por extensión. El conjunto es no vacío porque contiene $(M,f)$. En una cadena, las funciones coinciden donde se superponen sus dominios. La unión es, por tanto, una función bien definida. Su dominio es un subespacio: dos vectores de la unión pertenecen juntos al dominio mayor de dos miembros comparables, y sus combinaciones lineales también. Esa misma observación verifica linealidad; cada valor sigue dominado por $p$ y la función sigue extendiendo a $f$. La unión es una cota superior de la cadena. La cadena vacía tiene, por ejemplo, la cota $(M,f)$.

El lema de Zorn proporciona un par maximal $(N,g)$. Si $N\ne X$, el paso de una dimensión construye una extensión propia dominada por $p$, contradicción. Por tanto $N=X$ y $F=g$.
\end{proof}

\begin{teorema}[Extensión real y compleja sin aumentar la norma]\label{v3:hb-norma}
Sea $M$ un subespacio de un espacio normado real o complejo $X$. Cada $f\in M^*$ admite una extensión $F\in X^*$ con $\norm F=\norm f$. No se exige que $M$ sea cerrado ni que $X$ sea completo.
\end{teorema}
\begin{proof}
En el caso real, aplique el teorema anterior a $p(x)=\norm f\norm x$. La cota $F(x)\leq\norm f\norm x$, aplicada también a $-x$, da $|F(x)|\leq\norm f\norm x$. Así $\norm F\leq\norm f$, mientras que restringir a $M$ da la desigualdad contraria.

En el caso complejo, considere los espacios reales subyacentes y extienda $\operatorname{Re}f$ a una funcional real $G$ con $G(x)\leq\norm f\norm x$. Defina
\[
 F(x)=G(x)-iG(ix).
\]
Es aditiva y real homogénea. Además $F(ix)=G(ix)+iG(x)=iF(x)$, de modo que es compleja lineal. Para $m\in M$, $G(m)=\operatorname{Re}f(m)$ y $G(im)=-\operatorname{Im}f(m)$; por tanto $F(m)=f(m)$. Si $F(x)\ne0$, elija $|\lambda|=1$ con $\lambda F(x)=|F(x)|$. Entonces
\[
 |F(x)|=\operatorname{Re}F(\lambda x)=G(\lambda x)\leq\norm f\norm x.
\]
El caso $F(x)=0$ es inmediato. Esto prueba $\norm F\leq\norm f$; la restricción vuelve a dar igualdad. La lectura atribuye esta versión compleja a Bohnenblust, Sobczyk y Suhomlinov; aquí se distingue el argumento real del paso adicional complejo.
\end{proof}
El operador de restricción $X^*\to M^*$ es lineal, contractivo y sobreyectivo. La cota contractiva se debe a que el supremo se toma sobre una bola menor. Puede tener núcleo no trivial: en $\R^2$, la segunda coordenada se anula en el eje horizontal.

\section{Norma, distancia y bidual}\label{lec:v3-02}
\fuentelectura{02}{Primeras aplicaciones del teorema de Hahn--Banach}
\begin{proposicion}[Funcionales normantes]\label{v3:normantes}
Para cada $x\ne0$ existe $\varphi\in X^*$ con $\norm\varphi=1$ y $\varphi(x)=\norm x$. En consecuencia,
\[
 \norm x=\sup_{\norm\varphi\leq1}|\varphi(x)|.
\]
El dual separa puntos y la aplicación canónica $J_X:X\to X^{**}$, $J_Xx(\varphi)=\varphi(x)$, es una isometría lineal.
\end{proposicion}
\begin{proof}
En $\F x$ defina $f(\alpha x)=\alpha\norm x$. Tiene norma uno; extiéndala por Hahn--Banach. Toda funcional de norma a lo sumo uno satisface $|\varphi(x)|\leq\norm x$, y la construida alcanza esa cota. Para $x=0$ la fórmula es directa. Si $u\ne v$, una funcional normante de $u-v$ satisface $\varphi(u)-\varphi(v)=\norm{u-v}\ne0$. La evaluación $J_Xx$ es lineal sobre $X^*$ y acotada por $\norm x$; la fórmula del supremo demuestra $\norm{J_Xx}=\norm x$. Su linealidad respecto de $x$ se verifica evaluando en cualquier $\varphi$.
\end{proof}
\index{bidual}\index{funcional!normante}
No se afirma que toda funcional alcance su norma. Para cada vector se construye una funcional adecuada; intercambiar esas dos elecciones produciría una afirmación distinta y generalmente falsa.

\begin{proposicion}[Medir la distancia a un subespacio]\label{v3:distancia-dual}
Si $M\leq X$ y $d=\dist(x,M)>0$, existe $\varphi\in X^*$ con $\norm\varphi=1$, $\varphi|_M=0$ y $\varphi(x)=d$.
\end{proposicion}
\begin{proof}
Como $x\notin M$, cada vector de $M+\F x$ tiene una única representación $m+\alpha x$. Defina $f(m+\alpha x)=\alpha d$. Si $\alpha\ne0$,
\[
 |f(m+\alpha x)|=|\alpha|d\leq|\alpha|\norm{x+m/\alpha}=\norm{m+\alpha x};
\]
si $\alpha=0$, la desigualdad también vale. Así $\norm f\leq1$. Elija $m_n\in M$ con $\norm{x-m_n}\to d$. Entonces $d=|f(x-m_n)|\leq\norm f\norm{x-m_n}$ y, al pasar al límite y dividir por $d>0$, $\norm f\geq1$. Extender $f$ conserva su norma y sus valores.
\end{proof}

\begin{definicion}
Para $S\subseteq X$ y $A\subseteq X^*$ definimos el anulador y preanulador
\[
 S^\perp=\{\varphi\in X^*: \varphi(s)=0\ \forall s\in S\},\qquad
 {}^\perp A=\{x\in X:\varphi(x)=0\ \forall\varphi\in A\}.
\]
Aquí $\perp$ no denota necesariamente ortogonalidad de Hilbert; indica anulación bajo el apareamiento entre un espacio y su dual.
\end{definicion}
\index{anulador}\index{preanulador}
Son subespacios cerrados: cada uno es una intersección de núcleos de funciones lineales continuas, en el segundo caso evaluaciones. La continuidad también da $M^\perp=(\cl M)^\perp$. Para todo subespacio $M$,
\begin{equation}\label{v3:doble-anulador}
 {}^\perp(M^\perp)=\cl M.
\end{equation}
Una inclusión es inmediata. Si $x\notin\cl M$, la proposición de distancia proporciona una funcional que se anula en $M$ y no en $x$, lo que excluye a $x$ del preanulador. Esta es también la demostración completa de la proposición8 que la lectura deja como ejercicio.

\begin{proposicion}[Sistema biortogonal finito]\label{v3:biortogonal}
Dados $x_1,\ldots,x_m$ linealmente independientes, existen $\varphi_1,\ldots,\varphi_m\in X^*$ con $\varphi_i(x_j)=\delta_{ij}$. Para $x$ en su espacio generado, $x=\sum_j\varphi_j(x)x_j$.
\end{proposicion}
\begin{proof}
En $E=\spanop(x_1,\ldots,x_m)$, las funcionales coordenada son continuas por equivalencia de normas en dimensión finita. Extiéndalas por Hahn--Banach. La fórmula se verifica escribiendo $x=\sum_ja_jx_j$ y aplicando cada coordenada. Alternativamente, se puede aplicar la proposición de distancia a $x_i$ y al espacio generado por los demás vectores, y dividir la funcional resultante por la distancia positiva; ese factor es necesario para obtener valor uno.
\end{proof}

\section{Ejercicios propuestos: extensión algebraica y sublinealidad}
\begin{ejercicio}[Lectura1, ejercicio1]\label{ej:v3-l01-1}
Si $p,q$ son sublineales reales y $\alpha>0$, pruebe que $p+q$ y $\alpha p$ son sublineales.
\end{ejercicio}

\begin{ejercicio}[Lectura1, ejercicio2]\label{ej:v3-l01-2}
Sea $(p_n)$ creciente y sublineal, con límite real $p(x)$ para cada $x$. Pruebe que $p$ es sublineal.
\end{ejercicio}

\begin{ejercicio}[Lectura1, ejercicio3]\label{ej:v3-l01-3}
Dados espacios vectoriales reales $X,Y$, un subespacio $M$ y una aplicación lineal $T:M\to Y$, construya una extensión lineal a $X$.
\end{ejercicio}

\begin{ejercicio}[Lectura1, ejercicio4]\label{ej:v3-l01-4}
Construya una solución de $f(x+y)=f(x)+f(y)$ en $\R$ que no tenga la forma $f(x)=cx$.
\end{ejercicio}


\section{Límites de Banach y convergencia casi uniforme de promedios}
\index{Banach, limite de@Banach, límite de}\index{casi convergencia}
La funcional límite en $c\subseteq\ell^\infty$ tiene norma uno: $|\lim x_n|\leq\norm x_\infty$, y la sucesión constante uno alcanza la cota. Hahn--Banach produce extensiones de norma uno a $\ell^\infty$. Sin embargo, \emph{esas extensiones no son automáticamente invariantes por desplazamientos}. La denominación de la lectura necesita reforzarse antes de resolver su ejercicio5.

\begin{definicion}[Convención corregida]
Un límite de Banach real es una funcional lineal $B:\ell^\infty(\R)\to\R$ positiva, con $B(\mathbf1)=1$ e invariante por el desplazamiento izquierdo $S(x_1,x_2,\ldots)=(x_2,x_3,\ldots)$. Su extensión compleja se define por $B(x+iy)=B(x)+iB(y)$ para sucesiones reales $x,y$.
\end{definicion}
La positividad implica $|B(x)|\leq\norm x_\infty$ para $x$ real. Además $B(e_1)=B(Se_1)=0$ y, por desplazamiento, $B(e_n)=0$ para todo $n$; por continuidad, $B$ se anula en $c_0$. Por tanto extiende el límite usual. También $B(Rx)=B(x)$ para el desplazamiento derecho $R$, pues $SR=I$. En el caso complejo, elija una fase $\lambda$ que haga real no negativo $\lambda B(z)$; entonces $|B(z)|=B(\operatorname{Re}(\lambda z))\leq\norm z_\infty$. La extensión tiene norma uno.

\begin{lema}\label{v3:banach-limites}
Para $x\in\ell^\infty(\R)$ defina
\[
 p(x)=\lim_{k\to\infty}\sup_{n\geq1}\frac1k\sum_{j=0}^{k-1}x_{n+j},\qquad q(x)=-p(-x).
\]
Estos límites existen. La función $p$ es sublineal y los valores posibles de $B(x)$, al variar los límites de Banach, forman exactamente $[q(x),p(x)]$.
\end{lema}
\begin{proof}
Ponga $a_k=\sup_n\sum_{j=0}^{k-1}x_{n+j}$. Separar un bloque en dos da $a_{k+l}\leq a_k+a_l$. Para un $m$ fijo escriba $k=tm+r$, $0\leq r<m$. Entonces $a_k\leq ta_m+r\norm x_\infty$. Al dividir por $k$ y tomar límite superior obtenemos $\limsup a_k/k\leq a_m/m$. Como cada $a_k/k$ es al menos $\inf_m a_m/m$, existe el límite y es ese ínfimo. Es finito, entre $-\norm x_\infty$ y $\norm x_\infty$.

Cada promedio es lineal; tomar supremos y después límites da $p(x+y)\leq p(x)+p(y)$ y homogeneidad positiva. También $p(x+c\mathbf1)=p(x)+c$. Los promedios de $Sx-x$ son $(x_{n+k}-x_n)/k$, de módulo a lo sumo $2\norm x_\infty/k$; por ello $p(Sx-x)=p(x-Sx)=0$.

Hahn--Banach aplicado a $c\mathbf1\mapsto c$, dominado por $p$, da una funcional $B\leq p$. Aplicando la dominación a $-x$ se obtiene $q(x)\leq B(x)$. Si $x\geq0$, $q(x)\geq0$ y $B$ es positiva. Las dos cotas sobre $Sx-x$ dan invariancia. Así existen límites de Banach.

Para alcanzar un valor $c\in[q(x),p(x)]$, defina primero $f(tx+s\mathbf1)=tc+s$. Si $x$ no es constante, la representación es única; si es constante, el intervalo contiene solo esa constante y la definición sigue siendo consistente. Para $t\geq0$, $tc+s\leq tp(x)+s=p(tx+s\mathbf1)$. Para $t<0$, $c\geq-p(-x)$ da la misma dominación. Hahn--Banach extiende $f$ dominada por $p$, y el argumento precedente produce un límite de Banach con $B(x)=c$.

Finalmente, todo límite de Banach satisface $B(x)=B(k^{-1}\sum_{j=0}^{k-1}S^jx)$. La positividad lo coloca entre el ínfimo y el supremo de las entradas de ese promedio. Al pasar $k\to\infty$, $q(x)\leq B(x)\leq p(x)$.
\end{proof}

\begin{ejercicio}[Lectura1, ejercicio5, incisos(i)--(ii)]\label{ej:v3-l01-5}
Una sucesión es casi convergente si todos sus límites de Banach tienen el mismo valor. (i) Estudie la sucesión $(1,0,1,0,\ldots)$. (ii) Pruebe que $x$ es casi convergente a $\alpha$ si y solo si sus promedios $k^{-1}\sum_{j=0}^{k-1}x_{n+j}$ convergen a $\alpha$ uniformemente respecto de $n$. Distinga la definición corregida de la extensión arbitraria utilizada inicialmente en la lectura.
\end{ejercicio}


\section{Ejercicios propuestos: anulación y cocientes}
\begin{ejercicio}[Lectura2, ejercicio1]\label{ej:v3-l02-1}
Si $X^*$ es separable, demuestre que $X$ es separable.
\end{ejercicio}

\begin{ejercicio}[Lectura2, ejercicio2, incisos(i)--(iv)]\label{ej:v3-l02-2}
Pruebe las identidades y contenciones: (i) $\{0\}^\perp=X^*$ y ${}^\perp X^*=\{0\}$; (ii) $M_1\subseteq M_2$ implica $M_1^\perp\supseteq M_2^\perp$; (iii) $(S\cup T)^\perp=S^\perp\cap T^\perp$ y $S^\perp+T^\perp\subseteq(S\cap T)^\perp$; (iv) $S\subseteq{}^\perp A$ equivale a $\varphi(s)=0$ para todo $s\in S$, $\varphi\in A$, y a $A\subseteq S^\perp$.
\end{ejercicio}

\begin{ejercicio}[Lectura2, ejercicio3: examen del enunciado]\label{ej:v3-l02-3}
El original pide, para $M_1\cap M_2=\{0\}$ finito dimensionales, funcionales $\varphi_i$ que se anulen en $M_i$ y sean no nulas en todo vector del otro subespacio. Determine cuándo es posible y formule una versión válida para dimensiones arbitrarias.
\end{ejercicio}

\begin{ejercicio}[Lectura2, ejercicio4]\label{ej:v3-l02-4}
Para $x\in X$, pruebe que $\{\varphi\in X^*: \varphi(x)=\norm x\}$ es no vacío y convexo.
\end{ejercicio}

\begin{ejercicio}[Lectura2, ejercicio5]\label{ej:v3-l02-5}
Si $M$ es un subespacio cerrado de $X$, pruebe que
\[
 \Phi:M^\perp\to(X/M)^*,\quad\Phi(\varphi)(x+M)=\varphi(x),\qquad
 \Psi:X^*/M^\perp\to M^*,\quad\Psi(\varphi+M^\perp)=\varphi|_M
\]
son isomorfismos. Precise sus normas.
\end{ejercicio}


\section{Del álgebra a la geometría}
Hahn--Banach convierte un cálculo en un subespacio pequeño en una funcional sobre todo el espacio. De ahí proceden la fórmula dual de la norma, la recuperación de subespacios cerrados mediante anuladores y las identificaciones de duales de cocientes. La extensión no añade propiedades que no se hayan impuesto: el caso de los límites de Banach distingue con precisión extensión, positividad e invariancia. En el capítulo siguiente usaremos funciones sublineales para producir hiperplanos que separan conjuntos convexos.

\chapter{Convexidad, separación y funciones sublineales}\label{cap:v3-03}
La separación geométrica necesita transformar un conjunto en una función sublineal. La funcional de Minkowski realiza ese paso, pero solo conserva la simetría cuando el conjunto la posee. Desarrollaremos las pruebas de separación antes del teorema del rango cerrado, para que este último pueda utilizar una herramienta ya demostrada. Cerraremos con una segunda construcción de extensiones por funcionales minimales y con la corrección de su fórmula auxiliar.

\section{La funcional de Minkowski}\label{lec:v3-05}
\fuentelectura{05}{Forma geométrica del teorema de Hahn--Banach}
Toda convexidad y toda desigualdad de esta sección se entienden sobre el cuerpo real. En un espacio complejo se pasa al espacio real subyacente y luego se representa una funcional real $g$ como la parte real de $\varphi(x)=g(x)-ig(ix)$. La comprobación de linealidad compleja es la misma del teorema~\ref{v3:hb-norma}.

\begin{proposicion}\label{v3:subnivel}
Para una funcional sublineal $p$ y $U_p=\{x:p(x)<1\}$, el conjunto $U_p$ es convexo y
\[
 tU_p=\{x:p(x)<t\}\quad(t>0),\qquad
 y+U_p=\{x:p(x-y)<1\}.
\]
Si $p\geq0$ y $f$ es lineal real no nula, entonces $f\leq p$ equivale a $f^{-1}(1)\cap U_p=\varnothing$. Si además $p$ es continua, $f\leq p$ implica que $f$ es continua.
\end{proposicion}
\begin{proof}
Para $x,y\in U_p$ y $0<t<1$, $p(tx+(1-t)y)\leq tp(x)+(1-t)p(y)<1$; en los extremos el resultado es inmediato. Las dos identidades se obtienen usando homogeneidad positiva y la definición de traslación.

Si $f\leq p$, un vector con $f(x)=1$ no puede satisfacer $p(x)<1$. Recíprocamente, si $f(x)>p(x)\geq0$, entonces $f(x)>0$ y $u=x/f(x)$ satisface $f(u)=1$, $p(u)<1$, contradicción. Para continuidad, dado $\eps>0$ elija $\delta>0$ tal que $\norm x<\delta$ implique simultáneamente $p(x)<\eps$ y $p(-x)<\eps$. Entonces $-p(-x)\leq f(x)\leq p(x)$ y $|f(x)|<\eps$. No se puede sustituir esta doble cota por $|f(x)|\leq p(x)$ cuando $p$ no es simétrica.
\end{proof}

\begin{teorema}[Funcional de calibración o de Minkowski]\label{v3:minkowski}
Sea $A$ convexo en un espacio normado real $X$, con $0\in\interior A$. Defina
\[
 p_A(x)=\inf\{t>0:x\in tA\}.
\]
Entonces $p_A$ es finita, no negativa, sublineal y continua. Se cumplen
\[
 \interior A=\{p_A<1\},\qquad \cl A=\{p_A\leq1\},\qquad
 \partial A=\{p_A=1\},\qquad X\setminus\cl A=\{p_A>1\}.
\]
La homogeneidad garantizada es $p_A(tx)=tp_A(x)$ para $t\geq0$; si $A=-A$, también $p_A(tx)=|t|p_A(x)$ para escalares reales.
\end{teorema}
\index{Minkowski, funcional de}\index{funcional!de calibracion@de calibración}
\begin{proof}
Elija $r>0$ con $B(0,r)\subseteq A$. Para $t>\norm x/r$, $x/t\in A$, así que $0\leq p_A(x)\leq\norm x/r$. Para $s>0$, el cambio $t=su$ en el ínfimo da $p_A(sx)=sp_A(x)$; para $s=0$ ambos miembros son cero.

Si $x\in aA$, $y\in bA$ con $a,b>0$, entonces
\[
 \frac{x+y}{a+b}=\frac a{a+b}\frac xa+\frac b{a+b}\frac yb\in A.
\]
Por ello $p_A(x+y)\leq a+b$. Elija $a,b$ en sus conjuntos definitorios a distancias arbitrariamente pequeñas de los ínfimos, para obtener subaditividad. Esta y la cota anterior dan
\[
 p_A(x)-p_A(y)\leq p_A(x-y)\leq\norm{x-y}/r,
\]
y al intercambiar $x,y$ se prueba continuidad Lipschitz.

Si $p_A(x)<1$, existe $0<t<1$ con $x=ta$, $a\in A$. Por convexidad,
\[
 x+(1-t)B(0,r)=ta+(1-t)B(0,r)\subseteq A,
\]
de modo que $x$ es interior. Si $x$ es interior y $x\ne0$, existe $s>1$ cercano a uno con $sx\in A$; luego $p_A(x)\leq1/s<1$. Para $x=0$ ya sabemos $p_A(0)=0$. Si $x\in\cl A$, la continuidad y $p_A\leq1$ en $A$ dan $p_A(x)\leq1$. Si $p_A(x)\leq1$, entonces $p_A(tx)<1$ para $0<t<1$; los puntos $tx\in\interior A$ convergen a $x$, así que $x\in\cl A$. Las fórmulas de frontera y exterior se siguen de las dos primeras. Finalmente $A=-A$ da $p_A(-x)=p_A(x)$ y con ello homogeneidad absoluta.
\end{proof}
\begin{observacion}[Una simetría que no puede suponerse]
Para $A=(-1,2)\subseteq\R$, $p_A(x)=x/2$ si $x\geq0$ y $p_A(x)=-x$ si $x<0$. Así $p_A(1)=1/2$ y $p_A(-1)=1$. La homogeneidad absoluta afirmada sin simetría en la lectura5 necesita corregirse; las demás caracterizaciones geométricas anteriores no requieren esa simetría.
\end{observacion}

\section{Separación por hiperplanos}
\begin{teorema}[Separación con interior]\label{v3:separacion-interior}
Sean $A,B$ convexos no vacíos disjuntos, con $\interior A\ne\varnothing$. Existe una funcional real continua no nula $f$ tal que
\[
 f(a)\leq f(b)\quad(a\in A,b\in B).
\]
Para $a\in\interior A$ la desigualdad es estricta.
\end{teorema}
\begin{proof}
El conjunto $C=\interior A-B$ es abierto y convexo, no vacío y no contiene cero. Elija $c_0\in C$ y ponga $U=C-c_0$, $z=-c_0$. Entonces $U$ es abierto, convexo, contiene cero y no contiene $z$. Por el teorema anterior $p_U(z)\geq1$. En $\R z$ defina $h(tz)=t$. Si $t\geq0$, $h(tz)=t\leq tp_U(z)$; si $t<0$, $h(tz)<0\leq p_U(tz)$. Hahn--Banach extiende $h$ a $f\leq p_U$. La continuidad se obtiene usando las cotas para $x$ y $-x$, y $f(z)=1$ asegura que $f\ne0$.

Para $c\in C$, $c-c_0\in U$ y $f(c)=f(c-c_0)+f(c_0)<1-1=0$. Por tanto $f(a)<f(b)$ si $a\in\interior A$, $b\in B$. Falta incluir los puntos restantes de $A$. Elija $a_0\in\interior A$ y $r>0$ con $B(a_0,r)\subseteq A$. Para $a\in A$ y $0<t\leq1$, el vector $a_t=(1-t)a+ta_0$ tiene la bola $B(a_t,tr)$ contenida en $A$; por tanto $a_t\in\interior A$. Como $a_t\to a$, la continuidad da $f(a)\leq f(b)$.
\end{proof}

\begin{teorema}[Separación estricta]\label{v3:separacion-estricta}
Si $A,B$ son convexos no vacíos disjuntos, $A$ es compacto y $B$ cerrado, existe una funcional real continua $f$ con
\[
 \sup_{a\in A}f(a)<\inf_{b\in B}f(b).
\]
En el caso complejo, las desigualdades se expresan mediante $\operatorname{Re}\varphi$ para una funcional compleja continua $\varphi$.
\end{teorema}
\begin{proof}
La función $a\mapsto\dist(a,B)$ es continua, positiva en $A$ y alcanza su mínimo $d>0$ en el compacto $A$. Elija $0<r<d$ y ponga $U=A+B(0,r)$. Es abierto, convexo y disjunto de $B$. El teorema anterior da $f\ne0$ con $f(u)\leq f(b)$ para $u\in U$, $b\in B$. Para $a\in A$, tome el supremo sobre $\norm h<r$ en $f(a+h)\leq f(b)$; resulta
\[
 f(a)+r\norm f\leq f(b).
\]
El supremo de $f(h)$ sobre esa bola es $r\norm f$ por la definición de la norma y la posibilidad de cambiar el signo de $h$. Tomando supremo en $a$ e ínfimo en $b$ se obtiene un salto de tamaño al menos $r\norm f>0$. La conversión compleja se explicó al inicio del capítulo.
\end{proof}
\index{separacion@separación!de conjuntos convexos}
\begin{corolario}[Separación absoluta de un conjunto balanceado]\label{v3:separacion-balanceada}
Sea $A\ne\varnothing$ convexo y balanceado en un espacio real o complejo: $\lambda A\subseteq A$ si $|\lambda|\leq1$. Si $x\notin\cl A$, existen $\varphi\in X^*$ y $\alpha>0$ tales que
\[
 |\varphi(a)|<\alpha<|\varphi(x)|\quad(a\in A).
\]
\end{corolario}
\begin{proof}
Separe estrictamente $\cl A$ del compacto $\{x\}$. En el caso complejo obtenga $\sup_A\operatorname{Re}\varphi<\operatorname{Re}\varphi(x)$. Como $A$ es balanceado, para cada $a$ todas sus rotaciones de módulo uno están en $A$, y $\sup_A\operatorname{Re}\varphi=\sup_A|\varphi|$. En el caso real basta cambiar el signo. Este supremo es no negativo porque $0\in A$. Elija $\alpha$ estrictamente entre el supremo y $\operatorname{Re}\varphi(x)\leq|\varphi(x)|$.
\end{proof}

\section{Ejercicios propuestos: calibración y convexidad}
\begin{ejercicio}[Lectura5, ejercicio1]\label{ej:v3-l05-1}
Calcule la funcional de calibración de la bola abierta unitaria de un espacio normado real.
\end{ejercicio}

\begin{ejercicio}[Lectura5, ejercicio2]\label{ej:v3-l05-2}
Demuestre que $f(x)=\tfrac12\norm x^2$ es convexa.
\end{ejercicio}

\begin{ejercicio}[Lectura5, ejercicio3, incisos(i)--(ii)]\label{ej:v3-l05-3}
Un espacio es estrictamente convexo si $\norm{u+v}<2$ para vectores unitarios \emph{distintos} $u,v$. (i) Pruebe la equivalencia con $\norm{(u+v)/2}<1$. (ii) Pruebe la equivalencia con la convexidad estricta de $x\mapsto\tfrac12\norm x^2$, también entendida para puntos distintos.
\end{ejercicio}


\begin{ejercicio}[Lectura5, ejercicio4: teorema del sándwich, incisos(i)--(iv)]\label{ej:v3-l05-4}
Sea $p:X\to\R$ sublineal, $S\subseteq X$ no vacío y $f:S\to\R$ con $f\leq p$ en $S$. Defina
\[
 g(x)=\inf_{s\in S,\,t\geq0}\{p(x+ts)-tf(s)\}.
\]
Pruebe: (i) $g$ es real, $g\leq p$ y el mismo ínfimo se obtiene con $t>0$; (ii) para una funcional lineal $h$, $h\leq g$ equivale a $h\leq p$ en $X$ y $h\geq f$ en $S$; (iii) si $S$ es convexo y $f$ cóncava, $g$ es sublineal; (iv) bajo esas hipótesis existe una funcional lineal entre las dos cotas. Examine también $S=\varnothing$.
\end{ejercicio}


\begin{ejercicio}[Lectura5, ejercicio5: separación de un abierto]\label{ej:v3-l05-5}
Examine si en la separación estricta puede sustituirse ``$A$ compacto'' por ``$A$ abierto'' manteniendo la misma conclusión.
\end{ejercicio}

\begin{ejercicio}[Lectura5, ejercicio6]\label{ej:v3-l05-6}
Si $H$ es un hiperplano afín cerrado real, pruebe que $H=f^{-1}(\alpha)$ para alguna $f\in X^*\setminus\{0\}$ y $\alpha\in\R$. Describa los lados y sus cierres.
\end{ejercicio}

\begin{ejercicio}[Lectura5, ejercicio7: teorema del soporte]\label{ej:v3-l05-7}
El original supone que $A$ es conexo, con interior no vacío, y pide un hiperplano cerrado por cada $x\notin\interior A$ que deje a $A$ en un semiespacio cerrado. Examine esa hipótesis y pruebe la versión válida.
\end{ejercicio}

\begin{ejercicio}[Lectura5, ejercicio8]\label{ej:v3-l05-8}
Demuestre que un convexo cerrado de un espacio normado real no nulo es la intersección de los semiespacios cerrados que lo contienen.
\end{ejercicio}


\section{Funcionales minimales y una segunda extensión}\label{lec:v3-06}
\fuentelectura{06}{Más sobre funcionales sublineales}
Las funcionales sublineales reales forman un cono ordenado por la comparación punto a punto. Una funcional $p$ es minimal si la única sublineal $q\leq p$ es $p$ misma.
\begin{proposicion}\label{v3:l06-p1}
Para una funcional sublineal $p$ son equivalentes: (i) $p$ es lineal; (ii) $p(x)+p(-x)=0$ para todo $x$; (iii) $p(x)+p(-x)\leq0$ para todo $x$.
\end{proposicion}
\begin{proof}
La linealidad implica(ii), y(ii) implica(iii). Siempre $0=p(0)\leq p(x)+p(-x)$, así que(iii) implica(ii). Si(ii) vale, entonces
\[
 p(x)\leq p(x+y)+p(-y)=p(x+y)-p(y),
\]
lo cual, junto con subaditividad, da aditividad. La homogeneidad positiva y $p(-x)=-p(x)$ dan homogeneidad para todos los escalares reales. Así(ii) implica(i), completando también la prueba propuesta en la lectura.
\end{proof}

\begin{observacion}[Corrección de la funcional auxiliar]
La fórmula original toma el ínfimo simultáneamente sobre $t\geq0$ y $w\in X$. Sin embargo,
\[
 \inf_{t\geq0,\,w\in X}\{p(x+tw)-tp(w)\}=-p(-x).
\]
La desigualdad $\geq$ procede de subaditividad y la igualdad se alcanza tomando $t=1$, $w=-x$. Para $p(x)=|x|$ en $\R$ el resultado es $-|x|$, que no es subaditivo: su valor en $1+(-1)$ es cero, mayor que $-2$. El teorema auxiliar original es falso. La corrección que conserva el método consiste en fijar la dirección $w$ antes de tomar el ínfimo.
\end{observacion}
\begin{lema}\label{v3:auxiliar-fija}
Para una dirección fija $w\in X$, defina
\[
 q_w(x)=\inf_{t\geq0}\{p(x+tw)-tp(w)\}.
\]
Entonces $q_w$ es sublineal, $-p(-x)\leq q_w(x)\leq p(x)$ y $q_w(-w)=-p(w)$.
\end{lema}
\begin{proof}
La cota inferior se obtiene de $tp(w)\leq p(x+tw)+p(-x)$, y la superior tomando $t=0$. En particular la función es real. El cambio $t=\alpha s$ demuestra homogeneidad positiva; para $x=0$ todos los términos del ínfimo son cero. Para $s,t\geq0$,
\[
 q_w(x+y)\leq p(x+y+(s+t)w)-(s+t)p(w)
 \leq[p(x+sw)-sp(w)]+[p(y+tw)-tp(w)].
\]
Los parámetros $s,t$ son independientes, así que elegir aproximantes a ambos ínfimos prueba subaditividad. Para $x=-w$, la cota inferior da $q_w(-w)\geq-p(w)$, y tomar $t=1$ da la desigualdad inversa.
\end{proof}
\begin{teorema}\label{v3:minimales}
Las funcionales sublineales minimales son exactamente las lineales. Toda funcional sublineal domina alguna funcional lineal.
\end{teorema}
\begin{proof}
Si $p$ es lineal y $q\leq p$ es sublineal, entonces $p(x)=-p(-x)\leq-q(-x)\leq q(x)\leq p(x)$, por lo que $q=p$. Si $p$ es minimal, el lema da $p=q_w$ para cada dirección $w$. Evaluar en $-w$ produce $p(-w)=-p(w)$, y la proposición anterior demuestra linealidad.

Para la existencia, considere $\mathcal P=\{q:q\text{ sublineal},q\leq p\}$, no vacío porque contiene $p$. Si $\mathcal C$ es una cadena no vacía, cada $q\in\mathcal C$ satisface $-p(-x)\leq-q(-x)\leq q(x)$; por tanto $h(x)=\inf_{q\in\mathcal C}q(x)$ es real. Es positivamente homogénea y $h(0)=0$. Dados $x,y$ y $\eps>0$, elija $q_1,q_2$ con $q_1(x)<h(x)+\eps/2$, $q_2(y)<h(y)+\eps/2$. Como la familia es una cadena, el menor de $q_1,q_2$ pertenece a ella y sirve simultáneamente para las dos cotas. Su subaditividad da $h(x+y)\leq h(x)+h(y)+\eps$. Así $h$ es sublineal y es una cota inferior en $\mathcal P$. La cadena vacía tiene la cota $p$. Zorn, aplicado al orden inverso, da un elemento minimal de $\mathcal P$; también es minimal entre todas las sublineales porque cualquier funcional menor sigue estando dominada por $p$. Es lineal por lo ya probado.
\end{proof}

\begin{proposicion}[Segunda prueba de Hahn--Banach]\label{v3:hb-minimal}
Si $f:M\to\R$ es lineal y $f\leq p$ en $M$, la función
\[
 q(x)=\inf_{m\in M}\{p(x-m)+f(m)\}
\]
es sublineal, satisface $q\leq p$ en $X$ y $q=f$ en $M$. Una funcional lineal dominada por $q$ extiende a $f$.
\end{proposicion}
\begin{proof}
La desigualdad $-f(m)=f(-m)\leq p(-m)\leq p(x-m)+p(-x)$ da $q(x)\geq-p(-x)$. Tomar $m=0$ da $q\leq p$. Si $x\in M$, entonces $f(x)=f(x-m)+f(m)\leq p(x-m)+f(m)$ y tomar $m=x$ da igualdad $q(x)=f(x)$.

La homogeneidad se obtiene cambiando $m=\alpha n$ para $\alpha>0$, y $q(0)=f(0)=0$. Para $x,y$ elija aproximantes $m,n\in M$ a los dos ínfimos; use
\[
 p(x+y-m-n)+f(m+n)\leq p(x-m)+f(m)+p(y-n)+f(n)
\]
y haga tender a cero el error de aproximación. Así $q$ es sublineal. El teorema de minimalidad produce una funcional lineal $F\leq q$. En $M$, $F\leq f$, y aplicar esa desigualdad a $-m$ da $F\geq f$. Por tanto $F|_M=f$ y $F\leq p$ en todo $X$.
\end{proof}

\section{Alcance de la separación}
La separación débil entre conjuntos no equivale a una distancia positiva entre sus imágenes. La compacidad de uno y la cerradura del otro producen esa distancia; la apertura produce desigualdades estrictas punto a punto. La funcional de Minkowski reúne estas ideas y permite verificar cada hipótesis. Su construcción, y la corrección de la auxiliar en dirección fija, muestran cómo una fórmula de ínfimos puede conservar unas propiedades y perder otras si se cambian sus parámetros.

\chapter{Aplicación abierta, gráfica cerrada y dualidad de operadores}\label{cap:v3-04}
La continuidad de un operador suele verificarse mediante una desigualdad. En este capítulo obtenemos desigualdades a partir de información cualitativa: suprayectividad, cierre de una gráfica o cierre de una imagen. La completitud permite pasar de aproximaciones sucesivas a soluciones exactas. Después estudiaremos cómo esas propiedades se reflejan en el operador sobre los espacios duales.

\section{De las aproximaciones a la aplicación abierta}\label{lec:v3-03}
\fuentelectura{3}{Los teoremas del mapeo abierto y de la gráfica cerrada}
Recordemos el teorema de Baire: un espacio métrico completo no vacío no puede ser unión numerable de cerrados de interior vacío. Una prueba breve aclara su uso: si todos los cerrados $F_n$ tuviesen interior vacío, dentro de cualquier bola abierta podríamos elegir bolas cerradas $\cl B(x_n,r_n)$, con $0<r_n<2^{-n}$, anidadas y disjuntas de $F_n$. Los centros forman una sucesión de Cauchy. Su límite pertenece a todas las bolas cerradas y a ninguno de los $F_n$, contradicción.

\begin{lema}[Levantamiento por aproximaciones]\label{v3:levantamiento}
Sean $X$ Banach, $Y$ normado y $T\in\Bcal(X,Y)$. Si
\[B_Y(0,c)\subset\cl{T(B_X(0,1))}\qquad(c>0),\]
entonces $B_Y(0,c/2)\subset T(B_X(0,1))$.
\end{lema}
\begin{proof}
Por homogeneidad, $B_Y(0,cr)\subset\cl{T(B_X(0,r))}$ para todo $r>0$. Dado $\norm y<c/2$, elijamos $x_1$ con $\norm{x_1}<1/2$ y $\norm{y-Tx_1}<c/4$. Por inducción elegimos $\norm{x_j}<2^{-j}$ de modo que
\[\norm{y-T\sum_{j=1}^n x_j}<c2^{-n-1}.
\]
La serie $x=\sum_jx_j$ converge en $X$. Además,
$\norm x\le\norm{x_1}+\sum_{j\ge2}2^{-j}<1$, y la continuidad de $T$ da $Tx=y$ al tomar límites. El control estricto de la primera norma asegura que llegamos a la bola abierta.
\end{proof}

\begin{teorema}[Aplicación abierta e inversa acotada]\label{v3:abierta}\index{aplicación abierta}\index{inversa acotada}
Una aplicación lineal continua y suprayectiva entre espacios de Banach es abierta. Si además es inyectiva, su inversa es continua.
\end{teorema}
\begin{proof}
Sea $T:X\to Y$ suprayectiva. Si $Y=\{0\}$ el resultado es inmediato. En otro caso,
$Y=\bigcup_{n\ge1}\cl{T(nB_X)}$. Por Baire, alguna de esas cerraduras contiene una bola $B(y_0,r)$. Al restar dos puntos de esta bola obtenemos
\[B(0,2r)\subset\cl{T(nB_X)}-\cl{T(nB_X)}
\subset\cl{T(2nB_X)}.
\]
La segunda inclusión se justifica aproximando cada uno de los dos puntos por imágenes y restando las aproximaciones. Por homogeneidad,
$B(0,r/n)\subset\cl{T(B_X)}$. El lema anterior produce $d>0$ tal que $B(0,d)\subset T(B_X)$. Si $U\subset X$ es abierto, $Tx\in T(U)$ y $x+\eps B_X\subset U$, entonces
\[Tx+\eps d B_Y\subset T(x+\eps B_X)\subset T(U).
\]
Así todo punto de $T(U)$ es interior. Cuando $T$ es biyectiva, para cada abierto $U\subset X$ tenemos $(T^{-1})^{-1}(U)=T(U)$ abierto: ésta es exactamente la continuidad de $T^{-1}$.
\end{proof}

\begin{teorema}[Acotación uniforme]\label{v3:acotacion-uniforme}\index{acotación uniforme}
Sea $X$ Banach y sea $(T_i)_{i\in I}\subset\Bcal(X,Y)$, con $Y$ normado. Si $\sup_i\norm{T_i x}<\infty$ para cada $x\in X$, entonces $\sup_i\norm{T_i}<\infty$.
\end{teorema}
\begin{proof}
Los conjuntos $E_n=\{x:\sup_i\norm{T_ix}\le n\}$ son cerrados, por ser intersecciones de cerrados, y cubren $X$. Por Baire, $B(x_0,r)\subset E_n$ para ciertos $n,r>0$. Si $\norm h<r$, tanto $x_0+h$ como $x_0$ están en $E_n$; por tanto $\norm{T_i h}\le2n$ para todo $i$. Escalando $h$ y tomando supremos se obtiene $\norm{T_i}\le2n/r$. La completitud sólo se usó en el dominio.
\end{proof}

\section{Gráficas y operadores cerrados}
\begin{proposicion}\label{v3:producto-banach}
La fórmula $\norm{(x,y)}=\norm x+\norm y$ define una norma en $X\times Y$. Este producto es Banach si y sólo si ambos factores lo son.
\end{proposicion}
\begin{proof}
Positividad, homogeneidad y desigualdad triangular se siguen sumando las propiedades de las dos normas. Una sucesión es de Cauchy en el producto si y sólo si cada coordenada lo es: una dirección usa $\norm{x_n-x_m}\le\norm{(x_n-x_m,y_n-y_m)}$ y la otra suma las dos estimaciones. Si ambos factores son completos, los límites coordenados dan el límite en el producto. Recíprocamente, una sucesión de Cauchy $(x_n)$ se inserta como $(x_n,0)$; el límite en el producto tiene segunda coordenada cero, de modo que $x_n$ converge en $X$. Lo mismo vale para $Y$.
\end{proof}

\begin{definicion}\index{operador!cerrado}\index{gráfica cerrada}
Una aplicación lineal $T:D(T)\subset X\to Y$ es \emph{cerrada} si su gráfica $G(T)=\{(x,Tx):x\in D(T)\}$ es cerrada en $X\times Y$. Equivalentemente, $x_n\to x$ y $Tx_n\to y$ implican $x\in D(T)$ y $Tx=y$. Es \emph{cerrable} si la cerradura de su gráfica es la gráfica de un operador, llamado $\cl T$.
\end{definicion}
En espacios métricos, ser cerrado equivale a contener los límites de todas las sucesiones convergentes de puntos del conjunto; esto explica la equivalencia anterior. Todo operador continuo de dominio completo $X$ tiene gráfica cerrada, pues $Tx_n\to Tx$ y los límites son únicos. El recíproco necesita completitud.

\begin{ejemplo}[Derivar con la norma uniforme]
En $X=C^1[0,1]$ con norma uniforme, $Df=f'$ define un operador cerrado hacia $C[0,1]$. En efecto, si $f_n\to f$ y $f_n'\to g$ uniformemente, la identidad
\[f_n(t)-f_n(0)=\int_0^t f_n'(s)\,ds\]
pasa al límite, pues el error de las integrales es a lo sumo $\norm{f_n'-g}_\infty$. Resulta $f(t)-f(0)=\int_0^t g(s)\,ds$, luego $f'=g$. Sin embargo, $f_n(t)=\sin(nt)/n$ tiende uniformemente a cero y $\norm{Df_n}_\infty=1$. Por tanto $D$ no es continuo. Aquí $C^1[0,1]$ con norma uniforme no es completo.
\end{ejemplo}

\begin{teorema}[Gráfica cerrada]\label{v3:grafica-cerrada}
Si $X,Y$ son Banach y $T:X\to Y$ es lineal con gráfica cerrada, entonces $T$ es acotado.
\end{teorema}
\begin{proof}
$G(T)$ es un subespacio cerrado del Banach $X\times Y$, luego es Banach. Las proyecciones $P_X:G(T)\to X$, $P_Y:G(T)\to Y$ son lineales y de norma a lo sumo uno. La primera es biyectiva: a cada $x$ le corresponde exactamente $(x,Tx)$. Esto no requiere que $T$ sea inyectiva. Por el teorema de la inversa acotada, $P_X^{-1}$ es continua; entonces $T=P_YP_X^{-1}$ es continua. La hipótesis de inyectividad y el tercer espacio que aparecían en el enunciado original son innecesarios.
\end{proof}

\begin{proposicion}[Criterio de cerrabilidad]\label{v3:cerrabilidad}
Un operador lineal $S:D\subset X\to Z$ es cerrable si y sólo si $x_n\in D$, $x_n\to0$ y $Sx_n\to z$ implican $z=0$.
\end{proposicion}
\begin{proof}
Si la cerradura es una gráfica, no puede contener $(0,z)$ salvo para $z=0$. Recíprocamente, la cerradura de un subespacio es un subespacio. Si contiene $(x,z)$ y $(x,w)$, contiene $(0,z-w)$, que se aproxima por una sucesión $(x_n,Sx_n)$. La condición implica $z=w$. Por ello asigna a cada primera coordenada un único valor; la linealidad de la gráfica da la linealidad del operador correspondiente.
\end{proof}

\begin{teorema}[Estimación de Hörmander]\label{v3:hormander}\index{Hörmander, estimación de}
Sean $X,Y,Z$ Banach, $D$ un subespacio de $X$, $T:D\to Y$ cerrado y $S:D\to Z$ cerrable. Existe $C>0$ tal que
\[\norm{Sx}\le C(\norm x+\norm{Tx})\qquad(x\in D).
\]
\end{teorema}
\begin{proof}
$G(T)$ es Banach. Definamos $U:G(T)\to Z$ por $U(x,Tx)=Sx$. Para probar que es cerrado, supongamos $(x_n,Tx_n)\to(x,Tx)$ en $G(T)$ y $Sx_n\to z$. Entonces $x_n-x\to0$ y $S(x_n-x)\to z-Sx$. El criterio de cerrabilidad implica $z=Sx$. Por el teorema de la gráfica cerrada, $U$ es acotado y
$\norm{Sx}=\norm{U(x,Tx)}\le C\norm{(x,Tx)}$. La lectura original considera el caso particular $D=X$; trabajar en la norma de la gráfica explica la forma de la estimación y conserva ese caso.
\end{proof}

\section{Ejercicios sobre aplicaciones abiertas y gráficas}
\begin{ejercicio}\label{ej:v3-l03-1}
Demuestre que la inversa de una aplicación lineal continua biyectiva entre espacios de Banach es continua.
\end{ejercicio}


\begin{ejercicio}\label{ej:v3-l03-2}
Dos normas $\norm\cdot,\norm\cdot'$ que hacen Banach al mismo espacio son equivalentes si y sólo si existe $C>0$ con $\norm x'\le C\norm x$ para todo $x$.
\end{ejercicio}


\begin{ejercicio}\label{ej:v3-l03-3}
No existe una sucesión compleja $(z_n)$ tal que, para toda sucesión $(x_n)$, $\sum_n|x_n|<\infty$ equivalga a que $(x_nz_n)$ sea acotada.
\end{ejercicio}


\begin{ejercicio}\label{ej:v3-l03-4}
Si $T:X\to Y$ es continua e inyectiva entre espacios de Banach, pruebe que $T^{-1}:\ran T\to X$ es continua si y sólo si $\ran T$ es cerrada.
\end{ejercicio}


\begin{ejercicio}\label{ej:v3-l03-5}
La inversa de una aplicación lineal biyectiva con gráfica cerrada también tiene gráfica cerrada.
\end{ejercicio}


\begin{ejercicio}\label{ej:v3-l03-6}
Si $T$ tiene gráfica cerrada, $x_n\to x$, $y_n\to x$ y ambas sucesiones de imágenes convergen, sus límites son iguales.
\end{ejercicio}


\begin{ejercicio}\label{ej:v3-l03-7}
El núcleo de una aplicación lineal con gráfica cerrada es cerrado.
\end{ejercicio}


\begin{ejercicio}\label{ej:v3-l03-8}
Sean $S,T:X\to Y$ lineales entre espacios normados. (i) Si $S$ es cerrado y $T$ continuo, $S+T$ es cerrado. (ii) ¿Basta que ambos sean cerrados?
\end{ejercicio}


\section{El adjunto sobre el dual}\label{lec:v3-04}
\fuentelectura{4}{Transformaciones adjuntas}
El apareamiento de $X$ y $X^*$ es $\langle x,\phi\rangle=\phi(x)$. Es bilineal incluso sobre $\C$: la linealidad en $x$ es la de $\phi$, y la linealidad en $\phi$ es la de las operaciones puntuales en el dual. No es un producto interno.

\begin{proposicion}[Adjunto de Banach]\label{v3:adjunto-banach}\index{adjunto!de Banach}
Para $T\in\Bcal(X,Y)$ existe un único operador $T':Y^*\to X^*$ tal que
\[\langle Tx,\psi\rangle=\langle x,T'\psi\rangle.
\]
Se define por $T'\psi=\psi\circ T$ y satisface
\[(\alpha S+\beta T)'=\alpha S'+\beta T',\qquad(UT)'=T'U',\qquad\norm{T'}=\norm T.
\]
\end{proposicion}
\begin{proof}
La composición es lineal y $|(T'\psi)(x)|\le\norm\psi\norm T\norm x$, por lo cual pertenece al dual. La aplicación $\psi\mapsto T'\psi$ es lineal por evaluación en cada $x$, y $\norm{T'}\le\norm T$. La identidad de apareamientos determina el valor de $T'\psi$ en todo $x$, lo que prueba unicidad. Las dos reglas algebraicas se comprueban evaluando: por ejemplo, $(UT)'\chi(x)=\chi(UTx)=(T'U'\chi)(x)$. Para la desigualdad inversa, el teorema de Hahn--Banach da
\[\norm{Tx}=\sup_{\norm\psi\le1}|\psi(Tx)|
\le\norm{T'}\norm x.
\]
Tomando el supremo en la bola de $X$ resulta $\norm T\le\norm{T'}$.
\end{proof}
Con producto interno lineal en la primera entrada, $J_Hy=\ip{\,\cdot\,}{y}$ es una \emph{isometría conjugado-lineal} de $H$ en $H^*$. Para operadores entre Hilbert se cumple
\[T'J_K=J_HT^*.
\]
En efecto, ambas expresiones evaluadas en $x$ dan $\ip{Tx}{y}=\ip{x}{T^*y}$. Esto relaciona ambos adjuntos, pero explica por qué el adjunto de Banach depende linealmente de $T$ y el de Hilbert depende conjugado-linealmente. Usaremos $T'$ y $T^*$ para mantener visible esta distinción.

\begin{teorema}[Schauder]\label{v3:schauder}\index{Schauder, teorema de}
Entre espacios de Banach, $T$ es compacto si y sólo si $T'$ es compacto.
\end{teorema}
\begin{proof}
Supongamos primero que $T$ es compacto. Dado $\eps>0$, escojamos $x_1,\ldots,x_m$ en la bola unidad de $X$ cuyas imágenes forman una $\eps/4$-red de $T(B_X)$. La aplicación
\[A:Y^*\to\F^m,\qquad A\psi=(\psi(Tx_1),\ldots,\psi(Tx_m))
\]
es acotada con la norma máximo, pues $\norm{A\psi}_\infty\le\norm T\norm\psi$. Su imagen de la bola unidad es acotada en dimensión finita, por tanto totalmente acotada. Elijamos $\psi_1,\ldots,\psi_k\in B_{Y^*}$ de modo que los $A\psi_j$ formen una $\eps/4$-red de esa imagen. Para $\psi\in B_{Y^*}$ elegimos $j$ con $\norm{A\psi-A\psi_j}_\infty<\eps/4$. Para cada $x\in B_X$ elegimos $i$ con $\norm{Tx-Tx_i}<\eps/4$; entonces
\begin{align*}
|\psi(Tx)-\psi_j(Tx)|
&\le|\psi(Tx-Tx_i)|+|\psi(Tx_i)-\psi_j(Tx_i)|+|\psi_j(Tx_i-Tx)|\\
&<3\eps/4.
\end{align*}
Al tomar el supremo en $x$, las imágenes $T'\psi_j$ dan una $\eps$-red. El dual $X^*$ es completo, de modo que la cerradura de esta imagen totalmente acotada es compacta.

Para el recíproco, lo demostrado aplicado a $T'$ hace compacto a $T''$. Las inserciones canónicas satisfacen $T''J_X=J_YT$: evaluando en $\psi\in Y^*$ ambos lados valen $\psi(Tx)$. Así $J_YT(B_X)$ es relativamente compacto en $Y^{**}$. Como $Y$ es Banach, $J_Y(Y)$ es cerrado y $J_Y$ es una isometría; el compacto que contiene la cerradura de esta imagen se transporta de vuelta a $Y$. Por tanto $T$ es compacto.
\end{proof}

\section{Cocientes y rango cerrado}
\begin{lema}\label{v3:cociente-banach}
Si $X$ es Banach y $M\subset X$ es un subespacio cerrado, el cociente con norma $\norm{x+M}=\inf_{m\in M}\norm{x-m}$ es Banach.
\end{lema}
\begin{proof}
La fórmula define una norma: si el ínfimo es cero, $x\in\cl M=M$, y las demás propiedades se deducen de las de $\norm\cdot$. Sea $(u_n)$ de Cauchy en el cociente. Elijamos una subsucesión $(u_{n_j})$ con $\norm{u_{n_{j+1}}-u_{n_j}}<2^{-j}$. Podemos representar cada diferencia por $x_j\in X$ con $\norm{x_j}<2^{1-j}$. Fijado un representante $x_0$ de $u_{n_1}$, la serie $x_0+\sum_jx_j$ converge en $X$. Sus sumas parciales representan $u_{n_j}$, por lo cual esta subsucesión converge en el cociente. Una sucesión de Cauchy con una subsucesión convergente converge al mismo límite: se aplica la desigualdad triangular pasando por un término de esa subsucesión.
\end{proof}

\begin{teorema}[Rango cerrado]\label{v3:rango-cerrado}\index{rango cerrado}
Si $T\in\Bcal(X,Y)$ y $X,Y$ son Banach, entonces $\ran T$ es cerrada si y sólo si $\ran T'$ es cerrada. En tal caso
\[\ran T'=(\ker T)^\perp.
\]
\end{teorema}
\begin{proof}
Si $\ran T$ es cerrada, el operador inducido $\widetilde T:X/\ker T\to\ran T$, $\widetilde T(x+\ker T)=Tx$, es una biyección continua entre Banach. Su inversa es acotada, luego
\[\norm{x+\ker T}\le C\norm{Tx}.
\]
Sea $\phi\in(\ker T)^\perp$. La fórmula $\chi(Tx)=\phi(x)$ está bien definida porque $Tx=Tx'$ implica $x-x'\in\ker T$. Además, para todo $m\in\ker T$,
\[|\chi(Tx)|=|\phi(x-m)|\le\norm\phi\norm{x-m}.
\]
Tomando ínfimos resulta $|\chi(Tx)|\le C\norm\phi\norm{Tx}$. Hahn--Banach extiende $\chi$ a $\psi\in Y^*$ y entonces $T'\psi=\phi$. La inclusión opuesta es inmediata: $T'\psi$ se anula en $\ker T$. El anulador es cerrado, de modo que $\ran T'$ es cerrada.

Supongamos ahora cerrada $\ran T'$. Sea $Y_0=\cl{\ran T}$ y sea $S:X\to Y_0$ el mismo operador con codominio reducido. Su imagen es densa, por lo que $S'$ es inyectiva: una funcional que se anula en una parte densa se anula en todo $Y_0$. Por Hahn--Banach, toda funcional en $Y_0$ se extiende a $Y$, y por restricción se obtiene la otra inclusión; así $\ran S'=\ran T'$. La inversa acotada aplicada a $S':Y_0^*\to\ran T'$ produce $c>0$ con
\[\norm{S'\phi}\ge c\norm\phi\qquad(\phi\in Y_0^*).
\]
Pongamos $K=\cl{S(B_X)}$, convexo y balanceado. Si $\norm y<c$ y $y\notin K$, la separación balanceada del teorema~\ref{v3:separacion-balanceada} da una funcional no nula $\phi$ tal que
\[\sup_{z\in K}|\phi(z)|<|\phi(y)|.
\]
Pero el supremo de la izquierda es $\norm{S'\phi}$, mientras que $|\phi(y)|\le\norm\phi\norm y<c\norm\phi\le\norm{S'\phi}$, contradicción. Así $cB_{Y_0}\subset\cl{S(B_X)}$. Por el lema~\ref{v3:levantamiento}, $S(B_X)$ contiene una bola en $Y_0$. Escalando cualquier vector de $Y_0$ hasta esa bola concluimos que $S$ es suprayectiva. Por tanto $\ran T=Y_0$ es cerrada. El uso de la cerradura de $S(B_X)$ evita la separación estricta injustificada de un punto de su frontera.
\end{proof}

\section{Ejercicios sobre adjuntos y anuladores}
\begin{ejercicio}\label{ej:v3-l04-1}
Para $T\in\Bcal(X,Y)$ entre Banach son equivalentes: (i) el isomorfismo canónico $X/\ker T\to\ran T$ es topológico; (ii) $\ran T$ es cerrada; (iii) existe $C>0$ con $\norm{x+\ker T}\le C\norm{Tx}$.
\end{ejercicio}


\begin{ejercicio}\label{ej:v3-l04-2}
Examine la afirmación original: si $A$ es convexo y balanceado y $x\notin A$, existen $\phi\in X^*$ y $\alpha>0$ con $|\phi(a)|<\alpha<|\phi(x)|$ para todo $a\in A$.
\end{ejercicio}


\begin{ejercicio}\label{ej:v3-l04-3}
Para $M\le X$ y $T\in\Bcal(X)$, estudie el mapeo de $(TM)^\perp$ en $M^\perp$ inducido por $T'$, y deduzca la invariancia del anulador si $M$ es $T$-invariante.
\end{ejercicio}


\begin{ejercicio}\label{ej:v3-l04-4}
Si $T\in\Bcal(X,Y)$ tiene rango finito, pruebe que $T'$ tiene rango finito y la misma dimensión de imagen.
\end{ejercicio}


\begin{ejercicio}\label{ej:v3-l04-5}
Pruebe: (i) $T'$ es inyectiva si y sólo si $\ran T$ es densa; (ii) si $\ran T'$ es densa en $X^*$, entonces $T$ es inyectiva.
\end{ejercicio}


\begin{ejercicio}\label{ej:v3-l04-6}
Demuestre las cuatro relaciones
\[\ker T'=(\ran T)^\perp,\qquad \ker T={}^\perp(\ran T'),\qquad
\cl{\ran T}={}^\perp(\ker T'),\qquad \ran T'\subset(\ker T)^\perp.
\]
\end{ejercicio}


\begin{ejercicio}\label{ej:v3-l04-7}
Explique con precisión la afirmación $T''|_X=T$.
\end{ejercicio}


Los teoremas de este capítulo convierten propiedades topológicas en estimaciones uniformes. La dualidad añade una segunda perspectiva: densidad e inyectividad se intercambian, la compacidad se preserva y el rango cerrado puede detectarse en el adjunto. Estas herramientas permitirán distinguir los tipos de espectro sin confundir una inversa algebraica con una inversa acotada.

\parte{Espectro y clases de operadores}{La invertibilidad organiza el espectro. Los ejemplos de multiplicación y desplazamiento muestran sus distintas partes; la compacidad y la geometría de Hilbert conducen a descomposiciones y normas de operadores.}
\chapter{Espectro, resolvente y modelos concretos}\label{cap:v3-05}
En dimensión finita, resolver $(T-\lambda I)x=y$ equivale a evitar los valores propios. En dimensión infinita aparecen dos dificultades adicionales: la imagen puede ser densa sin ser todo el espacio, o puede dejar fuera un subespacio detectable por funcionales. La teoría espectral organiza estas posibilidades. En todo el capítulo $X$ será un espacio de Banach complejo \emph{no nulo}; escribiremos $\lambda$ por $\lambda I$ cuando no haya ambigüedad.

\section{Invertibilidad y tipos de espectro}\label{lec:v3-07}
\fuentelectura{7}{El espectro de un operador}
\begin{lema}[Serie de Neumann]\label{v3:neumann}\index{Neumann, serie de}
Si $A\in\Bcal(X)$ y $\norm A<1$, entonces
\[(I-A)^{-1}=\sum_{n=0}^{\infty}A^n,\qquad
\norm{(I-A)^{-1}}\le\frac1{1-\norm A}.
\]
El conjunto de operadores invertibles es abierto. Si $T$ es invertible y $\norm{S-T}<\norm{T^{-1}}^{-1}$, entonces $S$ es invertible.
\end{lema}
\begin{proof}
La serie converge en el Banach $\Bcal(X)$ porque $\sum\norm{A^n}\le\sum\norm A^n<\infty$. Para $P_n=\sum_{k=0}^nA^k$ se tiene $(I-A)P_n=P_n(I-A)=I-A^{n+1}$. La continuidad del producto y $A^{n+1}\to0$ prueban ambas identidades de inversa; la suma geométrica da la cota. Para la última afirmación, escribimos
\[S=T[I+T^{-1}(S-T)]
\]
y aplicamos lo anterior al segundo factor. El producto de dos operadores invertibles tiene inversa en orden inverso: $(ST)(T^{-1}S^{-1})=I=(T^{-1}S^{-1})(ST)$. Esta verificación también resuelve la demostración propuesta en la proposición 2 original. Finalmente, $1=\norm I\le\norm T\norm{T^{-1}}$.
\end{proof}

\begin{definicion}\index{espectro}\index{resolvente}\index{espectro!puntual}\index{espectro!continuo}\index{espectro!residual}
El \emph{resolvente} de $T$ es $\rho(T)=\{\lambda:T-\lambda I\text{ es invertible}\}$ y su \emph{espectro} es $\sigma(T)=\C\setminus\rho(T)$. Usaremos $r(T)$ para el radio espectral, evitando confundirlo con el conjunto $\rho(T)$. El espectro se divide en tres partes disjuntas:
\begin{itemize}
\item $\sigma_p(T)$: $T-\lambda I$ no es inyectivo; sus elementos son valores propios.
\item $\sigma_c(T)$: $T-\lambda I$ es inyectivo, tiene imagen densa y no es suprayectivo.
\item $\sigma_r(T)$: $T-\lambda I$ es inyectivo y su imagen no es densa.
\end{itemize}
\end{definicion}
Si $T-\lambda I$ fuese biyectivo, su inversa sería acotada por aplicación abierta. Por eso las tres alternativas anteriores agotan el espectro. En dimensión finita inyectividad equivale a suprayectividad; allí todo el espectro es puntual.

\begin{ejemplo}[Multiplicación por la coordenada]\label{v3:mult-coordenada}
En $L^1[0,1]$, $Mf(t)=tf(t)$ tiene norma uno y
\[\sigma(M)=\sigma_c(M)=[0,1].
\]
La cota $\norm{Mf}_1\le\norm f_1$ es inmediata; los indicadores normalizados de intervalos próximos a uno muestran que la norma es uno. Si $\lambda\notin[0,1]$, la función $(t-\lambda)^{-1}$ es acotada y multiplicar por ella da la inversa. Para cualquier $\lambda$, la igualdad $(t-\lambda)f(t)=0$ casi en todas partes implica $f=0$, pues el posible conjunto excepcional $\{\lambda\}\cap[0,1]$ tiene medida cero. Si $\lambda\in[0,1]$, para $g\in L^1$ definimos
\[f_\delta(t)=\frac{g(t)}{t-\lambda}\boldsymbol1_{\{|t-\lambda|\ge\delta\}}.
\]
Es integrable porque $|f_\delta|\le|g|/\delta$, y $(M-\lambda)f_\delta\to g$ en $L^1$ por convergencia dominada. La imagen es densa. No contiene la función constante uno: su única preimagen formal sería $1/(t-\lambda)$, no integrable incluso si $\lambda$ es un extremo. Así la imagen es propia. Esta prueba desarrolla la aproximación propuesta en la lectura y evita los errores de signos de sus intervalos de corte.
\end{ejemplo}

\begin{ejemplo}[Corrimiento derecho en $\ell^\infty$]\label{v3:shift-infinito}\index{corrimiento!derecho}
Sea $R(x_1,x_2,\ldots)=(0,x_1,x_2,\ldots)$. Entonces $\norm R=1$ y
\[\sigma(R)=\sigma_r(R)=\{\lambda:|\lambda|\le1\}.
\]
El operador $R-\lambda$ es inyectivo: para $\lambda=0$ esto es directo; para $\lambda\ne0$, la primera coordenada y luego las siguientes fuerzan sucesivamente $x_n=0$. Fuera del disco unidad es invertible por Neumann. Si $|\lambda|<1$, la funcional no nula
\[\phi(y)=\sum_{n\ge1}\lambda^{n-1}y_n
\]
es continua y se anula en $\ran(R-\lambda)$, como se comprueba desplazando el índice de una serie absolutamente convergente. Por tanto esa imagen no es densa.

Si $|\lambda|=1$, la ecuación $(R-\lambda)x=y$ determina por inducción
\[x_n=-\lambda^{-n}\sum_{k=1}^n\lambda^{k-1}y_k.
\]
Tomemos $y_k=\lambda^{1-k}$. Si $\norm z_\infty<1/2$ y $(R-\lambda)x=y+z$, entonces
\[|x_n|=\left|n+\sum_{k=1}^n\lambda^{k-1}z_k\right|
\ge n-n\norm z_\infty>n/2.
\]
No existe tal $x\in\ell^\infty$, y la bola $B(y,1/2)$ no intersecta la imagen. Esto prueba también en la frontera el carácter residual. Los exponentes y el signo de esta recurrencia corrigen las fórmulas de la lectura original.
\end{ejemplo}

\section{Compacidad del espectro y cálculo polinómico}\label{lec:v3-08}
\fuentelectura{8}{Propiedades del espectro de un operador}
\begin{teorema}\label{v3:espectro-compacto}
El espectro de todo $T\in\Bcal(X)$ es un compacto no vacío, contenido en $\{z:|z|\le\norm T\}$.
\end{teorema}
\begin{proof}
La apertura de los invertibles implica que $\rho(T)$ es abierto. Si $|z|>\norm T$, la factorización $zI-T=z(I-T/z)$ y Neumann dan invertibilidad y
\[\norm{(zI-T)^{-1}}\le\frac1{|z|-\norm T}.
\]
Así el espectro es cerrado y acotado, por tanto compacto. Para la no vacuidad, supongamos que fuera vacío y pongamos $R(z)=(zI-T)^{-1}$ para todo $z$. La identidad de inversas da
\[R(z)-R(w)=(w-z)R(z)R(w).
\]
Además, la expansión de Neumann alrededor de cualquier $w$ prueba que $R(z)\to R(w)$ en norma. Dividiendo la identidad anterior por $z-w$ obtenemos $R'(w)=-R(w)^2$. Para cada $\Phi\in\Bcal(X)^*$, la función escalar $f(z)=\Phi(R(z))$ es entera. La cota anterior la acota fuera de un disco, y la continuidad la acota dentro. Por Liouville es constante, y como $f(z)\to0$ cuando $|z|\to\infty$, es idénticamente cero. Hahn--Banach, aplicado al espacio $\Bcal(X)$, implica $R(z)=0$, imposible porque $R(z)(zI-T)=I\ne0$. La hipótesis $X\ne\{0\}$ es la que asegura esta última desigualdad.
\end{proof}
La demostración sólo usa el teorema escalar de Liouville, que se deduce de la fórmula de Cauchy: para una función entera acotada por $M$, $|f'(z)|\le M/R$ en circunferencias de radio arbitrario $R$, luego $f'=0$.

\begin{lema}\label{v3:producto-conmutativo}
Un producto finito de operadores que conmutan dos a dos es invertible si y sólo si cada factor lo es.
\end{lema}
\begin{proof}
Basta tratar dos factores $A,B$. Una dirección usa $(AB)^{-1}=B^{-1}A^{-1}$. Si $AB$ es invertible, $A$ y $B$ conmutan con $AB$, luego también con $(AB)^{-1}$: de $AC=CA$ se deduce $C^{-1}A=AC^{-1}$ multiplicando por $C^{-1}$ a ambos lados. Así $B(AB)^{-1}$ es a la vez inversa derecha e izquierda de $A$; análogamente $A(AB)^{-1}$ lo es de $B$. La inducción da el caso finito.
\end{proof}

\begin{teorema}[Aplicación espectral polinómica]\label{v3:espectral-polinomios}\index{aplicación espectral}
Para todo polinomio complejo $p$, $\sigma(p(T))=p(\sigma(T))$.
\end{teorema}
\begin{proof}
Si $p=c$ es constante, $p(T)=cI$ tiene espectro $\{c\}$, igual a $p(\sigma(T))$ por no vacuidad. Si $p$ tiene grado $m\ge1$, para cada $\mu$ el teorema fundamental del álgebra da
\[p(z)-\mu=a\prod_{j=1}^m(z-\lambda_j),\qquad a\ne0.
\]
Sustituyendo $T$, los factores conmutan. El lema anterior dice que $p(T)-\mu$ es invertible exactamente cuando todos los $T-\lambda_j$ lo son, es decir, cuando ningún punto de $\sigma(T)$ satisface $p(\lambda)=\mu$. Esto equivale a $\mu\notin p(\sigma(T))$ y prueba la igualdad.
\end{proof}

\begin{teorema}[Espectro y dualidad]\label{v3:espectro-adjunto}
Para el adjunto de Banach se cumple
\[\sigma(T')=\sigma(T),\qquad
\sigma_r(T)\subset\sigma_p(T')\subset\sigma_p(T)\cup\sigma_r(T),\qquad
\sigma_c(T')\subset\sigma_c(T).
\]
\end{teorema}
\begin{proof}
Si $A$ es invertible, $A'$ tiene inversa $(A^{-1})'$ por la regla de composición. Si $A'$ es invertible, su imagen densa implica que $A$ es inyectiva, y su núcleo cero implica que $\ran A$ es densa, por el ejercicio~\ref{ej:v3-l04-5}. Además $\ran A'$ es cerrada, así que $\ran A$ es cerrada por el teorema~\ref{v3:rango-cerrado}. Es, pues, todo $X$. Aplicando esto a $A=T-\lambda$ y usando $A'=T'-\lambda$ se obtiene la igualdad de espectros.

Si $\lambda\in\sigma_r(T)$, la imagen no densa de $T-\lambda$ tiene un anulador no nulo; éste es $\ker(T'-\lambda)$, luego $\lambda\in\sigma_p(T')$. A la inversa, un vector no nulo de ese núcleo demuestra que la imagen de $T-\lambda$ no es densa. Por tanto $\lambda$ no puede estar en $\sigma_c(T)$ ni en el resolvente, y pertenece a $\sigma_p(T)\cup\sigma_r(T)$. Finalmente, si $\lambda\in\sigma_c(T')$, la inyectividad de $T'-\lambda$ implica densidad de la imagen de $T-\lambda$, y la densidad de la imagen del primero implica inyectividad del segundo. La igualdad de espectros excluye la suprayectividad, de modo que $\lambda\in\sigma_c(T)$. No se ha confundido el anulador cero con densidad en norma de un subespacio del dual: ese paso, empleado sin hipótesis en la prueba original, se sustituye aquí por rango cerrado.
\end{proof}

\section{Ejercicios de la lectura 7}
\begin{ejercicio}\label{ej:v3-l07-1}
Sea $(a_n)$ acotada y $Dx=(a_nx_n)$ en $\ell^1$. Calcule su espectro.
\end{ejercicio}


\begin{ejercicio}\label{ej:v3-l07-2}
Demuestre que el corrimiento derecho bilateral en $\ell^2(\Z)$ es invertible y determine su inversa.
\end{ejercicio}


\begin{ejercicio}\label{ej:v3-l07-3}
Para $f\in C[0,1]$, sea $Tg(t)=f(1-t)g(t)$ en $L^1[0,1]$. Pruebe que $\sigma(T)=f([0,1])$ y determine los valores propios.
\end{ejercicio}


\begin{ejercicio}\label{ej:v3-l07-4}
Los vectores propios asociados a valores propios \emph{distintos} $\lambda_1,\ldots,\lambda_m$ son linealmente independientes.
\end{ejercicio}


\begin{ejercicio}\label{ej:v3-l07-5}
Si $T\in\Bcal(X)$ es invertible, pruebe $\sigma(T^{-1})=\{\lambda^{-1}:\lambda\in\sigma(T)\}$.
\end{ejercicio}


\begin{ejercicio}\label{ej:v3-l07-6}
Para el adjunto de Hilbert, pruebe $\sigma(T^*)=\{\cl\lambda:\lambda\in\sigma(T)\}$.
\end{ejercicio}


\begin{ejercicio}\label{ej:v3-l07-7}
Para $L(x_1,x_2,\ldots)=(x_2,x_3,\ldots)$ en $\ell^1$, pruebe (i) $\sigma_p(L)=\{|z|<1\}$; (ii) $\sigma_c(L)=\{|z|=1\}$; (iii) $\sigma_r(L)=\varnothing$.
\end{ejercicio}


\section{Ejercicios de la lectura 8}
\begin{ejercicio}\label{ej:v3-l08-1}
Demuestre que un producto finito de operadores que conmutan dos a dos es invertible si y sólo si cada factor lo es.
\end{ejercicio}


\begin{ejercicio}\label{ej:v3-l08-2}
Pruebe: (i) $T^2=T$ implica $\sigma(T)\subset\{0,1\}$; (ii) $T$ nilpotente implica $\sigma(T)=\{0\}$.
\end{ejercicio}


\begin{ejercicio}\label{ej:v3-l08-3}
Si $T_n\to T$ en norma, $\lambda_n\to\lambda$ y $\lambda_n\in\sigma(T_n)$, pruebe $\lambda\in\sigma(T)$.
\end{ejercicio}


\begin{ejercicio}\label{ej:v3-l08-4}
Para los corrimientos unilaterales $L,R$ en $\ell^\infty$, demuestre $\sigma(LR)\ne\sigma(RL)$.
\end{ejercicio}


\begin{ejercicio}\label{ej:v3-l08-5}
Si $(S\times T)(x,y)=(Sx,Ty)$, pruebe $\sigma(S\times T)=\sigma(S)\cup\sigma(T)$.
\end{ejercicio}


\begin{ejercicio}\label{ej:v3-l08-6}
Si $T_n\to T$ en norma, para todo $\eps>0$ existe $N$ con $\sigma(T_n)\subset\sigma(T)+B(0,\eps)$ cuando $n\ge N$.
\end{ejercicio}


\begin{ejercicio}\label{ej:v3-l08-7}
Todo compacto no vacío $K\subset\C$ es el espectro de un operador acotado en un espacio de Banach.
\end{ejercicio}


\begin{ejercicio}\label{ej:v3-l08-8}
Los operadores acotados inyectivos con imagen cerrada forman un abierto de $\Bcal(X)$.
\end{ejercicio}


\section{Radio espectral y crecimiento de potencias}\label{lec:v3-09}
\fuentelectura{9}{El radio espectral}
\begin{definicion}\index{radio espectral}
Definimos $r(T)=\max\{|\lambda|:\lambda\in\sigma(T)\}$. El máximo existe por compacidad y $0\le r(T)\le\norm T$. Un operador con $r(T)=0$ se llama \emph{cuasinilpotente}.
\end{definicion}
Un nilpotente no nulo muestra que la desigualdad con la norma puede ser estricta: en los polinomios de grado a lo sumo $m$, la derivada tiene potencia $m+1$ igual a cero, aunque no es el operador cero si $m\ge1$.

\begin{teorema}[Fórmula de Gelfand]\label{v3:gelfand-radio}\index{Gelfand!fórmula del radio}
Para $T\in\Bcal(X)$,
\[r(T)=\lim_{n\to\infty}\norm{T^n}^{1/n}.
\]
\end{teorema}
\begin{proof}
Aplicación espectral da $\lambda^n\in\sigma(T^n)$ para $\lambda\in\sigma(T)$; luego $r(T)\le\norm{T^n}^{1/n}$ para todo $n$. Falta la cota para el límite superior.

Sea $D=\{z:|z|<1/r(T)\}$, interpretado como $\C$ si $r(T)=0$. Para $z\ne0$ en $D$, $1/z\notin\sigma(T)$ y $I-zT$ es invertible. En $z=0$ también lo es. El mismo argumento local de Neumann usado para el resolvente muestra que $z\mapsto(I-zT)^{-1}$ es holomorfa en norma en $D$. Para cada $\Phi\in\Bcal(X)^*$, la función escalar
\[g_\Phi(z)=\Phi((I-zT)^{-1})
\]
es holomorfa allí. Cerca de cero, Neumann da $g_\Phi(z)=\sum_{n\ge0}\Phi(T^n)z^n$. Por unicidad de los coeficientes de Taylor y el teorema de expansión de una función holomorfa en un disco, esta serie converge para todo $z\in D$. Este último hecho se obtiene de la fórmula de Cauchy en cualquier circunferencia compactamente contenida en $D$; no se presupone convergencia de una serie de operadores fuera del disco $|z|<1/\norm T$.

Fijemos $z\in D\setminus\{0\}$. La convergencia escalar implica que $(z^n\Phi(T^n))_n$ es acotada para cada $\Phi$. En el Banach $\Bcal(X)^*$ consideramos las funcionales $F_n(\Phi)=\Phi(z^nT^n)$, cuya norma es $\norm{z^nT^n}$ por la inserción isométrica en el bidual. Acotación uniforme proporciona $C_z<\infty$ con $|z|^n\norm{T^n}\le C_z$. Por tanto
\[\limsup_n\norm{T^n}^{1/n}\le |z|^{-1}.
\]
Si $r(T)>0$, hacemos $|z|\uparrow1/r(T)$; si $r(T)=0$, hacemos $|z|\to\infty$. En ambos casos el límite superior es a lo sumo $r(T)$, y la cota inicial concluye. Aquí las funcionales actúan sobre $\Bcal(X)$, no sobre $X$; además el argumento correcto es $I-zT$, dos precisiones necesarias en la prueba original.
\end{proof}

\begin{ejemplo}[Cuasinilpotente que no es nilpotente]
En $\ell^\infty$, sea $Tx=(0,x_1,x_2/2,x_3/3,\ldots)$. Por inducción,
\[(T^nx)_k=\begin{cases}0,&k\le n,\\
\displaystyle\frac{x_{k-n}}{(k-1)(k-2)\cdots(k-n)},&k>n.
\end{cases}
\]
El denominador mínimo es $n!$, alcanzado cuando $k=n+1$. Así $\norm{T^n}\le1/n!$ y la igualdad se alcanza en $e_1$, pues $T^ne_1=e_{n+1}/n!$. Además $(n!)^{1/n}\to\infty$: al menos $\lfloor n/2\rfloor$ factores son mayores o iguales que $n/2$. La fórmula de Gelfand da $r(T)=0$, de modo que $\sigma(T)=\{0\}$. El operador es inyectivo y su imagen tiene primera coordenada cero, por tanto cero es residual. Ninguna potencia es cero, como muestra la misma fórmula en $e_1$.
\end{ejemplo}

\section{Ejercicios de la lectura 9}
\begin{ejercicio}\label{ej:v3-l09-1}
Si $ST=TS$, pruebe $r(ST)\le r(S)r(T)$ y $r(S+T)\le r(S)+r(T)$.
\end{ejercicio}


\begin{ejercicio}\label{ej:v3-l09-2}
Pruebe $r(T^m)=r(T)^m$ para todo entero $m\ge1$.
\end{ejercicio}


\begin{ejercicio}\label{ej:v3-l09-3}
Para una matriz compleja $A$ con la norma de operador euclídea, pruebe $\norm A^2=r(A^*A)$.
\end{ejercicio}


\begin{ejercicio}\label{ej:v3-l09-4}
Demuestre que $T\mapsto r(T)$ es semicontinua superior para la norma de operadores.
\end{ejercicio}


\section{Ampliación: multiplicadores y corrimientos bilaterales}
La siguiente formulación desarrolla el modelo presentado por \textbf{Gerardo Ramos Vázquez}, \emph{Los operadores de multiplicación en espacios de medida $\sigma$-finita} \cite{ramos2018}, alojado en el sitio de Egor Maximenko. La atribución corresponde al autor de esa presentación; las demostraciones y ejercicios siguientes son desarrollos de esta edición.

\begin{teorema}\label{v3:multiplicadores}\index{operador!de multiplicación}\index{rango esencial}
Sea $(\Omega,\mu)$ un espacio de medida $\sigma$-finito con $L^2(\mu)\ne\{0\}$ y $a\in L^\infty(\mu)$. El operador $M_af=af$ satisface
\[\norm{M_a}=\norm a_\infty,\quad M_a^*=M_{\cl a},\quad
\sigma(M_a)=\operatorname{essran}(a),
\]
donde $\lambda\in\operatorname{essran}(a)$ significa que $\mu\{|a-\lambda|<\eps\}>0$ para todo $\eps>0$. Sus valores propios son los $\lambda$ con $\mu\{a=\lambda\}>0$; el resto del espectro es continuo y el residual es vacío.
\end{teorema}
\begin{proof}
La desigualdad $\norm{af}_2\le\norm a_\infty\norm f_2$ da una cota. Si $0<c<\norm a_\infty$, el conjunto $\{|a|>c\}$ tiene medida positiva. Por $\sigma$-finitud contiene un subconjunto $E$ de medida positiva finita. El vector $f=\boldsymbol1_E/\sqrt{\mu(E)}$ tiene norma uno y $\norm{M_af}_2\ge c$. Al hacer $c\uparrow\norm a_\infty$ se obtiene la igualdad; si la norma es cero no hace falta este paso. La fórmula del adjunto se deduce de
$\int af\cl g\,d\mu=\int f\cl{\cl a g}\,d\mu$.

Si $\lambda$ no está en el rango esencial, $|a-\lambda|\ge\eps$ casi en todas partes para algún $\eps>0$, y $M_{1/(a-\lambda)}$ es una inversa acotada. Si está, los conjuntos $\{|a-\lambda|<1/n\}$ contienen subconjuntos $E_n$ de medida positiva finita. Sus indicadores normalizados $f_n$ satisfacen $\norm{(M_a-\lambda)f_n}_2\le1/n$, lo que excluye una inversa acotada. Esto prueba la igualdad de espectros.

La ecuación $(a-\lambda)f=0$ equivale a que $f$ esté soportada, salvo conjuntos nulos, en $\{a=\lambda\}$. Por $\sigma$-finitud existe allí un indicador integrable no nulo exactamente cuando ese conjunto tiene medida positiva. Si su medida es cero, el operador es inyectivo; además $\ker(M_a-\lambda)^*=\ker M_{\cl a-\cl\lambda}=0$, de modo que su imagen es densa. Si $\lambda$ está en el espectro, esa imagen no es todo el espacio y el punto es continuo. Por ello nunca aparece espectro residual.
\end{proof}

\begin{ejercicio}[Complemento: norma y espectro de una proyección]\label{ej:v3-ampl-multiplicador}
Sea $E\subset\Omega$ medible y suponga que $E$ y su complemento tienen medida positiva. Determine norma, espectro y espacios propios de $M_{\boldsymbol1_E}$ en $L^2(\mu)$.
\end{ejercicio}


\begin{ejercicio}[Complemento: espectro del corrimiento bilateral]\label{ej:v3-ampl-shift}
Para $R$ en $\ell^2(\Z)$ del ejercicio~\ref{ej:v3-l07-2}, pruebe $\sigma(R)=\sigma_c(R)=\{|z|=1\}$.
\end{ejercicio}


Los modelos de multiplicación y corrimiento muestran que conocer el conjunto espectral no basta para conocer el operador: el mismo disco puede tener espectro puntual o residual según el espacio y la dirección del corrimiento. La fórmula de Gelfand añade una información cuantitativa, el crecimiento asintótico de las potencias. En el siguiente capítulo la compacidad reducirá gran parte de la teoría espectral a espacios propios de dimensión finita.

\chapter{Operadores compactos y descomposición espectral}\label{cap:v3-06}
La compacidad conserva una parte del comportamiento de las matrices: fuera de cero, el espectro está formado por valores propios de multiplicidad finita. Para operadores normales en Hilbert, los espacios propios son ortogonales y permiten reconstruir el operador mediante una serie convergente en norma. Desarrollaremos primero la alternativa de Fredholm en Banach y después la geometría que hace posible esa reconstrucción.

\section{Espectro de un compacto en Banach}\label{lec:v3-10}
\fuentelectura{10}{El espectro de un operador compacto}
\begin{lema}[Riesz]\label{v3:riesz-lema}
Si $M$ es un subespacio cerrado propio de un espacio normado $X$ y $0<\theta<1$, existe $x\in X$ con $\norm x=1$ y $\dist(x,M)>\theta$.
\end{lema}
\begin{proof}
Elijamos $y\notin M$ y $d=\dist(y,M)>0$. Por la definición de ínfimo, existe $m\in M$ con $\norm{y-m}<d/\theta$. Para $x=(y-m)/\norm{y-m}$ y $z\in M$,
\[\norm{x-z}=\frac{\norm{y-(m+\norm{y-m}z)}}{\norm{y-m}}
\ge\frac d{\norm{y-m}}>\theta.
\]
\end{proof}
Aplicándolo sucesivamente a spans finitos se obtiene, en un espacio de dimensión infinita, una sucesión unitaria cuyos términos distan entre sí más de $1/2$. Por tanto la bola unidad no es relativamente compacta y la identidad no es compacta.

\begin{proposicion}\label{v3:compacto-cero}
Si $X$ es un Banach complejo de dimensión infinita y $K\in\Kcal(X)$, entonces $0\in\sigma(K)$. Para cada $\eps>0$ hay sólo un número finito de valores propios \emph{distintos} de módulo al menos $\eps$. Cada espacio propio correspondiente a un valor propio no nulo tiene dimensión finita.
\end{proposicion}
\begin{proof}
Si $K$ fuera invertible, $I=K^{-1}K$ sería compacto, contradiciendo el lema anterior. Supongamos que existen valores propios distintos $\lambda_n$ con $|\lambda_n|\ge\eps$, y escojamos vectores propios $v_n$. Son linealmente independientes por el ejercicio~\ref{ej:v3-l07-4}. Sea $M_n=\spanop(v_1,\ldots,v_n)$. Por Riesz, podemos elegir $x_n\in M_n$ unitario con $\dist(x_n,M_{n-1})>1/2$ para $n\ge2$, y cualquier $x_1$ unitario en $M_1$. Como $(K-\lambda_n)M_n\subset M_{n-1}$,
\[K(x_n/\lambda_n)=x_n+z_n,\qquad z_n\in M_{n-1}.
\]
Si $n>m$, entonces $x_m+z_m-z_n\in M_{n-1}$, y
$\norm{K(x_n/\lambda_n)-K(x_m/\lambda_m)}>1/2$. Esto contradice compacidad porque $\norm{x_n/\lambda_n}\le1/\eps$.

Si $E_\lambda=\ker(K-\lambda)$ fuese infinito-dimensional para $\lambda\ne0$, su bola unidad no sería relativamente compacta, mientras que $K|_{E_\lambda}=\lambda I$ la haría relativamente compacta. Por tanto su dimensión es finita. El requisito $\lambda\ne0$ falta en el corolario original y es indispensable: el operador cero tiene núcleo infinito-dimensional.
\end{proof}

\begin{proposicion}\label{v3:compacto-rango}
Si $K$ es compacto y $\lambda\ne0$, la imagen de $A=K-\lambda I$ es cerrada y existe $C>0$ tal que
\[\dist(x,\ker A)\le C\norm{Ax}\qquad(x\in X).
\]
\end{proposicion}
\begin{proof}
Si no existiese esa constante, en el cociente $X/\ker A$ encontraríamos clases de norma uno cuyas imágenes bajo $A$ tienden a cero. Elijamos representantes $x_n$ con $\dist(x_n,\ker A)=1$, $\norm{x_n}\le2$ y $Ax_n\to0$. Por compacidad, pasando a una subsucesión, $Kx_n\to y$. Como $\lambda x_n=Kx_n-Ax_n$, se tiene $x_n\to x=y/\lambda$. La continuidad de $A$ implica $Ax=0$, pero la continuidad de la función distancia da $\dist(x,\ker A)=1$, contradicción.

Para ver el cierre de la imagen, sea $Ax_n\to y$. La cota permite reemplazar cada $x_n$ por $u_n$ en la misma clase con $\norm{u_n}\le C\norm{Ax_n}+1$, de modo que $(u_n)$ es acotada y $Au_n=Ax_n$. Alguna subsucesión satisface $Ku_n\to z$; entonces $u_n\to(z-y)/\lambda=:u$, y por continuidad $Au=y$. Así $y\in\ran A$.
\end{proof}

\begin{proposicion}\label{v3:compacto-sobreyectivo}
Si $K$ es compacto, $\lambda\ne0$ y $K-\lambda I$ es suprayectivo, entonces es inyectivo.
\end{proposicion}
\begin{proof}
Supongamos $A=K-\lambda I$ suprayectivo y elijamos $x_0\ne0$ en $\ker A$. Podemos construir $x_{j+1}$ con $Ax_{j+1}=x_j$. Entonces $A^jx_j=x_0\ne0$ y $A^{j+1}x_j=0$; los subespacios cerrados $N_j=\ker A^j$ forman una cadena estrictamente creciente. Por Riesz elegimos $u_j\in N_j$ unitario con $\dist(u_j,N_{j-1})>1/2$. Para $j>i$, $Au_j-Au_i-\lambda u_i\in N_{j-1}$, por lo cual
\[\norm{Ku_j-Ku_i}
=\norm{\lambda u_j+(Au_j-Au_i-\lambda u_i)}>|\lambda|/2.
\]
Las imágenes de la sucesión acotada $(u_j)$ no tienen subsucesión de Cauchy, contradicción.
\end{proof}

\begin{teorema}[Alternativa de Fredholm]\label{v3:fredholm}\index{Fredholm, alternativa de}
Para $K\in\Kcal(X)$ y $\lambda\ne0$, $K-\lambda I$ es inyectivo si y sólo si es suprayectivo. Además
\[\ker(K-\lambda I)\ne0\quad\Longleftrightarrow\quad\ker(K'-\lambda I)\ne0.
\]
Por tanto todo punto no nulo del espectro de $K$ es un valor propio.
\end{teorema}
\begin{proof}
Ya probamos que suprayectividad implica inyectividad. Sea ahora $A=K-\lambda I$ inyectivo. Por rango cerrado y la identidad de anuladores,
\[\ran A'=(\ker A)^\perp=X^*.
\]
Por Schauder, $K'$ es compacto; aplicando a $K'-\lambda I=A'$ la proposición anterior, su suprayectividad implica que es inyectivo. La identidad $\ker A'=(\ran A)^\perp$ da densidad de $\ran A$, y ésta es cerrada por la proposición~\ref{v3:compacto-rango}; así es todo $X$. Hemos probado la equivalencia.

Si $\ker A\ne0$, $A$ no es suprayectivo. Como su imagen es cerrada propia, Hahn--Banach da un anulador no nulo, o sea $\ker A'\ne0$. Si $\ker A'\ne0$ pero $\ker A=0$, el primer párrafo haría invertible a $A$, y por tanto a $A'$, contradicción. Finalmente, si $\lambda\ne0$ no es un valor propio, $A$ es biyectivo y su inversa es acotada; así $\lambda$ está en el resolvente.
\end{proof}
En consecuencia, los puntos espectrales no nulos son a lo sumo numerables y su único posible punto de acumulación es cero: para cada $m$, sólo hay un número finito de ellos fuera de $B(0,1/m)$. Si son infinitos, una enumeración sin repeticiones converge a cero. La imagen de $K-\lambda I$ también admite la descripción de compatibilidad
\[y\in\ran(K-\lambda I)\quad\Longleftrightarrow\quad
\phi(y)=0\text{ para toda }\phi\in\ker(K'-\lambda I),
\]
porque está cerrada y es el preanulador de ese núcleo. Esto concreta las condiciones necesarias para resolver una ecuación no invertible.

\section{Ejercicios de Fredholm}
\begin{ejercicio}\label{ej:v3-l10-1}
La lectura afirma, para $S$ acotado y $K$ compacto, que (i) $\dim\ker(S+K)<\infty$ y (ii) $\codim\ran(S+K)<\infty$. Examine el enunciado y pruebe una versión válida.
\end{ejercicio}
\begin{solucion}
Ambos incisos son falsos sin una hipótesis sobre $S$: en un espacio infinito-dimensional, $S=K=0$ tiene núcleo de dimensión infinita e imagen $\{0\}$ de codimensión infinita. Una versión suficiente y natural exige que $S$ sea invertible. Entonces
\[S+K=S(I+S^{-1}K),
\]
donde $C=S^{-1}K$ es compacto. (i) El núcleo de $I+C$ es el espacio propio de $C$ para $-1$, de dimensión finita por la proposición~\ref{v3:compacto-cero}. Es también el núcleo de $S+K$. (ii) La imagen de $I+C$ es cerrada y su anulador es $\ker(I+C')$, de dimensión finita. Para un subespacio cerrado $M$, $(X/M)^*\cong M^\perp$; si ese dual es finito-dimensional, $X/M$ también lo es, pues un conjunto linealmente independiente de $n$ vectores admite $n$ funcionales biortogonales por Hahn--Banach. Por tanto la codimensión es finita. La biyección $S$ preserva codimensión y cierre, dando el inciso (ii) para $S+K$.
\end{solucion}

\begin{ejercicio}\label{ej:v3-l10-2}
Si $K$ es compacto y $r(K)>0$, existe un valor propio de módulo $r(K)$.
\end{ejercicio}
\begin{solucion}
La compacidad no vacía de $\sigma(K)$ asegura que el máximo del módulo se alcanza en algún $\lambda\in\sigma(K)$. Como $|\lambda|=r(K)>0$, el teorema~\ref{v3:fredholm} afirma que ese punto es un valor propio.
\end{solucion}

\begin{ejercicio}[Segunda alternativa de Fredholm]\label{ej:v3-l10-3}
Para $\lambda\ne0$, demuestre que o bien $(K-\lambda)x=y$ tiene una única solución para cada $y$, o bien $K'\phi=\lambda\phi$ tiene una solución no nula.
\end{ejercicio}
\begin{solucion}
Si $\ker(K'-\lambda)=0$, el teorema de Fredholm da $\ker(K-\lambda)=0$ y suprayectividad de $K-\lambda$. Existe entonces una única solución para cada dato. Si ese núcleo del adjunto no es cero, obtenemos la segunda alternativa. No pueden ocurrir ambas: si $K-\lambda$ es invertible, también lo es su adjunto, cuyo núcleo debe ser cero. Para un dato particular en la segunda alternativa, la condición de solvencia es que lo anulen todas las soluciones de la ecuación adjunta, como muestra la fórmula previa a los ejercicios.
\end{solucion}

\begin{ejercicio}[Tercera alternativa de Fredholm]\label{ej:v3-l10-4}
Para $\lambda\ne0$, pruebe $\dim\ker(K-\lambda)=\dim\ker(K'-\lambda)$.
\end{ejercicio}
\begin{solucion}
Pongamos $A=K-\lambda$, $n=\dim\ker A$, $m=\dim\ker A'$, ambos finitos. Elijamos una base $x_1,\ldots,x_n$ de $\ker A$ y funcionales continuas $f_i$ con $f_i(x_j)=\delta_{ij}$. Sea $\psi_1,\ldots,\psi_m$ una base de $\ker A'$. Existen $y_1,\ldots,y_m\in X$ con $\psi_i(y_j)=\delta_{ij}$: el mapeo $x\mapsto(\psi_i(x))$ es suprayectivo sobre $\C^m$, pues de lo contrario una combinación lineal no nula de sus coordenadas se anularía en su imagen y contradice la independencia de las $\psi_i$.

Supongamos $n<m$ y definamos $F x=\sum_{i=1}^nf_i(x)y_i$. Si $(A+F)x=0$, aplicar $\psi_j$ para $j\le n$ da $f_j(x)=0$, luego $Ax=0$. Escribiendo $x$ en la base del núcleo se obtiene $x=0$. Pero $\psi_{n+1}$ se anula en la imagen de $A+F$, así que este operador no es suprayectivo. Contradice Fredholm, porque $A+F=(K+F)-\lambda I$ y $K+F$ es compacto.

Supongamos ahora $n>m$ y definamos $F x=\sum_{i=1}^mf_i(x)y_i$. Este operador $A+F$ no es inyectivo, pues anula $x_{m+1}$. Sin embargo es suprayectivo. Dado $z\in X$, el vector $z_0=z-\sum_{i=1}^m\psi_i(z)y_i$ es anulado por todas las $\psi_i$, luego pertenece a $\ran A$. Escojamos $u$ con $Au=z_0$ y pongamos
\[x=u+\sum_{i=1}^m[\psi_i(z)-f_i(u)]x_i.
\]
Entonces $Ax=z_0$ y $f_i(x)=\psi_i(z)$, por lo que $(A+F)x=z$. Esto vuelve a contradecir Fredholm. Los argumentos incluyen los casos $n=0$ o $m=0$ interpretando las sumas vacías como cero. Por tanto $n=m$. En la perturbación finita se necesitan los vectores $y_i$ que detectan el conúcleo, no los $x_i$ del núcleo usados en la primera sugerencia original.
\end{solucion}

\section{La geometría de un compacto normal}\label{lec:v3-11}
\fuentelectura{11}{El teorema espectral}
El corrimiento ponderado $Tx=(0,x_1,x_2/2,\ldots)$ del capítulo anterior es compacto: truncando después de la coordenada $N+1$ se obtiene un operador de rango finito $T_N$ y $\norm{T-T_N}\le1/(N+1)$. Un límite en norma de compactos es compacto, pues para cada tolerancia una aproximación compacta proporciona una red finita de la imagen de la bola. Este operador no tiene valores propios y su espectro es $\{0\}$. Por ello necesitaremos normalidad para una descomposición en vectores propios.

\begin{lema}\label{v3:normal-nucleos}
Si $T$ es normal, $\norm{Tx}=\norm{T^*x}$ y
\[\ker(T-\lambda I)=\ker(T^*-\cl\lambda I).
\]
Los espacios propios de valores distintos son ortogonales, y cada uno es invariante tanto bajo $T$ como bajo $T^*$.
\end{lema}
\begin{proof}
Por adjunción, $\norm{Tx}^2=\ip{T^*Tx}{x}$ y $\norm{T^*x}^2=\ip{TT^*x}{x}$; normalidad da igualdad. También $T-\lambda I$ es normal, y aplicar esa igualdad a él prueba la identidad de núcleos. Si $Tv=\lambda v$ y $Tw=\mu w$, entonces $T^*w=\cl\mu w$ y
\[\lambda\ip vw=\ip{Tv}w=\ip v{T^*w}=\mu\ip vw.
\]
Para $\lambda\ne\mu$, esto fuerza $\ip vw=0$. La identidad de núcleos da además las dos invariancias. Obsérvese el conjugado: un valor propio de $T$ corresponde a su conjugado para $T^*$.
\end{proof}

\begin{lema}[Norma de un autoadjunto]\label{v3:norma-cuadratica}
Para $A=A^*$ acotado en un Hilbert,
\[\norm A=\sup_{\norm x=1}|\ip{Ax}x|.
\]
\end{lema}
\begin{proof}
Denotemos por $c$ el supremo. Cauchy--Schwarz da $c\le\norm A$. La forma cuadrática $q(x)=\ip{Ax}x$ es real y satisface $|q(x)|\le c\norm x^2$. Por expansión,
\[4\Re\ip{Au}v=q(u+v)-q(u-v).
\]
Si $u,v$ son unitarios, paralelogramo da
\[|\Re\ip{Au}v|\le\frac c4(\norm{u+v}^2+\norm{u-v}^2)=c.
\]
Multiplicando $u$ por un escalar de módulo uno que haga real no negativa a $\ip{Au}v$, obtenemos $|\ip{Au}v|\le c$. Tomando el supremo en los vectores unitarios $v$ y después $u$, resulta $\norm A\le c$. La resta de las dos formas cuadráticas, y no su suma, es la identidad necesaria en este paso.
\end{proof}

\begin{proposicion}[Subsucesiones débiles en Hilbert]\label{v3:hilbert-subsucesion}
Toda sucesión acotada $(x_n)$ en un Hilbert tiene una subsucesión $(x_{n_j})$ y un vector $x$ tales que $\ip{x_{n_j}}y\to\ip xy$ para todo $y$. Si $\norm{x_n}\le1$, entonces $\norm x\le1$.
\end{proposicion}
\begin{proof}
Sea $M=\cl{\spanop\{x_n:n\ge1\}}$, que es separable porque las combinaciones finitas con coeficientes racionales complejos forman un denso numerable. Aplicando Gram--Schmidt a una sucesión densa se obtiene una base ortonormal finita o numerable $(e_k)$ de $M$. Cada sucesión escalar $\ip{x_n}{e_k}$ es acotada. La selección sucesiva de subsucesiones y una diagonal produce una subsucesión en la que todas las coordenadas convergen, digamos a $a_k$. Si $\norm{x_n}\le C$, Bessel y el paso al límite en sumas finitas dan $\sum_{k=1}^m|a_k|^2\le C^2$ para todo $m$. Así $x=\sum_ka_ke_k$ existe y $\norm x\le C$. Para un vector en el span finito de los $e_k$, la convergencia de apareamientos es inmediata. Para un $y\in M$ arbitrario, aproximamos por ese span; el error es a lo sumo $2C$ veces la distancia de la aproximación. Si $y\in H$, lo descomponemos en $P_My+(I-P_M)y$; la segunda parte es ortogonal a todos los $x_n$ y a $x$. Esto prueba la afirmación sin presuponer separabilidad de $H$ ni los teoremas posteriores de compacidad débil. El caso de una sucesión unitaria sólo garantiza que el límite esté en la bola \emph{cerrada}.
\end{proof}

\begin{teorema}\label{v3:normal-valor-propio}
Todo operador compacto normal \emph{no nulo} tiene un valor propio no nulo. Si es autoadjunto, uno de $\norm T$ y $-\norm T$ es un valor propio.
\end{teorema}
\begin{proof}
Primero sea $T=T^*\ne0$. Por el lema de la norma, existe una sucesión unitaria $x_n$ con $|\ip{Tx_n}{x_n}|\to\norm T$. Pasando a una subsucesión, los valores reales convergen a $\lambda=\norm T$ o $\lambda=-\norm T$. Entonces
\[\norm{(T-\lambda)x_n}^2
=\norm{Tx_n}^2-2\lambda\ip{Tx_n}{x_n}+\lambda^2
\le\norm T^2-2\lambda\ip{Tx_n}{x_n}+\lambda^2\longrightarrow0.
\]
Por compacidad, una subsucesión de $Tx_n$ converge a $y$; como $\lambda\ne0$, esa misma subsucesión de $x_n$ converge a $x=y/\lambda$. Tiene norma uno y $(T-\lambda)x=0$.

Para $T$ normal, escribamos $T=A+iB$, donde $A=(T+T^*)/2$, $B=(T-T^*)/(2i)$ son compactos autoadjuntos y conmutan. Al menos uno es no nulo. Si $B=0$, el caso anterior termina la prueba. Si $B\ne0$, tiene un valor propio real $\beta\ne0$. Su espacio propio $E_\beta$ es finito-dimensional y es $A$-invariante porque $AB=BA$. La restricción autoadjunta de $A$ a él tiene un vector propio $x\ne0$, $Ax=\alpha x$, por el teorema espectral finito. Entonces $Tx=(\alpha+i\beta)x$, con $\alpha+i\beta\ne0$. La hipótesis $T\ne0$, ausente en el enunciado original, es indispensable.
\end{proof}
Si $T$ es compacto, todos los puntos no nulos de $\sigma(T^*T)$ son valores propios; para un vector propio unitario, su valor es $\ip{T^*Tx}x=\norm{Tx}^2\ge0$. Como cero también es no negativo, $\sigma(T^*T)\subset[0,\infty)$. De igual modo, para $T$ compacto autoadjunto, sus valores propios son reales porque $\ip{Tx}x$ es real, y $\sigma(T)\subset\R$. Esto prueba la proposición 2 de la lectura.

\begin{teorema}[Espectral para compactos normales]\label{v3:espectral-compacto}\index{teorema espectral!compactos normales}
Sea $T$ compacto normal en $H$. Enumeremos sus valores propios distintos no nulos como $(\lambda_j)$, una lista finita o numerable, y sea $P_j$ la proyección ortogonal sobre $E_{\lambda_j}$. Entonces
\[H=\ker T\oplus\bigoplus_j E_{\lambda_j},\qquad
Tx=\sum_j\lambda_jP_jx.
\]
La suma de espacios es una suma ortogonal \emph{cerrada}. En el caso infinito, $\lambda_j\to0$ y la serie de operadores converge en norma, con
\[\norm{T-\sum_{j=1}^NP_j\lambda_j}=\sup_{j>N}|\lambda_j|.
\]
\end{teorema}
\begin{proof}
Los espacios propios son ortogonales y reducen a $T$ por el lema~\ref{v3:normal-nucleos}. Sea $M$ la cerradura de su suma. Tanto $M$ como $M^\perp$ son invariantes bajo $T,T^*$: para $x\in M^\perp$, $y\in M$, se tiene $\ip{Tx}y=\ip x{T^*y}=0$, y análogamente para $T^*x$. La restricción a $M^\perp$ es compacta normal. Si fuera no nula, el teorema anterior produciría allí un vector propio de valor no nulo, que pertenecería a $M$ y a $M^\perp$, contradicción. Por tanto $M^\perp\subset\ker T$. Recíprocamente, el núcleo es ortogonal a cada espacio propio no nulo por el lema, así que $\ker T\subset M^\perp$. Esto da la descomposición.

Para $x=x_0+\sum_jP_jx$, la continuidad de $T$ da la fórmula de la serie. Por ortogonalidad,
\[\norm{\sum_{j>N}\lambda_jP_jx}^2
=\sum_{j>N}|\lambda_j|^2\norm{P_jx}^2
\le\bigl(\sup_{j>N}|\lambda_j|\bigr)^2\norm x^2.
\]
Evaluar en un vector unitario de cada espacio propio restante da la cota inversa para la norma. La proposición~\ref{v3:compacto-cero} implica $\lambda_j\to0$ si la lista es infinita, y con ello la convergencia en norma. Si no hay valores propios no nulos, el argumento demuestra $T=0$.
\end{proof}
La descomposición autoadjunta y el control exacto del resto amplían las notas de Maximenko \cite{maxcompactos}. La prueba anterior incluye el caso normal complejo, la multiplicidad y la posible dimensión no numerable del núcleo, sin exigir una base numerable de todo $H$.

\section{Ejercicios de descomposición espectral}
\begin{ejercicio}\label{ej:v3-l11-1}
Si $a_n\to0$, pruebe que $D(x_n)=(a_nx_n)$ en $\ell^\infty$ es compacto y $\sigma(D)=\{a_n:n\ge1\}\cup\{0\}$.
\end{ejercicio}
\begin{solucion}
Sea $D_N$ la misma diagonal con ceros después de $N$. Tiene rango finito y $\norm{D-D_N}=\sup_{n>N}|a_n|\to0$, así que $D$ es compacto. Fuera de la cerradura de los $a_n$, la inversa diagonal está acotada; cada $a_n$ es un valor propio. Cero pertenece al espectro porque el espacio es infinito-dimensional y $D$ es compacto. Como la única posible acumulación de la sucesión es cero, su cerradura es exactamente el conjunto indicado. También puede comprobarse directamente que los vectores $e_n$ impiden una inversa acotada en cero si ningún $a_n$ se anula.
\end{solucion}

\begin{ejercicio}\label{ej:v3-l11-2}
Para un operador normal, demuestre que los espacios propios de dos valores distintos son ortogonales.
\end{ejercicio}
\begin{solucion}
Si $Tv=\lambda v$, $Tw=\mu w$, normalidad implica $T^*w=\cl\mu w$ mediante $\norm{(T^*-\cl\mu)w}=\norm{(T-\mu)w}=0$. Por tanto $\lambda\ip vw=\ip{Tv}w=\ip v{T^*w}=\mu\ip vw$. Si $\lambda\ne\mu$, el producto interno se anula. Esto desarrolla todos los pasos del lema~\ref{v3:normal-nucleos} en el caso pedido.
\end{solucion}

\begin{ejercicio}\label{ej:v3-l11-3}
Para $T\in\Bcal(H)$ son equivalentes: (i) $T$ compacto; (ii) $T^*T$ compacto; (iii) $T^*$ compacto.
\end{ejercicio}
\begin{solucion}
(i)$\Rightarrow$(ii) usa que componer un compacto con un acotado da un compacto. (i)$\Leftrightarrow$(iii) se deduce de Schauder y la identificación de Riesz, o de aplicar una dirección a $T^*$ y usar $T^{**}=T$. Para (ii)$\Rightarrow$(i), sea $(x_n)$ acotada, $\norm{x_n}\le C$. Elegimos una subsucesión para la que $T^*Tx_n$ converge. Para dos términos de ella, poniendo $u=x_n-x_m$,
\[\norm{Tu}^2=\ip{T^*Tu}u\le\norm{T^*Tu}\norm u\le2C\norm{T^*Tx_n-T^*Tx_m}.
\]
Así $(Tx_n)$ es de Cauchy y converge por completitud de $H$. Éste es el criterio secuencial de compacidad del operador.
\end{solucion}

\begin{ejercicio}\label{ej:v3-l11-4}
Si $(e_n)$ es ortonormal y $T$ compacto, demuestre $Te_n\to0$ en norma.
\end{ejercicio}
\begin{solucion}
Bessel implica $\ip{e_n}z\to0$ para todo $z\in H$, porque sus módulos cuadrados forman una serie convergente. Así $\ip{Te_n}y=\ip{e_n}{T^*y}\to0$. Si las normas $\norm{Te_n}$ no tendieran a cero, habría una subsucesión con normas al menos $\eps>0$; por compacidad, una subsucesión adicional convergería en norma a $v$. Sus apareamientos convergen tanto a $\ip vy$ como a cero para todo $y$, luego $v=0$, contradiciendo la cota de normas.
\end{solucion}

\begin{ejercicio}[Volterra]\label{ej:v3-l11-5}\index{Volterra, operador de}
En $L^2[0,1]$ sea $(Vf)(x)=\int_0^xf(t)\,dt$. Pruebe (i) que es compacto; (ii) que $(V^*g)(t)=\int_t^1g(x)\,dx$; (iii) que $M_a=\{f:f=0\text{ casi en todas partes en }[0,a]\}$ es $V$-invariante para $0\le a\le1$.
\end{ejercicio}
\begin{solucion}
(i) Cauchy--Schwarz da $|Vf(x)|\le\norm f_2$ y
\[|Vf(x)-Vf(y)|\le|x-y|^{1/2}\norm f_2.
\]
La imagen de la bola unidad es uniformemente acotada y equicontinua en $C[0,1]$. Por Arzelà--Ascoli su cerradura allí es compacta; la inclusión continua $C[0,1]\hookrightarrow L^2[0,1]$ preserva compacidad, de modo que $V$ es compacto.

(ii) Para $f,g\in L^2$, la integral doble de $|f(t)g(x)|$ sobre el triángulo $0\le t\le x\le1$ está acotada por $\norm f_1\norm g_1\le\norm f_2\norm g_2$, así que Fubini es aplicable. Entonces
\[\ip{Vf}g=\int_0^1\int_0^x f(t)\cl{g(x)}\,dt\,dx
=\int_0^1 f(t)\overline{\int_t^1g(x)\,dx}\,dt.
\]
Esto identifica el adjunto, cuya fórmula define un operador acotado por la misma estimación de Cauchy--Schwarz.

(iii) Si $f=0$ casi en todas partes en $[0,a]$, para todo $x\le a$ la integral $\int_0^xf$ es cero. Por tanto $Vf\in M_a$. Estos subespacios son cerrados porque son núcleos del multiplicador continuo por $\boldsymbol1_{[0,a]}$. La clasificación de \emph{todos} los subespacios invariantes de Volterra, mencionada sin prueba en la lectura, es un resultado adicional que no se usa aquí; el ejercicio sólo solicita esta invariancia.
\end{solucion}

El teorema espectral convierte un compacto normal en una diagonal ortogonal y controla el error al truncarla. Para compactos generales esa diagonalización falla, como mostró el corrimiento ponderado. El sustituto correcto será la descomposición en valores singulares, que conserva dos familias ortonormales y permite medir cuán rápidamente decae la acción del operador.

\chapter{Valores singulares, traza y clases de Schatten}\label{cap:v3-07}
La norma de operador mide la mayor expansión posible, pero no describe cuántas direcciones contribuyen a ella. Los valores singulares organizan esas expansiones en una sucesión. Su suma y la suma de sus cuadrados producen dos clases fundamentales de operadores: los de clase de traza y los de Hilbert--Schmidt. Probaremos también la completitud de todas las clases de Schatten, desarrollando el ejercicio más extenso de la lectura.

\section{Raíz positiva y descomposición singular}\label{lec:v3-12}
\fuentelectura{12}{Operadores de Hilbert--Schmidt y de clase de traza}
Diremos que $A\in\Bcal(H)$ es positivo si es autoadjunto y $\ip{Ax}x\ge0$ para todo $x$. En particular $T^*T$ es positivo para todo $T$ acotado.

\begin{proposicion}\label{v3:raiz-compacta}
Si $T$ es compacto, existe un único operador positivo $|T|$ tal que $|T|^2=T^*T$. Es compacto y satisface $\norm{|T|x}=\norm{Tx}$.
\end{proposicion}
\begin{proof}
El operador $T^*T$ es compacto positivo. Por el teorema espectral, sus valores propios no nulos, repetidos según multiplicidad, son $a_n>0$, con vectores propios ortonormales $f_n$, y
\[T^*Tx=\sum_n a_n\ip x{f_n}f_n.
\]
La lista puede ser finita; en el caso infinito $a_n\to0$. Definimos
\[|T|x=\sum_n\sqrt{a_n}\ip x{f_n}f_n,
\]
con valor cero en el núcleo. La serie define un operador positivo y las truncaciones convergen en norma porque el supremo de los coeficientes restantes tiende a cero. Es, pues, compacto y su cuadrado es $T^*T$. Además $\norm{|T|x}^2=\ip{T^*Tx}x=\norm{Tx}^2$.

Si $B$ es otra raíz positiva, $B$ conmuta con $T^*T=B^2$. Por ello preserva cada espacio propio $E_a$ para $a>0$, y su restricción positiva a ese espacio finito tiene cuadrado $aI$. El teorema espectral finito obliga a que todos sus valores propios sean $\sqrt a$, luego $B|_{E_a}=\sqrt a I$. En el núcleo de $T^*T$, $\norm{Bx}^2=\ip{B^2x}x=0$, de modo que $B$ se anula. La descomposición ortogonal demuestra $B=|T|$. La unicidad se afirma entre raíces \emph{positivas}; otras raíces pueden existir.
\end{proof}

\begin{definicion}\index{valor singular}
Los valores propios positivos de $|T|$, repetidos según multiplicidad y ordenados de manera no creciente, se llaman \emph{valores singulares} de $T$ y se denotan $s_n(T)$. Si la lista es finita, la completamos con ceros cuando una fórmula lo requiera.
\end{definicion}

\begin{teorema}[Descomposición singular y polar]\label{v3:singulares}\index{descomposición!singular}\index{descomposición!polar}\index{isometría parcial}
Existen familias ortonormales $(f_n),(g_n)$ tales que, con $s_n=s_n(T)>0$,
\begin{align*}
Tx&=\sum_ns_n\ip x{f_n}g_n,& |T|x&=\sum_ns_n\ip x{f_n}f_n,\\
T^*x&=\sum_ns_n\ip x{g_n}f_n,& |T^*|x&=\sum_ns_n\ip x{g_n}g_n.
\end{align*}
Las series de operadores convergen en norma. Además $T=U|T|$, donde $U$ es una isometría de $\cl{\ran|T|}=(\ker T)^\perp$ sobre $\cl{\ran T}$ y vale cero en $\ker T$.
\end{teorema}
\begin{proof}
Escojamos los $f_n$ como base ortonormal de los espacios propios positivos de $|T|$ y definamos $g_n=Tf_n/s_n$. Entonces
\[\ip{g_n}{g_m}=\frac{\ip{T^*Tf_n}{f_m}}{s_ns_m}=\delta_{nm}.
\]
El núcleo de $|T|$ coincide con el de $T$ por la igualdad de normas anterior. Aplicando $T$ a la descomposición de $x$ en ese núcleo y los $f_n$ obtenemos la primera fórmula; la norma de su cola es $\sup_{n>N}s_n$, por ortogonalidad. Para el adjunto, evaluar los apareamientos en la serie da la tercera fórmula. De ella se deduce $TT^*g_n=s_n^2g_n$ y que $TT^*$ se anula en el complemento ortogonal de los $g_n$; la unicidad de la raíz positiva da la cuarta.

La aplicación definida inicialmente por $U(|T|x)=Tx$ está bien definida porque $\ker|T|=\ker T$, y preserva normas. Se extiende por continuidad a la cerradura de esa imagen, con imagen $\cl{\ran T}$, y se define como cero en su complemento ortogonal. Esto prueba la descomposición polar. En general $U$ es una \emph{isometría parcial}, no un unitario como afirma la lectura: el corrimiento compacto ponderado inyectivo en $\ell^2$ tiene parte polar igual al corrimiento unilateral, que no es suprayectivo.
\end{proof}

\section{Hilbert--Schmidt y la independencia de la base}
Una suma de números no negativos sobre un conjunto arbitrario de índices se define como el supremo de sus sumas finitas. Si es finita, sólo un número a lo sumo numerable de términos es no nulo: para cada $m$, sólo puede haber finitos términos mayores que $1/m$.

\begin{proposicion}\label{v3:hilbert-schmidt}\index{operador!de Hilbert--Schmidt}
Para cualquier $T\in\Bcal(H)$, el número
\[\norm T_{\mathrm{HS}}^2=\sum_{e\in\mathcal E}\norm{Te}^2\in[0,\infty]
\]
no depende de la base ortonormal $\mathcal E$. Además $\norm T_{\mathrm{HS}}=\norm{T^*}_{\mathrm{HS}}$. Si es finito, $T$ es compacto y
\[\norm T\le\norm T_{\mathrm{HS}},\qquad
\norm T_{\mathrm{HS}}^2=\sum_ns_n(T)^2.
\]
\end{proposicion}
\begin{proof}
Para bases ortonormales $\mathcal E,\mathcal F$, Parseval y el intercambio de sumas no negativas dan
\[\sum_{e\in\mathcal E}\norm{Te}^2
=\sum_{e\in\mathcal E}\sum_{f\in\mathcal F}|\ip{Te}f|^2
=\sum_{f\in\mathcal F}\sum_{e\in\mathcal E}|\ip e{T^*f}|^2
=\sum_{f\in\mathcal F}\norm{T^*f}^2.
\]
Tomando primero $\mathcal F=\mathcal E$ obtenemos igualdad para $T$ y $T^*$ en una base; combinándola con la identidad para dos bases obtenemos independencia. Parseval proviene de la expansión ortogonal de los vectores: en particular $\ip xy=\sum_{e\in\mathcal E}\ip xe\ip ey$, con convergencia absoluta por Cauchy--Schwarz.

Para $x=\sum_e\ip xe\,e$ y $y$ unitario, Cauchy--Schwarz da
\[|\ip{Tx}y|\le\Bigl(\sum_e|\ip xe|^2\Bigr)^{1/2}
\Bigl(\sum_e|\ip{Te}y|^2\Bigr)^{1/2}
\le\norm x\norm T_{\mathrm{HS}}.
\]
Así $\norm T\le\norm T_{\mathrm{HS}}$. Si esta última norma es finita, para cada $\eps>0$ existe un conjunto finito $F\subset\mathcal E$ con $\sum_{e\notin F}\norm{Te}^2<\eps^2$. La misma estimación aplicada a $T(I-P_F)$ da $\norm{T-TP_F}<\eps$. Como $TP_F$ tiene rango finito, $T$ es compacto. Finalmente, extendemos los vectores singulares $f_n$ a una base con vectores del núcleo; en ella $\norm{Tf_n}^2=s_n^2$ y los demás términos son cero, dando la fórmula de valores singulares.
\end{proof}

\section{Una fórmula variacional para las clases de Schatten}
Para $1\le p<\infty$, definimos
\[\mathcal S_p(H)=\{T\in\Kcal(H):\sum_ns_n(T)^p<\infty\},\qquad
\norm T_p=\Bigl(\sum_ns_n(T)^p\Bigr)^{1/p}.
\]
Escribimos $\mathcal S_2=\mathrm{HS}(H)$ y $\mathcal S_1=\mathrm{TR}(H)$; la norma de esta última también se denota $\norm T_{\mathrm{tr}}$. La siguiente prueba establece de una vez la desigualdad triangular y la completitud, que no se deducen simplemente de la desigualdad triangular para sucesiones, porque los valores singulares de una suma no se suman término a término.

\begin{lema}[Fórmula variacional]\label{v3:schatten-variacional}
Sea $T$ compacto, $1\le p<\infty$ y $q$ el exponente conjugado, con $q=\infty$ si $p=1$. Entonces, admitiendo el valor $+\infty$,
\[\norm T_p=
\sup\left\{\left|\sum_{i=1}^m b_i\ip{Te_i}{h_i}\right|:
\begin{array}{l}(e_i)_{i=1}^m,(h_i)_{i=1}^m\text{ ortonormales},\\
b_i\ge0,\ \norm b_{\ell^q}\le1,\ m<\infty
\end{array}\right\}.
\]
\end{lema}
\begin{proof}
Use la descomposición $Tx=\sum_js_j\ip x{f_j}g_j$. Para las familias finitas del enunciado pongamos
$a_{ij}=|\ip{e_i}{f_j}|\,|\ip{g_j}{h_i}|$. Por Cauchy--Schwarz y Bessel,
\[\sum_ja_{ij}\le1\quad\text{para cada }i,\qquad
\sum_ia_{ij}\le1\quad\text{para cada }j.
\]
Entonces el módulo de la suma está acotado por $\sum_{i,j}b_is_ja_{ij}$. Si $1<p<\infty$, Hölder para esta suma con pesos no negativos da
\[\sum_{i,j}b_is_ja_{ij}
\le\Bigl(\sum_{i,j}b_i^qa_{ij}\Bigr)^{1/q}
\Bigl(\sum_{i,j}s_j^pa_{ij}\Bigr)^{1/p}
\le\norm b_q\norm s_p\le\norm T_p.
\]
Para $p=1$ se usa directamente $b_i\le1$ y $\sum_i a_{ij}\le1$, obteniendo la misma cota. Las sumas infinitas se justifican por límites monótonos de sumas finitas.

Para la cota opuesta, escoja $e_i=f_i$, $h_i=g_i$ para los primeros $m$ valores singulares positivos. Si $p=1$, tome $b_i=1$; obtiene $\sum_{i=1}^ms_i$. Si $p>1$, tome
$b_i=s_i^{p-1}/(\sum_{j=1}^ms_j^p)^{1/q}$, cuando el denominador no sea cero. Entonces $\norm b_q=1$ y la suma es $(\sum_{i=1}^ms_i^p)^{1/p}$. Al aumentar $m$ se alcanza el valor requerido. Si $T=0$, ambas partes son cero.
\end{proof}

\begin{teorema}\label{v3:schatten-banach}\index{Schatten, clases de}
Para cada $1\le p<\infty$, $\mathcal S_p(H)$ es un espacio de Banach con la norma $\norm\cdot_p$, y $\norm T\le\norm T_p$. En particular $\norm\cdot_{\mathrm{HS}}$ y $\norm\cdot_{\mathrm{tr}}$ son normas.
\end{teorema}
\begin{proof}
Los compactos forman un espacio vectorial. La fórmula variacional expresa $\norm T_p$ como un supremo de módulos de funcionales lineales en $T$; por ello es homogénea y subaditiva. Si $S,T\in\mathcal S_p$, esa subaditividad hace finita la norma de $S+T$, así que la clase es un espacio vectorial. El caso $m=1$, $b_1=1$ muestra $\norm T\le\norm T_p$, y de aquí se deduce que sólo el operador cero tiene norma cero.

Sea $(T_n)$ de Cauchy en $\norm\cdot_p$. También lo es en norma de operador, luego converge a un operador $T$; como los compactos son cerrados en esa norma, $T$ es compacto. Para cada funcional finita $F$ del lema, $F(T_n-T_m)\to F(T_n-T)$ cuando $m\to\infty$. Dado $\eps>0$, si $n,m\ge N$ tenemos $|F(T_n-T_m)|\le\eps$, por tanto $|F(T_n-T)|\le\eps$ para $n\ge N$. Tomando el supremo en todas las elecciones finitas, $\norm{T_n-T}_p\le\eps$. En particular $T-T_N\in\mathcal S_p$, y entonces $T\in\mathcal S_p$. Esto prueba la completitud sin suponer de antemano continuidad de cada valor singular.
\end{proof}

\section{Traza y productos de operadores}
Para $u,v\in H$ escribiremos $u\otimes v$ para el operador de rango uno $x\mapsto\ip xv u$. Su único valor singular no nulo es $\norm u\norm v$ si ambos vectores son no nulos: el cuadrado del módulo es el operador $\norm u^2(v\otimes v)$. Por tanto $\norm{u\otimes v}_p=\norm u\norm v$ para todo $p$.

\begin{proposicion}\label{v3:traza-variacional}
Para un compacto $T$,
\[\norm T_{\mathrm{tr}}=
\sup_{\mathcal E,\mathcal F}\sum_{i\in I}|\ip{Te_i}{h_i}|,
\]
donde se consideran bases ortonormales $\mathcal E=(e_i)_{i\in I}$, $\mathcal F=(h_i)_{i\in I}$ con un mismo conjunto de índices. Si $T$ es de clase de traza, la serie
\[\tr T=\sum_{e\in\mathcal E}\ip{Te}e
\]
converge absolutamente y es independiente de la base. Para una descomposición singular,
\[\tr T=\sum_ns_n\ip{g_n}{f_n},\qquad |\tr T|\le\norm T_{\mathrm{tr}}.
\]
\end{proposicion}
\begin{proof}
La cota superior para la suma de módulos sigue del lema~\ref{v3:schatten-variacional} con $p=1$, absorbiendo las fases en los $h_i$ y tomando supremos sobre subconjuntos finitos. Para la cota inferior, fijado $N$ extendemos las familias \emph{finitas} $f_1,\ldots,f_N$ y $g_1,\ldots,g_N$ a bases completas con índices coincidentes en esos primeros lugares; sus complementos tienen la misma dimensión porque se quitó la misma cantidad finita de vectores. La suma es al menos $\sum_{n=1}^Ns_n$. Hacer $N\to\infty$ da la igualdad. Este argumento no necesita extender simultáneamente dos familias infinitas cuyos complementos podrían tener dimensiones distintas.

Para una base $\mathcal E$ y $T$ de traza, la descomposición singular da
\[\sum_e\ip{Te}e=\sum_e\sum_ns_n\ip e{f_n}\ip{g_n}e.
\]
La suma de módulos de los términos dobles es a lo sumo
\[\sum_ns_n\Bigl(\sum_e|\ip e{f_n}|^2\Bigr)^{1/2}
\Bigl(\sum_e|\ip{g_n}e|^2\Bigr)^{1/2}=\sum_ns_n<\infty.
\]
Podemos, pues, intercambiar las sumas. Parseval da $\sum_e\ip e{f_n}\ip{g_n}e=\ip{g_n}{f_n}$, y obtenemos una expresión independiente de $\mathcal E$. La cota de la traza sigue tomando módulos. La linealidad de la traza se obtiene sumando en una misma base las series absolutamente convergentes.
\end{proof}
Si $T$ es además normal, su descomposición espectral y una base de vectores propios dan
\[\tr T=\sum_{\lambda\in\sigma(T)\setminus\{0\}}\lambda\dim E_\lambda,
\]
con convergencia absoluta, pues los valores singulares son los módulos de esos valores propios con sus multiplicidades.

\begin{observacion}[Una igualdad de imágenes que debe corregirse]
La lectura afirma $\ran\sqrt{|T|}=\ran|T|$. En general sólo coinciden sus \emph{cerraduras}, ambas iguales a $(\ker T)^\perp$. Por ejemplo, en $\ell^2$, sea $|T|e_n=n^{-2}e_n$. El vector $y_n=n^{-2}$ pertenece a $\ran\sqrt{|T|}$ porque su preimagen es $(1/n)\in\ell^2$, pero no a $\ran|T|$, pues la única preimagen formal sería $(1,1,\ldots)$. La demostración de independencia de la traza dada arriba evita esta afirmación falsa. Sí es válida la factorización $T=(U\sqrt{|T|})\sqrt{|T|}$, con ambos factores de Hilbert--Schmidt cuando $T$ es de clase de traza, porque los cuadrados de sus valores singulares suman $\sum_ns_n(T)$.
\end{observacion}

\begin{proposicion}\label{v3:ideales-hs-traza}
$\mathcal S_2$ y $\mathcal S_1$ son ideales bilaterales de $\Bcal(H)$ y satisfacen, para $A$ acotado,
\[\norm{AT}_{\mathrm{HS}},\norm{TA}_{\mathrm{HS}}\le\norm A\norm T_{\mathrm{HS}},
\qquad
\norm{AT}_{\mathrm{tr}},\norm{TA}_{\mathrm{tr}}\le\norm A\norm T_{\mathrm{tr}}.
\]
Si $S,T$ son Hilbert--Schmidt, $ST$ es de traza y $\norm{ST}_{\mathrm{tr}}\le\norm S_{\mathrm{HS}}\norm T_{\mathrm{HS}}$.
\end{proposicion}
\begin{proof}
Para Hilbert--Schmidt, en cualquier base,
$\sum_e\norm{ATe}^2\le\norm A^2\sum_e\norm{Te}^2$. La cota para $TA$ se obtiene aplicando la anterior a $A^*T^*$ y usando la igualdad de normas Hilbert--Schmidt de un operador y su adjunto.

Para traza, la descomposición singular $T=\sum_ns_n(g_n\otimes f_n)$ converge en norma de traza: la norma de la cola es $\sum_{n>N}s_n$. Como $\norm{(Ag_n)\otimes f_n}_{\mathrm{tr}}\le\norm A$, la serie $\sum_ns_n(Ag_n\otimes f_n)$ converge en el Banach $\mathcal S_1$ y coincide en norma de operador con $AT$. Esto da la primera cota. De igual modo $TA=\sum_ns_n(g_n\otimes A^*f_n)$ da la segunda.

Finalmente, $ST$ es compacto. Para familias ortonormales finitas $(e_i),(h_i)$,
\[\sum_i|\ip{STe_i}{h_i}|=\sum_i|\ip{Te_i}{S^*h_i}|
\le\Bigl(\sum_i\norm{Te_i}^2\Bigr)^{1/2}
\Bigl(\sum_i\norm{S^*h_i}^2\Bigr)^{1/2}
\le\norm T_{\mathrm{HS}}\norm S_{\mathrm{HS}}.
\]
La fórmula variacional para $p=1$ implica que $ST$ es de traza y da la cota.
\end{proof}

\section{Ejercicios propuestos}
\begin{ejercicio}\label{ej:v3-l12-1}
Demuestre que $\norm\cdot_{\mathrm{HS}}$ y $\norm\cdot_{\mathrm{tr}}$ son normas en sus respectivas clases.
\end{ejercicio}


\begin{ejercicio}\label{ej:v3-l12-2}
Todo operador de Hilbert--Schmidt es compacto.
\end{ejercicio}


\begin{ejercicio}\label{ej:v3-l12-3}
Pruebe $\norm T\le\norm T_{\mathrm{HS}}$.
\end{ejercicio}


\begin{ejercicio}\label{ej:v3-l12-4}
Demuestre que las clases de Hilbert--Schmidt y de traza son ideales bilaterales en $\Bcal(H)$.
\end{ejercicio}


\begin{ejercicio}\label{ej:v3-l12-5}
Pruebe: (i) si $T$ es de traza, $T^*$ también y $\tr(T^*)=\overline{\tr T}$; (ii) para $S,T$ Hilbert--Schmidt, $\tr(ST)=\tr(TS)$; (iii) $\langle S,T\rangle_{\mathrm{HS}}=\tr(ST^*)$ es un producto interno que hace Hilbert a $\mathrm{HS}(H)$.
\end{ejercicio}


\begin{ejercicio}[Todas las clases de Schatten]\label{ej:v3-l12-6}
Para $1\le p<\infty$, pruebe que $\mathcal S_p(H)$ es un espacio vectorial complejo, que $\norm\cdot_p$ es una norma y que ese espacio es completo.
\end{ejercicio}


\section{Ampliación: núcleos integrales y prueba de Schur}\index{Schur, prueba de}
El siguiente criterio desarrolla la nota de Egor Maximenko \emph{La prueba de Schur simplificada para operadores integrales} \cite{maxschur}. Se distinguen los espacios de entrada y salida para evitar invertirlos en la fórmula.

\begin{teorema}[Schur, versión no ponderada]\label{v3:schur}
Sean $(X,\mu),(Y,\nu)$ espacios de medida $\sigma$-finitos y $K:X\times Y\to\C$ medible. Supongamos
\[\esssup_{x\in X}\int_Y|K(x,y)|\,d\nu(y)\le C_1<\infty,
\quad
\esssup_{y\in Y}\int_X|K(x,y)|\,d\mu(x)\le C_2<\infty.
\]
Entonces $Af(x)=\int_YK(x,y)f(y)\,d\nu(y)$ define un operador $L^2(Y,\nu)\to L^2(X,\mu)$ y $\norm A\le\sqrt{C_1C_2}$.
\end{teorema}
\begin{proof}
Tonelli da
\[\int_X\int_Y|K(x,y)||f(y)|^2\,d\nu(y)\,d\mu(x)
\le C_2\norm f_2^2.
\]
Así la integral interior es finita para casi todo $x$. Cauchy--Schwarz aplicado a $|K|^{1/2}$ y $|K|^{1/2}|f|$ implica allí integrabilidad absoluta de $Kf$ y
\[|Af(x)|^2\le\left(\int_Y|K(x,y)|\,d\nu(y)\right)
\left(\int_Y|K(x,y)||f(y)|^2\,d\nu(y)\right).
\]
Integramos y usamos las dos cotas en su orden:
\[\norm{Af}_2^2\le C_1\int_Y|f(y)|^2\left(\int_X|K(x,y)|\,d\mu(x)\right)d\nu(y)
\le C_1C_2\norm f_2^2.
\]
Esto prueba existencia casi en todas partes, independencia del representante y acotación. No debe introducirse un segundo factor $C_1C_2$ al reescribir la integral intermedia, como ocurre en una línea de la nota complementaria.
\end{proof}

\begin{proposicion}\label{v3:nucleo-hs}
Si $K\in L^2([0,1]^2)$, el operador $A_Kf(x)=\int_0^1K(x,y)f(y)\,dy$ es Hilbert--Schmidt y
\[\norm{A_K}_{\mathrm{HS}}=\norm K_{L^2([0,1]^2)}.
\]
\end{proposicion}
\begin{proof}
Por Cauchy--Schwarz en $y$ y después integración en $x$, $\norm{A_Kf}_2\le\norm K_2\norm f_2$. Sea $(e_n)$ una base ortonormal de $L^2[0,1]$. Para casi todo $x$, $K(x,\cdot)\in L^2$ por Fubini y
\[\sum_n\left|\int_0^1K(x,y)e_n(y)\,dy\right|^2
=\norm{K(x,\cdot)}_2^2,
\]
porque $(\cl{e_n})$ también es una base ortonormal y se aplica Parseval. Tonelli permite integrar esta suma no negativa:
\[\sum_n\norm{A_Ke_n}_2^2=\int_0^1\norm{K(x,\cdot)}_2^2\,dx=\norm K_2^2.
\]
La proposición~\ref{v3:hilbert-schmidt} da la conclusión y también la compacidad.
\end{proof}

\begin{ejercicio}[Complemento: núcleo acotado en medida finita]\label{ej:v3-ampl-schur}
Si $\mu(X),\nu(Y)<\infty$ y $|K(x,y)|\le M$ casi en todas partes, pruebe $\norm A\le M\sqrt{\mu(X)\nu(Y)}$.
\end{ejercicio}


\begin{ejercicio}[Complemento: tamaño Hilbert--Schmidt de Volterra]\label{ej:v3-ampl-volterra-hs}
Calcule la norma Hilbert--Schmidt de $V$ en $L^2[0,1]$ y determine si $V^*V$ es de traza.
\end{ejercicio}


Las clases de Schatten proporcionan normas que distinguen diferentes grados de compacidad. La traza recupera una suma de valores propios para compactos normales, mientras que la clase Hilbert--Schmidt convierte el espacio de operadores en un Hilbert. Los criterios integrales conectan estas construcciones con ejemplos calculables; más adelante, las topologías débiles permitirán estudiar límites que no convergen en ninguna de estas normas.

\parte{Topologías débiles y reflexividad}{Las redes permiten estudiar topologías que no se describen mediante una sola métrica. Banach--Alaoglu, Eberlein--Šmulian y la reflexividad precisan cuándo la acotación se convierte en compacidad débil.}
\chapter{Topologías débiles, redes y convergencia}\label{cap:v3-08}
Una topología débil observa vectores mediante funcionales: exige convergencia de todas las mediciones escalares, pero permite perder control sobre la norma. Las redes sirven para describir esa topología cuando las sucesiones no bastan. Desarrollaremos cuidadosamente la diferencia entre ambas, porque varias afirmaciones válidas para sucesiones dejan de serlo para redes.

\section{Topologías generadas por aplicaciones}\label{lec:v3-13}
\fuentelectura{13}{La topología débil}
\begin{definicion}\index{topología!inicial}
Sea $E$ un conjunto y $f_i:E\to Y_i$ una familia de aplicaciones a espacios topológicos. La \emph{topología inicial} de esta familia es la topología con menos abiertos que hace continuas a todas las $f_i$.
\end{definicion}
\begin{proposicion}\label{v3:topologia-inicial}
Los conjuntos $f_i^{-1}(U)$, con $U\subset Y_i$ abierto, forman una subbase para la topología inicial. Las intersecciones finitas de esos conjuntos, incluyendo la intersección vacía $E$, forman una base. Si los $Y_i$ son Hausdorff y la familia separa puntos, la topología inicial es Hausdorff.
\end{proposicion}
\begin{proof}
La familia de intersecciones finitas cubre $E$ porque incluye $E$, y la intersección de dos de sus elementos sigue siendo una intersección finita. Por tanto sus uniones forman una topología. En ella son abiertos todos los $f_i^{-1}(U)$, así que cada $f_i$ es continua. Cualquier otra topología que haga continuas a todas las aplicaciones debe contener esos conjuntos, sus intersecciones finitas y sus uniones; por ello contiene a la construida. Si $x\ne y$, existe $i$ con $f_i(x)\ne f_i(y)$. En $Y_i$ hay abiertos disjuntos alrededor de esos dos valores; sus preimágenes son abiertos disjuntos que separan $x$ e $y$.
\end{proof}
La topología producto en $\prod_iY_i$ es, por definición, la topología inicial de las proyecciones coordenadas. Sus abiertos básicos sólo imponen condiciones en un número finito de coordenadas, aunque el producto tenga infinitos factores.

\section{Topologías débil y débil estrella}
\begin{definicion}\label{v3:topologia-debil}\index{topología!débil}
En un espacio normado $X$, la topología \emph{débil} $\sigma(X,X^*)$ es la topología inicial de las funcionales continuas $X^*$. En el dual $X^*$, la topología \emph{débil estrella} $\sigma(X^*,X)$ es la topología inicial de las evaluaciones $\phi\mapsto\phi(x)$, con $x\in X$.
\end{definicion}
\label{v3:topologia-debil-estrella}\index{topología!débil estrella}
Hahn--Banach asegura que $X^*$ separa puntos de $X$, y dos funcionales distintas difieren en algún vector. Por tanto ambas topologías son Hausdorff. La topología de la norma en $X$ es más fina que la débil, porque todas las funcionales usadas ya son continuas en norma. En $X^*$, la topología de norma es más fina que la débil $\sigma(X^*,X^{**})$, y ésta es más fina que la débil estrella, pues las evaluaciones $J_Xx$ son sólo algunas de las funcionales de $X^{**}$.

\begin{proposicion}\label{v3:vecindades-debiles}
Una base de vecindades débiles de cero es
\[D(F,\eps)=\{x:|\phi(x)|<\eps\text{ para toda }\phi\in F\},
\qquad F\subset X^*\text{ finito},\quad\eps>0.
\]
Una base en un punto $x_0$ se obtiene trasladando estos conjuntos por $x_0$. En el dual, las vecindades débiles estrella se obtienen intercambiando los papeles de vectores y funcionales.
\end{proposicion}
\begin{proof}
Una vecindad inicial de cero contiene una intersección finita $\bigcap_{j=1}^m\phi_j^{-1}(U_j)$ con $0\in U_j$ abierto en $\F$. Cada $U_j$ contiene una bola escalar de radio $\eps_j>0$. Tomando $\eps=\min_j\eps_j$ obtenemos un $D(F,\eps)$ contenido en la vecindad. Cada uno de estos conjuntos es débilmente abierto por su definición. La igualdad $\phi(x-x_0)=\phi(x)-\phi(x_0)$ demuestra la afirmación sobre traslaciones. La misma prueba con evaluaciones da la versión débil estrella. Sobre $\C$ se usan discos $|z|<\eps$, no intervalos reales.
\end{proof}
Los conjuntos $D(F,\eps)$ son convexos y balanceados. La suma es continua en estas topologías porque $D(F,\eps/2)+D(F,\eps/2)\subset D(F,\eps)$; la multiplicación escalar también lo es, al estimar cada una de las finitísimas funcionales. Así obtenemos espacios vectoriales topológicos localmente convexos.

\section{Ejercicios de topologías iniciales}
\begin{ejercicio}\label{ej:v3-l13-1}
Suponga que una topología $\tau$ en $E$ tiene la propiedad: para toda aplicación $g:Z\to E$ desde un espacio topológico, $g$ es continua si y sólo si cada $f_i\circ g$ lo es. Pruebe que $\tau$ es la topología inicial de las $f_i$.
\end{ejercicio}
\begin{solucion}
Tomando $g$ igual a la identidad de $(E,\tau)$ en sí mismo, la hipótesis implica que cada $f_i$ es continua para $\tau$. Por minimalidad, la topología inicial $\tau_0$ está contenida en $\tau$. Ahora tome $Z=(E,\tau_0)$ y $g$ la identidad hacia $(E,\tau)$. Todas las composiciones $f_i\circ g$ son continuas por construcción, así que $g$ es continua y $\tau\subset\tau_0$. Las dos inclusiones dan igualdad. A la inversa, en la topología inicial, la continuidad de todas las composiciones hace abiertas las preimágenes bajo $g$ de los abiertos subbásicos y, por uniones e intersecciones finitas, de todos los abiertos.
\end{solucion}

\begin{ejercicio}\label{ej:v3-l13-2}
Si la topología inicial de una familia es Hausdorff, demuestre que la familia separa puntos.
\end{ejercicio}
\begin{solucion}
Si dos puntos distintos $x,y$ tuvieran $f_i(x)=f_i(y)$ para todo $i$, pertenecerían exactamente a los mismos abiertos subbásicos. Lo mismo ocurriría para intersecciones finitas y uniones, por tanto para todos los abiertos. No podrían tener vecindades abiertas disjuntas, contradiciendo Hausdorff. Esta dirección ni siquiera requiere que los espacios de llegada sean Hausdorff.
\end{solucion}

\begin{ejercicio}\label{ej:v3-l13-3}
Demuestre que los conjuntos $D^*(S,1)=\{\phi:|\phi(x)|<1\text{ para todo }x\in S\}$, con $S\subset X$ finito, forman una base de vecindades débiles estrella de cero.
\end{ejercicio}
\begin{solucion}
La proposición~\ref{v3:vecindades-debiles} da primero los conjuntos con tolerancia arbitraria $\eps>0$. Pero
\[D^*(\{x_1,\ldots,x_m\},\eps)
=D^*(\{x_1/\eps,\ldots,x_m/\eps\},1),
\]
por linealidad de cada $\phi$. Por tanto fijar el radio en uno no reduce la familia de vecindades necesarias. Cada conjunto indicado es abierto por ser una intersección finita de preimágenes de discos bajo evaluaciones.
\end{solucion}

\begin{ejercicio}\label{ej:v3-l13-4}
En dimensión infinita, las topologías débil y de norma no coinciden.
\end{ejercicio}
\begin{solucion}
Toda vecindad débil de cero contiene un $D(F,\eps)$, donde $F=\{\phi_1,\ldots,\phi_m\}$. El mapeo lineal $x\mapsto(\phi_1(x),\ldots,\phi_m(x))$ desde un espacio infinito-dimensional a $\F^m$ tiene núcleo no nulo. Si $u\ne0$ pertenece a él, toda la recta $\{tu:t\in\F\}$ está en $D(F,\eps)$, que por ello no es acotado. La bola unidad abierta en norma no puede contener ninguna vecindad débil de cero; por tanto no es débilmente abierta. En dimensión finita las coordenadas de una base son funcionales continuas y las dos topologías sí coinciden.
\end{solucion}

\section{Redes, subredes y puntos de acumulación}\label{lec:v3-14}
\fuentelectura{14}{Redes}
\begin{definicion}\label{v3:redes}\index{red}\index{subred}
Un conjunto dirigido es un conjunto no vacío $A$ parcialmente ordenado en el que cualesquiera dos elementos tienen una cota superior común. Una \emph{red} es una aplicación $\alpha\mapsto x_\alpha$ desde un conjunto dirigido a un conjunto $E$. Está \emph{eventualmente} en $B\subset E$ si existe $\alpha_0$ con $x_\alpha\in B$ para todo $\alpha\ge\alpha_0$. Está \emph{frecuentemente} en $B$ si para todo $\alpha_0$ existe $\alpha\ge\alpha_0$ con $x_\alpha\in B$.
\end{definicion}
Ejemplos de conjuntos dirigidos son $\N$ con su orden usual, los enteros positivos ordenados por divisibilidad (el producto es una cota común) y las vecindades de un punto ordenadas por \emph{inclusión inversa} (la intersección es una cota común). Una sucesión es una red, pero el conjunto de índices de una red no tiene por qué ser numerable.

Un conjunto de índices es \emph{cofinal} si contiene algún índice posterior a cada índice dado. Un conjunto superior no vacío contiene una cola y permite describir la eventualidad. Es necesario excluir el conjunto vacío al usar esa descripción: de lo contrario cualquier red estaría ``eventualmente'' en cualquier conjunto por vacuidad, un defecto de la definición original.

La red converge a $x$ si está eventualmente en cada vecindad de $x$; $x$ es un punto de acumulación de la red si ésta está frecuentemente en cada vecindad suya. Una \emph{subred} tiene la forma $y_\beta=x_{a(\beta)}$, donde $a:B\to A$ satisface: para todo $\alpha_0\in A$ existe $\beta_0\in B$ tal que $a(\beta)\ge\alpha_0$ cuando $\beta\ge\beta_0$. No exigimos que $a$ sea inyectiva ni monótona. Toda propiedad eventual de una red pasa a sus subredes, directamente por esta condición.

\begin{lema}\label{v3:subred-familia}
Sea $\mathcal F$ una familia no vacía de subconjuntos tal que, dados $F,G\in\mathcal F$, existe $H\in\mathcal F$ con $H\subset F\cap G$. Si una red está frecuentemente en cada elemento de $\mathcal F$, tiene una subred que está eventualmente en todos ellos.
\end{lema}
\begin{proof}
Considere $B=\{(\alpha,F):F\in\mathcal F,\ x_\alpha\in F\}$, ordenado por
\[(\alpha,F)\preceq(\beta,G)\quad\Longleftrightarrow\quad
\alpha\le\beta\text{ y }G\subset F.
\]
Es no vacío por frecuencia. Dados dos pares, elija $H$ contenido en la intersección de sus conjuntos y un índice posterior a los dos índices. Por frecuencia en $H$, hay un índice aún posterior $\gamma$ con $x_\gamma\in H$; el par $(\gamma,H)$ es una cota común. La proyección $a(\alpha,F)=\alpha$ satisface la condición de subred: para cualquier $\alpha_0$, la frecuencia permite hallar un par con primer índice posterior a $\alpha_0$, y todos los pares posteriores conservan esa propiedad. Finalmente, fijado $F\in\mathcal F$, hay un par $(\alpha,F)$; después de él los conjuntos del segundo componente están contenidos en $F$, y sus valores de la red pertenecen a $F$.
\end{proof}

\begin{proposicion}\label{v3:acumulacion-subred}
Un punto es de acumulación de una red si y sólo si alguna subred converge a él.
\end{proposicion}
\begin{proof}
Si es de acumulación, el lema anterior se aplica a sus vecindades, cerradas bajo intersección finita, y produce una subred eventualmente en cada una. Recíprocamente, si $x_{a(\beta)}\to x$, dados una vecindad $U$ de $x$ y un índice $\alpha_0$, la condición de subred asegura $a(\beta)\ge\alpha_0$ eventualmente y la convergencia asegura $x_{a(\beta)}\in U$ eventualmente. Una cota común de ambos umbrales da un índice con las dos propiedades. Así la red original está frecuentemente en $U$.
\end{proof}

\begin{teorema}\label{v3:compacto-redes}
Un espacio topológico es compacto si y sólo si toda red tiene una subred convergente.
\end{teorema}
\begin{proof}
Si el espacio es compacto y $(x_\alpha)$ es una red, las cerraduras de sus colas $C_\alpha=\cl{\{x_\beta:\beta\ge\alpha\}}$ tienen la propiedad de intersección finita: una cola posterior a finitísimos índices está contenida en todas ellas. La compacidad implica $\bigcap_\alpha C_\alpha\ne\varnothing$. Un punto de esta intersección tiene toda vecindad intersectando cada cola, por tanto es de acumulación y admite una subred convergente.

En sentido contrario, sea $\mathcal F$ una familia de cerrados con la propiedad de intersección finita. Su familia $\mathcal G$ de intersecciones finitas no vacías está dirigida por inclusión \emph{inversa}. Elijamos $x_G\in G$ para cada $G\in\mathcal G$. Por hipótesis una subred converge a algún $x$. Para cada $F\in\mathcal F$, la red original está eventualmente en $F$, y también la subred. Un límite de una red eventualmente contenida en un cerrado pertenece a ese cerrado: si no, el complemento abierto contendría eventualmente la red y contradiría la eventual pertenencia al cerrado. Así $x\in\bigcap\mathcal F$, y el criterio de intersección finita prueba compacidad. Las colas deben cerrarse en la primera dirección, y el orden de la segunda es inverso, precisiones que no deben perderse en la notación.
\end{proof}

\section{Filtros y redes universales}\index{filtro}\index{ultrafiltro}\index{red!universal}
Un \emph{filtro} $\mathcal U$ en $E$ es una familia no vacía de subconjuntos que no contiene $\varnothing$, es cerrada bajo intersecciones finitas y contiene todo superconjunto de cualquiera de sus miembros. Es un \emph{ultrafiltro} si para todo $B\subset E$, uno de $B$ y $E\setminus B$ pertenece a $\mathcal U$. No pueden pertenecer ambos, porque su intersección es vacía.

\begin{lema}\label{v3:ultrafiltro-existencia}
Todo filtro está contenido en un ultrafiltro. Toda red tiene una subred universal, es decir, una subred que para cada $B\subset E$ está eventualmente en $B$ o eventualmente en su complemento.
\end{lema}
\begin{proof}
Ordenemos por inclusión los filtros que contienen al filtro dado. La unión de una cadena es un filtro: para dos miembros se elige un filtro de la cadena que contenga a ambos, y allí se toman sus intersecciones. Zorn da uno maximal $\mathcal U$. Si ni $B$ ni su complemento pertenecieran a él, todo $U\in\mathcal U$ intersectaría $B$; de lo contrario $U\subset E\setminus B$ forzaría que el complemento perteneciera al filtro. Las intersecciones $U\cap B$ generarían entonces un filtro propio estrictamente mayor, contradicción. Así el maximal es ultrafiltro.

Los conjuntos en los que una red está eventualmente forman un filtro: la intersección de dos eventualidades se obtiene tomando una cota común de sus umbrales; el vacío no puede contener una cola. Extiéndalo a un ultrafiltro $\mathcal U$. La red está frecuentemente en cada $U\in\mathcal U$, pues toda imagen de una cola pertenece al filtro inicial y su intersección con $U$ no puede ser vacía. El lema~\ref{v3:subred-familia} produce una subred eventualmente en todo $U\in\mathcal U$. Para cada $B$, el ultrafiltro contiene $B$ o su complemento; por tanto la subred es universal.
\end{proof}
Una red universal que tiene un punto de acumulación converge a él: cada vecindad es visitada frecuentemente, así que la universalidad impide que la red esté eventualmente en su complemento y obliga a la eventualidad dentro de la vecindad.

\section{Ejercicios de redes y filtros}
\begin{ejercicio}\label{ej:v3-l14-1}
En un espacio Hausdorff, una red convergente tiene límite único.
\end{ejercicio}
\begin{solucion}
Si convergiera a $x\ne y$, elija vecindades abiertas disjuntas $U,V$ de esos puntos. La red estaría eventualmente en cada una, por tanto eventualmente en su intersección vacía al tomar una cota común de los dos umbrales, imposible. También vale el recíproco mencionado en la lectura: si $x\ne y$ no admiten vecindades disjuntas, dirija los pares $(U,V)$ de vecindades por inclusión inversa y elija $z_{U,V}\in U\cap V$. Esa red converge simultáneamente a $x$ e $y$, contradiciendo unicidad de límites.
\end{solucion}

\begin{ejercicio}\label{ej:v3-l14-2}
Para $A\subset E$, pruebe: (i) $x$ es interior a $A$ si y sólo si ninguna red en $E\setminus A$ converge a $x$; (ii) $x\in\cl A$ si y sólo si alguna red en $A$ converge a $x$.
\end{ejercicio}
\begin{solucion}
(ii) Si una red en $A$ converge a $x$, toda vecindad de $x$ contiene algún término y, por tanto, intersecta $A$. Recíprocamente, si toda vecindad $U$ de $x$ intersecta $A$, elija $x_U\in A\cap U$ y ordene las vecindades por inclusión inversa. Dada una vecindad $V$, después del índice $V$ todos los $U$ están contenidos en $V$, luego $x_U\in V$. La red converge a $x$.

(i) Si $x$ tiene una vecindad contenida en $A$, ninguna red íntegramente en el complemento puede estar eventualmente en ella. Si $x$ no es interior, toda vecindad intersecta el complemento, o sea $x\in\cl{E\setminus A}$. El inciso (ii) construye entonces una red en ese complemento convergente a $x$.
\end{solucion}

\begin{ejercicio}\label{ej:v3-l14-3}
Una aplicación $f:E\to Y$ es continua si y sólo si preserva los límites de todas las redes.
\end{ejercicio}
\begin{solucion}
Si es continua y $x_\alpha\to x$, la preimagen de cualquier vecindad de $f(x)$ es una vecindad de $x$, de modo que $f(x_\alpha)\to f(x)$. Si no es continua en $x$, existe una vecindad abierta $V$ de $f(x)$ tal que ninguna vecindad $U$ de $x$ está contenida en $f^{-1}(V)$. Elija $x_U\in U\setminus f^{-1}(V)$; la red dirigida por inclusión inversa converge a $x$, pero sus imágenes nunca pertenecen a $V$, así que no convergen a $f(x)$.
\end{solucion}

\begin{ejercicio}\label{ej:v3-l14-4}
Un espacio topológico es compacto si y sólo si toda red universal converge.
\end{ejercicio}
\begin{solucion}
En un compacto, cualquier red tiene un punto de acumulación por el teorema~\ref{v3:compacto-redes}; si es universal, la observación posterior al lema~\ref{v3:ultrafiltro-existencia} muestra que converge a ese punto. Recíprocamente, cada red tiene una subred universal por ese lema. Si todas las universales convergen, toda red tiene una subred convergente y el teorema~\ref{v3:compacto-redes} da compacidad. Así se justifica también la existencia de subredes universales, que la lectura sólo menciona.
\end{solucion}

\begin{ejercicio}\label{ej:v3-l14-5}
Pruebe que los conjuntos en los que una red está eventualmente forman un filtro. Recíprocamente, dado un filtro $\mathcal F$, pruebe que elegir $x_F\in F$ define una red.
\end{ejercicio}
\begin{solucion}
El conjunto total pertenece a la familia de eventualidades y el vacío no. Si la red está eventualmente en $F$ y $F\subset G$, está eventualmente en $G$. Si está eventualmente en $F$ y $G$, después de una cota común de los umbrales está en $F\cap G$. Éstos son los axiomas del filtro. A la inversa, ordene $\mathcal F$ por inclusión inversa. Es un conjunto dirigido porque $F\cap G\in\mathcal F$ y está contenido en ambos. Cada miembro es no vacío, por lo que puede elegirse $x_F\in F$; la aplicación $F\mapsto x_F$ es una red. Su filtro de eventualidades contiene a $\mathcal F$, aunque para elecciones arbitrarias no tiene por qué ser exactamente igual a él.
\end{solucion}

\begin{ejercicio}\label{ej:v3-l14-6}
Un filtro converge a $x$ cuando contiene todas las vecindades de $x$. Pruebe que una red converge a $x$ si y sólo si su filtro de eventualidades converge a $x$.
\end{ejercicio}
\begin{solucion}
Por definición, una vecindad $U$ pertenece al filtro de eventualidades exactamente cuando la red está eventualmente en $U$. Exigir esto para todas las vecindades de $x$ es precisamente exigir convergencia de la red a $x$. La frase ``converge a $x$'' completa la omisión verbal del enunciado original.
\end{solucion}

\begin{ejercicio}\label{ej:v3-l14-7}
Si $y_\beta=x_{a(\beta)}$ es una subred, su filtro de eventualidades contiene al de la red original.
\end{ejercicio}
\begin{solucion}
Si $x_\alpha\in F$ para todos los $\alpha\ge\alpha_0$, la condición de subred proporciona $\beta_0$ con $a(\beta)\ge\alpha_0$ para $\beta\ge\beta_0$. Para esos índices, $y_\beta=x_{a(\beta)}\in F$. Así toda eventualidad de la red original es una eventualidad de la subred, que es exactamente la inclusión de filtros pedida.
\end{solucion}

\begin{ejercicio}\label{ej:v3-l14-8}
Una red es universal si y sólo si su filtro asociado es un ultrafiltro. Deduzca: (i) todo filtro está contenido en un ultrafiltro; (ii) un espacio es compacto si y sólo si todo ultrafiltro converge.
\end{ejercicio}
\begin{solucion}
La equivalencia traduce literalmente la definición: para cada $B$, la red es eventualmente en $B$ o en su complemento si y sólo si uno de ellos pertenece al filtro de eventualidades. Para (i), el argumento de Zorn del lema~\ref{v3:ultrafiltro-existencia} da una extensión maximal y demuestra que decide cada subconjunto; ésa es la extensión requerida. También puede construirse la red seleccionada del ejercicio~\ref{ej:v3-l14-5} y tomar una subred universal: su ultrafiltro de eventualidades contiene al filtro inicial por el ejercicio~\ref{ej:v3-l14-7}.

(ii) En un compacto, sea $\mathcal U$ un ultrafiltro y elijamos la red $x_U\in U$ ordenada por inclusión inversa. Su filtro de eventualidades contiene a $\mathcal U$ y, por maximalidad del ultrafiltro, es igual a él. La red es universal y converge por el ejercicio~\ref{ej:v3-l14-4}; el ejercicio~\ref{ej:v3-l14-6} da convergencia de $\mathcal U$. Recíprocamente, si todo ultrafiltro converge, el filtro asociado a cualquier red universal converge, así que esa red converge; el ejercicio~\ref{ej:v3-l14-4} da compacidad.
\end{solucion}

\section{Convergencia débil y el papel de la acotación}\label{lec:v3-15}
\fuentelectura{15}{Convergencia débil}
\begin{proposicion}\label{v3:convergencia-debil}
Una red $x_\alpha$ converge débilmente a $x$, escrito $x_\alpha\rightharpoonup x$, si y sólo si $\phi(x_\alpha)\to\phi(x)$ para toda $\phi\in X^*$. Una red $\phi_\alpha$ converge débil estrella a $\phi$ si y sólo si $\phi_\alpha(x)\to\phi(x)$ para todo $x\in X$.
\end{proposicion}
\begin{proof}
La convergencia en la topología inicial implica convergencia de todas las aplicaciones que la generan, por continuidad. Para el recíproco, una vecindad básica sólo impone finitísimas desigualdades. La convergencia de cada coordenada da un umbral para cada una, y la dirección del conjunto de índices permite tomar una cota común para todos ellos, por inducción finita. Después de esa cota se cumplen simultáneamente las desigualdades y la red está en la vecindad. El mismo razonamiento vale para las evaluaciones del dual.
\end{proof}
La convergencia en norma implica la débil porque $|\phi(x_\alpha-x)|\le\norm\phi\norm{x_\alpha-x}$. La recíproca falla: en $L^1[0,2\pi]$, $f_n(t)=1+\cos(nt)$ converge débilmente a uno por Riemann--Lebesgue, probado en el ejercicio~\ref{ej:v3-l15-1}, y por la identificación $L^1{}^*=L^\infty$. En efecto, cada $g\in L^\infty$ pertenece también a $L^1$ en este intervalo finito y $\int g\cos(nt)\to0$. Sin embargo
\[\norm{f_n-1}_1=\frac1n\int_0^{2\pi n}|\cos u|\,du=4,
\]
pues hay $n$ períodos de integral cuatro. Por tanto no hay convergencia en norma.

\begin{proposicion}\label{v3:debil-acotacion-secuencias}
Una \emph{sucesión} débilmente convergente en un espacio normado es acotada en norma. Una sucesión débil estrella convergente en $X^*$ es acotada si $X$ es Banach. Para cualquier \emph{red} débilmente convergente,
\[\norm x\le\liminf_\alpha\norm{x_\alpha},
\]
y para cualquier red débil estrella convergente, $\norm\phi\le\liminf_\alpha\norm{\phi_\alpha}$.
\end{proposicion}
\begin{proof}
Si $x_n\rightharpoonup x$, para cada $\psi\in X^*$ la sucesión escalar $\psi(x_n)$ converge y es acotada. Aplicamos acotación uniforme a las evaluaciones $J_Xx_n$ sobre el Banach $X^*$; sus normas son $\norm{x_n}$, por Hahn--Banach. En el caso débil estrella aplicamos acotación uniforme directamente a las funcionales $\phi_n$ sobre el Banach $X$.

Para la semicontinuidad con redes, fijemos $\psi\in X^*$ de norma a lo sumo uno. Dado $\eps>0$, eventualmente $|\psi(x)|\le|\psi(x_\alpha)|+\eps\le\norm{x_\alpha}+\eps$. Esto implica $|\psi(x)|\le\liminf_\alpha\norm{x_\alpha}+\eps$, donde $\liminf_\alpha a_\alpha=\sup_{\alpha_0}\inf_{\alpha\ge\alpha_0}a_\alpha$. Haciendo $\eps\downarrow0$ y tomando el supremo en $\psi$, obtenemos la desigualdad. La versión débil estrella usa de igual modo $|\phi(x)|=\lim|\phi_\alpha(x)|$ para cada $\norm x\le1$, y toma el supremo en $x$.
\end{proof}

\begin{ejemplo}[Una red débilmente convergente sin cota eventual]\label{v3:red-debil-no-acotada}
En un espacio normado infinito-dimensional $X$, dirijamos los pares $(F,n)$, donde $F\subset X^*$ es finito y $n\ge1$, mediante inclusión de $F$ y orden usual de $n$. Para cada $F$, la intersección $\bigcap_{\phi\in F}\ker\phi$ contiene un vector unitario $u_F$, por el argumento del ejercicio~\ref{ej:v3-l13-4}. Definamos $x_{F,n}=nu_F$. Para cada $\phi$, después de incluirla en $F$ se tiene $\phi(x_{F,n})=0$; así $x_{F,n}\rightharpoonup0$. Pero $\norm{x_{F,n}}=n\to\infty$, de modo que ni siquiera existe una cola acotada en norma. Esto refuta el inciso (i) de la proposición 4 original para redes y la necesidad de acotación global del teorema 5 original.

Incluso una red escalar convergente puede tener un conjunto total de valores no acotado: en $\N^2$ con orden producto, sea $a_{m,n}=n$ si $m=1$ y cero si $m\ge2$. Converge a cero, pero sus valores incluyen todos los enteros positivos. Por tanto el argumento ``converge, luego su rango es acotado'', válido para sucesiones, no se traslada a redes.
\end{ejemplo}

\begin{proposicion}[Comprobar convergencia en un conjunto total]\label{v3:total-debil}
Sea $D\subset X^*$ con span denso en norma. Para una red \emph{acotada en norma}, $x_\alpha\rightharpoonup x$ equivale a $\phi(x_\alpha)\to\phi(x)$ para cada $\phi\in D$. Para sucesiones, la convergencia débil equivale a la conjunción de acotación en norma y convergencia en $D$. Análogamente, en una red normoacotada de $X^*$ basta comprobar la convergencia débil estrella en un conjunto cuyo span sea normodenso en $X$.
\end{proposicion}
\begin{proof}
La necesidad de las evaluaciones es inmediata; para sucesiones, la necesidad de acotación es la proposición anterior. Para la suficiencia, sea $C=\sup_\alpha\norm{x_\alpha-x}<\infty$. Si $C=0$ no hay nada que probar. Dada $\psi\in X^*$ y $\eps>0$, elija $\phi\in\spanop D$ con $\norm{\psi-\phi}<\eps/(2C)$. La convergencia para cada elemento de $D$ implica la de cada combinación finita, de modo que eventualmente $|\phi(x_\alpha-x)|<\eps/2$. Entonces
\[|\psi(x_\alpha-x)|\le\norm{\psi-\phi}\norm{x_\alpha-x}+|\phi(x_\alpha-x)|<\eps.
\]
La prueba débil estrella intercambia vectores y funcionales, con una cota uniforme de las normas de las funcionales. También basta una cota en alguna cola, pues la convergencia sólo depende de las colas.
\end{proof}
Si ``base'' significa base de Hamel de $X^*$, cada funcional ya es una combinación \emph{finita} de elementos de la base; en ese caso la convergencia de las evaluaciones basta para redes sin ninguna acotación. La formulación con conjunto total normodenso cubre también las bases topológicas y hace explícito dónde se necesita la estimación uniforme.

\begin{proposicion}\label{v3:operador-debil-continuo}
Para una aplicación lineal $T:X\to Y$ entre Banach son equivalentes: (i) $T$ es acotada; (ii) es continua de $\sigma(X,X^*)$ a $\sigma(Y,Y^*)$; (iii) lleva redes débilmente convergentes a cero en redes débilmente convergentes a cero.
\end{proposicion}
\begin{proof}
(i)$\Rightarrow$(ii): para $\psi\in Y^*$, $\psi\circ T\in X^*$; la propiedad universal de la topología inicial da continuidad. (ii)$\Rightarrow$(iii) es conservación de límites. Para (iii)$\Rightarrow$(i), si $x_n\to0$ en norma, también converge débilmente a cero; por (iii), $Tx_n\rightharpoonup0$. Es una \emph{sucesión}, por lo que sus normas están acotadas. El ejercicio~\ref{ej:v3-l15-3}, demostrado abajo, implica que $T$ es continuo. Este argumento sólo utiliza la acotación secuencial válida, sin recurrir a la afirmación falsa para redes.
\end{proof}

\section{Ejercicios de convergencia débil}
\begin{ejercicio}[Riemann--Lebesgue]\label{ej:v3-l15-1}
Para $f\in L^1[0,2\pi]$, pruebe $\int_0^{2\pi}f(t)\cos(nt)\,dt\to0$.
\end{ejercicio}
\begin{solucion}
Las funciones $\cos(nt)/\sqrt\pi$, $n\ge1$, son ortonormales en $L^2[0,2\pi]$: la identidad $2\cos(nt)\cos(mt)=\cos((n-m)t)+\cos((n+m)t)$ y la integral de un coseno de frecuencia entera no nula dan ortogonalidad y norma cuadrada $\pi$. Si $g\in L^2$, Bessel implica que los coeficientes $\int g(t)\cos(nt)\,dt$ tienden a cero. Dado $f\in L^1$ y $\eps>0$, elija una función continua $g$ con $\norm{f-g}_1<\eps/2$; esa densidad se obtiene aproximando funciones simples por continuas mediante regularidad de la medida de Lebesgue. Como $g\in L^2$, para $n$ grande su coeficiente tiene módulo menor que $\eps/2$. Entonces
\[\left|\int f\cos(nt)\right|\le\norm{f-g}_1+\left|\int g\cos(nt)\right|<\eps.
\]
El factor correcto es $\cos(nt)$; el enunciado original imprime $\cos(t)$, cuya integral no depende de $n$ y en general no es cero, como muestra $f(t)=\cos t$.
\end{solucion}

\begin{ejercicio}\label{ej:v3-l15-2}
Si $K$ es compacto Hausdorff y $f_n,f\in C(K)$, pruebe que $f_n\rightharpoonup f$ si y sólo si $(f_n)$ es uniformemente acotada y $f_n(x)\to f(x)$ para cada $x\in K$.
\end{ejercicio}
\begin{solucion}
Si hay convergencia débil, la proposición~\ref{v3:debil-acotacion-secuencias} da $\sup_n\norm{f_n}_\infty<\infty$. Las evaluaciones $\delta_x:g\mapsto g(x)$ son funcionales continuas, así que $f_n(x)\to f(x)$.

Recíprocamente, por el teorema de representación de Riesz para $C(K)$, toda funcional es integración respecto de una medida de Borel regular finita real o compleja $\mu$. Si $\norm{f_n}_\infty\le C$, también $|f(x)|\le C$ por el límite puntual. Así $|f_n-f|\le2C$, integrable respecto de la variación total finita $|\mu|$. Convergencia dominada da
\[\left|\int(f_n-f)\,d\mu\right|\le\int|f_n-f|\,d|\mu|\longrightarrow0.
\]
Por tanto todas las funcionales convergen y la sucesión converge débilmente. Se usa aquí la representación de medidas desarrollada en el volumen II; también puede consultarse \cite{rudin1991}. La hipótesis de que el límite puntual $f$ sea continuo es parte del enunciado y no se deduce de la mera acotación.
\end{solucion}

\begin{ejercicio}\label{ej:v3-l15-3}
Una aplicación lineal entre espacios normados es continua si y sólo si lleva toda sucesión que converge a cero en una sucesión acotada.
\end{ejercicio}
\begin{solucion}
Si $T$ es continua, $x_n\to0$ implica $Tx_n\to0$, y toda sucesión convergente es acotada. Si $T$ no es continua, no existe una cota global de operador. Para cada $n$ escoja $u_n$ unitario con $\norm{Tu_n}>n^3$. Entonces $x_n=u_n/n$ tiene norma $1/n\to0$, pero $\norm{Tx_n}>n^2$, que no es acotada. Esto prueba el contrapositivo sin requerir completitud.
\end{solucion}

\begin{ejercicio}\label{ej:v3-l15-4}
Todo conjunto convexo cerrado en norma es débilmente cerrado.
\end{ejercicio}
\begin{solucion}
Si $C$ es vacío o todo el espacio, es inmediato. Para $x\notin C$, el teorema de separación estricta entre el compacto $\{x\}$ y el convexo cerrado $C$ da una funcional continua real $f$ y un real $a$ con $f(x)>a\ge\sup_{c\in C}f(c)$, cambiando el signo si fuera necesario. En un espacio complejo, $f=\Re\phi$ para una funcional compleja continua $\phi$, por la conversión probada en el capítulo~\ref{cap:v3-03}. El conjunto $\{y:\Re\phi(y)>a\}$, o $\{y:f(y)>a\}$ en el caso real, es débilmente abierto, contiene a $x$ y no intersecta $C$. Por tanto el complemento de $C$ es débilmente abierto. No se necesita que una bola cerrada sea compacta en dimensión infinita: separar el punto de $C$ evita esa confusión.
\end{solucion}

Las redes recuperan la relación entre cierre, continuidad y compacidad en topologías que no son métricas. A la vez exigen respetar la distinción entre una cola y el conjunto completo de valores, y entre acotación secuencial y acotación de redes. Con estas precisiones, podemos abordar la compacidad débil estrella de la bola del dual y los criterios de reflexividad del capítulo siguiente.

\chapter{Compacidad débil y reflexividad}\label{cap:v3-09}
Una bola cerrada de dimensión infinita deja de ser compacta en norma, pero puede recuperar compacidad cuando sólo exigimos convergencia de las mediciones lineales. Este capítulo explica tres pasos distintos: Alaoglu da compacidad en un dual; Eberlein--Šmulian permite reconocer compacidad débil mediante sucesiones; la reflexividad determina cuándo el espacio original ocupa todo su bidual. No deben confundirse las topologías ni sustituirse redes por sucesiones sin justificarlo.

Escribimos $B_X=\{x\in X:\norm x\leq1\}$ y $J_X:X\to X^{**}$ para la isometría canónica $J_Xx(\varphi)=\varphi(x)$. La norma de $J_Xx$ es $\norm x$ por Hahn--Banach. En particular, la topología débil de $X$ coincide, bajo $J_X$, con la topología débil* de $X^{**}$ restringida a $J_X(X)$: ambas comparan exactamente los números $\varphi(x)$, $\varphi\in X^*$.

\section{El teorema de Banach--Alaoglu}\label{lec:v3-16}
\fuentelectura{16}{El teorema de Banach--Alaoglu}
\index{Banach--Alaoglu, teorema de}\index{compacidad!débil*}
\begin{lema}[Tychonoff]\label{v3:tychonoff}
El producto de una familia de espacios compactos es compacto para la topología producto.
\end{lema}
\begin{proof}
Si un factor es vacío, el producto es vacío y compacto. En otro caso el producto es no vacío por elección. Sea $\mathcal U$ un ultrafiltro sobre el producto $P=\prod_{i\in I}K_i$. Para cada proyección $p_i$, la familia $\mathcal U_i=\{A\subset K_i:p_i^{-1}(A)\in\mathcal U\}$ es un ultrafiltro: las propiedades de intersección y monotonía pasan a las preimágenes, y para cada $A$ el ultrafiltro $\mathcal U$ contiene exactamente una de $p_i^{-1}(A)$ y su complemento. Por compacidad de $K_i$ y el ejercicio \ref{ej:v3-l14-8}, $\mathcal U_i$ converge a algún $x_i\in K_i$. Elija tales puntos y ponga $x=(x_i)$. Toda vecindad básica de $x$ contiene una intersección finita de conjuntos $p_i^{-1}(V_i)$, con $V_i$ vecindades de $x_i$. Cada uno pertenece a $\mathcal U$ por convergencia de $\mathcal U_i$, y su intersección también pertenece a $\mathcal U$. Así $\mathcal U$ converge a $x$. El mismo criterio de ultrafiltros demuestra que $P$ es compacto. La existencia de ultrafiltros y la equivalencia utilizada se desarrollaron en el capítulo anterior; la finitud de las restricciones en una vecindad básica es el paso decisivo.
\end{proof}

\begin{teorema}[Banach--Alaoglu]\label{v3:alaoglu}
Para todo espacio normado $X$, la bola $B_{X^*}$ es compacta en $\sigma(X^*,X)$, la topología débil*. No se requiere que $X$ sea completo.
\end{teorema}
\begin{proof}
Para cada $x\in X$, sea $D_x=\{z\in\F:|z|\leq\norm x\}$, un compacto. Tychonoff afirma que $P=\prod_{x\in X}D_x$ es compacto en la topología producto. Sus puntos son funciones $u:X\to\F$ con $|u(x)|\leq\norm x$. Dentro de $P$, las ecuaciones
\[
 u(x+y)-u(x)-u(y)=0,\qquad u(\lambda x)-\lambda u(x)=0
\]
definen cerrados, porque cada expresión depende continuamente de un número finito de coordenadas. Su intersección para todos $x,y,\lambda$ es exactamente el conjunto de las funciones lineales de norma a lo sumo uno: la cota que define $P$ asegura su continuidad. Por tanto esta intersección es $B_{X^*}$ y es compacta. La topología producto restringida a ella es la de convergencia puntual sobre $X$, es decir, la débil*. Esto completa la identificación topológica que necesita la prueba.
\end{proof}
\begin{corolario}\label{v3:alaoglu-cor}
Toda bola cerrada de $X^*$ es débil*-compacta. Todo conjunto acotado en norma y \emph{débil*-cerrado} es débil*-compacto.
\end{corolario}
\begin{proof}
Las traslaciones y las multiplicaciones por escalares no nulos son homeomorfismos para la topología débil*; el caso de radio cero es un singleton. Un conjunto acotado está contenido en una bola cerrada suficientemente grande, y un cerrado de un compacto es compacto.
\end{proof}
\begin{observacion}[Qué significa «cerrado»]
El corolario no vale si sólo se exige cerradura en norma. En $X=\ell^2$, los funcionales coordenados $\varphi_n(x)=x_n$ tienen norma uno y convergen débil* a cero, pues $x_n\to0$ para todo $x\in\ell^2$. La esfera unitaria de $X^*$ es normacerrada y acotada, pero no débil*-cerrada; tampoco puede ser débil*-compacta en el espacio Hausdorff $X^*$. Así se precisa el segundo inciso del corolario de la lectura 16.
\end{observacion}

\begin{proposicion}[Acotación débil*]\label{v3:acotacion-debil-estrella}
Si $X$ es Banach, un conjunto $A\subset X^*$ es acotado en norma si y sólo si para toda vecindad débil* $U$ de cero existe $r>0$ tal que $A\subset tU$ para todo $t\geq r$.
\end{proposicion}
\begin{proof}
Si $\norm\varphi\leq M$ en $A$ y $U$ contiene una vecindad básica definida por $|\varphi(x_j)|<\eps$ para $j=1,\ldots,m$, basta tomar $r>M\max_j\norm{x_j}/\eps$. Entonces $|\varphi(x_j)/t|<\eps$ para $t\geq r$. Si $M=0$ o todos los $x_j$ son cero, cualquier $r>0$ sirve.

Recíprocamente, aplique la condición a $U_x=\{\varphi:|\varphi(x)|<1\}$ para cada $x$. Se obtiene $\sup_{\varphi\in A}|\varphi(x)|<\infty$. El principio de acotación uniforme, con dominio Banach $X$, da $\sup_{\varphi\in A}\norm\varphi<\infty$. Para $A=\varnothing$ la afirmación es inmediata.
\end{proof}

Una red $(\varphi_\alpha)$ es de Cauchy débil* si para cada $x\in X$ la red escalar $(\varphi_\alpha(x))$ es de Cauchy. La equivalencia con la definición por vecindades se debe a que sólo hay un número finito de mediciones en cada vecindad básica. Una sucesión de Cauchy débil* es acotada en norma cuando $X$ es Banach: cada sucesión escalar es acotada, y se aplica acotación uniforme. Para redes, esta última conclusión \emph{no} es válida en general.

\begin{ejemplo}[Una red convergente sin cola acotada]\label{v3:red-no-acotada}
Sea $X$ infinito dimensional. Dirijamos los pares $(F,m)$, donde $F\subset X$ es finito y $m\in\N$, por inclusión en la primera coordenada y orden usual en la segunda. Para cada $F$, Hahn--Banach proporciona $\psi_F\in X^*$ de norma uno que se anula en $\spanop F$: ese subespacio es finito dimensional, cerrado y propio. Definamos $\varphi_{F,m}=m\psi_F$. Dado $x$, todos los índices posteriores a $(\{x\},1)$ tienen $\varphi_{F,m}(x)=0$. Por tanto la red converge débil* a cero y es de Cauchy débil*. Sin embargo, cualquier cola contiene los pares $(F_0,m)$ para $m$ arbitrariamente grande, y $\norm{\varphi_{F_0,m}}=m$. No hay cola acotada.

El error de la observación de la lectura consiste en tratar los índices anteriores a un umbral de una red como si fueran un conjunto finito. Eso sí ocurre para una sucesión; no ocurre para una red arbitraria.
\end{ejemplo}

\begin{teorema}[Representación mediante funciones continuas]\label{v3:encaje-CK}
Todo espacio normado $X$ admite un encaje lineal isométrico en $C(K)$ para algún compacto Hausdorff $K$. Si $X$ es Banach, su imagen es cerrada.
\end{teorema}
\begin{proof}
Tomemos $K=B_{X^*}$ con su topología débil*, compacto por Alaoglu y Hausdorff porque las evaluaciones separan funcionales diferentes. Definamos $(Tx)(\varphi)=\varphi(x)$. La definición de la topología débil* asegura que $Tx$ es continua; la linealidad de cada $\varphi$ da la linealidad de $T$. Hahn--Banach implica
\[
 \norm{Tx}_\infty=\sup_{\norm\varphi\leq1}|\varphi(x)|=\norm x.
\]
Así $T$ es una isometría. Una imagen isométrica de un espacio completo es completa y, por tanto, cerrada en el espacio métrico $C(K)$.
\end{proof}

\begin{teorema}[Metrizabilidad de la bola dual]\label{v3:bola-dual-metrizable}
La bola $B_{X^*}$ con topología débil* es metrizable si y sólo si $X$ es separable.
\end{teorema}
\begin{proof}
Si $X$ es separable, elijamos una sucesión $(x_n)$ densa en $X$ y pongamos
\[
 d(\varphi,\psi)=\sum_{n=1}^{\infty}2^{-n}
          \min\{1,|\varphi(x_n)-\psi(x_n)|\}.
\]
La serie converge uniformemente. Es una métrica: positividad y simetría son claras; la desigualdad triangular usa $\min(1,a+b)\leq\min(1,a)+\min(1,b)$; y $d=0$ obliga a que los funcionales coincidan en un denso y, por continuidad, en todo $X$.

La convergencia débil* implica convergencia para $d$: se elige primero una cola de la serie pequeña y después se controlan las finitas coordenadas iniciales. Recíprocamente, la convergencia para $d$ controla cada $x_n$. Para un $x$ arbitrario y $\varphi,\psi\in B_{X^*}$,
\[
 |(\varphi-\psi)(x)|\leq2\norm{x-x_n}
                 +|(\varphi-\psi)(x_n)|.
\]
La densidad da primero un $n$ adecuado; el segundo término se controla con $d$. Aplicado a un número finito de $x$, esto demuestra igualdad de las dos topologías, para redes y no sólo para sucesiones.

En sentido inverso, supongamos que una métrica induce la topología débil* en la bola. Cada vecindad métrica $U_n$ de cero de radio $1/n$ contiene una vecindad básica relativa definida por un conjunto finito $S_n\subset X$. Si un funcional en la bola se anula en $S=\bigcup_nS_n$, pertenece a todos los $U_n$ y es cero. Normalizar muestra que todo funcional continuo que se anula en $S$ es cero. Por Hahn--Banach, el subespacio $\overline{\spanop S}$ debe ser todo $X$: si fuera propio, existiría un funcional no nulo que lo anulara. Las combinaciones finitas de $S$ con coeficientes racionales o racionales complejos forman un conjunto numerable denso. Así $X$ es separable.
\end{proof}

\section{Ejercicios de Alaoglu y metrizabilidad}
\begin{ejercicio}[Lectura 16, ejercicio 1; formulación precisa]\label{ej:v3-l16-1}
Sean $\tau$ y $\tau'$ topologías en un mismo conjunto, con $(K,\tau)$ compacto y $(K,\tau')$ Hausdorff. Si $\tau'\subseteq\tau$, pruebe mediante redes que $\tau'=\tau$.
\end{ejercicio}


\begin{ejercicio}[Lectura 16, ejercicio 2]\label{ej:v3-l16-2}
Sea $(f_n:X\to Y_n)$ una familia numerable que separa puntos y suponga que cada $Y_n$ es métrico. Construya una métrica en $X$ que induzca la topología inicial de estas aplicaciones. Deduzca que un compacto con una familia numerable de funciones \emph{continuas} complejas que separa puntos es metrizable.
\end{ejercicio}


\begin{ejercicio}[Lectura 16, ejercicio 3]\label{ej:v3-l16-3}
Si $X$ es Banach, pruebe que toda sucesión de Cauchy débil* en $X^*$ converge débil*.
\end{ejercicio}


\begin{ejercicio}[Lectura 16, ejercicio 4]\label{ej:v3-l16-4}
Sean $1<p<\infty$, $1/p+1/q=1$, y $(u_\alpha)$ una red en $\ell^q=(\ell^p)^*$. Pruebe: (i) la convergencia débil* a $u$ implica convergencia de cada coordenada; (ii) si la red es acotada en norma y converge coordenada a coordenada a $u\in\ell^q$, converge débil* a $u$.
\end{ejercicio}


\begin{ejercicio}[Lectura 16, ejercicio 5]\label{ej:v3-l16-5}
Demuestre las dos afirmaciones anteriores para $\ell^1=(c_0)^*$ y $\ell^\infty=(\ell^1)^*$, indicando en cada caso el predual.
\end{ejercicio}


\begin{ejercicio}[Lectura 16, ejercicio 6]\label{ej:v3-l16-6}
Pruebe que todo espacio normado separable se encaja lineal e isométricamente (i) en $C(\mathcal C)$, donde $\mathcal C$ es el conjunto de Cantor; (ii) en $C([0,1])$.
\end{ejercicio}


\section{Familias normantes y Eberlein--Šmulian}\label{lec:v3-17}
\fuentelectura{17}{El teorema de Eberlein--Šmulian}
\index{familia!normante}\index{Eberlein--Smulian@Eberlein--Šmulian, teorema de}
\begin{definicion}
Un conjunto $A\subset X^*$ es $\alpha$-normante para un subespacio $Y\subset X$, donde $0<\alpha\le1$, si
\[
 \alpha\norm y\leq\sup_{\varphi\in A\setminus\{0\}}
       \frac{|\varphi(y)|}{\norm\varphi}\qquad(y\in Y).
\]
Cuando los funcionales tienen norma uno, el denominador desaparece. La palabra «normaliza» de la lectura expresa esta misma propiedad. Para $\alpha>0$, una familia normante separa puntos de $Y$.
\end{definicion}
\begin{lema}[Selección desde una familia normante]\label{v3:seleccion-normante}
Sea $Y\ne\{0\}$ un subespacio de $X$ y sea $A\subset X^*$ una familia $1$-normante de funcionales de norma uno. Si $0<\eps<1/2$ y $D$ es una $\eps$-red de la esfera unitaria de $Y$, elija para cada $v\in D$ un $\varphi_v\in A$ con $|\varphi_v(v)|>1-\eps$. Entonces $\{\varphi_v:v\in D\}$ es $(1-2\eps)$-normante para $Y$.
\end{lema}
\begin{proof}
La elección es posible por $1$-normación. Dado $y\in Y$ unitario, elija $v\in D$ con $\norm{y-v}<\eps$. Como $\norm{\varphi_v}=1$,
\[
 |\varphi_v(y)|\ge|\varphi_v(v)|-|\varphi_v(y-v)|>1-2\eps.
\]
Tomar el supremo y usar homogeneidad prueba la afirmación para todo $y$. Se precisa $\eps<1/2$ para que la constante de normación sea positiva; el valor absoluto corrige la comparación del original en el caso complejo.
\end{proof}

\begin{lema}\label{v3:normantes-finitos}
Todo espacio normado separable tiene una familia numerable de funcionales de norma uno que es $1$-normante. Si $W\subset X^{**}$ es finito dimensional y $0<\alpha<1$, existe una familia \emph{finita} $A\subset B_{X^*}$ tal que
\[
 \alpha\norm\Phi\leq\max_{\varphi\in A}|\Phi(\varphi)|\qquad(\Phi\in W).
\]
\end{lema}
\begin{proof}
En un espacio separable no nulo tomemos $(y_n)$ densa en su esfera unitaria. Por Hahn--Banach, existen $\psi_n$ de norma uno con $\psi_n(y_n)=1$. Si $y$ es unitario y $\norm{y-y_n}<\eps$, entonces
$|\psi_n(y)|\geq1-\eps$. Por tanto el supremo de esas mediciones es uno; por homogeneidad la familia es $1$-normante. Si se trabaja en un subespacio separable de un espacio mayor, Hahn--Banach extiende los funcionales sin aumentar la norma. El caso del espacio cero es vacío de contenido.

Para la segunda afirmación, si $W=\{0\}$ tomamos $A=\{0\}$, que satisface la desigualdad escrita. Supongamos $W\ne\{0\}$ y elijamos $\delta>0$ con $1-2\delta>\alpha$ y una $\delta$-red finita $w_1,\ldots,w_m$ en la esfera de $W$, compacta por dimensión finita. De la definición de norma en $X^{**}$, para cada $w_j$ existe $\varphi_j\in B_{X^*}$ con $|w_j(\varphi_j)|>1-\delta$. Para $w$ unitario elijamos $j$ con $\norm{w-w_j}<\delta$; entonces
\[
 |w(\varphi_j)|\geq|w_j(\varphi_j)|-|(w-w_j)(\varphi_j)|>1-2\delta>\alpha.
\]
Homogeneidad termina la prueba. Podemos normalizar cada $\varphi_j$ a norma uno, lo que sólo aumenta sus valores absolutos. Es imprescindible que $W$ sea finito dimensional; el corolario 5 de la lectura omite esa hipótesis. También hay que usar valor absoluto en el caso complejo, en vez de comparar un número complejo con $1-\delta$.
\end{proof}

\begin{teorema}[Eberlein--Šmulian]\label{v3:eberlein}
En un espacio de Banach $X$, un conjunto $K$ es débilmente compacto si y sólo si toda sucesión en $K$ tiene una subsucesión que converge débilmente a un elemento de $K$.
\end{teorema}
\begin{estrategia}
En una dirección, una sucesión genera un subespacio separable y la compacidad vuelve metrizable la parte relevante. En la otra, usamos Alaoglu en el bidual y una elección inductiva de un número finito de mediciones en cada etapa. Esas mediciones obligarán al límite de una subsucesión a coincidir con un punto dado de la clausura débil*.
\end{estrategia}
\begin{proof}
Supongamos primero que $K$ es débilmente compacto y sea $(x_n)\subset K$. El espacio $Y=\overline{\spanop\{x_n:n\geq1\}}$ es separable y cerrado en norma. Hahn--Banach muestra que es débilmente cerrado: si $x\notin Y$, un funcional que se anula en $Y$ y no en $x$ separa ese punto de $Y$. Así $K\cap Y$ es débilmente compacto.

El lema anterior proporciona una familia numerable de funcionales de $X^*$ que separa puntos de $Y$. Restringidos a $K\cap Y$ son continuos para la topología débil. Por el ejercicio \ref{ej:v3-l16-2}, esta topología en $K\cap Y$ es metrizable. Un compacto métrico es secuencialmente compacto, luego $(x_n)$ tiene una subsucesión convergente débilmente a un punto de $K\cap Y$.

Supongamos ahora la condición secuencial. Para cada $\varphi\in X^*$, el conjunto $\varphi(K)$ es acotado: si no lo fuera, podríamos escoger $x_n\in K$ con $|\varphi(x_n)|>n$, y ninguna subsucesión sería débilmente convergente. Aplicando acotación uniforme a los funcionales $J_Xx$ sobre el Banach $X^*$, obtenemos $\sup_{x\in K}\norm x<\infty$. Reescalamos para suponer $K\subset B_X$; el caso $K=\varnothing$ es inmediato.

Identifiquemos $X$ con $J_X(X)$ dentro de $X^{**}$. Por Alaoglu, basta probar que $K$ es débil*-cerrado en $B_{X^{**}}$. Sea $\Phi$ un punto de su clausura débil*. Fijemos $0<\alpha<1$. Construiremos $x_n\in K$ y familias finitas crecientes $A_n\subset B_{X^*}$ tales que $A_n$ sea $\alpha$-normante para
$W_n=\spanop\{\Phi,x_1,\ldots,x_n\}$ y
\[
 |x_n(\varphi)-\Phi(\varphi)|<1/n\qquad(\varphi\in A_{n-1}).
\]
El lema proporciona $A_0$ para $W_0=\spanop\{\Phi\}$. Dado $A_{n-1}$, la condición de clausura permite escoger $x_n\in K$ en la vecindad de \emph{$\Phi$} definida por esas finitas desigualdades. Añadimos a $A_{n-1}$ una familia finita normante para $W_n$ y obtenemos $A_n$. En la lectura falta trasladar una de estas vecindades desde cero hasta $\Phi$; la elección debe satisfacer exactamente las desigualdades escritas.

Sea $A=\bigcup_nA_n$ y $W=\overline{\bigcup_nW_n}$ en norma del bidual. La familia $A$ es $\alpha$-normante en cada $W_n$ y, por continuidad de las mediciones y la cota uniforme $\norm\varphi\leq1$, también en $W$. Por hipótesis, alguna subsucesión $x_{n_j}$ converge débilmente a $x\in K$. Como el subespacio $Y=\overline{\spanop\{x_n\}}\subset X$ es débilmente cerrado, $x\in Y$; luego, en el bidual, $x\in W$. También $\Phi\in W$.

Para cada $\varphi\in A$, a partir de cierto índice pertenece a todos los $A_n$. Así $x_n(\varphi)\to\Phi(\varphi)$, mientras que la convergencia débil de la subsucesión da $x_{n_j}(\varphi)\to x(\varphi)$. Por tanto $(x-\Phi)(\varphi)=0$ para todo $\varphi\in A$. La propiedad normante implica $\alpha\norm{x-\Phi}=0$, es decir, $\Phi=x\in K$. Hemos probado cerradura débil* y, con ella, compacidad débil.
\end{proof}

\section{Ejercicios sobre compacidad y sucesiones}
\begin{ejercicio}[Lectura 17, ejercicio 1]\label{ej:v3-l17-1}
En $(\ell^\infty)^*$ considere $\varphi_n(x)=x_n$. Pruebe que pertenecen a la bola unitaria y que ninguna subsucesión converge débil*. Concluya que compacidad débil* no implica compacidad secuencial débil*.
\end{ejercicio}


\begin{ejercicio}[Lectura 17, ejercicio 2: lema de Mazur]\label{ej:v3-l17-2}\label{v3:mazur}
Si $x_n\rightharpoonup x$ en un espacio de Banach, pruebe que $x$ pertenece al cierre en norma de la envolvente convexa de los $x_n$. Más precisamente, encuentre $y_n\in\co\{x_k:k\geq n\}$ con $\norm{y_n-x}\to0$.
\end{ejercicio}


\begin{ejercicio}[Lectura 17, ejercicio 3]\label{ej:v3-l17-3}
Pruebe: (i) $X$ admite una familia numerable $\alpha$-normante para algún $0<\alpha<1$ si y sólo si existe $T\in\Bcal(X,\ell^\infty)$ y $\beta>0$ con $\norm{Tx}_\infty\geq\beta\norm x$; (ii) admite una familia numerable $1$-normante si y sólo si se encaja lineal e isométricamente en $\ell^\infty$. Deduzca que todo Banach separable se encaja como subespacio cerrado de $\ell^\infty$.
\end{ejercicio}


\begin{ejercicio}[Lectura 17, ejercicio 4]\label{ej:v3-l17-4}
Para un subespacio $Y\subset X$, pruebe la equivalencia entre: (i) $Y$ separa puntos de $X^*$ mediante evaluación; (ii) $J_X(Y)\subset X^{**}$ es $\alpha$-normante para $X^*$ para algún $0<\alpha<1$; (iii) $\overline Y=X$.
\end{ejercicio}


\begin{ejercicio}[Lectura 17, ejercicio 5]\label{ej:v3-l17-5}
En $\ell^2$ sea $v_n=n^{1/4}e_n$. Pruebe: (i) cero es punto de acumulación débil del conjunto de términos; (ii) ninguna subsucesión converge débilmente. Explique por qué no hay contradicción con Eberlein--Šmulian.
\end{ejercicio}



\begin{ejemplo}[Un producto compacto que no es secuencialmente compacto]
El producto $[0,1]^{[0,1]}$ es compacto por Tychonoff. Sin embargo, las funciones $f_n(t)=nt-\lfloor nt\rfloor$ no tienen subsucesión convergente en ese producto. En efecto, si $f_{n_j}(t)$ convergiera para todo $t$, también lo haría $u_j(t)=e^{2\pi i n_jt}$. La convergencia dominada aplicada a $|u_j-u|^2\leq4$ daría convergencia en $L^2[0,1]$. Pero para índices distintos,
\[
 \norm{u_j-u_k}_2^2=2-2\operatorname{Re}\int_0^1e^{2\pi i(n_j-n_k)t}\,dt=2,
\]
pues la integral de esa exponencial es cero. La sucesión no es de Cauchy en $L^2$, contradicción. Este argumento desarrolla el ejemplo de partes fraccionarias de la lectura y muestra por qué no se puede extender Eberlein--Šmulian a cualquier topología compacta.
\end{ejemplo}

\section{Reflexividad y compacidad de la bola}\label{lec:v3-18}
\fuentelectura{18}{Caracterización de espacios reflexivos}
\index{reflexividad}\index{bidual}\index{Goldstine, teorema de}
\begin{definicion}
Un espacio de Banach $X$ es \emph{reflexivo} si la isometría canónica $J_X:X\to X^{**}$ es sobreyectiva. La palabra «canónica» es esencial: la existencia de alguna isometría sobreyectiva entre $X$ y $X^{**}$ no basta para concluir que $J_X$ lo sea.
\end{definicion}

\begin{teorema}[Goldstine]\label{v3:goldstine}
Para todo espacio normado $X$, la clausura débil* de $J_X(B_X)$ en $X^{**}$ es $B_{X^{**}}$.
\end{teorema}
\begin{proof}
La inclusión de la clausura en $B_{X^{**}}$ se sigue de que esta bola es débil*-cerrada: está definida por las desigualdades $|\Phi(\varphi)|\leq\norm\varphi$, para $\varphi\in X^*$. Para la otra inclusión, fijemos $\Phi\in B_{X^{**}}$, funcionales $\varphi_1,\ldots,\varphi_m$ y $\eps>0$. Debemos encontrar $x\in B_X$ tal que $|\varphi_j(x)-\Phi(\varphi_j)|<\eps$ para cada $j$.

Sea $T:X\to\F^m$, $Tx=(\varphi_j(x))_{j=1}^m$, y sea $y=(\Phi(\varphi_j))_{j=1}^m$. Si $y$ no perteneciera al cerrado convexo $\overline{T(B_X)}$, el teorema de separación en dimensión finita daría escalares $a_j$ tales que
\[
 \operatorname{Re}\sum_j a_j\Phi(\varphi_j)
 >\sup_{x\in B_X}\operatorname{Re}\sum_j a_j\varphi_j(x).
\]
En el caso complejo se considera $\C^m$ como espacio real; todo funcional real en él tiene esa forma. Con $\varphi=\sum_j a_j\varphi_j$, el miembro derecho es $\norm\varphi$. Para comprobarlo, multiplique un vector casi maximizante por un escalar de módulo uno que haga real no negativo el valor del funcional. El miembro izquierdo es $\operatorname{Re}\Phi(\varphi)\leq|\Phi(\varphi)|\leq\norm\varphi$, una contradicción. Así $y\in\overline{T(B_X)}$, que proporciona el vector $x$ requerido. Hemos aproximado toda vecindad básica débil*, no sólo una colección particular de funcionales.
\end{proof}

\begin{teorema}\label{v3:reflexividad-bola}
Para un espacio de Banach $X$ son equivalentes:
\begin{enumerate}
\item $X$ es reflexivo;
\item $B_X$ es débilmente compacta;
\item toda sucesión acotada en $X$ tiene una subsucesión débilmente convergente a un elemento de $X$.
\end{enumerate}
\end{teorema}
\begin{proof}
Si $J_X$ es sobreyectiva, identifica la bola $B_X$, con su topología débil, con $B_{X^{**}}$, con su topología débil*. Esta última es compacta por Alaoglu aplicado al espacio $X^*$.

Si $B_X$ es débilmente compacta, $J_X(B_X)$ es débil*-compacto y, por ser Hausdorff la topología débil*, es cerrado en $X^{**}$. Goldstine obliga entonces a que $J_X(B_X)=B_{X^{**}}$. Multiplicar por escalares muestra $J_X(X)=X^{**}$.

La equivalencia entre (ii) y (iii) resulta de Eberlein--Šmulian, con dos precisiones. Si (ii) vale, una sucesión acotada está contenida en $rB_X$ para algún $r>0$, compacto débil. Si (iii) vale, toda sucesión en $B_X$ tiene límite débil subsecuencial en esa misma bola: de $x_n\rightharpoonup x$ y $\norm{x_n}\leq1$ se deduce $|\varphi(x)|\leq\norm\varphi$ para cada $\varphi$, y Hahn--Banach da $\norm x\leq1$. Así $B_X$ es secuencialmente compacta débil y se aplica Eberlein--Šmulian.
\end{proof}

\begin{proposicion}[Subespacios cerrados]\label{v3:subespacio-reflexivo}
Todo subespacio cerrado $Y$ de un espacio de Banach reflexivo $X$ es reflexivo.
\end{proposicion}
\begin{proof}
Por Hahn--Banach, todo funcional continuo de $Y$ se extiende a $X$. Por ello, la topología débil de $Y$ coincide con la restricción de la topología débil de $X$. Además $Y$ es débilmente cerrado: si $x\notin Y$, Hahn--Banach da un funcional que se anula en $Y$ y no en $x$. Así $B_Y=Y\cap B_X$ es cerrado en el compacto débil $B_X$. La caracterización anterior, aplicada al espacio de Banach $Y$, demuestra la afirmación.
\end{proof}

\subsection{Ejemplos y límites de la reflexividad}
Todo espacio de dimensión finita es reflexivo: $J_X$ es inyectiva y dominio y codominio tienen la misma dimensión. Un espacio de Hilbert también lo es, o bien por la representación de Riesz, o bien por el siguiente argumento secuencial que evita convenciones sobre conjugación. Una sucesión acotada $(x_n)$ está contenida en el subespacio cerrado separable $Y$ que genera. Eligiendo una base ortonormal numerable de $Y$ y extrayendo sucesivamente subsucesiones, obtenemos una subsucesión cuyas coordenadas convergen. Si su cota es $M$, la suma de cuadrados de cualquier número finito de coordenadas límite es a lo sumo $M^2$; por tanto esas coordenadas representan un vector $y\in Y$. Aproximar un vector de prueba por combinaciones finitas de la base y usar la cota uniforme demuestra convergencia débil a $y$. La proyección ortogonal de un vector de $X$ sobre $Y$ muestra que la convergencia es también débil en $X$. Se aplica el teorema \ref{v3:reflexividad-bola}.

Para $1<p<\infty$, la dualidad de $L^p$ del volumen II muestra reflexividad cuando la medida es $\sigma$-finita: si $q=p/(p-1)$, las dos representaciones $(L^p)^*=L^q$ y $(L^q)^*=L^p$ identifican la evaluación canónica con la identidad. También se obtiene el resultado para una medida arbitraria. Dada una sucesión acotada $(f_n)$ de $L^p$, el conjunto
\[
 S=\bigcup_{n,k\geq1}\{x:|f_n(x)|>1/k\}
\]
es $\sigma$-finito, pues cada conjunto de la unión tiene medida a lo sumo $k^p\norm{f_n}_p^p$. Todos los $f_n$ se anulan fuera de $S$, salvo conjuntos nulos que se pueden eliminar al elegir representantes. En $L^p(S)$ existe una subsucesión débilmente convergente; la extensión por cero a $L^p(X)$ es lineal continua y preserva convergencia débil, porque todo funcional se compone con esa extensión. Otra aplicación del criterio secuencial termina la prueba.

Los casos siguientes explican por qué los extremos $p=1,\infty$ requieren cuidado.
\begin{itemize}
\item $c_0$ no es reflexivo: bajo las identificaciones $c_0^*=\ell^1$ y $(\ell^1)^*=\ell^\infty$, la imagen canónica es $c_0$, que no contiene la sucesión constante uno.
\item $\ell^1$ no es reflexivo. En cualquier subsucesión de $(e_n)$, un funcional de $\ell^\infty$ que tome alternadamente los valores $1,-1$ en sus índices impide convergencia débil.
\item $\ell^\infty$ no es reflexivo, porque contiene el subespacio cerrado $c_0$.
\item $L^1[0,1]$ no es reflexivo: $f_n=n\mathbf1_{[0,1/n]}$ tiene norma uno. Si una subsucesión convergiera débilmente a $f$, las pruebas contra indicadores medibles contenidos en $[\delta,1]$, para cada $\delta>0$, obligarían a que $f=0$ casi dondequiera en $(0,1]$. Pero probar contra la función constante uno daría $\int f=1$, contradicción.
\item $C[0,1]$ no es reflexivo: una subsucesión de $t\mapsto t^n$ sólo podría tener como límite puntual la función que vale cero en $[0,1)$ y uno en $1$, que no es continua. Las evaluaciones puntuales pertenecen al dual, de modo que tampoco hay subsucesión débilmente convergente en $C[0,1]$.
\item $C_0(\R)$ no es reflexivo: escoja una función continua de soporte compacto, norma uno, y trasládelas a intervalos disjuntos que escapen a infinito. El mapa que asigna a $(a_n)\in c_0$ la suma de esas funciones multiplicadas por $a_n$ es una isometría lineal de $c_0$ en $C_0(\R)$; la suma converge uniformemente y su imagen es cerrada. Se aplica la proposición \ref{v3:subespacio-reflexivo}.
\end{itemize}
Estas afirmaciones no dicen que \emph{todos} los espacios llamados $L^1$, $L^\infty$ o $C(K)$ sean no reflexivos: cuando tienen dimensión finita sí son reflexivos. En particular, importa el espacio de medida o el compacto de partida.

\subsection{Un espacio de sucesiones de James}
\index{James!espacio de}
Para una sucesión real $x\in c_0$ y una lista \emph{estrictamente creciente} de índices naturales $P=(n_1,\ldots,n_{2k+1})$, pongamos
\[
 p_P(x)=\left(x_{n_{2k+1}}^2+
       \sum_{j=1}^k(x_{n_{2j-1}}-x_{n_{2j}})^2\right)^{1/2},
 \qquad \norm x_J=\sup_P p_P(x).
\]
Se permite $k=0$. Definimos $J=\{x\in c_0:\norm x_J<\infty\}$. La condición de orden sobre los índices es indispensable y corrige una omisión del original: repetir arbitrariamente una pareja de coordenadas distintas haría infinito el supremo.

Cada $p_P$ es una seminorma, por la desigualdad triangular euclidiana; el supremo satisface la desigualdad triangular y $\norm x_\infty\leq\norm x_J$, tomando una lista de un solo índice. Esto prueba que $\norm\cdot_J$ es norma. Para demostrar completitud, sea $(x^r)$ de Cauchy en $J$. Es de Cauchy en $c_0$ y converge uniformemente a cierto $x\in c_0$. Si $\norm{x^r-x^s}_J\leq\eps$ para $r,s\geq r_0$, pasar al límite $s\to\infty$ en cada lista finita $P$ da $p_P(x^r-x)\leq\eps$. Tomar el supremo da $\norm{x^r-x}_J\leq\eps$, de donde $x\in J$ y convergencia en $J$.

El espacio $J$ no es reflexivo. Los vectores $s_N=e_1+\cdots+e_N$ tienen norma $J$ igual a uno: en una lista creciente, sus coordenadas son una cadena de unos seguida de ceros; hay a lo sumo una diferencia no nula entre parejas, y si la última coordenada es uno, todas las diferencias son cero. Además la lista $(1)$ da la cota inferior uno. Toda evaluación de coordenada es continua en $J$, por lo que un límite débil subsecuencial de $(s_N)$ tendría todas sus coordenadas iguales a uno. Tal sucesión no está en $c_0$, y por tanto no está en $J$. El criterio secuencial demuestra la falta de reflexividad.

El artículo de James de 1951 \cite{james1951} construye un espacio de Banach no reflexivo isométricamente isomorfo a su segundo dual. Es un resultado adicional sobre la estructura del bidual, más profundo que la completitud y la no reflexividad que acabamos de demostrar. Lo citamos como contexto histórico, sin usarlo en ninguna prueba ni confundir un isomorfismo abstracto con la aplicación canónica $J_X$.

\subsection{Existencia de puntos de mínima distancia}
\index{aproximación!mejor}\index{funcional!que alcanza su norma}
\begin{teorema}\label{v3:proximinal-reflexivo}
Sea $X$ un espacio de Banach reflexivo.
\begin{enumerate}
\item Si $A\subset X$ es no vacío, convexo y cerrado, para cada $x\in X$ existe $a\in A$ con $\norm{x-a}=d(x,A)$.
\item Todo $\varphi\in X^*$ alcanza su norma en $B_X$: existe $x\in B_X$ tal que $\varphi(x)=\norm\varphi$. Si $\varphi\ne0$, necesariamente $\norm x=1$.
\end{enumerate}
\end{teorema}
\begin{proof}
(i) Sea $d=d(x,A)$ y defina $A_n=A\cap(x+(d+1/n)B_X)$. Cada $A_n$ es no vacío por la definición de ínfimo. El conjunto $A$ es débilmente cerrado por separación de convexos; la bola trasladada es débilmente compacta por reflexividad. Por tanto $A_n$ es débilmente compacto. Los $A_n$ forman una familia decreciente de cerrados no vacíos del compacto $A_1$, y su intersección es no vacía por la propiedad de intersección finita. Todo $a$ en la intersección satisface $\norm{x-a}\leq d+1/n$ para cada $n$, luego $\norm{x-a}=d$. No se ha supuesto que $A$ sea acotado o compacto: lo que se compacta son los conjuntos truncados $A_n$.

(ii) La función débilmente continua $x\mapsto\operatorname{Re}\varphi(x)$ alcanza su máximo en el compacto $B_X$. Ese máximo es $\norm\varphi$, por cambio de fase de los vectores casi maximizantes. Si $\operatorname{Re}\varphi(x)=\norm\varphi$ y $|\varphi(x)|\leq\norm\varphi$, la parte imaginaria debe ser cero. Esto prueba $\varphi(x)=\norm\varphi$. Para $\varphi\ne0$, la desigualdad $|\varphi(x)|\leq\norm\varphi\norm x$ fuerza $\norm x=1$.
\end{proof}
La mejor aproximación puede no ser única. En $(\R^2,\norm\cdot_\infty)$, todos los puntos del segmento $A=\{(1,t):|t|\leq1\}$ están a distancia uno del origen. La unicidad del volumen I utilizaba la geometría estrictamente convexa del espacio de Hilbert, no sólo reflexividad.

\begin{proposicion}[Reducción a subespacios separables]\label{v3:reflexividad-separable}
Un espacio de Banach $X$ es reflexivo si y sólo si cada subespacio cerrado separable de $X$ es reflexivo.
\end{proposicion}
\begin{proof}
La necesidad es la proposición \ref{v3:subespacio-reflexivo}. Para la suficiencia, tome una sucesión acotada $(x_n)$ y el espacio $Y=\overline{\spanop\{x_n:n\in\N\}}$. Es separable, porque las combinaciones finitas con coeficientes racionales (racionales complejos, si procede) forman un conjunto numerable denso. Es cerrado, de modo que es Banach y por hipótesis reflexivo. Existe una subsucesión que converge débilmente en $Y$. También converge débilmente en $X$, pues todo funcional de $X$ se restringe a un funcional de $Y$. El criterio secuencial prueba reflexividad de $X$. Tomar la clausura en la definición de $Y$ evita aplicar un teorema de espacios de Banach a un subespacio que podría no ser completo.
\end{proof}

\begin{observacion}[El criterio profundo de James]\index{James!criterio de reflexividad}
El recíproco del inciso (ii) del teorema \ref{v3:proximinal-reflexivo} también es cierto: un espacio de Banach es reflexivo si y sólo si todo funcional continuo alcanza su norma \emph{en la bola unitaria cerrada}. Es el criterio de James; una referencia primaria es \cite{james1972}. Como en el curso original, se enuncia como panorama y no se usa como paso de una demostración posterior. Su prueba requiere un desarrollo adicional distinto de los argumentos de compacidad ya establecidos.

La restricción $x\in B_X$ no puede omitirse. Si $\varphi\ne0$, elija $u$ con $\varphi(u)\ne0$ y ponga $x=(\norm\varphi/\varphi(u))u$. Entonces $\varphi(x)=\norm\varphi$ en \emph{cualquier} espacio normado. Esa igualdad sin control de $\norm x$ no caracteriza reflexividad.
\end{observacion}

\subsection{Ejercicios de la lectura 18}
\begin{ejercicio}[Lectura 18, ejercicio 1]\label{ej:v3-l18-1}
Si $X$ es un espacio de Banach reflexivo y $Y$ es un subespacio cerrado, pruebe que $Y$ y $X/Y$ son reflexivos.
\end{ejercicio}


\begin{ejercicio}[Lectura 18, ejercicio 2]\label{ej:v3-l18-2}
Pruebe que un espacio de Banach $X$ es reflexivo si y sólo si $X^*$ lo es.
\end{ejercicio}


\begin{ejercicio}[Lectura 18, ejercicio 3; hipótesis precisadas]\label{ej:v3-l18-3}
Sea $X\ne\{0\}$ un espacio de Banach reflexivo y sea $T:X\to\F^n$ lineal \emph{continuo}. Pruebe que existe $x\in X$ con $\norm x=1$ y $\norm{Tx}=\norm T$.
\end{ejercicio}


\begin{ejercicio}[Lectura 18, ejercicio 4]\label{ej:v3-l18-4}
Para un subespacio $Y$ de un espacio de Banach $X$, pruebe la equivalencia entre: (i) $Y$ separa puntos de $X^*$; (ii) $J_X(Y)$ es $\alpha$-normante para $X^*$ para algún $0<\alpha<1$; (iii) $\overline Y=X$.
\end{ejercicio}


\begin{ejercicio}[Lectura 18, ejercicio 5]\label{ej:v3-l18-5}
Un operador acotado $T:X\to Y$ entre espacios de Banach se llama de Dunford--Pettis si transforma toda sucesión débilmente convergente en una sucesión convergente en norma. Pruebe que todo operador compacto es de Dunford--Pettis, y que el recíproco vale cuando $X$ es reflexivo.
\end{ejercicio}

\parte{Álgebras de Banach y funciones del espectro}{La composición de operadores es un producto. Estudiamos álgebras normadas, cocientes, derivadas, resolventes y cálculo holomorfo; los caracteres describen el caso conmutativo.}
\chapter{Álgebras normadas, ideales y adjunción de una identidad}\label{cap:v3-10}

En un espacio normado podemos sumar vectores y tomar límites. Un álgebra añade un producto entre sus elementos: componer operadores, multiplicar funciones o convolucionar señales son ejemplos de esta operación. La desigualdad que relaciona producto y norma permitirá pasar de las identidades algebraicas a las series convergentes. Estudiaremos también qué significa añadir una identidad a un álgebra que no la tiene y qué propiedades sobreviven al tomar un cociente.

Trabajaremos sobre $\F=\R$ o $\C$. Una identidad será bilateral y distinta de cero. Por tanto, cuando hablemos de un álgebra unitaria, excluiremos el álgebra nula. Un ideal maximal será siempre \emph{propio}: omitir esta condición haría maximal a toda el álgebra y destruiría los resultados sobre cocientes.

\section{El producto y su relación con la norma}\label{lec:v3-19}
\fuentelectura{19}{Álgebras normadas}

\begin{definicion}
Un álgebra asociativa $A$ sobre $\F$ es un espacio vectorial con un producto bilineal $A\times A\to A$, $(a,b)\mapsto ab$, tal que $(ab)c=a(bc)$. La bilinealidad comprende las dos distributividades y las identidades
\[
 (\alpha a)b=\alpha(ab)=a(\alpha b).
\]
Es conmutativa si $ab=ba$, y unitaria si existe $1\ne0$ con $1a=a1=a$ para todo $a$. Una \emph{norma de álgebra} satisface $\norm{ab}\leq\norm a\norm b$. Si el espacio normado es completo, hablamos de un álgebra de Banach.
\end{definicion}
\index{algebra@álgebra!normada}\index{algebra@álgebra!de Banach}

El nombre ``producto vectorial'' de la lectura significa aquí producto del álgebra; no se refiere al producto cruz de $\R^3$. No exigimos conmutatividad. Las matrices y los operadores proporcionan ejemplos fundamentales en que el orden de los factores importa.

\begin{proposicion}\label{v3:algebra-identidades}
En toda álgebra se tiene $0a=a0=0$, $(-a)b=a(-b)=-(ab)$ y $(-a)(-b)=ab$. La identidad bilateral, cuando existe, es única. Si $a,b$ conmutan y $n\geq1$, entonces
\[
 (a+b)^n=\sum_{k=0}^n\binom nk a^{n-k}b^k.
\]
En un álgebra normada, $\norm{a^n}\leq\norm a^n$ para $n\geq1$; si es unitaria, $\norm1\geq1$.
\end{proposicion}
\begin{proof}
La igualdad $0a=(0+0)a=0a+0a$ permite cancelar un sumando y obtener $0a=0$; la otra igualdad se obtiene por la segunda distributividad. Como $(a+(-a))b=0$, los elementos $ab$ y $(-a)b$ son opuestos. Se razona igual en el otro factor, y aplicar dos veces la regla da $(-a)(-b)=ab$. Si $e,e'$ son identidades bilaterales, $e=ee'=e'$.

Para el binomio, el caso $n=1$ es la distributividad. Multiplique la fórmula de $n$ por $a+b$. La conmutación permite escribir $a^{n-k}b^ka=a^{n+1-k}b^k$; después se agrupan los términos de la misma potencia y se usa $\binom nk+\binom n{k-1}=\binom{n+1}k$. En un álgebra sin identidad, los términos extremos se interpretan como $a^n$ y $b^n$, sin introducir artificialmente una potencia de orden cero. Finalmente, la submultiplicatividad da la cota de las potencias por inducción. De $\norm1=\norm{1^2}\leq\norm1^2$ y $\norm1>0$ se obtiene $\norm1\geq1$.
\end{proof}

\begin{observacion}[Dos erratas algebraicas]
La regla de signos de la lectura contiene una letra ajena al producto. Además, en su fórmula del binomio debe decir $(a+b)^n$, no $(ab)^n$. La conmutación es necesaria para la fórmula usual: sin ella, $(a+b)^2=a^2+ab+ba+b^2$.
\end{observacion}

\begin{proposicion}\label{v3:producto-continuo}
El producto de un álgebra normada es conjuntamente continuo. Si $\sum_n a_n$ converge en $A$ y $b\in A$, entonces
\[
 \sum_n a_nb=\left(\sum_n a_n\right)b,
 \qquad \sum_n ba_n=b\left(\sum_n a_n\right).
\]
\end{proposicion}
\begin{proof}
Si $a_n\to a$ y $b_n\to b$, escribimos
\[
 a_nb_n-ab=a_n(b_n-b)+(a_n-a)b.
\]
Una sucesión convergente es acotada, de modo que $\norm{a_n}\leq M$ para algún $M$. Por tanto
\[
 \norm{a_nb_n-ab}\leq M\norm{b_n-b}+\norm b\norm{a_n-a}\longrightarrow0.
\]
Este criterio secuencial caracteriza continuidad en los espacios métricos considerados. Para las series, las sumas parciales de $\sum a_nb$ son $(\sum_{n\leq N}a_n)b$ por distributividad. Su límite es el producto del límite por $b$; el argumento del otro lado es idéntico, conservando el orden de los factores.
\end{proof}

\section{Ejemplos e ideales: hipótesis que importan}

\begin{definicion}
Una subálgebra es un subespacio estable bajo producto. Un subespacio $I\subset A$ es ideal izquierdo si $AI\subset I$, derecho si $IA\subset I$, y bilateral si cumple ambas condiciones. Escribimos $I\mathrel{\triangleleft}A$ para un ideal bilateral. No exigimos que una subálgebra contenga la identidad del álgebra ambiente.
\end{definicion}
\index{ideal!izquierdo}\index{ideal!bilateral}

Todo ideal es una subálgebra: en la condición de ideal puede elegirse también el factor exterior dentro del propio ideal. En un álgebra conmutativa, izquierda y derecha coinciden. El recíproco falla: las constantes forman una subálgebra de $C(X)$, pero, si existe una función no constante $f$, el producto $1f=f$ impide que esa subálgebra sea ideal.

Los espacios $C(X)$, para $X$ compacto Hausdorff no vacío, y $C_b(X)$, para $X$ arbitrario, son álgebras conmutativas de Banach con norma uniforme. El producto preserva continuidad y acotación, y
\[
 \norm{fg}_\infty\leq\norm f_\infty\norm g_\infty.
\]
La completitud es la del espacio de funciones uniformemente acotadas y la continuidad de los límites uniformes. En un espacio no compacto, $C(X)$ puede contener funciones no acotadas; la norma uniforme no es entonces una norma finita sobre todo $C(X)$.

Para un espacio localmente compacto Hausdorff, $C_0(X)$ consiste en las funciones continuas cuyos conjuntos $\{|f|\geq\eps\}$ son compactos para todo $\eps>0$. Es un subespacio cerrado de $C_b(X)$: si $f_n\to f$ uniformemente, el conjunto cerrado $\{|f|\geq\eps\}$ está contenido en el compacto $\{|f_n|\geq\eps/2\}$ para un $n$ suficientemente grande. Además, es un ideal de $C_b(X)$, pues $|fg|\geq\eps$ implica $|g|\geq\eps/\norm f_\infty$ cuando $f\ne0$.

Si $X$ no es compacto, $C_0(X)$ no tiene identidad. Para verlo, usamos la propiedad de funciones de corte de los espacios localmente compactos Hausdorff: para cada $x$ existe $u\in C_c(X)$ con $u(x)=1$. Una identidad $e$ cumpliría $eu=u$, luego $e(x)=1$ para todo $x$. Pero la constante $1$ no se anula en el infinito cuando $X$ no es compacto. La propiedad de corte se obtiene escogiendo una vecindad de clausura compacta de $x$, separando $x$ de su frontera mediante el lema de Urysohn en ese compacto Hausdorff, y prolongando por cero; la función se anula en una vecindad de la frontera, lo que justifica la continuidad de la prolongación.

\begin{observacion}[El álgebra ambiente no puede omitirse]
La afirmación de que $C_0(X)$ es ideal de todo $C(X)$ es falsa. En $X=\R$, tome $g(x)=1/(1+x^2)\in C_0(\R)$ y $f(x)=1+x^2\in C(\R)$. Su producto es $1\notin C_0(\R)$. La prueba de la lectura utiliza $\norm f_\infty$ y, por ello, demuestra la afirmación correcta para $C_b(X)$.
\end{observacion}

En un espacio de medida, $L^\infty$ es un álgebra de Banach con producto de clases de funciones. El producto está bien definido porque cambiar representantes en conjuntos nulos no cambia su clase; la desigualdad esencial $|fg|\leq\norm f_\infty\norm g_\infty$ prueba submultiplicatividad. Si $\mu(X)>0$, la clase de $1$ es una identidad de norma uno; si $\mu(X)=0$, el álgebra cociente es nula y constituye la excepción a nuestra convención sobre identidades.

Para un espacio de Banach $V$, $\Bcal(V,V)$ es un álgebra de Banach bajo composición, y $\Kcal(V,V)$ es un ideal bilateral cerrado. Si $T$ es compacto y $S$ acotado, una sucesión acotada admite una subsucesión en la que $Tx_n$ converge, y entonces $STx_n$ converge; para $TS$, se aplica compacidad a la sucesión acotada $Sx_n$. La cerradura en norma se prueba aproximando la imagen de la bola por una imagen totalmente acotada y usando completitud del codominio. En particular, las álgebras de operadores y de matrices no son generalmente conmutativas; la excepción no se limita a los cuaternios.

Finalmente, $C^n([a,b])$ es subálgebra de $C([a,b])$ por Leibniz. Si $n\geq1$ y $a<b$, no es ideal: contiene a $1$, pero no a la función continua $|t-(a+b)/2|$. Para $n=0$ coincide con toda el álgebra.

\section{Ejercicios sobre productos, inversos y cocientes}

\begin{ejercicio}[Lectura 19, ejercicio 1]\label{ej:v3-l19-1}
Estudie las siguientes álgebras normadas y justifique todos los productos:
\begin{enumerate}
\item $C^n([a,b])$ con producto puntual y la norma de la lectura 0, $\norm f_n=\sum_{k=0}^n\norm{f^{(k)}}_\infty$. Determine si esa norma es submultiplicativa y, cuando haga falta, corríjala mediante una norma equivalente.
\item $\ell^1(\Z)$ con $(x*y)_k=\sum_{j\in\Z}x_jy_{k-j}$.
\item $L^1(\R)$ con $(f*g)(x)=\int_\R f(y)g(x-y)\,dy$, definida casi en todas partes.
\end{enumerate}
\end{ejercicio}
\begin{solucion}
(i) La norma dada no es submultiplicativa para $n\geq2$. Ponga $c=(a+b)/2$, $r=(b-a)/2$ y $f(t)=g(t)=t-c$. Entonces $\norm f_n=\norm g_n=r+1$, mientras que
\[
 \norm{fg}_n=r^2+2r+2>(r+1)^2.
\]
La versión válida utiliza
\[
 \norm f_{n,!}=\sum_{k=0}^n\frac{\norm{f^{(k)}}_\infty}{k!}.
\]
Leibniz da
\[
 \frac{\norm{(fg)^{(k)}}_\infty}{k!}
 \leq\sum_{j=0}^k\frac{\norm{f^{(j)}}_\infty}{j!}
                       \frac{\norm{g^{(k-j)}}_\infty}{(k-j)!}.
\]
Al sumar en $k\leq n$, aparecen algunos de los términos no negativos del producto $\norm f_{n,!}\norm g_{n,!}$, y se obtiene la desigualdad requerida. Las normas son equivalentes porque $\norm f_n/n!\leq\norm f_{n,!}\leq\norm f_n$; para $n=0,1$ coinciden. Su completitud puede verse directamente: una sucesión de Cauchy tiene límites uniformes $g_k$ de sus derivadas de orden $k$, y pasar al límite en $f_j^{(k)}(t)-f_j^{(k)}(a)=\int_a^t f_j^{(k+1)}$ da $g_k'=g_{k+1}$. Así $g_0\in C^n$ y la convergencia ocurre en la norma indicada. El producto es asociativo y conmutativo punto por punto, y la constante $1$ es identidad.

(ii) Tonelli para series no negativas y el cambio de índice $l=k-j$ dan
\[
 \sum_{k\in\Z}\sum_{j\in\Z}|x_jy_{k-j}|
 =\sum_j|x_j|\sum_l|y_l|=\norm x_1\norm y_1.
\]
Cada coeficiente está definido por una serie absolutamente convergente, la sucesión resultante pertenece a $\ell^1$, y $\norm{x*y}_1\leq\norm x_1\norm y_1$. Bilinealidad y conmutatividad siguen de las series y del cambio $j\mapsto k-j$. Para asociatividad, ambas expresiones tienen coeficiente $\sum_{j,l}x_jy_lz_{k-j-l}$. La suma de los valores absolutos, también sobre $k$, es $\norm x_1\norm y_1\norm z_1$, por lo que Fubini justifica reagrupar. La completitud es la de $\ell^1$.

(iii) Elija representantes medibles. Tonelli e invariancia por traslación de Lebesgue implican
\[
 \int_\R\int_\R |f(y)g(x-y)|\,dy\,dx
 =\int_\R |f(y)|\left(\int_\R|g(x-y)|\,dx\right)dy
 =\norm f_1\norm g_1.
\]
La integral interior existe para casi todo $x$, define una función medible integrable y satisface $\norm{f*g}_1\leq\norm f_1\norm g_1$. El mismo cálculo aplicado a diferencias nulas casi en todas partes prueba independencia de representantes. Bilinealidad se hereda de la integral. La sustitución $u=x-y$ demuestra conmutatividad. Para tres funciones, la integral absoluta total de $f(s)g(t)h(x-s-t)$ sobre $\R^3$ vale $\norm f_1\norm g_1\norm h_1$; Fubini permite calcular las integrales en ambos órdenes y demuestra $(f*g)*h=f*(g*h)$ casi en todas partes. El resultado es un álgebra de Banach porque $L^1$ es completo. La ausencia de identidad se demostrará en el ejercicio~\ref{ej:v3-l20-5}.
\end{solucion}

\begin{ejercicio}[Lectura 19, ejercicio 2]\label{ej:v3-l19-2}
Sea $A$ un álgebra unitaria. Un inverso izquierdo de $a$ satisface $ba=1$; un inverso derecho satisface $ac=1$. Pruebe:
\begin{enumerate}
\item Si $a$ tiene inverso izquierdo pero no derecho, posee infinitos inversos izquierdos.
\item Si $b$ es inverso izquierdo y $c$ derecho, entonces $b=c$.
\item El inverso bilateral, cuando existe, es único y se denota $a^{-1}$.
\item Si $a,b$ son invertibles, $\alpha\ne0$ y $n\geq1$, también lo son $ab,\alpha a,a^n$, y sus inversos son $b^{-1}a^{-1},\alpha^{-1}a^{-1},(a^{-1})^n$.
\item Si $A$ es normada y $b$ es inverso por alguno de los lados, determine la cota inferior correcta para $\norm b$.
\end{enumerate}
\end{ejercicio}
\begin{solucion}
(i) Sea $ba=1$. Como $a$ no tiene inverso derecho, $q=1-ab\ne0$. Se tiene $qa=a-aba=0$. Por tanto $(b+tq)a=1$ para todo $t\in\F$. Si $s\ne t$, $(s-t)q\ne0$; los inversos son distintos y $\R,\C$ tienen infinitos escalares.

(ii) La asociatividad da $b=b(ac)=(ba)c=c$. (iii) Dos inversos bilaterales son, en particular, un inverso izquierdo y uno derecho, y el inciso anterior prueba unicidad. (iv) Multiplicar en ambos órdenes verifica, por ejemplo, $(ab)(b^{-1}a^{-1})=a1a^{-1}=1$ y $(b^{-1}a^{-1})(ab)=b^{-1}1b=1$. La compatibilidad con escalares verifica el segundo caso; el tercero se obtiene iterando el primero. Es indispensable $\alpha\ne0$, pues $0a=0$ no es invertible.

(v) En cualquiera de los dos órdenes, $\norm1\leq\norm a\norm b$, de modo que
\[
 \norm b\geq\frac{\norm1}{\norm a}\geq\frac1{\norm a}.
\]
La cota $\norm b\geq\norm a$ impresa en la lectura es falsa incluso en $\R$: tome $a=2$, $b=1/2$. En particular, $\norm{a^{-1}}\geq1/\norm a$ si $a$ es invertible.
\end{solucion}

\begin{ejercicio}[Lectura 19, ejercicio 3]\label{ej:v3-l19-3}
Para $a$ en un álgebra normada unitaria, defina $L_a(x)=ax$. Pruebe que $L_a$ es lineal continuo y que $a\mapsto L_a$ es lineal, multiplicativo e inyectivo. Deduzca una norma equivalente y submultiplicativa para la que la identidad tiene norma uno.
\end{ejercicio}
\begin{solucion}
Distributividad y compatibilidad escalar dan linealidad, y $\norm{L_ax}\leq\norm a\norm x$ da continuidad. Evaluando en cada $x$ se obtiene $L_{a+b}=L_a+L_b$, $L_{\alpha a}=\alpha L_a$ y $L_{ab}=L_aL_b$. Si $L_a=0$, evaluar en $1$ da $a=0$.

Defina $\norm a'=\norm{L_a}$. La inyectividad garantiza que se trata de una norma; la desigualdad de la composición la hace submultiplicativa. Además,
\[
 \norm a'\leq\norm a,
 \qquad \norm a=\norm{L_a1}\leq\norm a'\norm1.
\]
Las normas son equivalentes. Finalmente $L_1$ es la identidad del espacio normado no nulo $A$, cuya norma de operador es uno. Así $\norm1'=1$. La equivalencia preserva también completitud.
\end{solucion}

\begin{ejercicio}[Lectura 19, ejercicio 4]\label{ej:v3-l19-4}
Sea $I$ un ideal \emph{bilateral} cerrado de un álgebra normada $A$. Pruebe que $(a+I)(b+I)=ab+I$ define un producto de álgebra normada sobre $A/I$, con su norma cociente.
\end{ejercicio}
\begin{solucion}
Si se sustituyen $a,b$ por $a+i,b+j$, la diferencia de productos es $aj+ib+ij\in I$. Este cálculo utiliza ambos lados del ideal y demuestra que el producto está bien definido. Todas las identidades algebraicas descienden de las de $A$. Para $\eps>0$, elija representantes $a',b'$ de las clases con normas menores que $\norm{a+I}+\eps$ y $\norm{b+I}+\eps$. Entonces
\[
 \norm{ab+I}\leq\norm{a'b'}
 \leq(\norm{a+I}+\eps)(\norm{b+I}+\eps).
\]
Dejar $\eps\downarrow0$ da submultiplicatividad. La cerradura de $I$ garantiza que la norma cociente separa clases: norma cero significa pertenecer a $\cl I=I$. Si $A$ es Banach, $A/I$ es Banach por el teorema de completitud de cocientes. Si $I=A$, se obtiene el álgebra nula; si $A$ es unitaria e $I$ propio, $1+I$ es una identidad no nula.
\end{solucion}

\begin{ejercicio}[Lectura 19, ejercicio 5]\label{ej:v3-l19-5}
Sea $T:A\to B$ lineal, continuo y multiplicativo; cuando se exige homomorfismo unital, suponga además $T(1_A)=1_B$. Pruebe que su núcleo es ideal bilateral y su imagen es subálgebra. Precise el significado del primer teorema de isomorfismo $A/\ker T\simeq T(A)$ y las hipótesis para que sea un isomorfismo topológico.
\end{ejercicio}
\begin{solucion}
Si $i\in\ker T$, entonces $T(ai)=T(a)T(i)=0$ y $T(ia)=0$. El núcleo es subespacio y, por continuidad, cerrado. Las imágenes de sumas, escalares y productos muestran que $T(A)$ es subálgebra. La aplicación
\[
 \widetilde T(a+\ker T)=T(a)
\]
está bien definida, es sobreyectiva sobre $T(A)$ y es inyectiva porque $T(a)=0$ equivale a $a\in\ker T$. Es lineal y multiplicativa, y
$\norm{T(a)}\leq\norm T\inf_{i\in\ker T}\norm{a+i}=\norm T\norm{a+\ker T}$.

Esto prueba un isomorfismo \emph{algebraico} y continuidad hacia la imagen, pero no continuidad automática de la inversa. Como contraejemplo, incluya $\ell^1$ en $c_0$, ambos con producto puntual, y use sus normas usuales. Es una inclusión multiplicativa continua; las sucesiones con sus primeros $n$ términos iguales a $1/n$ tienen norma $\ell^1$ igual a $1$ y norma uniforme $1/n$. La inversa sobre la imagen no es continua. Si $A$ y $B$ son Banach y $T(A)$ es cerrado, tanto el cociente como la imagen son Banach y el teorema de la aplicación abierta aplicado a $\widetilde T$ hace continua su inversa. Estas hipótesis justifican la versión topológica, que no se deduce solamente de la identidad algebraica.
\end{solucion}

\begin{ejercicio}[Lectura 19, ejercicio 6: variables y cerradura]\label{ej:v3-l19-6}
Sea $X$ localmente compacto Hausdorff y $E\subset X$. Defina
\[
 I(E)=\{f:f|_E=0\}.
\]
Pruebe que es un ideal en $C_b(X)$, con norma uniforme, y en $C(X)$, con la topología de convergencia uniforme en compactos. Examine la afirmación de que ese ideal es cerrado si y solo si $E$ es cerrado.
\end{ejercicio}
\begin{solucion}
La notación de la lectura usa la misma letra para el conjunto de puntos y el espacio de funciones, por lo que hemos distinguido $E$ e $I(E)$. Las funciones que se anulan en $E$ forman un subespacio, y $(fg)(x)=f(x)g(x)=0$ en $E$ si $g$ se anula allí. Esto prueba la propiedad de ideal.

Para cada $x$, la evaluación $f\mapsto f(x)$ es continua en ambas topologías: en norma uniforme, $|f(x)|\leq\norm f_\infty$; en la otra, se utiliza el compacto $\{x\}$. Así
$I(E)=\bigcap_{x\in E}\ker(f\mapsto f(x))$ es siempre cerrado. Además, una función continua que se anula en $E$ se anula en $\cl E$, y por ello $I(E)=I(\cl E)$. No puede recuperarse la condición de que $E$ sea cerrado mirando solamente el ideal. Por ejemplo, $E=(0,1)\subset\R$ no es cerrado y su ideal es cerrado. Para los conjuntos cerrados sí puede recuperarse el conjunto como los ceros comunes de su ideal: si $x\notin E$, una función de corte soportada en $X\setminus E$ separa $x$ de $E$. Esto explica la correspondencia geométrica correcta sin afirmar el falso ``si y solo si''.
\end{solucion}

\section{Adjuntar una identidad sin alterar el álgebra original}\label{lec:v3-20}
\fuentelectura{20}{Extensión de álgebras de Banach}

Sea $A$ un álgebra de Banach, inicialmente sin identidad. Construimos
\[
 A^{\sim}=A\oplus\F,\qquad
 (a,\alpha)(b,\beta)=(ab+\alpha b+\beta a,\alpha\beta),
 \qquad \norm{(a,\alpha)}_1=\norm a+|\alpha|.
\]
La expresión imita el producto formal $(a+\alpha1)(b+\beta1)$. Esa motivación sugiere la fórmula, pero sus propiedades deben verificarse dentro del espacio construido.

\begin{teorema}\label{v3:unitizacion}
Con estas operaciones, $A^{\sim}$ es un álgebra de Banach unitaria, con identidad $(0,1)$ de norma uno. La aplicación $J(a)=(a,0)$ es un homomorfismo isométrico y su imagen es un ideal bilateral cerrado maximal.
\end{teorema}
\begin{proof}
La suma de las dos normas es positiva y se anula únicamente en $(0,0)$; homogeneidad y desigualdad triangular se obtienen sumando las correspondientes desigualdades de cada coordenada. Si $(a_n,\alpha_n)$ es de Cauchy, ambas coordenadas lo son, luego convergen a $a\in A$ y $\alpha\in\F$. La suma de sus errores demuestra convergencia en la norma de $A^{\sim}$. Equivalentemente, una serie absolutamente convergente converge coordenada a coordenada; sus sumas parciales son $\sum_{k=1}^n a_k$, no las colas $\sum_{k=n}^\infty a_k$ que aparecen por errata en la lectura.

La primera coordenada de cualquiera de las dos asociaciones de un producto de tres factores es
\[
 abc+\alpha bc+\beta ac+\gamma ab
       +\alpha\beta c+\alpha\gamma b+\beta\gamma a;
\]
la segunda es $\alpha\beta\gamma$. La asociatividad de $A$ y de los escalares prueba la de $A^{\sim}$. Cada término de la fórmula del producto es lineal respecto de cada par de coordenadas, por lo que se cumplen las dos distributividades y la compatibilidad escalar. Multiplicar por $(0,1)$ en ambos órdenes devuelve $(a,\alpha)$. Además,
\[
 \norm{(a,\alpha)(b,\beta)}_1
 \leq\norm a\norm b+|\alpha|\norm b+|\beta|\norm a+|\alpha\beta|
 =\norm{(a,\alpha)}_1\norm{(b,\beta)}_1.
\]
La norma de la identidad es $1$.

Las identidades $J(a+b)=J(a)+J(b)$, $J(ab)=J(a)J(b)$ y $\norm{J(a)}_1=\norm a$ verifican todas las propiedades de $J$. La proyección $q(a,\alpha)=\alpha$ es un homomorfismo continuo sobre $\F$, cuyo núcleo es $J(A)$; por tanto ese núcleo es cerrado y bilateral. Para comprobar maximalidad sin invertir flechas de cocientes, sea $I$ un ideal que contiene a $J(A)$ y contiene también algún $(a,\alpha)$ con $\alpha\ne0$. Restar $(a,0)$ y dividir por $\alpha$ muestra que $(0,1)\in I$, de modo que $I=A^{\sim}$. No hay ideal intermedio. Equivalentemente, $A^{\sim}/J(A)\simeq\F$ tiene dimensión uno.
\end{proof}

\begin{ejemplo}[Sucesiones convergentes]
Sea $c$ el álgebra de sucesiones convergentes, con norma uniforme y producto puntual. Si $\ell=\lim x_n$, defina $T(x)=(x-\ell1,\ell)\in c_0^{\sim}$. El límite es lineal, así que $T$ es lineal; su inversa es $(a,\alpha)\mapsto a+\alpha1$. Las cotas
\[
 \norm{T(x)}_1\leq\norm{x-\ell1}_\infty+|\ell|
 \leq3\norm x_\infty,
 \qquad \norm{a+\alpha1}_\infty\leq\norm a_\infty+|\alpha|
\]
hacen continuas ambas aplicaciones. Si $\ell=\lim x_n$ y $m=\lim y_n$, la primera coordenada del producto es
\[
 (x-\ell1)(y-m1)+\ell(y-m1)+m(x-\ell1)=xy-\ell m1,
\]
y la segunda es $\ell m=\lim x_ny_n$. Por ello $T(xy)=T(x)T(y)$. Es un isomorfismo de Banach, generalmente no isométrico para la norma suma: $x=(-1,1,1,\ldots)$ tiene norma uniforme uno y $\norm{T(x)}_1=3$. La inclusión de $c_0$ en $c$, en cambio, sí es isométrica.
\end{ejemplo}

\section{Una herramienta para ideales cerrados}

\begin{lema}[Serie de Neumann]\label{v3:neumann-basico}
En un álgebra de Banach unitaria, si $\norm a<1$, entonces $1-a$ es invertible y
\[
 (1-a)^{-1}=\sum_{n=0}^\infty a^n.
\]
\end{lema}
\begin{proof}
El término $n=0$ es la identidad, de norma finita; la cola converge absolutamente porque $\norm{a^n}\leq\norm a^n$ para $n\geq1$. Si $s_N=1+a+\cdots+a^N$, la distributividad produce $(1-a)s_N=s_N(1-a)=1-a^{N+1}$. Como $\norm{a^{N+1}}\to0$, la continuidad del producto permite pasar al límite en ambas igualdades. La suma de la serie es, pues, un inverso por ambos lados.
\end{proof}

\section{Ejercicios sobre ideales maximales y existencia de identidad}

\begin{ejercicio}[Lectura 20, ejercicio 1]\label{ej:v3-l20-1}
Pruebe que la clausura de un ideal izquierdo, derecho o bilateral de un álgebra normada es un ideal del mismo tipo. Examine la conclusión de que todo ideal maximal es cerrado y demuestre la versión válida para álgebras de Banach unitarias.
\end{ejercicio}
\begin{solucion}
El cierre de un subespacio es subespacio por continuidad de las operaciones vectoriales. Si $i_n\in I$ converge a $i$ y $a\in A$, entonces $ai_n\to ai$ y $i_na\to ia$. Cada condición de ideal que valía para $I$ pasa así a $\cl I$.

Supongamos ahora $A$ Banach unitaria e $I$ maximal propio. Si $\cl I=A$, existe $i\in I$ con $\norm{1-i}<1$. El lema~\ref{v3:neumann-basico} hace invertible a $i$; un ideal que contiene un invertible contiene a $1$ y es todo $A$, contradicción. Así $\cl I$ es propio. La maximalidad y $I\subset\cl I$ implican $I=\cl I$. El argumento funciona también para ideales maximales izquierdos o derechos, multiplicando el inverso por el lado correspondiente.

La completitud no puede suprimirse. En $\C[z]$ con la norma uniforme del disco unidad, $I=(z-2)$ es maximal porque es el núcleo de la evaluación sobreyectiva en $2$. Sin embargo,
\[
 p_N(z)=-\frac12\sum_{k=0}^N(z/2)^k,
 \qquad (z-2)p_N(z)=1-(z/2)^{N+1}\longrightarrow1
\]
uniformemente en el disco. La clausura del ideal contiene a $1$, luego a toda el álgebra, aunque $I$ es propio. Por tanto el enunciado original para cualquier álgebra normada requiere la corrección indicada.
\end{solucion}

\begin{ejercicio}[Lectura 20, ejercicio 2]\label{ej:v3-l20-2}
Pruebe que un álgebra conmutativa unitaria no nula tiene ideales maximales.
\end{ejercicio}
\begin{solucion}
Ordene por inclusión los ideales propios. La familia no es vacía, pues contiene a $\{0\}$. La unión de una cadena es subespacio e ideal: dados finitos elementos de la unión, alguno de los miembros de la cadena los contiene simultáneamente y permite realizar las operaciones. Esa unión sigue siendo propia, pues si contuviera a $1$, alguno de los ideales de la cadena ya contendría a $1$. Por Zorn existe un elemento maximal. Para la cadena vacía sirve $\{0\}$ como cota superior. Este argumento no presupone norma ni completitud.
\end{solucion}

\begin{ejercicio}[Lectura 20, ejercicio 3: el campo cociente]\label{ej:v3-l20-3}
Sea $A$ conmutativa unitaria sobre $\F$ e $I$ un ideal propio. Determine la relación entre maximalidad de $I$ y el cociente $A/I$. Examine la afirmación de la lectura $I$ maximal si y solo si $A/I\simeq\F$.
\end{ejercicio}
\begin{solucion}
La afirmación algebraica general es: $I$ es maximal si y solo si $A/I$ es un \emph{campo}, que puede ser una extensión propia de $\F$. Si $I$ es maximal y $a\notin I$, el ideal $I+Aa$ contiene estrictamente a $I$, luego es todo $A$. Así $1=i+ba$ para ciertos $i\in I$, $b\in A$, y $(b+I)(a+I)=1+I$. Toda clase no nula tiene inversa. Recíprocamente, si el cociente es campo y un ideal $J$ contiene estrictamente a $I$, elija $a\in J\setminus I$. La inversa de su clase proporciona $b$ con $ba-1\in I\subset J$, de donde $1\in J$ y $J=A$.

Si $A/I\simeq\F$, ciertamente es campo y $I$ es maximal. El recíproco, tal como está impreso, es falso: considere $A=\C$ como álgebra real e $I=\{0\}$. Es maximal, pero $\C$ no es isomorfo a $\R$ como álgebra real. Incluso sobre $\C$, un álgebra puramente algebraica puede ser $\C(t)$, cuyo ideal cero es maximal y cuyo campo no es $\C$. Para un álgebra de Banach \emph{compleja} unitaria, el cociente por un maximal es Banach por el ejercicio anterior y el teorema de Gelfand--Mazur, demostrado en el capítulo~\ref{cap:v3-11}, identifica ese campo con $\C$. Esas son las hipótesis analíticas adicionales que hacen válida la conclusión escalar.
\end{solucion}

\begin{ejercicio}[Lectura 20, ejercicio 4]\label{ej:v3-l20-4}
Demuestre que $e_0=(\delta_{0,n})_{n\in\Z}$ es la identidad de $\ell^1(\Z)$ con el producto de Cauchy.
\end{ejercicio}
\begin{solucion}
La sucesión pertenece a $\ell^1$ y tiene norma uno. Para cada $k$,
\[
 (e_0*x)_k=\sum_j\delta_{0,j}x_{k-j}=x_k,
 \qquad (x*e_0)_k=\sum_jx_j\delta_{0,k-j}=x_k.
\]
Cada suma tiene un único término no nulo. Las dos igualdades verifican una identidad bilateral.
\end{solucion}

\begin{ejercicio}[Lectura 20, ejercicio 5]\label{ej:v3-l20-5}
Pruebe que $L^1(\R)$ con convolución no tiene identidad. Si existiera $e\in L^1$ con $e*f=f$ casi en todas partes para toda $f\in L^1$, pruebe que la medida $\mu_e(E)=\int_E e\,d\lambda$ asignaría masa uno a todo intervalo abierto acotado que contiene al origen, y obtenga una contradicción con su continuidad absoluta.
\end{ejercicio}
\begin{solucion}
Sea $f_n=(n/2)\mathbf1_{(-1/n,1/n)}$, cuya integral es uno. Fije $a<0<b$. Para $n$ suficientemente grande, $\int_a^b f_n=1$, y la identidad supuesta da $\int_a^b(e*f_n)=1$. Fubini, justificado por $\norm e_1\norm{f_n}_1<\infty$, permite escribir
\[
 1=\int_\R e(y)h_n(y)\,dy,
 \qquad h_n(y)=\frac n2\lambda\big((a,b)\cap(y-1/n,y+1/n)\big).
\]
Se tiene $0\leq h_n\leq1$. Para $y\in(a,b)$, la pequeña ventana termina contenida en el intervalo y $h_n(y)=1$; para $y\notin[a,b]$, termina disjunta y $h_n(y)=0$. Solo los dos extremos pueden tener otro límite. Convergencia dominada con dominante $|e|$ demuestra
$1=\int_a^b e=\mu_e((a,b))$.

En particular, $\mu_e((-1/n,1/n))=1$. Como esos intervalos decrecen a $\{0\}$ y $e$ es integrable, convergencia dominada vuelve a dar $\mu_e(\{0\})=1$. Pero $\lambda(\{0\})=0$ y una medida de densidad integrable satisface $\mu_e(\{0\})=\int_{\{0\}}e\,d\lambda=0$. Esta contradicción prueba la ausencia de identidad. El obstáculo es concreto: la identidad formal de la convolución es una masa puntual, que no tiene densidad en $L^1(\R)$.
\end{solucion}

La norma controla productos y series, mientras que los ideales permiten conservar parte del álgebra al pasar a un cociente. La unitización añade exactamente una dirección escalar. Estas construcciones preparan la pregunta espectral: para qué valores de un escalar deja de ser invertible un elemento menos ese escalar por la identidad.

\chapter{Diferenciación, espectro y cálculo holomorfo}\label{cap:v3-11}

La inversión de matrices es una operación no lineal. Sin embargo, lejos de las matrices singulares, una perturbación pequeña admite una expansión controlada. El mismo mecanismo funciona en cualquier álgebra de Banach: la serie geométrica proporciona primero una derivada, después una función resolvente y finalmente una manera de evaluar funciones holomorfas en elementos del álgebra. Desarrollaremos esta cadena sin confundir el producto no conmutativo con el producto de números.

\section{La derivada como aproximación lineal}\label{lec:v3-21}
\fuentelectura{21}{La derivada de Fréchet}

\begin{definicion}
Sean $X,Y$ espacios normados sobre $\F$, $U\subset X$ abierto y $a\in U$. Una aplicación $f:U\to Y$ es \emph{diferenciable en el sentido de Fréchet} en $a$ si existe $T\in\Lcal(X,Y)$ tal que
\[
 f(a+h)=f(a)+Th+r(h),\qquad
 \frac{\norm{r(h)}}{\norm h}\longrightarrow0\quad(h\to0,\ h\ne0).
\]
Escribimos $Df(a)=T$. Cuando $\F=\C$, la aplicación $T$ debe ser compleja lineal; la diferenciabilidad real es una condición diferente.
\end{definicion}
\index{derivada de Fréchet}

El cociente controla el error simultáneamente en todas las direcciones. Las derivadas direccionales aisladas no proporcionan este control uniforme. La apertura de $U$ garantiza que todos los incrementos suficientemente pequeños están permitidos.

\begin{proposicion}[Reglas elementales]\label{v3:reglas-frechet}
La derivada, cuando existe, es única. La diferenciabilidad implica continuidad; la suma y la multiplicación por escalares se diferencian término a término. Si $f:U\to Y$ es diferenciable en $a$ y $g:V\to Z$ lo es en $f(a)$, con $f(U)\subset V$, entonces
\[
 D(g\circ f)(a)=Dg(f(a))\,Df(a).
\]
\end{proposicion}
\begin{proof}
Si $T,S$ son derivadas, sustituyendo $h=tv$ y dividiendo entre $|t|$ obtenemos $\norm{(T-S)v}=0$ para cada $v$. Si $r(h)=o(\norm h)$, para $h$ pequeño se tiene $\norm{f(a+h)-f(a)}\le(\norm T+1)\norm h$, lo cual prueba continuidad incluso cuando $T=0$. Las reglas de linealidad se siguen de que una combinación lineal de errores $o(\norm h)$ vuelve a ser de ese orden.

Para la composición pongamos $k=f(a+h)-f(a)=Th+r(h)$. Entonces $\norm k\le C\norm h$ y
\[
 g(f(a)+k)-g(f(a))=Sk+s(k)=STh+Sr(h)+s(k).
\]
El primer resto es $o(\norm h)$. Para el segundo, si $k\ne0$,
$\norm{s(k)}/\norm h=(\norm{s(k)}/\norm k)(\norm k/\norm h)\to0$; si $k=0$, el resto es cero. Esta observación evita dividir entre un incremento que puede anularse.
\end{proof}

\begin{ejemplo}
Una aplicación lineal acotada $T$ tiene derivada $T$ en todos los puntos. En un espacio de Hilbert real, para $b(x,y)=\ip{x}{y}$,
\[
 Db(x,y)[h,k]=\ip h y+\ip x k,
\]
pues el resto $\ip h k$ tiene módulo a lo sumo $\norm h\norm k$. Usamos en el producto la norma $\norm{(h,k)}=\norm h+\norm k$. El caso complejo exige distinguir las dos nociones de derivada; véase el ejercicio~\ref{ej:v3-l21-2}.
\end{ejemplo}

\begin{proposicion}
Si $f:X\to Y$ es diferenciable y $Df(x)=0$ para todo $x$, entonces $f$ es constante. Más generalmente, basta que el dominio sea un abierto conexo.
\end{proposicion}
\begin{proof}
Fijemos un segmento $[a,b]$ contenido en el dominio y $\varphi\in Y^*$. La función $t\mapsto\varphi(f(a+t(b-a)))$ tiene derivada cero en $0<t<1$. El teorema del valor medio, aplicado a sus partes real e imaginaria si es necesario, da $\varphi(f(b)-f(a))=0$. Hahn--Banach separa puntos y por ello $f(a)=f(b)$. En un abierto conexo, el conjunto de puntos que pueden unirse a $a$ por una poligonal contenida en el abierto es abierto; su complemento también lo es, pues cada punto tiene una bola convexa contenida en el dominio. La conexidad implica que ese conjunto es todo el dominio. Aplicamos el argumento en cada segmento de la poligonal.
\end{proof}

\subsection{Producto, potencias e inversión}
En esta sección $A$ es un álgebra de Banach unitaria no nula. Podemos normalizar $\norm1=1$ mediante la norma equivalente del ejercicio~\ref{ej:v3-l19-3}; esto no modifica diferenciabilidad, espectro ni límites de raíces de normas de potencias.

\begin{proposicion}\label{v3:derivada-producto}
La multiplicación $m:A\times A\to A$ es diferenciable y
\[
 Dm(x,y)[h,k]=hy+xk.
\]
\end{proposicion}
\begin{proof}
Al desarrollar $(x+h)(y+k)$ obtenemos $xy+hy+xk+hk$. La parte propuesta es lineal y acotada, y $\norm{hk}\le\norm h\norm k\le\norm{(h,k)}^2/4$. El orden $hy$ es esencial: en un álgebra no conmutativa no puede sustituirse por $yh$.
\end{proof}

\begin{lema}[Límite submultiplicativo]\label{v3:limite-raices}
Para cada $x\in A$ existe
\[
 \rho_0(x):=\lim_{n\to\infty}\norm{x^n}^{1/n}
 =\inf_{n\ge1}\norm{x^n}^{1/n}\le\norm x.
\]
Si $\rho_0(x)<1$, la serie $\sum_{n\ge0}x^n$ converge en norma y es la inversa de $1-x$.
\end{lema}
\begin{proof}
Si alguna potencia es cero, todas las posteriores lo son y el resultado es inmediato. En otro caso fijemos $N$. Para $n=qN+r$, $0\le r<N$, se tiene
$\norm{x^n}\le\norm{x^N}^{q}C_N$, donde $C_N=\max_{0\le r<N}\norm{x^r}$. Por tanto
$\limsup_n\norm{x^n}^{1/n}\le\norm{x^N}^{1/N}$. Tomando el ínfimo en $N$ y comparando con el límite inferior queda probada la igualdad. Si $\rho_0(x)<s<1$, para $n$ grande $\norm{x^n}\le s^n$, de modo que la serie converge absolutamente. Multiplicar sus sumas parciales por $1-x$ a derecha e izquierda da $1-x^{n+1}$; la continuidad del producto permite pasar al límite.
\end{proof}

\begin{teorema}[Apertura del grupo de invertibles]\label{v3:derivada-inversa}
El conjunto $G(A)$ es un grupo multiplicativo abierto. Si $x\in G(A)$ y $\norm h<\norm{x^{-1}}^{-1}$, entonces $x+h$ es invertible. La aplicación $\operatorname{inv}:G(A)\to G(A)$ es diferenciable y
\[
 D\operatorname{inv}(x)[h]=-x^{-1}hx^{-1}.
\]
\end{teorema}
\begin{proof}
Las identidades $(xy)^{-1}=y^{-1}x^{-1}$ y $(x^{-1})^{-1}=x$ prueban la afirmación algebraica. Como $x+h=x(1+x^{-1}h)$, la serie de Neumann da
\[
 (x+h)^{-1}=\sum_{n=0}^{\infty}(-x^{-1}h)^nx^{-1}
 =x^{-1}-x^{-1}hx^{-1}+R(h),
\]
con
\[
 \norm{R(h)}\le
 \frac{\norm{x^{-1}}^3\norm h^2}{1-\norm{x^{-1}}\norm h}.
\]
La cota prueba el resto requerido y la apertura. También demuestra continuidad de la inversión. En particular, la derivada depende continuamente de $x$, porque los factores $x^{-1}$ lo hacen.
\end{proof}

\section{Ejercicios sobre diferenciación y exponenciales}

\begin{ejercicio}\label{ej:v3-l21-1}
Demuestre la unicidad de la derivada de Fréchet. Explique por qué la hipótesis de que el dominio sea abierto es necesaria para esa afirmación.
\end{ejercicio}


\begin{ejercicio}\label{ej:v3-l21-2}
Determine si el producto interno de un espacio de Hilbert complejo, visto como aplicación $H\times H\to\C$, es diferenciable.
\end{ejercicio}


\begin{ejercicio}\label{ej:v3-l21-3}
Sean $x,y\in A$. Pruebe: (i) la serie $e^x=\sum_{n\ge0}x^n/n!$ converge absolutamente; (ii) si $xy=yx$, entonces $e^{x+y}=e^xe^y$; (iii) $e^x$ es invertible y su inversa es $e^{-x}$; (iv) $e^x=\lim_n(1+x/n)^n$; (v) la fórmula de producto de Trotter
\[
 e^{x+y}=\lim_{n\to\infty}(e^{x/n}e^{y/n})^n
\]
es válida sin suponer que $x$ e $y$ conmuten.
\end{ejercicio}


\begin{ejercicio}\label{ej:v3-l21-4}
Pruebe que $t\mapsto e^{tx}$ es diferenciable y su derivada es $xe^{tx}=e^{tx}x$.
\end{ejercicio}


\begin{ejercicio}\label{ej:v3-l21-5}
Sea $x\in\partial G(A)$. Pruebe: (i) para cada $y\in G(A)$,
$\norm{y^{-1}}\ge1/\norm{x-y}$; (ii) $x$ es un divisor topológico de cero por ambos lados, es decir, existen $u_n$ con $\norm{u_n}=1$ y $xu_n,u_nx\to0$; (iii) $x$ no tiene inversa izquierda ni derecha.
\end{ejercicio}


\section{El espectro en un álgebra de Banach}\label{lec:v3-22}
\fuentelectura{22}{El espectro de un elemento en un álgebra}

Desde ahora el campo es $\C$, salvo indicación contraria. Es indispensable esta hipótesis: una matriz real de rotación puede tener espectro real vacío.

\begin{definicion}
Para $x\in A$, su \emph{espectro} y su \emph{conjunto resolvente} son
\[
 \sigma_A(x)=\{\lambda\in\C:\lambda1-x\notin G(A)\},
 \qquad \varrho_A(x)=\C\setminus\sigma_A(x).
\]
Escribimos $R_x(\lambda)=(\lambda1-x)^{-1}$ para $\lambda\in\varrho_A(x)$.
\end{definicion}
\index{espectro!en un álgebra de Banach}

\begin{ejemplo}
Si $A=C(K)$, con $K$ compacto Hausdorff no vacío, entonces $\sigma(f)=f(K)$. En efecto, si $\lambda\notin f(K)$, la función $t\mapsto1/(\lambda-f(t))$ es continua y es la inversa de $\lambda-f$. Si $f(t_0)=\lambda$, evaluar cualquier supuesta identidad $(\lambda-f)g=1$ en $t_0$ da $0=1$.
\end{ejemplo}

\subsection{Integración vectorial y el puente con análisis complejo}
Para una función continua $F:[a,b]\to A$, las sumas de Riemann forman una red de Cauchy por continuidad uniforme; la completitud de $A$ define $\int_a^bF(t)\,dt$. La desigualdad triangular pasa al límite y da $\norm{\int F}\le\int\norm F$. Una curva a trozos $C^1$ define $\int_\gamma F(z)\,dz=\int F(\gamma(t))\gamma'(t)\,dt$. Toda aplicación lineal acotada conmuta con esta integral, porque conmuta con las sumas finitas y con sus límites. Las integrales dobles de funciones continuas sobre productos compactos pueden intercambiarse: las dos iteraciones son límites de las mismas sumas rectangulares.

Recordemos la forma precisa de los resultados escalares que utilizaremos. Si $f$ es holomorfa en un abierto $U$ y $\Gamma$ es un ciclo cerrado a trozos suave cuya región interior está contenida en $U$, entonces $\int_\Gamma f=0$. Para $w$ fuera de $\Gamma$,
\[
 \frac1{2\pi i}\int_\Gamma\frac{f(z)}{z-w}\,dz
 =\operatorname{Ind}(\Gamma,w)f(w)
\]
si $w\in U$; el miembro izquierdo es cero cuando el integrando es holomorfo en toda la región encerrada. En particular, para una curva simple orientada positivamente el índice interior vale uno. La simple holomorfía cerca de un punto, sin control de la región encerrada, no basta para el teorema integral.

Para recordar por qué no se necesita continuidad de la derivada en estas fórmulas, el argumento de Goursat subdivide un triángulo en cuatro y conserva uno cuya integral tiene al menos un cuarto del módulo original. Los triángulos anidados convergen a $z_0$; la expansión $f(z)=f(z_0)+f'(z_0)(z-z_0)+o(|z-z_0|)$ y la anulación de las integrales de las dos primeras partes dan, tras multiplicar por $4^n$, una cota que tiende a cero para la integral inicial. La cancelación de aristas interiores extiende la conclusión a regiones trianguladas. Al quitar un disco pequeño alrededor de $w$, se aplica el teorema integral a $f(z)/(z-w)$ en la región restante; en la circunferencia pequeña, separar $f(z)=f(w)+(f(z)-f(w))$ deja $2\pi i f(w)$ y un error acotado por $2\pi\sup|f(z)-f(w)|\to0$. La misma partición, contando orientaciones, da la versión para ciclos. Aproximar las fronteras a trozos suaves por polígonos dentro del abierto no cambia las integrales en el límite, por continuidad uniforme en una vecindad compacta de la curva.

Estas fórmulas valen también para funciones holomorfas con valores en $A$: aplíquense a $\varphi\circ F$ para cada $\varphi\in A^*$ y úsese Hahn--Banach para recuperar la identidad vectorial. Del mismo modo, una función entera vectorial acotada es constante. En el caso escalar esto se obtiene de la fórmula de Cauchy para la derivada, $|f'(w)|\le M/R$, dejando $R\to\infty$; la versión vectorial se obtiene por separación mediante funcionales.

\begin{teorema}[Compacidad y no vacuidad]\label{v3:espectro-algebra-compacto}
El espectro de cada elemento de un álgebra de Banach compleja unitaria no nula es compacto y no vacío; además $\sigma(x)\subset\{\lambda:|\lambda|\le\norm x\}$.
\end{teorema}
\begin{proof}
La apertura de $G(A)$ implica que el resolvente es abierto. Para $|\lambda|>\norm x$, la serie de Neumann da
\[
 R_x(\lambda)=\sum_{n=0}^{\infty}\frac{x^n}{\lambda^{n+1}},
 \qquad\norm{R_x(\lambda)}\le\frac1{|\lambda|-\norm x}.
\]
Por tanto el espectro es cerrado y acotado. La regla de la cadena y el teorema~\ref{v3:derivada-inversa} muestran que $R_x$ es holomorfa en el resolvente, con derivada $-R_x(\lambda)^2$. Si el espectro fuera vacío, esta función sería entera. La cota anterior la acota fuera de un disco y la continuidad la acota dentro. Liouville y el límite $R_x(\lambda)\to0$ en infinito implicarían $R_x\equiv0$, imposible porque $(\lambda1-x)R_x(\lambda)=1\ne0$.
\end{proof}

\begin{teorema}[Fórmula del radio espectral de Gelfand]\label{v3:radio-gelfand}
Para cada $x\in A$,
\[
 r(x):=\max_{\lambda\in\sigma(x)}|\lambda|
 =\lim_{n\to\infty}\norm{x^n}^{1/n}.
\]
\end{teorema}
\begin{proof}
Si $|\lambda|>\rho_0(x)$, el lema~\ref{v3:limite-raices} aplicado a $x/\lambda$ prueba que $\lambda1-x$ es invertible. Luego $r(x)\le\rho_0(x)$. Para la desigualdad opuesta fijemos $s>r(x)$. Sobre una circunferencia de radio $R>\norm x$, la expansión geométrica converge uniformemente y la integración término a término da
\[
 x^n=\frac1{2\pi i}\int_{|z|=R}z^nR_x(z)\,dz.
\]
El integrando es holomorfo en el anillo $s\le|z|\le R$; el teorema integral de Cauchy permite reemplazar $R$ por $s$. Por ello
\[
 \norm{x^n}\le s^{n+1}\max_{|z|=s}\norm{R_x(z)}.
\]
Tomar raíces $n$-ésimas y pasar al límite da $\rho_0(x)\le s$. Dejando $s\downarrow r(x)$, incluso cuando $r(x)=0$, obtenemos la igualdad.
\end{proof}

\begin{proposicion}[Operaciones elementales con el espectro]\label{v3:espectro-operaciones}
Para $x,y\in A$ y $\alpha\in\C$:
\[
 \sigma(x+\alpha1)=\sigma(x)+\alpha,\qquad
 \sigma(\alpha x)=\alpha\sigma(x).
\]
Si $x$ es invertible, $\sigma(x^{-1})=\{\lambda^{-1}:\lambda\in\sigma(x)\}$. Además,
\[
 \sigma(xy)\setminus\{0\}=\sigma(yx)\setminus\{0\}.
\]
\end{proposicion}
\begin{proof}
Las dos primeras fórmulas resultan de restar el parámetro espectral y factorizar el escalar; para $\alpha=0$, ambos lados son $\{0\}$ porque el espectro de $x$ es no vacío y $\sigma(0)=\{0\}$. Para $\mu\ne0$,
$x^{-1}-\mu1=-\mu x^{-1}(x-\mu^{-1}1)$; ambos factores externos son invertibles. El cero no está en ninguno de los espectros relevantes.

Si $\lambda\ne0$ y $xy-\lambda1$ es invertible, la identidad
\[
 (yx-\lambda1)^{-1}=\lambda^{-1}\bigl(y(xy-\lambda1)^{-1}x-1\bigr)
\]
se verifica multiplicando por ambos lados: use $(yx-\lambda1)y=y(xy-\lambda1)$ y $x(yx-\lambda1)=(xy-\lambda1)x$. Intercambiar $x,y$ completa la prueba. La formulación mediante eliminación del cero expresa de forma simétrica exactamente qué parte de los dos espectros coincide.
\end{proof}

\begin{corolario}[Gelfand--Mazur]\label{v3:gelfand-mazur}
Toda álgebra de Banach compleja de división es isomorfa a $\C$ mediante $\lambda\mapsto\lambda1$; con $\norm1=1$, el isomorfismo es isométrico.
\end{corolario}
\begin{proof}
Para $x\in A$ elija $\lambda\in\sigma(x)$. Como todo elemento no nulo es invertible, necesariamente $x-\lambda1=0$. La aplicación indicada es entonces sobreyectiva, además de lineal, multiplicativa e inyectiva.
\end{proof}

\begin{corolario}
Si un álgebra de Banach compleja unitaria $A$ no es isomorfa a $\C$, tiene un divisor topológico de cero no nulo. Si además es conmutativa, tiene un ideal propio no nulo.
\end{corolario}
\begin{proof}
Elija $x$ que no sea múltiplo de $1$. El compacto no vacío $\sigma(x)$ tiene frontera no vacía: un subconjunto de $\C$ con frontera vacía sería simultáneamente abierto y cerrado, y por conexidad sería vacío o todo $\C$, alternativas imposibles. Para $\lambda\in\partial\sigma(x)$ hay $\lambda_n\notin\sigma(x)$ con $\lambda_n\to\lambda$. Así $x-\lambda1$ es un punto no invertible adherente a $G(A)$; el ejercicio~\ref{ej:v3-l21-5} prueba que es divisor topológico de cero. Es no nulo por la elección de $x$.

En el caso conmutativo, si los únicos ideales fueran $0,A$, para cada $x\ne0$ el ideal $xA$ sería $A$ y existiría $y$ con $xy=1$. La conmutatividad haría que $y$ fuera inversa bilateral y Gelfand--Mazur daría $A\simeq\C$, contradicción.
\end{proof}

\begin{observacion}[El ambiente importa]\label{v3:fallo-permanencia}
Si $B\subset A$ es una subálgebra cerrada con la misma identidad y $x\in B$, siempre $\sigma_A(x)\subset\sigma_B(x)$: una inversa en $B$ también lo es en $A$. La igualdad afirmada sin condiciones adicionales en el corolario 9 de la lectura 22 es falsa.

En efecto, tome $A=C(\mathbb T)$ y sea $B$ la clausura uniforme en $\mathbb T$ de los polinomios en $z$ con potencias no negativas. Es una subálgebra cerrada unitaria. Para $u(z)=z$, $\sigma_A(u)=\mathbb T$. En cambio $\sigma_B(u)=\overline{\mathbb D}$. Para justificarlo sin presuponer una representación de $B$, note que el principio del máximo da $|p(w)|\le\norm p_{\mathbb T}$ cuando $|w|\le1$. Por ello la evaluación $p\mapsto p(w)$ se extiende a un homomorfismo $\chi_w:B\to\C$. Si $u-w$ tuviera inversa en $B$, aplicar $\chi_w$ daría $0=1$. Luego todo el disco cerrado pertenece al espectro; la inclusión opuesta se sigue de $\norm u=1$. Más adelante probaremos que la igualdad sí se cumple para subálgebras $C^*$ unitarias.
\end{observacion}

\begin{teorema}[Semicontinuidad superior del espectro]\label{v3:semicontinuidad-espectro}
Si $x_n\to x$ en norma, para cada $\eps>0$ existe $N$ tal que
$\sigma(x_n)\subset\sigma(x)+B(0,\eps)$ cuando $n\ge N$.
\end{teorema}
\begin{proof}
Sea $U=\sigma(x)+B(0,\eps)$. El resolvente $R_x$ está uniformemente acotado en $\C\setminus U$: fuera de un disco grande esto se sigue de la cota geométrica y dentro del disco, de continuidad sobre el compacto que queda al quitar $U$. Sea $M$ una cota positiva. Si $\norm{y-x}<1/M$ y $\lambda\notin U$, entonces
\[
 \lambda1-y=(\lambda1-x)\bigl(1-R_x(\lambda)(y-x)\bigr)
\]
es invertible por Neumann. Por tanto $\sigma(y)\subset U$, y finalmente se aplica la convergencia $x_n\to x$.
\end{proof}

\section{Ejercicios sobre el espectro}

\begin{ejercicio}\label{ej:v3-l22-1}
Si $p$ es un polinomio complejo, pruebe $\sigma(p(x))=p(\sigma(x))$.
\end{ejercicio}


\begin{ejercicio}\label{ej:v3-l22-2}
Para $x\in A$, $y\in B$, pruebe $\sigma_{A\times B}((x,y))=\sigma_A(x)\cup\sigma_B(y)$.
\end{ejercicio}


\begin{ejercicio}\label{ej:v3-l22-3}
Pruebe que $\lambda\mapsto(x-\lambda1)^{-1}$ es diferenciable en el resolvente y que su derivada de Fréchet es continua.
\end{ejercicio}


\begin{ejercicio}\label{ej:v3-l22-4}
En $C^1([0,1])$ con $\norm f=\norm f_\infty+\norm{f'}_\infty$, calcule el radio espectral de $f(t)=t$.
\end{ejercicio}


\begin{ejercicio}\label{ej:v3-l22-5}
Pruebe $r(xy)=r(yx)$.
\end{ejercicio}


\begin{ejercicio}\label{ej:v3-l22-6}
Si $A$ es conmutativa, pruebe
\[
 \sigma(x+y)\subset\sigma(x)+\sigma(y),\qquad
 \sigma(xy)\subset\sigma(x)\sigma(y),
\]
y dé ejemplos de contenciones propias.
\end{ejercicio}


\section{Evaluar funciones holomorfas en un elemento}\label{lec:v3-23}
\fuentelectura{23}{El cálculo funcional holomorfo}

Una función holomorfa cerca del espectro contiene suficiente información para definir $f(x)$. No se requiere que esté definida en todo $\C$, ni que su dominio sea simplemente conexo. Esta precisión permite, por ejemplo, usar $f(z)=1/z$ cuando el espectro evita cero.

\begin{definicion}
Sea $K=\sigma(x)$. Denotamos por $\mathcal O(K)$ el álgebra de \emph{gérmenes holomorfos cerca de $K$}: dos funciones definidas en abiertos que contienen a $K$ representan el mismo germen si coinciden en alguna vecindad de $K$. Para $f$ holomorfa en $U\supset K$, elija una región poligonal acotada $\Omega$ con
$K\subset\Omega\subset\overline\Omega\subset U$ y $\partial\Omega\cap K=\varnothing$, y oriente positivamente su frontera completa $\Gamma$ (las fronteras de los agujeros tienen orientación negativa). Definimos
\[
 f(x)=\frac1{2\pi i}\int_\Gamma f(z)R_x(z)\,dz.
\]
\end{definicion}
\index{cálculo funcional!holomorfo}

Tal región se construye con una cuadrícula suficientemente fina: tome una unión finita de cuadrados que cubra $K$, añada una capa de cuadrados vecinos y conserve su clausura dentro de $U$. La frontera queda separada de $K$. Puede haber varias componentes y agujeros; la orientación completa asegura índice uno sobre $K$ y cero fuera de $\overline\Omega$. La integral existe porque el integrando es continuo en una curva compacta. Dos elecciones dan el mismo valor: se elige una tercera región más pequeña alrededor de $K$, contenida en ambas, y se aplica Cauchy a la diferencia de fronteras, donde el resolvente es holomorfo. Este argumento prueba también que sólo importa el germen de $f$.

\begin{teorema}[Cálculo holomorfo]\label{v3:calculo-holomorfo}
La aplicación $H_x:\mathcal O(\sigma(x))\to A$, $f\mapsto f(x)$, es un homomorfismo unital. Envía $z$ a $x$, es compatible con la composición y satisface:
\begin{enumerate}[label=(\roman*)]
\item $(\alpha f+\beta g)(x)=\alpha f(x)+\beta g(x)$ y $(fg)(x)=f(x)g(x)$;
\item si $g$ es holomorfa cerca de $f(\sigma(x))$, entonces $(g\circ f)(x)=g(f(x))$;
\item si $f_n\to f$ uniformemente en compactos de un mismo abierto que contiene a $\sigma(x)$, entonces $f_n(x)\to f(x)$ en norma.
\end{enumerate}
Además,
\[
 \sigma(f(x))=f(\sigma(x)).
\]
\end{teorema}
\begin{proof}
Para $1$ y $z$, podemos integrar sobre un círculo grande, pues ambas funciones son enteras. La expansión geométrica de $R_x$ y las integrales de los monomios dan $1(x)=1$ y $z(x)=x$. La linealidad es inmediata. Para el producto elija dos fronteras anidadas $\Gamma_1,\Gamma_2$, la primera dentro de la segunda, ambas en un dominio común de $f,g$. La identidad resolvente
\[
 R_x(z)R_x(w)=\frac{R_x(w)-R_x(z)}{z-w}
\]
se verifica multiplicando por $(z1-x)(w1-x)$. Integrando primero en $z$ en el término con $R_x(w)$ se obtiene cero, porque $w$ está fuera de la región interior de $\Gamma_1$. Integrando primero en $w$ en el término con $R_x(z)$, la fórmula de Cauchy da $g(z)$. Fubini para funciones continuas sobre las curvas justifica el cambio de orden. Queda
\[
 f(x)g(x)=\frac1{2\pi i}\int_{\Gamma_1}f(z)g(z)R_x(z)\,dz=(fg)(x).
\]

Probemos ahora el mapeo espectral, antes de usarlo para la composición. Si $\mu\notin f(\sigma(x))$, la función $1/(f-\mu)$ es holomorfa en una vecindad del espectro. El homomorfismo ya probado produce una inversa de $f(x)-\mu1$. Así $\sigma(f(x))\subset f(\sigma(x))$. Recíprocamente, para $\lambda\in\sigma(x)$, defina
\[
 h(z)=\frac{f(z)-f(\lambda)}{z-\lambda},\quad h(\lambda)=f'(\lambda).
\]
La singularidad es removible, y $f(x)-f(\lambda)1=(x-\lambda1)h(x)$. Los dos factores conmutan, pues pertenecen a la imagen de un álgebra conmutativa. Si el producto fuera invertible, también lo sería $x-\lambda1$, contradicción. Se obtiene la inclusión inversa.

Para la composición, tome una frontera $\Delta$ en el dominio de $g$ que encierre $f(\sigma(x))$. Se puede reducir la vecindad de $\sigma(x)$ para que $f$ de su clausura esté dentro de esa región y lejos de $\Delta$. Para $w\in\Delta$, el homomorfismo muestra
\[
 (w1-f(x))^{-1}=\frac1{2\pi i}\int_\Gamma
 \frac{R_x(z)}{w-f(z)}\,dz.
\]
Sustituya esta fórmula en la integral que define $g(f(x))$. Fubini y Cauchy en la variable $w$ convierten la integral interior en $g(f(z))$, y queda exactamente $(g\circ f)(x)$.

Por último, sobre una frontera fija común,
\[
 \norm{f_n(x)-f(x)}\le
 \frac{\operatorname{long}(\Gamma)}{2\pi}
 \max_{z\in\Gamma}\norm{R_x(z)}\,
 \max_{z\in\Gamma}|f_n(z)-f(z)|\longrightarrow0.
\]
La estimación sólo necesita convergencia uniforme en la curva elegida; la convergencia uniforme en compactos la garantiza.
\end{proof}

\begin{proposicion}[Derivada de una función de un elemento]\label{v3:derivada-calculo}
Si $f$ es holomorfa en $U$, la aplicación $x\mapsto f(x)$ es diferenciable en el abierto $\{x:\sigma(x)\subset U\}$. Su derivada es
\[
 Df_A(x)[h]=\frac1{2\pi i}\int_\Gamma f(z)R_x(z)hR_x(z)\,dz.
\]
\end{proposicion}
\begin{proof}
La apertura del dominio se sigue del teorema~\ref{v3:semicontinuidad-espectro}. Fije $\Gamma$ y $M=\max_\Gamma\norm{R_x(z)}$. Para $M\norm h<1/2$, la serie geométrica uniformemente convergente sobre $\Gamma$ da
\[
 R_{x+h}(z)=R_x(z)+R_x(z)hR_x(z)+E_z(h),
 \qquad \norm{E_z(h)}\le2M^3\norm h^2.
\]
Integre multiplicando por $f(z)$. La parte lineal es acotada por
$\operatorname{long}(\Gamma)\max_\Gamma|f|M^2\norm h/(2\pi)$; el resto es $O(\norm h^2)$ uniformemente, y por tanto es $o(\norm h)$. Que la integral lineal tienda a cero no habría bastado por sí solo: es esta cota del resto la que demuestra diferenciabilidad. En general la derivada no es $f'(x)h$, salvo cuando $h$ conmuta con $x$.
\end{proof}

\section{Ejercicios sobre el cálculo holomorfo}

\begin{ejercicio}\label{ej:v3-l23-1}
Si $T\in\Lcal(V)$ y $Tv=\lambda v$ con $v\ne0$, pruebe que $f(T)v=f(\lambda)v$ para toda $f$ holomorfa cerca de $\sigma(T)$.
\end{ejercicio}


\begin{ejercicio}\label{ej:v3-l23-2}
Si $f$ toma el valor cero en algún punto de $\sigma(x)$, pruebe que $f(x)$ no es invertible.
\end{ejercicio}


\begin{ejercicio}\label{ej:v3-l23-3}
Sean $A,B$ álgebras de Banach complejas unitarias y $T:A\to B$ un homomorfismo \emph{continuo y unital}. Pruebe $f(Tx)=T(f(x))$ para $f$ holomorfa cerca de $\sigma_A(x)$.
\end{ejercicio}


\begin{ejercicio}\label{ej:v3-l23-4}
Si $e^x=1$, demuestre $\sigma(x)\subset2\pi i\Z$.
\end{ejercicio}


\begin{ejercicio}\label{ej:v3-l23-5}
Sea $f(z)=\sum_{n\ge0}a_nz^n$ con radio de convergencia $R>r(x)$. Pruebe $f(x)=\sum_{n\ge0}a_nx^n$.
\end{ejercicio}


\begin{ejercicio}\label{ej:v3-l23-6}
Para $T\in\Lcal(V)$, formule y pruebe la compatibilidad del cálculo holomorfo con el adjunto en el dual. Distinga esta operación del adjunto de Hilbert.
\end{ejercicio}


\begin{ejercicio}\label{ej:v3-l23-7}
Sea $I$ un ideal de $A$, no necesariamente cerrado, $x\in I$, y sea $f$ holomorfa en una vecindad de $\sigma(x)\cup\{0\}$, con $f(0)=0$. Pruebe $f(x)\in I$.
\end{ejercicio}


La función resolvente ha convertido una cuestión algebraica ---la existencia de inversas--- en una herramienta analítica. En un álgebra conmutativa veremos que el espectro puede describirse mediante homomorfismos escalares. En las álgebras $C^*$, la involución permitirá extender después el cálculo de funciones holomorfas a todas las funciones continuas del espectro. Para otras presentaciones y desarrollos del cálculo holomorfo pueden consultarse \cite{conway1990,rudin1991,murphy1990}.

\chapter{Caracteres y transformada de Gelfand}\label{cap:v3-12}

Un carácter convierte cada elemento de un álgebra conmutativa en un número, respetando sumas, productos e identidad. Reunir todos esos modos de evaluación produce un espacio compacto; cada elemento del álgebra se convierte entonces en una función continua sobre ese espacio. La construcción no siempre conserva toda la información: su núcleo medirá precisamente lo que los caracteres no pueden detectar.

\section{Caracteres e ideales maximales}\label{lec:v3-24}
\fuentelectura{24}{La transformada de Gelfand}

En este capítulo $A$ es un álgebra de Banach compleja, conmutativa, unitaria y no nula, con $\norm1=1$.

\begin{definicion}
Un \emph{carácter} de $A$ es un homomorfismo complejo lineal y multiplicativo $\varphi:A\to\C$ que no es idénticamente cero. El conjunto de caracteres se denota $\Delta(A)$.
\end{definicion}
\index{carácter}\index{espacio de ideales maximales}

La exclusión del homomorfismo cero es indispensable. Si $\varphi(a)\ne0$, la identidad $\varphi(1)\varphi(a)=\varphi(a)$ da $\varphi(1)=1$. No necesitamos suponer continuidad: será consecuencia de la estructura algebraica.

\begin{proposicion}\label{v3:caracteres-maximales}
Cada carácter satisface $\varphi(a)\in\sigma(a)$ y $\norm\varphi=1$. Su núcleo es un ideal maximal cerrado. Recíprocamente, cada ideal maximal $M$ es el núcleo de un único carácter, determinado por
\[
 \varphi_M(a)=\lambda\quad\Longleftrightarrow\quad a-\lambda1\in M.
\]
Además,
\[
 \sigma(a)=\{\varphi(a):\varphi\in\Delta(A)\}.
\]
\end{proposicion}
\begin{proof}
Si $a-\varphi(a)1$ fuera invertible, aplicar $\varphi$ a su producto con la inversa daría $0=1$. Así $|\varphi(a)|\le r(a)\le\norm a$, y $\norm\varphi\le1$; evaluar en $1$ da la igualdad. El núcleo es cerrado por continuidad y es ideal por multiplicatividad. Como $\varphi$ es sobreyectiva ($\varphi(\lambda1)=\lambda$), el cociente por su núcleo es $\C$ y el núcleo es maximal.

Para la recíproca, un ideal maximal de un álgebra de Banach unitaria es cerrado, y $A/M$ es un álgebra de Banach de división. Gelfand--Mazur identifica cada clase con un único $\lambda[1]$. La unicidad se sigue también de que $\lambda1\in M$ con $\lambda\ne0$ implicaría $1\in M$. Sumar y multiplicar clases prueba que $\varphi_M$ es lineal, multiplicativa y unital; su núcleo es exactamente $M$. Cualquier carácter con ese núcleo debe asignar a $a$ el mismo escalar, lo cual prueba la unicidad.

Ya se probó que todos los valores de caracteres pertenecen al espectro. Si $\lambda\in\sigma(a)$, el ideal $(a-\lambda1)A$ es propio, pues de otro modo $a-\lambda1$ tendría inversa en el álgebra conmutativa. Está contenido en un ideal maximal $M$ y entonces $\varphi_M(a)=\lambda$. Esto prueba la inclusión restante y, en particular, la existencia de caracteres.
\end{proof}

\begin{ejemplo}[Series absolutamente convergentes]
En $\ell^1(\N_0)$ con el producto de Cauchy, para $|z|\le1$ defina
\[
 \varphi_z(a)=\sum_{n=0}^{\infty}a_nz^n.
\]
La serie converge absolutamente y $|\varphi_z(a)|\le\norm a_1$, por lo que la aplicación es lineal continua. Además $\varphi_z(e_0)=1$. La doble serie $\sum_{j,k}|a_jb_kz^{j+k}|$ es finita; reagruparla por $n=j+k$ demuestra $\varphi_z(a*b)=\varphi_z(a)\varphi_z(b)$. Más adelante verificaremos que éstos son todos los caracteres.
\end{ejemplo}

\section{La topología del espacio de caracteres}

Dotamos a $\Delta(A)\subset A^*$ de la topología débil estrella: una red $\varphi_\alpha$ converge a $\varphi$ si y sólo si $\varphi_\alpha(a)\to\varphi(a)$ para cada $a\in A$. Esta topología registra la convergencia de todas las evaluaciones, sin exigir convergencia en norma de los funcionales.

\begin{teorema}\label{v3:compacto-caracteres}
El espacio $\Delta(A)$ es compacto Hausdorff. Si $A$ es separable, es metrizable.
\end{teorema}
\begin{proof}
La topología débil estrella es Hausdorff porque dos funcionales distintos difieren en algún elemento de $A$. La bola dual cerrada unitaria es compacta por Banach--Alaoglu. Dentro de ella,
\[
 \Delta(A)=\{\varphi:\varphi(1)=1,\ \varphi(ab)=\varphi(a)\varphi(b)\text{ para todo }a,b\in A\}
\]
es cerrado: cada igualdad se conserva bajo convergencia puntual de funcionales. En particular, la ecuación $\varphi(1)=1$ impide que una red de caracteres tenga límite cero; verificar únicamente la multiplicatividad no bastaría.

Si $\{a_n\}$ es denso en la bola unitaria de $A$, la fórmula
\[
 d(\varphi,\psi)=\sum_{n=1}^{\infty}2^{-n}|\varphi(a_n)-\psi(a_n)|
\]
define una métrica en la bola dual: la serie converge pues cada diferencia es a lo sumo $2$, y la densidad garantiza que distancia cero implica igualdad. La convergencia débil estrella implica convergencia para esta métrica al controlar primero una suma finita y luego la cola geométrica. Recíprocamente, la convergencia de las evaluaciones en $a_n$, junto con la cota uniforme de las normas, permite aproximar cualquier $a$ por un múltiplo de algún $a_n$ y obtener convergencia en $a$. Así la métrica induce exactamente la topología deseada en $\Delta(A)$.
\end{proof}

\section{La representación por funciones continuas}

\begin{definicion}
Para $a\in A$, su \emph{transformada de Gelfand} es la función
\[
 \widehat a:\Delta(A)\to\C,\qquad \widehat a(\varphi)=\varphi(a).
\]
La aplicación de Gelfand es $\mathcal G:A\to C(\Delta(A))$, $a\mapsto\widehat a$.
\end{definicion}
\index{transformada de Gelfand}

\begin{teorema}\label{v3:transformada-gelfand}
La aplicación $\mathcal G$ es un homomorfismo unital continuo, y
\[
 \widehat a(\Delta(A))=\sigma(a),\qquad
 \norm{\widehat a}_\infty=r(a)\le\norm a.
\]
Si $f$ es holomorfa cerca de $\sigma(a)$, entonces
\[
 \widehat{f(a)}=f\circ\widehat a.
\]
\end{teorema}
\begin{proof}
La función $\widehat a$ es continua por la definición de la topología débil estrella: es precisamente una de sus evaluaciones generadoras. La desigualdad $|\varphi(a)-\psi(a)|\le\norm{\varphi-\psi}\norm a$ sólo probaría continuidad para la topología de la norma, que es más fuerte y no resuelve esta cuestión. La linealidad y la multiplicatividad de cada carácter dan, punto por punto,
$\widehat{\alpha a+\beta b}=\alpha\widehat a+\beta\widehat b$, $\widehat{ab}=\widehat a\,\widehat b$ y $\widehat1=1$. La descripción del espectro de la proposición~\ref{v3:caracteres-maximales} da las igualdades de imagen y norma, de donde resulta continuidad de $\mathcal G$.

Finalmente, un carácter conmuta con la integral vectorial por ser lineal continuo y envía $(z1-a)^{-1}$ a $(z-\varphi(a))^{-1}$. La fórmula de Cauchy implica
\[
 \varphi(f(a))=\frac1{2\pi i}\int_\Gamma
 \frac{f(z)}{z-\varphi(a)}\,dz=f(\varphi(a)).
\]
Esto vale para cada carácter, que es exactamente la última identidad.
\end{proof}

\section{Ejercicios sobre el radical y la representación}

\begin{ejercicio}\label{ej:v3-l24-1}
Un elemento $a$ es \emph{cuasinilpotente} si $\sigma(a)=\{0\}$. (i) Pruebe que
\[
 J=\{a\in A:ab\text{ es cuasinilpotente para todo }b\in A\}
\]
es un ideal cerrado, llamado radical de Jacobson, y que es la intersección de todos los ideales maximales. (ii) Pruebe que todo elemento cuasinilpotente es divisor topológico de cero. (iii) Pruebe $\ker\mathcal G=J$. Verifique además que cuasinilpotente no implica nilpotente para el operador de desplazamiento ponderado
\[
 T(x_1,x_2,\ldots)=(0,x_1,x_2/2,x_3/3,\ldots)
 \quad\text{en }\ell^\infty.
\]
\end{ejercicio}


\begin{ejercicio}\label{ej:v3-l24-2}
Pruebe que todo carácter de $\ell^1(\N_0)$ con convolución de Cauchy es $\varphi_z$ para un único $z$ con $|z|\le1$.
\end{ejercicio}


\begin{ejercicio}\label{ej:v3-l24-3}
Demuestre que todo carácter es una aplicación abierta.
\end{ejercicio}


\begin{ejercicio}\label{ej:v3-l24-4}
Pruebe $\sigma(a+b)\subset\sigma(a)+\sigma(b)$ y $\sigma(ab)\subset\sigma(a)\sigma(b)$, y dé ejemplos en que ambas inclusiones sean propias.
\end{ejercicio}


\begin{ejercicio}\label{ej:v3-l24-5}
Pruebe que $a$ es invertible si y sólo si $\widehat a$ es invertible en $C(\Delta(A))$.
\end{ejercicio}


\begin{ejercicio}\label{ej:v3-l24-6}
Pruebe $r(a)=\norm{\widehat a}_\infty$. Deduzca que $\mathcal G$ es isométrica si y sólo si $\norm{a^2}=\norm a^2$ para cada $a\in A$.
\end{ejercicio}


El radical explica por qué una representación por funciones escalares puede perder información. La identidad $C^*$ eliminará esa pérdida en el caso conmutativo y permitirá identificar el álgebra completa con todas las funciones continuas de su espacio de caracteres. La exposición se complementa con \cite{maxgelfand,murphy1990,rudin1991}.

\parte{Álgebras C* y representaciones}{La involución relaciona norma, espectro y positividad. El cálculo continuo produce raíces y descomposiciones; los estados dan lugar a representaciones mediante la construcción GNS.}
\chapter{Álgebras \texorpdfstring{$C^*$}{C*} y cálculo funcional continuo}\label{cap:v3-13}

El adjunto de un operador no es una operación auxiliar: vincula su estructura algebraica con su norma. La identidad $\norm{a^*a}=\norm a^2$ recoge ese vínculo y tiene consecuencias muy fuertes: determina la norma, permite reconocer el espectro de elementos normales y convierte un álgebra conmutativa en un espacio de funciones continuas.

\section{Involución e identidad \texorpdfstring{$C^*$}{C*}}\label{lec:v3-25}
\fuentelectura{25}{Álgebras C*}

\begin{definicion}
Una involución en un álgebra compleja $A$ es una operación $a\mapsto a^*$ tal que
\[
 (\alpha a+\beta b)^*=\overline\alpha a^*+\overline\beta b^*,\quad
 (ab)^*=b^*a^*,\quad (a^*)^*=a.
\]
Una \emph{álgebra $C^*$} es un álgebra de Banach con involución que satisface $\norm{a^*a}=\norm a^2$ para todo $a$. No se exige identidad multiplicativa.
\end{definicion}
\index{algebra C*@$C^*$-álgebra}\index{involucion@involución}

La segunda propiedad se llama \emph{antimultiplicatividad}: invierte el orden del producto; no afirma $ab=-ba$. Las conjugaciones de los escalares son indispensables en la primera propiedad.

\begin{proposicion}\label{v3:involucion-isometrica}
La involución de un álgebra $C^*$ es isométrica. Si existe identidad, $1^*=1$ y $\norm1=1$ en el álgebra no nula. Además $(a^{-1})^*=(a^*)^{-1}$ para todo invertible.
\end{proposicion}
\begin{proof}
Para $a\ne0$, $\norm a^2=\norm{a^*a}\le\norm{a^*}\norm a$, de donde $\norm a\le\norm{a^*}$. Aplicar esto a $a^*$ da la desigualdad inversa. De $a1=a=1a$ y la sobreyectividad de la involución se obtiene que $1^*$ también es identidad, luego coincide con $1$. La identidad $C^*$ da $\norm1=\norm1^2$, y la norma no es cero. Aplicar la involución a $aa^{-1}=a^{-1}a=1$ prueba la fórmula de inversas.
\end{proof}

\begin{ejemplo}
$\C$, $C(K)$ para $K$ compacto Hausdorff, $C_b(X)$ y $C_0(X)$ para $X$ localmente compacto Hausdorff, y $L^\infty(\mu)$ son álgebras $C^*$ con conjugación puntual. La identidad de norma se reduce a $\norm{|f|^2}_\infty=\norm f_\infty^2$. En un espacio no compacto, $C(X)$ completo, sin restricción de acotación, no lleva en general una norma del supremo finita; ésa es la corrección necesaria al ejemplo original. En $\Lcal(H)$, el adjunto de Hilbert satisface $\norm{T^*T}=\norm T^2$ y proporciona el ejemplo no conmutativo fundamental.

El álgebra de convolución $L^1(\R)$ tiene involución $f^*(t)=\overline{f(-t)}$. Es conjugada lineal e involutiva, y
\[
 (f*g)^*(t)=\int_{\R}\overline{f(s)g(-t-s)}\,ds
 =\int_{\R}\overline{g(-u)}\,\overline{f(u-t)}\,du
 =(g^**f^*)(t)
\]
casi en todas partes, por el cambio $u=t+s$. No es un álgebra $C^*$ con la norma $L^1$. Para $f=\mathbf1_{[0,1]}-\mathbf1_{[1,2]}$, $\norm f_1^2=4$, mientras su autocorrelación $f^**f$ es par y vale $2-3t$ en $0\le t\le1$, $t-2$ en $1\le t\le2$ y cero para $t>2$. Integrando su valor absoluto se obtiene $\norm{f^**f}_1=8/3<4$.
\end{ejemplo}

\begin{definicion}
Un elemento es \emph{normal} si $a^*a=aa^*$, \emph{autoadjunto} si $a=a^*$, \emph{unitario} si $a^*a=aa^*=1$, \emph{positivo} si tiene la forma $b^*b$, y \emph{proyección} si $p=p^*=p^2$.
\end{definicion}

Todo autoadjunto y todo unitario es normal; todo $b^*b$ es autoadjunto y toda proyección es positiva porque $p=p^*p$. La equivalencia entre esta definición algebraica de positividad y la condición espectral $\sigma(a)\subset[0,\infty)$ se demostrará en el capítulo siguiente, sin utilizarla anticipadamente aquí.

\subsection{Norma y espectro de elementos normales}
Por ahora suponemos que $A$ es unitaria. La unitización que construiremos después extenderá las afirmaciones espectrales al caso no unitario.

\begin{proposicion}\label{v3:normal-radio-norma}
Se tiene $\sigma(a^*)=\overline{\sigma(a)}$. Si $a$ es normal, $r(a)=\norm a$. Si $u$ es unitario, $\sigma(u)\subset\mathbb T$; si $a$ es autoadjunto, $\sigma(a)\subset\R$.
\end{proposicion}
\begin{proof}
La fórmula de inversas muestra que $a-\lambda1$ es invertible si y sólo si $a^*-\overline\lambda1$ lo es. Para $a=a^*$, la identidad $C^*$ implica $\norm{a^2}=\norm a^2$, y al iterar se obtiene $\norm{a^{2^n}}=\norm a^{2^n}$. La fórmula de Gelfand da $r(a)=\norm a$. Si $a$ es normal,
\[
 \norm{a^2}^2=\norm{(a^*)^2a^2}
 =\norm{(a^*a)^2}=\norm{a^*a}^2=\norm a^4.
\]
Las potencias de $a$ siguen siendo normales, y la misma iteración demuestra $r(a)=\norm a$.

Si $u$ es unitario, $\norm u^2=\norm1=1$ y también $\norm{u^{-1}}=1$. Para $\lambda\in\sigma(u)$, las cotas espectrales de $u$ y $u^{-1}$ dan $|\lambda|\le1$ y $|\lambda|^{-1}\le1$, de donde $|\lambda|=1$. Si $a=a^*$, la convergencia en norma de la serie exponencial y la continuidad de la involución dan $(e^{ia})^*=e^{-ia}=(e^{ia})^{-1}$. Así $e^{ia}$ es unitario. El mapeo espectral implica $|e^{i\lambda}|=1$ para $\lambda\in\sigma(a)$; como ese módulo es $e^{-\operatorname{Im}\lambda}$, necesariamente $\operatorname{Im}\lambda=0$.
\end{proof}

\begin{teorema}[Fuglede]\label{v3:fuglede}
Si $a$ es normal y $ab=ba$ en un álgebra $C^*$ unitaria, entonces $a^*b=ba^*$.
\end{teorema}
\begin{proof}
La función $F(z)=e^{za^*}be^{-za^*}$ es entera con valores en $A$, como se ve multiplicando sus series localmente uniformes. Como $b$ conmuta con $a$, conmuta con todas sus potencias y exponenciales. Además $a$ conmuta con $a^*$. Por tanto
\[
 F(z)=e^{za^*-\overline z a}\,b\,e^{-za^*+\overline z a}.
\]
El exponente $za^*-\overline z a$ es antiautoadjunto: su adjunto es su negativo. Su exponencial es unitaria, porque la serie y la involución dan adjunto igual a la exponencial opuesta. En consecuencia $\norm{F(z)}\le\norm b$ para todo $z$. Liouville vectorial implica que $F$ es constante. Derivar en cero da $0=F'(0)=a^*b-ba^*$.

Todo elemento también admite una descomposición única $c=h+ik$ con $h,k$ autoadjuntos, dada por $h=(c+c^*)/2$ y $k=(c-c^*)/(2i)$. La normalidad de $c$ equivale a $hk=kh$, como se verifica desarrollando $c^*c-cc^*=2i(hk-kh)$. Estas identidades explican la separación en partes autoadjuntas empleada en la prueba original; la demostración anterior evita ambigüedades entre $z$ y $\overline z$.
\end{proof}

\section{Ejercicios sobre la estructura \texorpdfstring{$C^*$}{C*}}

\begin{ejercicio}\label{ej:v3-l25-1}
Si $a=a^*$, pruebe que $e^{ia}$ es unitario.
\end{ejercicio}
\begin{solucion}
La involución pasa al límite de la serie por ser isométrica, así $(e^{ia})^*=\sum(-ia)^n/n!=e^{-ia}$. Las exponenciales de $ia$ y $-ia$ multiplican a $1$ en ambos órdenes. Por tanto $(e^{ia})^*e^{ia}=e^{ia}(e^{ia})^*=1$. Si $A$ no es unitaria, la exponencial y la afirmación se entienden en su unitización.
\end{solucion}

\begin{ejercicio}\label{ej:v3-l25-2}
Para $a$ normal pruebe: (i) existe $\lambda\in\sigma(a)$ con $|\lambda|=\norm a$; (ii) si $a^n=0$ para algún $n\ge1$, entonces $a=0$.
\end{ejercicio}
\begin{solucion}
(i) El módulo alcanza su máximo en el espectro compacto no vacío, y ese máximo es $r(a)=\norm a$. (ii) La nilpotencia implica $r(a)=\lim_k\norm{a^k}^{1/k}=0$. La normalidad convierte esa igualdad en $\norm a=0$. Sin normalidad la conclusión es falsa: la matriz $\left(\begin{smallmatrix}0&1\\0&0\end{smallmatrix}\right)$ es nilpotente y no nula.
\end{solucion}

\begin{ejercicio}\label{ej:v3-l25-3}
Decimos que $a,b$ son unitariamente equivalentes si $b=uau^*$ para algún unitario $u$. Pruebe: (i) ésta es una relación de equivalencia; (ii) preserva normalidad, autoadjunción y positividad; (iii) $\sigma(uau^*)=\sigma(a)$.
\end{ejercicio}
\begin{solucion}
(i) La identidad prueba reflexividad; si $b=uau^*$, entonces $a=u^*bu$; y si $c=vbv^*$, entonces $c=(vu)a(vu)^*$, con $vu$ unitario. (ii) Para $b=uau^*$ se tiene $b^*=ua^*u^*$, $b^*b=ua^*au^*$ y $bb^*=uaa^*u^*$; las dos primeras propiedades se conservan. Si $a=c^*c$, entonces $b=(cu^*)^*(cu^*)$, y es positivo. (iii) $b-\lambda1=u(a-\lambda1)u^*$; ese elemento es invertible exactamente cuando lo es $a-\lambda1$, porque los factores externos son invertibles.
\end{solucion}

\begin{ejercicio}\label{ej:v3-l25-4}
Pruebe que los unitarios forman un subgrupo cerrado de $G(A)$. Determine si es necesariamente normal.
\end{ejercicio}
\begin{solucion}
El producto de unitarios es unitario y la inversa de $u$ es $u^*$. Si $u_n\to u$ en norma, la continuidad del producto y de la involución da $u^*u=uu^*=1$; el grupo unitario es incluso cerrado en $A$. No tiene por qué ser normal en $G(A)$. En $M_2(\C)$, tome
\[
 u=\begin{pmatrix}0&1\\1&0\end{pmatrix},\qquad
 s=\begin{pmatrix}2&0\\0&1\end{pmatrix}.
\]
Entonces $sus^{-1}=\left(\begin{smallmatrix}0&2\\1/2&0\end{smallmatrix}\right)$ y $(sus^{-1})^*(sus^{-1})=\operatorname{diag}(1/4,4)\ne I$.
\end{solucion}

\begin{ejercicio}\label{ej:v3-l25-5}
Pruebe que $A$ es conmutativa si y sólo si todos sus elementos son normales.
\end{ejercicio}
\begin{solucion}
La implicación directa es inmediata. Recíprocamente, si $h,k$ son autoadjuntos, la normalidad de $h+ik$ y el desarrollo al final de la prueba de Fuglede dan $hk=kh$. Cada elemento es combinación compleja lineal de dos autoadjuntos; distribuir los productos muestra entonces que cualesquiera dos elementos de $A$ conmutan.
\end{solucion}

\begin{ejercicio}\label{ej:v3-l25-6}
Pruebe que una misma álgebra con involución tiene a lo sumo una norma que la convierta en álgebra $C^*$.
\end{ejercicio}
\begin{solucion}
En cualquiera de esas normas, $a^*a$ es autoadjunto, de modo que
\[
 \norm a^2=\norm{a^*a}=r(a^*a).
\]
El miembro derecho sólo depende de la invertibilidad algebraica de $a^*a-\lambda1$, y no de la norma elegida. Por tanto las dos normas coinciden en cada $a$. En el caso no unitario se usa la unitización algebraica $A\oplus\C$, que es la misma para ambas normas; el teorema siguiente garantiza una norma $C^*$ en ella y permite aplicar el mismo argumento sin suponer de antemano que las normas originales sean equivalentes.
\end{solucion}

\begin{ejercicio}\label{ej:v3-l25-7}
Pruebe que un álgebra de Banach con involución es $C^*$ si y sólo si $\norm{a^*a}\ge\norm a^2$ para cada $a$.
\end{ejercicio}
\begin{solucion}
La necesidad es inmediata. Para la suficiencia, si $a\ne0$,
$\norm a^2\le\norm{a^*a}\le\norm{a^*}\norm a$, así $\norm a\le\norm{a^*}$. Sustituyendo $a$ por $a^*$ se obtiene la igualdad de normas. Por tanto $\norm{a^*a}\le\norm{a^*}\norm a=\norm a^2$, que junto con la hipótesis da la identidad $C^*$.
\end{solucion}

\begin{ejercicio}\label{ej:v3-l25-8}
Sea $(A_j)_{j\in J}$ una familia de álgebras $C^*$. Pruebe que son álgebras $C^*$, con operaciones coordenadas y norma del supremo: (i) $\prod_j^{\mathrm b}A_j=\{(a_j):\sup_j\norm{a_j}<\infty\}$; (ii) $\bigoplus_j^{c_0}A_j=\{(a_j):\text{para cada }\eps>0,\ \{j:\norm{a_j}\ge\eps\}\text{ es finito}\}$.
\end{ejercicio}
\begin{solucion}
(i) Las operaciones preservan la acotación y la norma es submultiplicativa. Una sucesión de Cauchy converge coordenada a coordenada en los espacios completos $A_j$; la cota de Cauchy uniforme pasa al límite y prueba convergencia en la norma del supremo hacia una familia acotada. La involución coordenada es isométrica y
$\norm{a^*a}=\sup_j\norm{a_j^*a_j}=\sup_j\norm{a_j}^2=\norm a^2$.

(ii) La condición dada implica acotación: fuera del conjunto finito correspondiente a $\eps=1$, las normas son menores que uno, y dentro hay un máximo finito. Sumas, productos y adjuntos conservan la condición; para productos, use $\norm{a_jb_j}\le\norm a\norm{b_j}$. Es un subespacio cerrado del producto acotado: si $a^{(n)}\to a$ uniformemente, elija $n$ con $\norm{a-a^{(n)}}<\eps/2$; entonces $\{j:\norm{a_j}\ge\eps\}\subset\{j:\norm{a_j^{(n)}}\ge\eps/2\}$, que es finito. Por tanto es completo y hereda la identidad $C^*$.
\end{solucion}

\section{Adjuntar una identidad compatible con el adjunto}\label{lec:v3-26}
\fuentelectura{26}{El cálculo funcional continuo}

Un $*$-homomorfismo es una aplicación compleja lineal, multiplicativa y que conserva la involución. Cuando se precisa preservación de la identidad se dice \emph{unital}. Más adelante veremos que entre álgebras $C^*$ la continuidad de un $*$-homomorfismo es automática.

\begin{teorema}[Unitización $C^*$]\label{v3:unitizacion-cstar}
Si $A$ es una álgebra $C^*$ no unitaria, su unitización algebraica $\widetilde A=A\oplus\C$, con
\[
 (a,\lambda)(b,\mu)=(ab+\lambda b+\mu a,\lambda\mu),\qquad
 (a,\lambda)^*=(a^*,\overline\lambda),
\]
es un álgebra $C^*$ unitaria con la norma
\[
 \norm{(a,\lambda)}_*=\sup_{\norm z\le1}\norm{az+\lambda z}.
\]
La inclusión de $A$ es isométrica y su imagen es un ideal maximal cerrado.
\end{teorema}
\begin{proof}
Las identidades de la involución se verifican coordenada a coordenada; en particular, al conjugar el producto su primera componente es $b^*a^*+\overline\lambda b^*+\overline\mu a^*$, la primera componente de $(b,\mu)^*(a,\lambda)^*$. Sea $L_{(a,\lambda)}z=az+\lambda z$. La norma propuesta es la norma de operador y $L_{xy}=L_xL_y$, de modo que es submultiplicativa.

Primero, para $a\in A$, $\norm{L_a}=\norm a$: la cota superior es inmediata, y si $a\ne0$ se prueba la inferior evaluando en $a^*/\norm a$, pues $\norm{aa^*}=\norm a^2$. Si $L_{(a,\lambda)}=0$ y $\lambda\ne0$, $-a/\lambda$ sería identidad izquierda de $A$. Una identidad izquierda $e$ en un álgebra con involución produce una identidad derecha $e^*$; como $ee^*=e^*$ por la primera propiedad y $ee^*=e$ por la segunda, ambas coinciden y $A$ sería unitaria. Por hipótesis esto es imposible. Si $\lambda=0$, $\norm{L_a}=\norm a$ da $a=0$. Así la norma es definida positiva.

Demostremos las cotas que aseguran completitud. Se tiene
\[
 |\lambda|\le\norm{L_{(a,\lambda)}},\qquad
 \norm a\le2\norm{L_{(a,\lambda)}}.
\]
Para la primera, si $\norm{L_a+\lambda I}<|\lambda|$, entonces $L_a=-\lambda[I-(L_a+\lambda I)/\lambda]$ sería invertible en $\Lcal(A)$ por Neumann. Tome $e\in A$ con $ae=a$. Para todo $b\in A$, $a(eb-b)=0$; la inyectividad de $L_a$ implica $eb=b$, contradiciendo que $A$ no tiene identidad. La segunda cota se sigue de $\norm a=\norm{L_a}\le\norm{L_a+\lambda I}+|\lambda|$. Por tanto
\[
 \norm{(a,\lambda)}_*\le\norm a+|\lambda|
 \le3\norm{(a,\lambda)}_*.
\]
La norma es equivalente a la norma suma completa, y por ello es completa.

Falta la identidad $C^*$. Para cualquier $v\in A$,
$\norm v=\sup_{\norm w\le1}\norm{wv}$, por el mismo argumento con $v^*/\norm v$. Si $x=(a,\lambda)$ y $z,w\in A$ tienen norma a lo sumo uno,
\[
 \norm{wx^*z}=\norm{z^*xw^*}\le\norm x_*.
\]
Tomar primero supremo en $w$ y luego en $z$ prueba $\norm{x^*}_*\le\norm x_*$; aplicar lo mismo a $x^*$ da igualdad. Además
\[
 \norm{xz}^2=\norm{z^*x^*xz}\le\norm{z^*}\norm{x^*xz}
 \le\norm{x^*x}_*\norm z^2.
\]
Tomar el supremo sobre $z$ da $\norm x_*^2\le\norm{x^*x}_*$. La desigualdad inversa resulta de submultiplicatividad e isometría de la involución. Finalmente, $A\oplus\{0\}$ es un ideal cerrado y el cociente, identificado con la segunda coordenada, es $\C$, luego es maximal.
\end{proof}

Para un álgebra no unitaria definimos el espectro de $a\in A$ dentro de $\widetilde A$. El cero siempre pertenece a él, porque la proyección escalar anula $a$ e impide que sea invertible. La norma original no cambia, de modo que las propiedades de elementos normales y autoadjuntos ya probadas se extienden a este caso.

\section{Representación de las álgebras conmutativas}

\begin{lema}[Permanencia espectral $C^*$]\label{v3:permanencia-cstar}
Si $B\subset A$ es una subálgebra $C^*$ cerrada que contiene la identidad de $A$, entonces $\sigma_B(b)=\sigma_A(b)$ para todo $b\in B$.
\end{lema}
\begin{proof}
Sólo hay que probar que un elemento de $B$ invertible en $A$ tiene su inversa en $B$. Primero suponga $b=b^*$ invertible en $A$. El espectro de $b$ en $A$ es un compacto real que evita cero. Elija $R>\norm b$. Por mapeo espectral polinomial y normalidad,
\[
 \norm{1-b^2/R^2}
 =\max_{t\in\sigma_A(b)}|1-t^2/R^2|<1.
\]
La serie de Neumann, cuyos términos pertenecen a $B$, muestra que $(b^2)^{-1}\in B$, y entonces $b^{-1}=b(b^2)^{-1}\in B$. Para $b$ general invertible, $b^*b$ es autoadjunto e invertible en $A$. El caso anterior da $(b^*b)^{-1}\in B$, y $b^{-1}=(b^*b)^{-1}b^*\in B$. No hemos usado positividad espectral de $b^*b$ en esta prueba, por lo que no hay circularidad con el capítulo siguiente.
\end{proof}

\begin{lema}[Stone--Weierstrass en compacto Hausdorff]\label{v3:stone-weierstrass}
Si $K$ es compacto Hausdorff y $\mathcal B\subset C(K,\C)$ es una subálgebra que contiene constantes, separa puntos y es estable bajo conjugación, entonces es uniformemente densa en $C(K,\C)$.
\end{lema}
\begin{proof}
Primero justificamos la herramienta de aproximación usada en la prueba. Para $h\in C([0,1],\R)$, los polinomios de Bernstein
\[
 B_nh(t)=\sum_{k=0}^nh(k/n)\binom nk t^k(1-t)^{n-k}
\]
convergen uniformemente a $h$. En efecto, los pesos son no negativos y suman uno; su promedio de $k/n$ es $t$ y su promedio de $(k/n-t)^2$ es $t(1-t)/n\le1/(4n)$, identidades que se obtienen diferenciando el binomio. Dado $\delta>0$, divida la suma del error entre $|k/n-t|<\delta$ y su complemento. La primera parte se acota por el módulo de continuidad $\omega_h(\delta)$; la suma de pesos de la segunda es a lo sumo $1/(4n\delta^2)$, de modo que
\[
 \norm{B_nh-h}_\infty\le\omega_h(\delta)+\frac{\norm h_\infty}{2n\delta^2}.
\]
Primero elija $\delta$ y después $n$. Por cambio afín, toda función continua real de un intervalo compacto, en particular $|t|$, admite aproximación polinomial uniforme.

Sea $\mathcal R$ la parte real de $\mathcal B$, es decir, sus elementos reales. Es una subálgebra real con constantes que separa puntos: si un elemento complejo distingue dos puntos, su parte real o su parte imaginaria también los distingue, y ambas pertenecen a $\mathcal B$ por estabilidad bajo conjugación. Su clausura uniforme $\overline{\mathcal R}$ es estable bajo valor absoluto: para $g$ en ella, aproxime $|t|$ por polinomios en un intervalo que contiene el rango de $g$, y evalúelos en $g$. Por tanto es estable bajo máximos y mínimos finitos, mediante $\max(g,h)=(g+h+|g-h|)/2$ y la fórmula análoga para el mínimo.

Fije $f\in C(K,\R)$ y $\eps>0$. Para $x,y\in K$ existe $a_{x,y}\in\mathcal R$ con $a_{x,y}(x)=f(x)$ y $a_{x,y}(y)=f(y)$: si $x\ne y$, una función que los separa se ajusta mediante una transformación afín; si coinciden, basta una constante. Para $x$ fijo, los abiertos $\{t:a_{x,y}(t)>f(t)-\eps\}$ cubren $K$ al variar $y$. Elija un subcubrimiento finito y sea $g_x$ el máximo de las correspondientes funciones. Entonces $g_x\in\overline{\mathcal R}$, $g_x>f-\eps$ en todo $K$ y $g_x(x)=f(x)$. Los abiertos $\{t:g_x(t)<f(t)+\eps\}$ cubren $K$ al variar $x$. Elija otro subcubrimiento finito y tome el mínimo $g$ de esas funciones $g_x$. Se tiene $f-\eps<g<f+\eps$ en todo $K$, y $g\in\overline{\mathcal R}$. Como esta clausura es cerrada y $\eps$ es arbitrario, $f\in\overline{\mathcal R}$. Finalmente aproxime por separado las partes real e imaginaria de una función compleja. Ningún paso utiliza una métrica en $K$.
\end{proof}

\begin{teorema}[Gelfand para álgebras $C^*$ conmutativas]\label{v3:gelfand-cstar}
Si $A$ es conmutativa y unitaria, $\mathcal G:A\to C(\Delta(A))$ es un $*$-isomorfismo isométrico sobreyectivo.
\end{teorema}
\begin{proof}
Todos los elementos son normales, por lo que $\norm{\widehat a}_\infty=r(a)=\norm a$. Así la imagen es cerrada: una sucesión de imágenes de Cauchy procede de una sucesión de Cauchy en el espacio completo $A$. Si $h=h^*$, cada carácter satisface $\varphi(h)\in\sigma(h)\subset\R$. Para $a=h+ik$ con $h,k$ autoadjuntos, se obtiene
$\varphi(a^*)=\varphi(h)-i\varphi(k)=\overline{\varphi(a)}$. Esto prueba que la imagen es estable bajo conjugación; la igualdad de conjuntos espectrales conjugados, por sí sola, no habría probado esta igualdad punto por punto.

La imagen contiene las constantes, porque $\widehat1=1$, y separa caracteres distintos, porque dos funcionales distintos difieren en algún $a$. El lema~\ref{v3:stone-weierstrass} implica que es densa, y al ser cerrada es todo $C(\Delta(A))$.
\end{proof}

\begin{proposicion}[Generadores y coordenadas]\label{v3:generador-espectro}
Si un álgebra de Banach conmutativa unitaria es la clausura de los polinomios en $a$, entonces $\varphi\mapsto\varphi(a)$ es un homeomorfismo de $\Delta(A)$ sobre $\sigma_A(a)$. Lo mismo ocurre si $A$ es $C^*$, conmutativa y unitaria, y está generada como álgebra $C^*$ unitaria por un elemento $a$.
\end{proposicion}
\begin{proof}
La aplicación es continua por la topología débil estrella y sobreyectiva por la descripción del espectro mediante caracteres. Si dos caracteres coinciden en $a$, coinciden en sus potencias, en todos los polinomios y, por continuidad, en su clausura. En el caso $C^*$ también coinciden en $a^*$ porque preservan conjugación, así que coinciden en todos los polinomios en $a,a^*$ y en su clausura. En ambos casos se obtiene una biyección continua de un compacto sobre un Hausdorff, cuya inversa es continua porque las imágenes de cerrados son compactas y por tanto cerradas.
\end{proof}

\begin{teorema}[Cálculo funcional continuo]\label{v3:calculo-continuo}
Sea $a$ normal en un álgebra $C^*$ unitaria $A$. Existe un único $*$-homomorfismo unital
\[
 C_a:C(\sigma_A(a))\longrightarrow A,\qquad f\longmapsto f(a),
\]
que envía la función coordenada $z$ a $a$. Es isométrico, su imagen es $C^*(1,a)$, extiende el cálculo holomorfo y cumple
\[
 \sigma_A(f(a))=f(\sigma_A(a)).
\]
\end{teorema}
\begin{proof}
La subálgebra $B=C^*(1,a)$, clausura de los polinomios en $a,a^*$, es conmutativa porque $a$ es normal. La permanencia espectral da $\sigma_B(a)=\sigma_A(a)$. Por la proposición anterior, $\tau:\Delta(B)\to\sigma_A(a)$, $\tau(\varphi)=\varphi(a)$, es un homeomorfismo. Defina $C_a(f)=\mathcal G_B^{-1}(f\circ\tau)$. Cada factor es un $*$-isomorfismo isométrico, por lo que se obtienen todas las propiedades algebraicas y de norma, incluida la imagen.

Si $f$ es holomorfa cerca del espectro, su integral calculada en $B$ coincide con la calculada en $A$: los resolventes pertenecen a $B$ por permanencia y son las mismas inversas. La transformada de Gelfand de esa integral es $f\circ\widehat a=f\circ\tau$, de modo que coincide con $C_a(f)$. Esta prueba evita necesitar una versión adicional del teorema de Runge. La unicidad se debe a que cualquier $*$-homomorfismo unital $F$ que envía $z$ a $a$ envía $\overline z$ a $a^*$ y coincide en los polinomios en $z,\overline z$; éstos son densos por Stone--Weierstrass. Para justificar el paso al límite sin suponer continuidad de antemano, una inversa se conserva bajo $F$, de modo que $\sigma(F(f))\subset\sigma(f)$. El elemento $F(f)$ es normal, porque $f$ conmuta con su conjugado y $F$ conserva productos y adjuntos. Por tanto $\norm{F(f)}=r(F(f))\le r(f)=\norm f_\infty$. Así $F$ es contractivo y la igualdad sobre el conjunto denso se extiende a todo el dominio.

Finalmente, la isomorfía con $C(\sigma(a))$ da el espectro de $f(a)$ dentro de $B$ como $f(\sigma(a))$. La permanencia espectral lo identifica con su espectro en $A$.
\end{proof}

\section{Ejercicios sobre funciones continuas y caracteres}

\begin{ejercicio}\label{ej:v3-l26-1}
Si $F:C(\sigma(a))\to A$ es un $*$-homomorfismo que extiende el cálculo holomorfo, pruebe $F=C_a$.
\end{ejercicio}
\begin{solucion}
La extensión implica $F(1)=1$ y $F(z)=a$; preservar el adjunto implica $F(\overline z)=a^*$. Por multiplicatividad y linealidad coincide con $C_a$ en cada polinomio en $z,\overline z$. Estos polinomios contienen las constantes, separan puntos del compacto $\sigma(a)$ y son estables bajo conjugación; Stone--Weierstrass los hace densos. La continuidad de ambas aplicaciones da la igualdad en toda función continua.
\end{solucion}

\begin{ejercicio}\label{ej:v3-l26-2}
Si $A$ es conmutativa, $C^*$ y no unitaria, pruebe que $\Delta(A)$ es localmente compacto Hausdorff y que $\mathcal G:A\to C_0(\Delta(A))$ es un $*$-isomorfismo isométrico.
\end{ejercicio}
\begin{solucion}
Sea $K=\Delta(\widetilde A)$. La proyección escalar $\chi_\infty(a,\lambda)=\lambda$ es un carácter. Cada carácter no nulo $\varphi$ de $A$ se extiende de manera única a $\widetilde\varphi(a,\lambda)=\varphi(a)+\lambda$. Recíprocamente, un carácter de $\widetilde A$ distinto de $\chi_\infty$ se restringe a un carácter no nulo de $A$. Las fórmulas de evaluación prueban que esta biyección identifica la topología débil estrella de $\Delta(A)$ con la de $K\setminus\{\chi_\infty\}$, que es localmente compacta Hausdorff.

La representación unitaria $\widetilde A\simeq C(K)$ lleva $A=\ker\chi_\infty$ al ideal de funciones que se anulan en $\chi_\infty$. Restringirlas a $K\setminus\{\chi_\infty\}$ identifica ese ideal con $C_0$ de dicho espacio: si $f(\chi_\infty)=0$, los conjuntos $\{|f|\ge\eps\}$ son compactos y evitan ese punto; recíprocamente, una función que se anula en infinito se extiende continuamente por cero, porque fuera de un compacto su módulo es menor que cualquier $\eps$. Todos estos mapas preservan productos, adjuntos y normas. El punto $\chi_\infty$ no es aislado: si lo fuera, su complemento sería compacto abierto y el indicador de ese complemento correspondería a una identidad de $A$, contradicción.
\end{solucion}

\begin{ejercicio}\label{ej:v3-l26-3}
Para un espacio compacto Hausdorff $X$, pruebe que $X$ es totalmente disconexo si y sólo si el subespacio generado por las proyecciones de $C(X)$ es denso.
\end{ejercicio}
\begin{solucion}
Una proyección de $C(X)$ toma sólo los valores $0,1$, por lo que es el indicador de un conjunto abierto y cerrado. Recordemos por qué en un compacto Hausdorff totalmente disconexo estos conjuntos separan puntos. Para $x\in X$, sea $Q_x$ la intersección de sus vecindades abiertas y cerradas. Afirmamos que $Q_x$ es conexo. Si se separara en dos compactos no vacíos $P,Q$, elija abiertos disjuntos $U,V$ que los contengan. Cada punto del compacto $X\setminus(U\cup V)$ es excluido por alguna vecindad abierta y cerrada de $x$; una intersección finita de ellas da un conjunto $C$ abierto y cerrado con $Q_x\subset C\subset U\cup V$. Los conjuntos $C\cap U$ y $C\cap V$ son abiertos y cerrados en $X$, y el que contiene a $x$ excluiría una parte de $Q_x$, contradicción. Así $Q_x$ es conexo y, por total desconexión, $Q_x=\{x\}$. Para $y\ne x$ existe entonces un conjunto abierto y cerrado que contiene a $x$ y no a $y$.

Los indicadores de esos conjuntos separan puntos. Su espacio generado es un álgebra (el producto de dos indicadores es el indicador de la intersección), contiene las constantes y es estable bajo conjugación; Stone--Weierstrass da densidad. Recíprocamente, una proyección es constante en cualquier subconjunto conexo, porque su imagen continua está en $\{0,1\}$. Si dos puntos distintos pertenecieran a un mismo subconjunto conexo, todas las combinaciones de proyecciones y sus límites uniformes tendrían iguales valores allí. Esto contradice que las funciones continuas de un compacto Hausdorff separan puntos. Por tanto toda componente conexa es un singleton.
\end{solucion}

\begin{ejercicio}\label{ej:v3-l26-4}
Si $A$ es una álgebra $C^*$ posiblemente no conmutativa y $a\in A$ es normal, construya su cálculo funcional continuo que extiende al holomorfo.
\end{ejercicio}
\begin{solucion}
En el caso unitario la construcción completa está en el teorema~\ref{v3:calculo-continuo}: se trabaja en $C^*(1,a)$, que es conmutativa, y la permanencia espectral asegura que se usa el espectro correcto. Si $A$ no tiene identidad, se realiza la construcción en $\widetilde A$. Para $f(0)=0$, el elemento $f(a)$ pertenece a $A$: la proyección escalar aplicada al cálculo continuo vale $f(0)$, primero para polinomios en $z,\overline z$ y luego por continuidad. Así el cálculo con valores en $A$ tiene como dominio natural $\{f\in C(\sigma_{\widetilde A}(a)):f(0)=0\}$. Las funciones con $f(0)\ne0$ dan en general elementos de la unitización. Esta precisión es necesaria si el enunciado no supone identidad.
\end{solucion}

\begin{ejercicio}\label{ej:v3-l26-5}
Sea $A$ una álgebra $C^*$ conmutativa unitaria generada por $a_1,\ldots,a_n$ (algebraicamente y luego por clausura, o como álgebra $C^*$). Pruebe que existe un compacto $X\subset\C^n$ y un $*$-isomorfismo $A\to C(X)$ que envía $a_k$ a la coordenada $z_k$.
\end{ejercicio}
\begin{solucion}
Defina $\tau:\Delta(A)\to\C^n$ por $\tau(\varphi)=(\varphi(a_1),\ldots,\varphi(a_n))$ y sea $X=\tau(\Delta(A))$. La continuidad de las coordenadas hace a $X$ compacto. Si dos caracteres coinciden en los generadores, coinciden en todos sus polinomios, en los adjuntos cuando éstos se incluyen, y por continuidad en todo $A$. Por ello $\tau$ es inyectiva y es un homeomorfismo sobre $X$. El isomorfismo buscado es $a\mapsto\widehat a\circ\tau^{-1}$. Evaluarlo en un generador da exactamente $z_k$. La hipótesis $C^*$ procede del contexto de la lista y no puede suprimirse: una álgebra de Banach conmutativa general puede tener radical no nulo y no ser isomorfa a un espacio $C(X)$.
\end{solucion}

\begin{ejercicio}\label{ej:v3-l26-6}
Sea $X$ compacto Hausdorff, $x_0\in X$, y sea $A$ el conjunto de sucesiones $(f_n)$ acotadas en $C(X)$ tales que $(f_n(x_0))$ converge. Pruebe que $A$ es una álgebra conmutativa unitaria y determine su espacio de caracteres.
\end{ejercicio}
\begin{solucion}
Con la norma $\norm f=\sup_n\norm{f_n}_\infty$, el producto acotado $B=\prod_{n\ge1}^{\mathrm b}C(X)$ es un álgebra $C^*$ conmutativa unitaria. La aplicación $q:B\to\ell^\infty$, $q(f)=(f_n(x_0))$, es un $*$-homomorfismo contractivo. Como el espacio $c$ de sucesiones convergentes es cerrado en $\ell^\infty$, $A=q^{-1}(c)$ es una subálgebra $C^*$ cerrada y contiene a la identidad. La cerradura de $c$ se verifica directamente: una aproximación uniforme por una sucesión convergente hace que la sucesión límite sea de Cauchy en sus índices.

Para describir todos los caracteres, incluidos los que no son evaluaciones en una coordenada finita, sea $K=\Delta(B)$ y defina $\gamma_n\in K$ por $\gamma_n(f)=f_n(x_0)$. Considere el compacto no vacío
\[
 F=\bigcap_{N\ge1}\overline{\{\gamma_n:n\ge N\}}^{\,K}.
\]
Es no vacío por la propiedad de intersección finita en el compacto $K$. Afirmamos que, bajo $B\simeq C(K)$, la subálgebra $A$ corresponde exactamente a las funciones continuas constantes sobre $F$. Si $f_n(x_0)\to L$, todo punto de adherencia de las colas tiene valor $L$ por continuidad. Recíprocamente, si $\widehat f=L$ en $F$ y la sucesión no convergiera a $L$, habría $\eps>0$ e infinitos índices con $|\gamma_n(f)-L|\ge\eps$. Las clausuras de las colas de esa subsucesión tendrían un punto común en $F$ donde, por cerradura de la desigualdad, el valor se alejaría de $L$ al menos $\eps$, contradicción.

Sea $K/F$ el espacio que identifica todos los puntos de $F$ en un único punto y deja distintos los demás. Es compacto Hausdorff: la relación de equivalencia es el conjunto cerrado diagonal unido a $F\times F$; alternativamente, la normalidad de $K$ separa el cerrado $F$ de cualquier punto exterior y permite construir vecindades disjuntas en el cociente. Las funciones continuas constantes en $F$ son exactamente las funciones que descienden a $K/F$. En consecuencia
\[
 \boxed{\ \Delta(A)\simeq K/F,\qquad K=\Delta\!\left(\prod_{n\ge1}^{\mathrm b}C(X)\right).\ }
\]
Esta descripción es concreta en términos de la compactificación de Stone--Čech: $B=C_b(\N\times X)$, con $\N$ discreto, así $K=\beta(\N\times X)$, y $F$ es el resto en infinito de la copia $\N\times\{x_0\}$. Por tanto se colapsa \emph{sólo ese resto} a un punto. El carácter del punto colapsado es $f\mapsto\lim_n f_n(x_0)$; cada $(n,x)$ da el carácter $f\mapsto f_n(x)$, y pueden existir otros caracteres en el resto no colapsado.

Para justificar la última advertencia, si $x_1\ne x_0$, elija $g\in C(X)$ con $g(x_0)=0$, $g(x_1)=1$. Toda sucesión escalar acotada $(c_n)$ produce $(c_ng)\in A$. Para un ultrafiltro libre $\mathcal U$ sobre $\N$, la fórmula $f\mapsto\lim_{n\to\mathcal U}f_n(x_1)$ es un carácter, porque los límites de ultrafiltro de sucesiones acotadas respetan sumas y productos. Distintos ultrafiltros se distinguen con sucesiones indicadoras y, por tanto, con elementos $(c_ng)$ de $A$. En cambio, si $X=\{x_0\}$, $A=c$ y el cociente se reduce a $\N\cup\{\infty\}$, la compactificación por un punto.
\end{solucion}

La norma ya permite evaluar cualquier función continua del espectro de un elemento normal. En particular, funciones reales no negativas y raíces cuadradas producirán una estructura de orden. Ese paso será la base para estudiar funcionales positivos y construir representaciones por operadores. Como apoyo adicional véanse \cite{maxcstar,maxcalculo,murphy1990,zhu1993}.

\chapter{Positividad, orden y estados}\label{cap:v3-14}

En una álgebra $C^*$, la expresión $a^*a$ desempeña el papel de un cuadrado no negativo. Demostrar que su espectro es realmente no negativo requiere un argumento: no basta trasladar sin prueba la intuición de los números reales. Una vez establecido ese hecho, el orden y los funcionales positivos permitirán observar los elementos del álgebra mediante valores escalares.

\section{Partes positiva y negativa}\label{lec:v3-27}
\fuentelectura{27}{Elementos positivos en álgebras C*}

En este capítulo $A$ es un álgebra $C^*$ unitaria no nula. Para evitar circularidad, distinguiremos al principio entre ser de la forma $b^*b$ y ser autoadjunto con espectro no negativo.

\begin{proposicion}\label{v3:partes-positivas}
Si $a=a^*$, existen únicos autoadjuntos $a_+,a_-$ con espectros contenidos en $[0,\infty)$ tales que
\[
 a=a_+-a_-,\qquad a_+a_-=a_-a_+=0.
\]
Están dados por $a_+=f(a)$ y $a_-=g(a)$, donde $f(t)=\max(t,0)$ y $g(t)=\max(-t,0)$.
\end{proposicion}
\begin{proof}
Las funciones $f,g$ son continuas y reales en el espectro real de $a$. El cálculo continuo conserva conjugación, productos y diferencias, de modo que sus imágenes son autoadjuntas, su diferencia es $a$ y su producto es cero. El mapeo espectral da los espectros no negativos.

Para la unicidad, suponga $a=u-v$, $uv=vu=0$, con $u,v$ autoadjuntos de espectro no negativo. El binomio da $a^n=u^n+(-v)^n$ para $n\ge1$, pues todos los términos mixtos contienen $uv$. Luego $p(a)=p(u)+p(-v)$ para polinomios con término constante cero. Elija un intervalo compacto real que contenga $\sigma(a),\sigma(u),\sigma(-v)$ y cero. Por Weierstrass, aproxime $f$ uniformemente por polinomios $q_n$ y ponga $p_n=q_n-q_n(0)$; entonces $p_n(0)=0$ y $p_n\to f$ uniformemente, ya que $f(0)=0$. Pasar al límite en la identidad polinomial da $a_+=f(u)+f(-v)=u+0$. Así $u=a_+$ y $v=a_-$.
\end{proof}

\begin{lema}\label{v3:cono-espectral}
Sea $P_{\mathrm{esp}}$ el conjunto de autoadjuntos con espectro no negativo. Para $h=h^*$ y $t\ge\norm h$,
\[
 h\in P_{\mathrm{esp}}\quad\Longleftrightarrow\quad\norm{t1-h}\le t.
\]
El conjunto $P_{\mathrm{esp}}$ es un cono convexo cerrado y $P_{\mathrm{esp}}\cap(-P_{\mathrm{esp}})=\{0\}$.
\end{lema}
\begin{proof}
Por normalidad y mapeo espectral, $\norm{t1-h}=\max_{\lambda\in\sigma(h)}|t-\lambda|$. Como $\sigma(h)\subset[-t,t]$, la cota por $t$ equivale a $\lambda\ge0$ para todos sus puntos. Si $h,k\in P_{\mathrm{esp}}$, con $s=\norm h$, $t=\norm k$,
\[
 \norm{(s+t)1-(h+k)}\le\norm{s1-h}+\norm{t1-k}\le s+t.
\]
Como $s+t\ge\norm{h+k}$, la equivalencia prueba que $h+k\in P_{\mathrm{esp}}$. Multiplicar por escalares no negativos conserva el espectro no negativo. Si $h_n\to h$ y todos están en el cono, $h=h^*$; elija $t$ mayor que todas las normas y pase al límite en $\norm{t1-h_n}\le t$. Finalmente, si $h$ y $-h$ están en el cono, $\sigma(h)\subset\{0\}$ y, por normalidad, $\norm h=0$.
\end{proof}

\begin{teorema}[Caracterizaciones de positividad]\label{v3:positividad-equivalencias}
Para $a\in A$ son equivalentes:
\begin{enumerate}[label=(\roman*)]
\item $a=b^*b$ para algún $b\in A$;
\item $a$ es autoadjunto y $\sigma(a)\subset[0,\infty)$;
\item $a=c^2$ para algún $c=c^*$;
\item $a=a^*$ y $\norm{t1-a}\le t$ para todo $t\ge\norm a$;
\item $a=a^*$ y esa desigualdad vale para algún $t\ge\norm a$.
\end{enumerate}
Si se supone de antemano que $a$ es normal, en (ii) basta la condición espectral: el cálculo continuo hace entonces $a=a^*$ cuando su espectro es real.
\end{teorema}
\begin{proof}
El lema anterior prueba (ii)$\Leftrightarrow$(iv)$\Leftrightarrow$(v). Si (ii) vale, el cálculo continuo aplicado a $\sqrt t$ produce un autoadjunto $c$ con $c^2=a$, y (iii) implica (i) porque $c^2=c^*c$.

Falta (i)$\Rightarrow$(ii), que es el paso sustancial. Suponga $a=b^*b$ y use su descomposición $a=a_+-a_-$. Ponga $y=ba_-$. Como $a_+a_-=0$,
\[
 y^*y=a_-aa_-=-a_-^3.
\]
Luego $\sigma(y^*y)\subset(-\infty,0]$. Los espectros de $y^*y$ y $yy^*$ coinciden fuera de cero, así también $\sigma(yy^*)\subset(-\infty,0]$. Escriba $y=u+iv$ con $u,v$ autoadjuntos. El desarrollo de ambos productos da
\[
 y^*y=2u^2+2v^2-yy^*.
\]
Los tres sumandos del miembro derecho pertenecen a $P_{\mathrm{esp}}$: los cuadrados por mapeo espectral y el último por la inclusión recién probada. El lema del cono, que se demostró sin asumir positividad de $b^*b$, implica $y^*y\in P_{\mathrm{esp}}$. Su espectro está por tanto en $\{0\}$; como es autoadjunto, $\norm{y^*y}=0$. Así $a_-^3=0$, y la normalidad de $a_-$ obliga $a_-=0$. Finalmente $a=a_+\in P_{\mathrm{esp}}$.
\end{proof}

\begin{corolario}[Raíz cuadrada positiva]\label{v3:raiz-positiva}
Cada $a\ge0$ tiene una única raíz cuadrada \emph{positiva}, denotada $a^{1/2}$. Está en $C^*(1,a)$ y conmuta con cualquier elemento que conmute con $a$ y $a^*$.
\end{corolario}
\begin{proof}
La existencia proviene de $t\mapsto\sqrt t$ y el mapeo espectral. Si $c\ge0$ y $c^2=a$, el cálculo continuo en $C^*(1,c)$ da $(c^2)^{1/2}=c$, pues $\sqrt{t^2}=t$ en $\sigma(c)\subset[0,\infty)$. El cálculo de $a$ es el mismo en la subálgebra y en $A$ por permanencia espectral, así $c=a^{1/2}$. La última afirmación se prueba aproximando uniformemente la raíz por polinomios en $a$. La unicidad entre \emph{todos} los autoadjuntos sería falsa: $1$ y $-1$ tienen ambos cuadrado $1$.
\end{proof}

\section{Orden de los elementos autoadjuntos}

Definimos $a\le b$ si $b-a\ge0$, para $a,b\in A_{\mathrm{sa}}=\{x:x=x^*\}$. El cono demostrado antes garantiza reflexividad, transitividad y antisimetría. Se tiene además $-\norm a1\le a\le\norm a1$ para todo autoadjunto, por mapeo espectral.

\begin{proposicion}\label{v3:propiedades-orden}
Si $a\le b$, entonces:
\begin{enumerate}[label=(\roman*)]
\item $a+c\le b+c$ para cada autoadjunto $c$;
\item $ta\le tb$ para $t\ge0$;
\item $c^*ac\le c^*bc$ para todo $c\in A$;
\item si $0\le a$, entonces $\norm a\le\norm b$;
\item si $0<a\le b$ y ambos son invertibles, entonces $b^{-1}\le a^{-1}$;
\item si $0\le a\le b$, entonces $a^{1/2}\le b^{1/2}$.
\end{enumerate}
\end{proposicion}
\begin{proof}
(i) y (ii) son las propiedades de traslación y escala del cono. Para (iii), escriba $b-a=d^*d$; entonces $c^*(b-a)c=(dc)^*(dc)\ge0$. Para (iv), $0\le a\le b\le\norm b1$. Por mapeo espectral en $C^*(1,a)$, la positividad de $\norm b1-a$ implica $\sigma(a)\subset[0,\norm b]$, y por normalidad $\norm a\le\norm b$.

Para (v), el cálculo continuo da que las raíces e inversas de elementos positivos invertibles son positivas. Congruencia por $a^{-1/2}$ produce $1\le d=a^{-1/2}ba^{-1/2}$. Su espectro está en $[1,\infty)$, de donde $d^{-1}\le1$. Como $d^{-1}=a^{1/2}b^{-1}a^{1/2}$, otra congruencia por $a^{-1/2}$ da $b^{-1}\le a^{-1}$. La hipótesis de positividad es necesaria: $a=-1\le1=b$, pero $b^{-1}=1\not\le-1=a^{-1}$.

Para (vi) probaremos primero: si $p,q\ge0$ y $p^2\le q^2$, entonces $p\le q$. Fije $t>0$ y escriba
\[
 (t1+q+p)(t1+q-p)=u+iv,
 \quad u=t^21+2tq+q^2-p^2\ge t^21,
\]
con $u,v$ autoadjuntos (son las partes real e imaginaria del producto). El elemento $u$ es positivo invertible. El autoadjunto $s=u^{-1/2}vu^{-1/2}$ tiene espectro real, por lo que $1+is$ es invertible. Así $u+iv=u^{1/2}(1+is)u^{1/2}$ es invertible. En consecuencia $t1+q-p$ tiene inversa izquierda. Un autoadjunto con inversa izquierda $c$ también tiene inversa derecha $c^*$; de $cB=1=Bc^*$ se deduce $c=cBc^*=c^*$, y por tanto es invertible. Hemos probado que $-t\notin\sigma(q-p)$ para todo $t>0$. Como $q-p$ es autoadjunto, su espectro es no negativo y $p\le q$. Aplique esto a $p=a^{1/2}$ y $q=b^{1/2}$.
\end{proof}

\subsection{Normas de homomorfismos y de cocientes}
\begin{proposicion}\label{v3:homomorfismo-cstar-norma}
Todo $*$-homomorfismo $\pi:A\to B$ entre álgebras $C^*$ unitarias es contractivo, aunque no se suponga continuo ni unital. Si es inyectivo, es isométrico. Para todo $a\in A$,
\[
 \norm{\pi(a)}=\inf_{k\in\ker\pi}\norm{a-k}.
\]
\end{proposicion}
\begin{proof}
Sea $p=\pi(1)$. Es una proyección y $\pi(A)\subset pBp$, una álgebra $C^*$ unitaria con identidad $p$. Si $p=0$, $\pi=0$ y las afirmaciones pertinentes son inmediatas. En otro caso, considere $\pi$ como aplicación unital hacia ese rincón. Una inversa se conserva bajo homomorfismos unitales, así $\sigma_{pBp}(\pi(a^*a))\subset\sigma_A(a^*a)$. Por normalidad,
\[
 \norm{\pi(a)}^2=\norm{\pi(a^*a)}
 =r_{pBp}(\pi(a^*a))\le r_A(a^*a)=\norm a^2.
\]
Esto prueba continuidad sin asumirla.

La compatibilidad $\pi(f(h))=f(\pi(h))$ para $h=h^*$ y $f$ continua en un intervalo que contiene ambos espectros se sigue de aproximar $f$ por polinomios y de la contractividad. Si $\pi$ es inyectivo y $\norm{\pi(h)}<\norm h$ para un positivo $h$, elija una función continua $f$ que se anule en $[0,\norm{\pi(h)}]$ y satisfaga $f(\norm h)=1$. Como $\norm h\in\sigma(h)$, $f(h)\ne0$, mientras $\pi(f(h))=f(\pi(h))=0$, contradicción. Aplicar este resultado a $h=a^*a$ demuestra isometría en todo $A$.

Para la fórmula del cociente, la contractividad da $\norm{\pi(a)}\le\norm{a-k}$ para cada $k\in\ker\pi$. Para la otra cota sea $m=\norm{\pi(a)}$, $\eps>0$, $h=a^*a$, y defina en $[0,\norm h]$
\[
 q(t)=\begin{cases}\min(1,(m+\eps)/\sqrt t),&t>0,\\1,&t=0.\end{cases}
\]
Es continua, y $q=1$ en $\sigma(\pi(h))\subset[0,m^2]$. Por tanto $\pi(a-aq(h))=0$. Además
\[
 \norm{aq(h)}^2=\norm{q(h)hq(h)}
 =\max_{t\in\sigma(h)}tq(t)^2\le(m+\eps)^2.
\]
Tomar $k=a-aq(h)$ demuestra que la distancia al núcleo es a lo sumo $m+\eps$. Finalmente deje $\eps\downarrow0$.
\end{proof}

\section{Ejercicios sobre raíces, orden y descomposición polar}

\begin{ejercicio}\label{ej:v3-l27-1}
Para $n\ge1$ y $a\ge0$, pruebe que existe un único $b\ge0$ con $b^n=a$.
\end{ejercicio}


\begin{ejercicio}\label{ej:v3-l27-2}
Pruebe $\norm a=\max(\norm{a_+},\norm{a_-})$ para $a=a^*$.
\end{ejercicio}


\begin{ejercicio}\label{ej:v3-l27-3}
Si $\pi:A\to B$ es un $*$-homomorfismo, pruebe $\norm{\pi(a)}=\inf\{\norm{a-b}:\pi(b)=0\}$.
\end{ejercicio}


\begin{ejercicio}\label{ej:v3-l27-4}
Pruebe que los positivos forman un cono: $a,b\ge0$ y $t\ge0$ implican $a+b,ta\ge0$.
\end{ejercicio}


\begin{ejercicio}\label{ej:v3-l27-5}
Pruebe que todo elemento de $A$ es combinación lineal de elementos unitarios.
\end{ejercicio}


\begin{ejercicio}\label{ej:v3-l27-6}
Para $a$ invertible defina $|a|=(a^*a)^{1/2}$. Pruebe que existe un unitario $u$ con $a=u|a|$.
\end{ejercicio}


\begin{ejercicio}\label{ej:v3-l27-7}
Pruebe que un $*$-homomorfismo preserva el orden.
\end{ejercicio}


\section{Funcionales positivos y estados}\label{lec:v3-28}
\fuentelectura{28}{El espacio de estados de un álgebra C*}

\begin{definicion}
Un funcional complejo lineal $\varphi$ es \emph{hermitiano} si $\varphi(a^*)=\overline{\varphi(a)}$, y \emph{positivo} si $\varphi(a)\ge0$ para cada $a\ge0$. Un \emph{estado} es un funcional positivo con $\varphi(1)=1$. Su conjunto se denota $S(A)$.
\end{definicion}
\index{estado}\index{funcional positivo}

Un funcional es hermitiano si y sólo si toma valores reales en los autoadjuntos: una implicación resulta de $h=h^*$; para la otra, escriba $a=h+ik$ y compare $a^*=h-ik$. Todo funcional positivo es hermitiano, porque cada autoadjunto es diferencia de sus partes positivas. La positividad implica además continuidad; no ocurre lo mismo con la sola condición hermitiana.

\begin{lema}[Cauchy--Schwarz para funcionales positivos]\label{v3:cs-positivo}
Si $\varphi$ es positivo, entonces
\[
 |\varphi(b^*a)|^2\le\varphi(a^*a)\varphi(b^*b).
\]
Además $\varphi$ es continuo y $\norm\varphi=\varphi(1)$.
\end{lema}
\begin{proof}
La forma $[a,b]=\varphi(b^*a)$ es lineal en la primera variable, conjugada lineal en la segunda y hermitiana. Si $[b,b]>0$, expanda $[a-\lambda b,a-\lambda b]\ge0$ y tome $\lambda=[a,b]/[b,b]$; queda $[a,a]-|[a,b]|^2/[b,b]\ge0$. Si $[b,b]=0$, la misma expansión, con un escalar de fase adecuada y módulo arbitrariamente grande, fuerza $[a,b]=0$. Esto demuestra la desigualdad incluso cuando la forma es degenerada.

Como $0\le a^*a\le\norm a^21$, la positividad da $\varphi(a^*a)\le\norm a^2\varphi(1)$. Aplicando Cauchy--Schwarz con $b=1$ se obtiene $|\varphi(a)|\le\varphi(1)\norm a$. Así es continuo, con norma a lo sumo $\varphi(1)$; evaluar en $1$ da la igualdad. Si $\varphi(1)=0$, la misma cota muestra que $\varphi=0$.
\end{proof}

\begin{proposicion}
Para un funcional hermitiano \emph{continuo},
\[
 \norm\varphi=\sup\{\varphi(h):h=h^*,\ \norm h\le1\}.
\]
\end{proposicion}
\begin{proof}
Si $\norm a\le1$, elija $\zeta\in\mathbb T$ con $\varphi(\zeta a)=|\varphi(a)|$. El autoadjunto $h=(\zeta a+(\zeta a)^*)/2$ tiene norma a lo sumo uno y $\varphi(h)=\operatorname{Re}\varphi(\zeta a)=|\varphi(a)|$. Tomar el supremo da una desigualdad; la otra es inmediata por la definición de norma.
\end{proof}

\begin{teorema}[Criterio de positividad]\label{v3:estado-norma}
Un funcional continuo $\varphi$ es positivo si y sólo si $\norm\varphi=\varphi(1)$.
\end{teorema}
\begin{proof}
Una dirección está en el lema anterior. Para la otra, el caso de norma cero es trivial; divida por la norma y suponga $\norm\varphi=\varphi(1)=1$. Si $h=h^*$ y $\varphi(h)=a+ib$, para todo $t\in\R$,
\[
 |1+it\varphi(h)|^2\le\norm{1+ith}^2
 =\norm{1+t^2h^2}\le1+t^2\norm h^2.
\]
El miembro izquierdo es $1-2tb+t^2(a^2+b^2)$. Al dividir por $t$ y dejar $t\to0$ por ambos signos se obtiene $b=0$. Así $\varphi$ es hermitiano. Si $h\ge0$ y $M=\norm h>0$, el espectro de $1-h/M$ está en $[0,1]$, por lo que $|1-\varphi(h)/M|\le1$. Como $\varphi(h)$ es real, esto implica $\varphi(h)\ge0$. El caso $h=0$ es inmediato.
\end{proof}

\begin{teorema}[Estados que realizan valores espectrales]\label{v3:estado-espectral}
Si $a\in A$ y $\lambda\in\sigma(a)$, existe $\varphi\in S(A)$ con $\varphi(a)=\lambda$. En particular $S(A)\ne\varnothing$.
\end{teorema}
\begin{proof}
En $E=\operatorname{span}\{1,a\}$ defina $\varphi_0(\alpha a+\beta1)=\alpha\lambda+\beta$. Está bien definida: si $\alpha a+\beta1=0$ y $\alpha\ne0$, entonces $a=-\beta/\alpha$ y su único valor espectral es ese escalar; si $\alpha=0$, también $\beta=0$. Por mapeo espectral afín, $\alpha\lambda+\beta\in\sigma(\alpha a+\beta1)$ y por tanto $|\varphi_0(\alpha a+\beta1)|\le\norm{\alpha a+\beta1}$. Como $\varphi_0(1)=1$, su norma es uno. Hahn--Banach complejo la extiende a un funcional de norma uno en $A$, todavía con valor uno en $1$. El criterio anterior hace que la extensión sea un estado.
\end{proof}

\begin{proposicion}\label{v3:estados-compactos-separan}
El conjunto $S(A)$ es convexo y compacto Hausdorff en la topología débil estrella. Además:
\[
 \begin{aligned}
 a=0&\ \Longleftrightarrow\ \varphi(a)=0\text{ para todo estado};\\
 a=a^*&\ \Longleftrightarrow\ \varphi(a)\in\R\text{ para todo estado};\\
 a\ge0&\ \Longleftrightarrow\ \varphi(a)\ge0\text{ para todo estado}.
 \end{aligned}
\]
\end{proposicion}
\begin{proof}
Las combinaciones convexas de estados conservan positividad y valor uno en la identidad. Todos tienen norma uno. Dentro de la bola dual cerrada, $S(A)$ está definido por las condiciones cerradas $\varphi(1)=1$ y $\varphi(h)\in[0,\infty)$ para cada $h\ge0$. Banach--Alaoglu da compacidad; la topología débil estrella es Hausdorff.

Si todos los estados anulan $a=h+ik$, con $h,k$ autoadjuntos, sus valores en $h,k$ son reales y necesariamente ambos cero. Por el teorema~\ref{v3:estado-espectral}, $\sigma(h)=\sigma(k)=\{0\}$; la igualdad de norma y radio para autoadjuntos da $h=k=0$. Si todos los valores de $a$ son reales, todos los estados anulan $a-a^*$, y el caso anterior implica autoadjunción. Finalmente, si todos los valores son no negativos, $a$ es autoadjunto y todo punto de su espectro es valor de un estado, de modo que el espectro es no negativo y $a\ge0$. Las implicaciones directas resultan de las definiciones.
\end{proof}

\section{Ejercicios sobre estados y estados puros}

\begin{ejercicio}\label{ej:v3-l28-1}
Examine la afirmación: ``todo funcional lineal hermitiano en un álgebra $C^*$ es continuo''. Dé una demostración si es verdadera o un contraejemplo si es falsa, e indique una hipótesis suficiente válida.
\end{ejercicio}


\begin{ejercicio}\label{ej:v3-l28-2}
Si $\norm\xi=1$ en un espacio de Hilbert $H$, pruebe que $T\mapsto\ip{T\xi}{\xi}$ es un estado de $\Lcal(H)$.
\end{ejercicio}


\begin{ejercicio}\label{ej:v3-l28-3}
Si $\varphi$ es un estado y $u$ unitario, pruebe que $a\mapsto\varphi(uau^*)$ es un estado.
\end{ejercicio}


\begin{ejercicio}\label{ej:v3-l28-4}
Si $a$ es normal, pruebe que existe un estado $\varphi$ con $|\varphi(a)|=\norm a$.
\end{ejercicio}


Los puntos extremos de $S(A)$ se llaman \emph{estados puros}; su conjunto es $P(A)$. El teorema de Krein--Milman, aplicado al espacio dual con topología débil estrella, dice que $S(A)$ es la clausura débil estrella de la envolvente convexa de $P(A)$. Esta aplicación usa la versión para espacios localmente convexos Hausdorff, no la compacidad en norma del dual.

\begin{ejercicio}\label{ej:v3-l28-5}
Sea $K$ un conjunto convexo compacto no vacío en un espacio localmente convexo Hausdorff y $\ell$ un funcional lineal continuo. Pruebe que $\operatorname{Re}\ell$ alcanza su máximo en algún punto extremo de $K$. En particular, vale para un compacto convexo de un espacio normado.
\end{ejercicio}


\begin{ejercicio}\label{ej:v3-l28-6}
Pruebe: (i) $a=0$ si y sólo si todos los estados puros lo anulan; (ii) $a=a^*$ si y sólo si todos los estados puros toman valores reales en $a$; (iii) $a\ge0$ si y sólo si todos los estados puros toman valores no negativos en $a$; (iv) si $a$ es normal, existe un estado puro $\varphi$ con $|\varphi(a)|=\norm a$.
\end{ejercicio}


\begin{ejercicio}\label{ej:v3-l28-7}
Pruebe que, si $A$ es \emph{conmutativa}, $P(A)=\Delta(A)$. Explique por qué la hipótesis de conmutatividad no puede omitirse.
\end{ejercicio}


Los estados detectan cero, autoadjunción y positividad, y los estados puros conservan ese poder de separación. En el siguiente capítulo construiremos, a partir de cada funcional positivo, un espacio de Hilbert donde el álgebra actúa por operadores y el funcional aparece como un producto interno. Véanse también \cite{murphy1990,zhu1993,conway1990}.

\chapter{La construcción GNS y representaciones fieles}\label{cap:v3-15}

Un estado asigna números a los elementos de un álgebra. La construcción de Gelfand--Naimark--Segal recupera, a partir de esos números, un espacio de Hilbert y una acción del álgebra. La idea es considerar los elementos del álgebra como vectores y declarar que su producto interno es el valor del estado en un producto adecuado. Primero se eliminan los vectores de longitud cero y después se completa el espacio.

\section{De un funcional positivo a un espacio de Hilbert}\label{lec:v3-29}
\fuentelectura{29}{La construcción GNS 1}

\begin{definicion}
Una representación de una álgebra $C^*$ $A$ es un $*$-homomorfismo $\pi:A\to\Lcal(H)$, donde $H$ es un espacio de Hilbert. Es \emph{fiel} si es inyectiva. Es \emph{cíclica} si existe $e\in H$ con $\overline{\pi(A)e}=H$.
\end{definicion}
\index{representación!cíclica}\index{construcción GNS}

Para un álgebra unitaria, una representación cíclica en un espacio no nulo es automáticamente unital: $\pi(1)$ actúa como identidad en el conjunto denso $\pi(A)e$ y por continuidad en todo $H$. Una representación general puede ser degenerada; lo señalaremos expresamente al estudiar sumas directas.

\begin{ejemplo}
Un carácter es una representación en $\C$. Si $X$ es compacto Hausdorff y $\mu$ es una medida de Borel positiva finita y regular, la multiplicación
\[
 M_f:L^2(\mu)\to L^2(\mu),\qquad M_fg=fg,
\]
define una representación de $C(X)$. En efecto $\norm{fg}_2\le\norm f_\infty\norm g_2$, $M_{fg}=M_fM_g$, $M_1=I$, y
$\ip{M_fg}{h}=\int fg\overline h\,d\mu=\ip g{M_{\overline f}h}$; luego $M_f^*=M_{\overline f}$. La conjugación es necesaria en la fórmula del adjunto.
\end{ejemplo}

Una forma sesquilineal no negativa $[\cdot,\cdot]$ es lineal en la primera variable, conjugada lineal en la segunda, satisface $[v,w]=\overline{[w,v]}$ y $[v,v]\ge0$. La demostración de Cauchy--Schwarz del lema~\ref{v3:cs-positivo} sólo utiliza estas propiedades, así vale también si existen vectores no nulos de longitud cero. En particular, $[v,v]=0$ implica $[v,w]=0$ para todo $w$.

\begin{teorema}[Gelfand--Naimark--Segal]\label{v3:gns}
Sea $\varphi$ un funcional positivo en una álgebra $C^*$ unitaria $A$. Existen un espacio de Hilbert $H_\varphi$, una representación cíclica $\pi_\varphi$ y un vector cíclico $e_\varphi$ tales que
\[
 \varphi(a)=\ip{\pi_\varphi(a)e_\varphi}{e_\varphi},
 \qquad \norm{e_\varphi}^2=\varphi(1),\qquad
 \norm{\pi_\varphi(a)}\le\norm a.
\]
\end{teorema}
\begin{proof}
La forma $[a,b]=\varphi(b^*a)$ es sesquilineal, hermitiana y no negativa por las propiedades de un funcional positivo. Defina
\[
 N_\varphi=\{a\in A:\varphi(a^*a)=0\}.
\]
Por Cauchy--Schwarz, todo elemento de $N_\varphi$ es ortogonal para esta forma a todos los elementos de $A$. Esto prueba que $N_\varphi$ es un subespacio: al expandir el cuadrado de una combinación lineal de dos elementos nulos, todos los términos son cero. Es cerrado porque $a\mapsto\varphi(a^*a)$ es continuo.

Además es ideal izquierdo. En efecto $c^*c\le\norm c^21$ implica
\[
 0\le a^*c^*ca\le\norm c^2a^*a.
\]
Si $a\in N_\varphi$, aplicar $\varphi$ da $\varphi((ca)^*(ca))=0$. En general no es ideal bilateral, por lo que el cociente siguiente es un espacio vectorial, no un álgebra cociente.

En $V_\varphi=A/N_\varphi$ defina
\[
 \ip{[a]}{[b]}_\varphi=\varphi(b^*a).
\]
Si se cambian los representantes, la diferencia es suma de términos con un factor nulo en una de las variables de la forma, y Cauchy--Schwarz los anula. Por ello está bien definido. Es positivo definido: $\norm{[a]}_\varphi^2=0$ equivale exactamente a $a\in N_\varphi$, es decir, $[a]=0$. Complete este espacio prehilbertiano para obtener $H_\varphi$; el producto interno se extiende por continuidad, pues Cauchy--Schwarz controla los límites de ambas variables.

Para $c\in A$ defina inicialmente $\pi_\varphi(c)[a]=[ca]$. Está bien definido porque $N_\varphi$ es ideal izquierdo. La desigualdad de orden anterior da
\[
 \norm{\pi_\varphi(c)[a]}_\varphi^2
 =\varphi(a^*c^*ca)\le\norm c^2\varphi(a^*a)
 =\norm c^2\norm{[a]}_\varphi^2.
\]
Así se extiende de forma única a un operador acotado de $H_\varphi$. Las identidades de linealidad y multiplicación se comprueban en las clases: $\pi_\varphi(cd)[a]=[cda]=\pi_\varphi(c)\pi_\varphi(d)[a]$, y se extienden por densidad. Para el adjunto,
\[
 \ip{\pi_\varphi(c)[a]}{[b]}_\varphi
 =\varphi(b^*ca)
 =\ip{[a]}{\pi_\varphi(c^*)[b]}_\varphi.
\]
Por densidad, $\pi_\varphi(c)^*=\pi_\varphi(c^*)$. Finalmente tome $e_\varphi=[1]$. Entonces $\pi_\varphi(A)e_\varphi=V_\varphi$ es denso, su norma cuadrada es $\varphi(1)$ y $\ip{\pi_\varphi(a)e_\varphi}{e_\varphi}=\varphi(a)$. Si $\varphi=0$, el espacio construido es el espacio cero, caso que también satisface las fórmulas.
\end{proof}

\begin{teorema}[Unicidad de la construcción]\label{v3:unicidad-gns}
Si $(\pi,H,e)$ es cíclica y $\varphi(a)=\ip{\pi(a)e}{e}$, existe un único unitario $U:H_\varphi\to H$ con $Ue_\varphi=e$ y $U\pi_\varphi(a)=\pi(a)U$ para todo $a$.
\end{teorema}
\begin{proof}
Defina $U_0[a]=\pi(a)e$. La identidad
\[
 \norm{\pi(a)e}^2=\ip{\pi(a^*a)e}{e}
 =\varphi(a^*a)=\norm{[a]}_\varphi^2
\]
prueba simultáneamente que está bien definido y que es isométrico. Su imagen es $\pi(A)e$, densa en $H$. La extensión isométrica a la completación tiene imagen cerrada (por completitud) y densa, luego es sobreyectiva y unitaria. En las clases, $U_0\pi_\varphi(c)[a]=\pi(ca)e=\pi(c)U_0[a]$; la continuidad da la identidad en todo el espacio. Para una representación cíclica $\pi(1)=I$, de modo que $Ue_\varphi=e$. Cualquier unitario con estas propiedades coincide en el conjunto denso de vectores $\pi_\varphi(a)e_\varphi$, lo que prueba unicidad.
\end{proof}

\section{Ejercicios sobre representaciones cíclicas}

\begin{ejercicio}\label{ej:v3-l29-1}
Si $\alpha>0$, pruebe que $\pi_\varphi$ y $\pi_{\alpha\varphi}$ son unitariamente equivalentes. Deduzca la reducción a estados para funcionales positivos no nulos.
\end{ejercicio}


\begin{ejercicio}\label{ej:v3-l29-2}
Pruebe que $(\pi_1\oplus\pi_2)(a)(\xi_1,\xi_2)=(\pi_1(a)\xi_1,\pi_2(a)\xi_2)$ es una representación. Si ambas representaciones son cíclicas, determine si su suma debe serlo.
\end{ejercicio}


\begin{ejercicio}\label{ej:v3-l29-3}
Sea $H$ separable de dimensión infinita, $(e_n)$ una base ortonormal y $\operatorname{BL}$ un límite de Banach positivo normalizado. Defina
\[
 \varphi(T)=\operatorname{BL}\bigl((\ip{Te_n}{e_n})_{n\ge1}\bigr).
\]
Pruebe: (i) $\varphi$ es un estado; (ii) $\ker\pi_\varphi=\Kcal(H)$.
\end{ejercicio}


\section{Sumas directas y representación universal}\label{lec:v3-30}
\fuentelectura{30}{La construcción GNS 2}

Para una familia $(H_j)_{j\in J}$, su suma hilbertiana es
\[
 \bigoplus_{j\in J}H_j
 =\{(\xi_j):\sum_j\norm{\xi_j}^2<\infty\},
 \quad\ip\xi\eta=\sum_j\ip{\xi_j}{\eta_j}.
\]
La suma de cantidades no negativas se define como el supremo de las sumas finitas. Una familia de cuadrados sumable tiene soporte a lo sumo numerable: para cada $n$ sólo hay finitos términos al menos $1/n$, y la unión de esos conjuntos contiene todo el soporte no nulo. Cauchy--Schwarz para sumas finitas y paso al límite dan convergencia absoluta del producto interno. Este espacio es la completación de las familias de soporte finito: truncar una familia sumable aproxima en norma, y una sucesión de Cauchy se resuelve coordenada a coordenada con control uniforme de las colas.

\begin{proposicion}
Si $\pi_j:A\to\Lcal(H_j)$ son representaciones, los operadores diagonales $\pi(a)(\xi_j)=(\pi_j(a)\xi_j)$ definen una representación $\pi=\bigoplus_j\pi_j$, y
\[
 \norm{\pi(a)}=\sup_j\norm{\pi_j(a)}\le\norm a.
\]
\end{proposicion}
\begin{proof}
La contractividad de cada $*$-homomorfismo da
$\sum_j\norm{\pi_j(a)\xi_j}^2\le\norm a^2\sum_j\norm{\xi_j}^2$, así el operador está bien definido y es acotado. Las identidades de linealidad y producto valen en cada coordenada. La identidad de adjuntos se obtiene sumando $\ip{\pi_j(a)\xi_j}{\eta_j}=\ip{\xi_j}{\pi_j(a^*)\eta_j}$, con series absolutamente convergentes. La igualdad de normas se prueba como en la suma de dos representaciones, ensayando vectores apoyados en una sola coordenada para obtener la cota inferior.
\end{proof}

\begin{teorema}[Descomposición en representaciones cíclicas]\label{v3:descomposicion-ciclica}
Toda representación unital de $A$ es unitariamente equivalente a una suma hilbertiana de representaciones cíclicas. Para una representación no unital, la afirmación se aplica a $\pi(1)H$ y queda además una representación cero en $(I-\pi(1))H$.
\end{teorema}
\begin{proof}
Si $\xi\ne0$, el espacio $H_\xi=\overline{\pi(A)\xi}$ es invariante bajo cada $\pi(a)$ y bajo su adjunto $\pi(a^*)$. Por ello su complemento ortogonal también es invariante: si $\eta\perp H_\xi$ y $v\in H_\xi$, $\ip{\pi(a)\eta}{v}=\ip\eta{\pi(a^*)v}=0$. Se dice que $H_\xi$ es \emph{reductor}.

Considere familias de vectores no nulos cuyos espacios cíclicos son mutuamente ortogonales para vectores distintos, ordenadas por inclusión. La unión de una cadena sigue teniendo esa propiedad, porque cualquier par de sus vectores aparece junto en un miembro de la cadena. Zorn produce una familia maximal $(\xi_j)$. Sea $K=\overline{\sum_jH_{\xi_j}}$. Si $K^\perp\ne0$, un vector $\eta\ne0$ allí generaría un espacio cíclico contenido en $K^\perp$, pues $K$ es reductor. Como $\pi(1)=I$, ese espacio contiene a $\eta$ y es no nulo, contradiciendo maximalidad. Así $K=H$.

La suma de vectores define una isometría de la suma algebraica de los $H_{\xi_j}$ sobre un subespacio denso de $H$, y se extiende a un unitario de la suma hilbertiana. Las restricciones $\pi_j(a)=\pi(a)|_{H_{\xi_j}}$ son representaciones porque los subespacios son reductores; el unitario entrelaza la suma con $\pi$. Para el caso general, $p=\pi(1)$ es una proyección y $\pi(a)=p\pi(a)p$. La restricción a $pH$ es unital y la de $(I-p)H$ es cero. Omitir ese sumando haría falso el enunciado para una representación cero en un espacio no nulo.
\end{proof}

\begin{teorema}[Gelfand--Naimark: representación fiel]\label{v3:gelfand-naimark}
Toda álgebra $C^*$ unitaria admite una representación isométrica. Si es separable, puede representarse isométricamente en un espacio de Hilbert separable.
\end{teorema}
\begin{proof}
Tome un subconjunto débil estrella denso $F\subset S(A)$; se puede tomar todo $S(A)$. Defina $\pi=\bigoplus_{\varphi\in F}\pi_\varphi$. La contractividad da $\norm{\pi(a)}\le\norm a$. Como $\norm{e_\varphi}=1$,
\[
 \norm{\pi(a)}^2\ge\norm{\pi_\varphi(a)}^2
 \ge\norm{\pi_\varphi(a)e_\varphi}^2=\varphi(a^*a).
\]
La segunda desigualdad tiene este sentido, inverso al que aparece en una línea de la lectura original. Un estado realiza el máximo espectral del positivo $a^*a$, de modo que $\sup_{\psi\in S(A)}\psi(a^*a)=\norm{a^*a}=\norm a^2$. La continuidad débil estrella de la evaluación y la densidad de $F$ hacen que el mismo supremo se obtenga sobre $F$. Así $\norm{\pi(a)}\ge\norm a$, y se obtiene igualdad.

Si $A$ es separable, la bola dual y su cerrado $S(A)$ son compactos metrizables en la topología débil estrella, con la métrica construida mediante una sucesión densa de elementos de $A$. Un compacto métrico es separable: elija, para cada $n$, una red finita de radio $1/n$ y reúna sus centros. Por tanto $F$ puede ser numerable. Cada $H_\varphi$ es separable porque el mapa continuo $a\mapsto[a]$ lleva un subconjunto numerable denso de $A$ a uno denso en $V_\varphi$, y éste es denso en $H_\varphi$. La suma \emph{numerable} de espacios de Hilbert separables es separable: use combinaciones finitas con coeficientes racionales complejos de sus respectivos conjuntos densos.
\end{proof}

Todo Hilbert separable de dimensión infinita es unitariamente isomorfo a $\ell^2$ mediante sus coordenadas respecto de una base ortonormal numerable. Si la representación anterior tiene dimensión finita, puede repetirse una cantidad numerable de veces para obtener una representación isométrica en un Hilbert isomorfo a $\ell^2$. Así toda álgebra $C^*$ unitaria separable no nula puede realizarse como una subálgebra de operadores de $\ell^2$.

\begin{ejemplo}[La construcción para una medida]
Para un estado $\varphi(f)=\int_Xf\,d\mu$ de $C(X)$, con $\mu$ probabilidad regular de Borel,
\[
 N_\varphi=\{f\in C(X):f=0\ \mu\text{-casi en todas partes}\},
 \qquad\norm{[f]}_\varphi^2=\int|f|^2\,d\mu.
\]
La completación es $L^2(\mu)$ y $\pi_\varphi(f)=M_f$, con vector cíclico $1$. Para justificar la densidad utilizada, las funciones simples son densas en $L^2$; por regularidad, cada indicador de Borel se aproxima en $L^2$ por una función continua entre cero y uno: elija compacto $K\subset E\subset U$ abierto con $\mu(U\setminus K)$ pequeño, y una función continua igual a uno en $K$ y cero fuera de $U$. Su error cuadrático está acotado por $\mu(U\setminus K)$. La representación universal de $C(X)$ reúne estas multiplicaciones para todos sus estados, es decir, para todas sus probabilidades regulares.
\end{ejemplo}

\section{Ejercicios sobre medidas y fidelidad}

\begin{ejercicio}\label{ej:v3-l30-1}
Sean $\mu,\nu$ medidas positivas finitas regulares de Borel sobre un compacto Hausdorff $X$. Pruebe: (i) $\pi_\mu\oplus\pi_\nu$ es cíclica si y sólo si $\mu\perp\nu$; (ii) si son mutuamente singulares, $\pi_\mu\oplus\pi_\nu$ es unitariamente equivalente a $\pi_{\mu+\nu}$.
\end{ejercicio}


\begin{ejercicio}\label{ej:v3-l30-2}
Si $A$ es separable y $(\varphi_n)$ es débil estrella densa en $S(A)$, defina $\varphi=\sum_{n\ge1}2^{-n}\varphi_n$. Pruebe que $\pi_\varphi$ es isométrica y deduzca que $A$ admite una representación cíclica isométrica.
\end{ejercicio}


\begin{ejercicio}\label{ej:v3-l30-3}
Si $\pi_\varphi(a)=0$ para todos los estados puros, pruebe $a=0$. Deduzca que la suma $\bigoplus_{\varphi\in P(A)}\pi_\varphi$ es isométrica.
\end{ejercicio}


Ya podemos estudiar cualquier álgebra $C^*$ concreta o abstracta mediante operadores en un espacio de Hilbert. Para reconocer qué propiedades dependen de cerrar una familia de operadores en norma y cuáles dependen de su acción sobre vectores, introduciremos ahora las topologías fuerte y débil del operador. Las construcciones de este capítulo se desarrollan también en \cite{murphy1990,zhu1993,conway1990}.

\parte{Álgebras de von Neumann}{Las topologías de operadores y los conmutantes permiten pasar del control en norma a cierres más amplios. Se desarrollan el doble conmutante, las medidas espectrales, el cálculo boreliano y el modelo conmutativo mediante funciones esencialmente acotadas.}
\chapter{Topologías de operadores y doble conmutante}\label{cap:v3-16}

La convergencia en norma controla simultáneamente todos los vectores unitarios. A menudo sólo se necesita controlar cada vector por separado, o incluso cada coeficiente de un producto interno. Estas exigencias definen topologías diferentes. El teorema del doble conmutante mostrará que el cierre en una de ellas puede reconocerse mediante una condición puramente algebraica de conmutación.

\section{Convergencia fuerte y débil del operador}\label{lec:v3-31}
\fuentelectura{31}{Las topologías del operador}

Para $\xi,\eta\in H$, las funciones $p_\xi(T)=\norm{T\xi}$ y $q_{\xi,\eta}(T)=|\ip{T\xi}{\eta}|$ son seminormas: la homogeneidad absoluta y la desigualdad triangular se siguen de las de la norma y del módulo. Ambas familias separan operadores, porque $T\xi\ne0$ puede detectarse tomando $\eta=T\xi$.

\begin{definicion}
La \emph{topología fuerte del operador} (TFO, o SOT) es la generada por las seminormas $p_\xi$; la \emph{topología débil del operador} (TDO, o WOT) es la generada por $q_{\xi,\eta}$. Para una red,
\[
 T_\alpha\xrightarrow{\mathrm{SOT}}T\iff T_\alpha\xi\to T\xi\quad(\xi\in H),
\]
\[
 T_\alpha\xrightarrow{\mathrm{WOT}}T\iff
 \ip{T_\alpha\xi}{\eta}\to\ip{T\xi}{\eta}\quad(\xi,\eta\in H).
\]
\end{definicion}
\index{topología!fuerte del operador}\index{topología!débil del operador}

Las desigualdades $\norm{(T_\alpha-T)\xi}\le\norm{T_\alpha-T}\norm\xi$ y $|\ip{(T_\alpha-T)\xi}{\eta}|\le\norm{(T_\alpha-T)\xi}\norm\eta$ dan
\[
 \text{convergencia en norma}\ \Longrightarrow\ \text{SOT}\ \Longrightarrow\ \text{WOT}.
\]
WOT es convergencia puntual débil porque, por Riesz, todos los funcionales continuos de $H$ son productos internos con un vector fijo. En dimensión finita, probar los coeficientes de una matriz basta para controlar su norma y las tres topologías coinciden. Los ejercicios con el desplazamiento mostrarán que en dimensión infinita las implicaciones son estrictas.

\begin{proposicion}\label{v3:operaciones-topologias}
La suma y la multiplicación por escalares son continuas en SOT y WOT. El producto de operadores es continuo en cada variable por separado en ambas topologías. La adjunción es continua en WOT. El producto es conjuntamente continuo en SOT para redes cuyos factores izquierdos están uniformemente acotados en norma.
\end{proposicion}
\begin{proof}
Las desigualdades de las seminormas prueban la continuidad de la suma. Para escalares $\lambda_\alpha\to\lambda$ y operadores $T_\alpha\to T$, escriba
$\lambda_\alpha T_\alpha-\lambda T=\lambda_\alpha(T_\alpha-T)+(\lambda_\alpha-\lambda)T$ y use que los escalares están eventualmente acotados. Para un operador fijo $S$, en SOT se tiene $\norm{S(T_\alpha-T)\xi}\le\norm S\norm{(T_\alpha-T)\xi}$ y $(T_\alpha-T)S\xi\to0$. En WOT se usan $\ip{S(T_\alpha-T)\xi}{\eta}=\ip{(T_\alpha-T)\xi}{S^*\eta}$ y la misma fórmula con $S\xi$. La identidad $\ip{T_\alpha^*\xi}{\eta}=\overline{\ip{T_\alpha\eta}{\xi}}$ prueba la continuidad del adjunto.

Finalmente, si $S_\alpha\to S$ y $T_\alpha\to T$ fuertemente, con $\norm{S_\alpha}\le M$, entonces
\[
 \norm{(S_\alpha T_\alpha-ST)\xi}
 \le M\norm{(T_\alpha-T)\xi}+\norm{(S_\alpha-S)T\xi}\to0.
\]
Una sucesión fuertemente convergente está uniformemente acotada por Banach--Steinhaus. Para redes arbitrarias no se debe suponer sin prueba una cota uniforme de todas sus normas; por eso se explicita aquí la hipótesis.
\end{proof}

\begin{teorema}[Dual continuo común]\label{v3:dual-sot-wot}
Para un funcional lineal $\Phi:\Lcal(H)\to\C$ son equivalentes: ser WOT-continuo; ser SOT-continuo; y admitir una expresión
\[
 \Phi(T)=\sum_{k=1}^n\ip{T\xi_k}{\eta_k}
\]
con un número finito de vectores $\xi_k,\eta_k\in H$.
\end{teorema}
\begin{proof}
Una suma finita de coeficientes es WOT-continua, y WOT-continuidad implica SOT-continuidad porque SOT es más fina. Suponga ahora continuidad fuerte. Existe una vecindad básica de cero, determinada por $\xi_1,\ldots,\xi_n$ y un $\delta>0$, en la que $|\Phi|<1$. Por homogeneidad se obtiene
\[
 |\Phi(T)|\le C\max_k\norm{T\xi_k}
 \le C\left(\sum_k\norm{T\xi_k}^2\right)^{1/2}
\]
para algún $C$. Si el máximo es cero, todos los múltiplos de $T$ están en la vecindad y necesariamente $\Phi(T)=0$; si es positivo, se escala $T$ para entrar en ella. Así se cubren ambos casos de la desigualdad.

En el subespacio $E=\{(T\xi_1,\ldots,T\xi_n):T\in\Lcal(H)\}\subset H^n$, defina $\ell(T\xi_1,\ldots,T\xi_n)=\Phi(T)$. La cota asegura que está bien definida y es continua. Se extiende a $\overline E$ por completitud de $\C$. Riesz representa esa extensión mediante un vector $(\eta_1,\ldots,\eta_n)\in\overline E\subset H^n$, y produce la expresión requerida.
\end{proof}

\section{Ejercicios sobre convergencia de operadores}

\begin{ejercicio}\label{ej:v3-l31-1}
Pruebe que $D(F,\eps)=\{T:\norm{T\xi}<\eps\text{ para todo }\xi\in F\}$, con $F\subset H$ finito y $\eps>0$, forma una base de vecindades fuertes de cero.
\end{ejercicio}
\begin{solucion}
Por definición de la topología generada por seminormas, una vecindad básica es una intersección finita $\bigcap_{j=1}^n\{T:p_{\xi_j}(T)<\eps_j\}$. El conjunto $D(\{\xi_1,\ldots,\xi_n\},\min_j\eps_j)$ está contenido en ella. A su vez, cada $D(F,\eps)$ es una intersección finita de esos abiertos. Por tanto ambas familias generan las mismas vecindades. Las vecindades en $T_0$ se obtienen por traslación.
\end{solucion}

\begin{ejercicio}\label{ej:v3-l31-2}
Si $T_\alpha^*T_\alpha\to0$ en WOT, pruebe $T_\alpha\to0$ en SOT. Si además $\sup_\alpha\norm{T_\alpha}<\infty$, pruebe $T_\alpha^*T_\alpha\to0$ en SOT.
\end{ejercicio}
\begin{solucion}
Para cada $\xi$, $\norm{T_\alpha\xi}^2=\ip{T_\alpha^*T_\alpha\xi}{\xi}\to0$. Bajo la cota uniforme $M$, se tiene $\norm{T_\alpha^*T_\alpha\xi}\le M\norm{T_\alpha\xi}\to0$, que prueba la segunda afirmación.
\end{solucion}

\begin{ejercicio}\label{ej:v3-l31-3}
Sea $L(x_1,x_2,\ldots)=(0,x_1,x_2,\ldots)$ en $\ell^2$. Pruebe: (i) $L^n\to0$ en WOT pero no en SOT; (ii) $(L^*)^n\to0$ en SOT pero no en norma.
\end{ejercicio}
\begin{solucion}
(ii) Primero, $\norm{(L^*)^nx}^2=\sum_{k>n}|x_k|^2\to0$. Sin embargo, $(L^*)^ne_{n+1}=e_1$, así cada operador tiene norma uno. (i) La identidad $\ip{L^nx}{y}=\ip x{(L^*)^ny}$ y Cauchy--Schwarz dan convergencia débil a cero. Como $L$ es isométrico, $\norm{L^nx}=\norm x$ para todo $n$; ningún vector no nulo converge a cero en norma bajo esas potencias. Este ejemplo también muestra que la adjunción no es SOT-continua: $(L^*)^n\to0$ fuertemente pero sus adjuntos $L^n$ no.
\end{solucion}

\begin{ejercicio}\label{ej:v3-l31-4}
Pruebe que la bola unitaria cerrada de $\Lcal(H)$ es WOT-compacta.
\end{ejercicio}
\begin{solucion}
Para cada par $\xi,\eta$, sea $D_{\xi,\eta}=\{z:|z|\le\norm\xi\norm\eta\}$. Tychonoff hace compacto el producto de todos estos discos. La aplicación $T\mapsto(\ip{T\xi}{\eta})_{\xi,\eta}$ identifica la bola con el conjunto de familias $b(\xi,\eta)$ que son sesquilineales y satisfacen $|b(\xi,\eta)|\le\norm\xi\norm\eta$. Las identidades de sesquilinealidad son condiciones cerradas coordenada a coordenada, de modo que ese conjunto es cerrado en el producto. Recíprocamente, Riesz aplicado a $\eta\mapsto\overline{b(\xi,\eta)}$, o a la forma sesquilineal acotada, produce un único operador $T$ con norma a lo sumo uno y esos coeficientes; la linealidad en $\xi$ asegura que el operador es lineal. La topología inducida del producto es exactamente WOT. Así la bola es compacta.
\end{solucion}

\begin{ejercicio}\label{ej:v3-l31-5}
Si $C\subset\Lcal(H)$ es convexo, pruebe que sus clausuras SOT y WOT coinciden.
\end{ejercicio}
\begin{solucion}
La clausura SOT está contenida en la WOT porque SOT es más fina. Si $T$ no pertenece a la clausura SOT de $C$, existen $\xi_1,\ldots,\xi_n$ y $\eps>0$ tales que ningún $S\in C$ satisface simultáneamente $\norm{(S-T)\xi_j}<\eps$. Defina $J(S)=(S\xi_1,\ldots,S\xi_n)\in H^n$. Entonces $J(T)$ está fuera de la clausura normada del conjunto convexo $J(C)$, pues su distancia a cada punto de éste es al menos $\eps$. La separación de Hahn--Banach en el espacio normado $H^n$ proporciona un funcional real continuo que separa estrictamente $J(T)$ de esa clausura. Es la parte real de un funcional complejo lineal continuo; por Riesz, éste tiene la forma $(v_1,\ldots,v_n)\mapsto\sum_j\ip{v_j}{\eta_j}$. Por tanto el separador compuesto con $J$ es la parte real de $\sum_j\ip{S\xi_j}{\eta_j}$, que es WOT-continuo. Su semiespacio abierto separa $T$ de $C$, así $T$ tampoco pertenece a la clausura WOT. Esto prueba la inclusión inversa usando sólo separación en espacios normados. El argumento funciona también para convexos reales, como el conjunto de autoadjuntos.
\end{solucion}

\begin{ejercicio}\label{ej:v3-l31-6}
Examine la afirmación: si $T_\alpha$ son normales y $T_\alpha\to T$ fuertemente, entonces $T_\alpha^*\to T^*$ fuertemente. Pruebe la versión correcta suponiendo además que $T$ es normal.
\end{ejercicio}
\begin{solucion}
La afirmación sin la hipótesis adicional es falsa. En $\ell^2$, sea $U_n$ el unitario que permuta cíclicamente $e_1,\ldots,e_n$, enviando $e_k$ a $e_{k+1}$ para $k<n$ y $e_n$ a $e_1$, y fija los demás vectores de la base. Para cada vector de soporte finito, $U_nx=Lx$ cuando $n$ es grande; la cota uniforme $\norm{U_n}=1$ extiende esta convergencia a todo $\ell^2$. Así $U_n\to L$ fuertemente. Pero $U_n^*e_1=e_n$ no converge en norma a $L^*e_1=0$.

Si $T$ también es normal, WOT-continuidad de la adjunción da $T_\alpha^*\xi\rightharpoonup T^*\xi$. La normalidad implica $\norm{T_\alpha^*\xi}=\norm{T_\alpha\xi}\to\norm{T\xi}=\norm{T^*\xi}$. En un Hilbert, convergencia débil más convergencia de normas implica convergencia en norma, como se ve al expandir
$\norm{T_\alpha^*\xi-T^*\xi}^2$. Esto prueba la versión corregida, también para redes. El desplazamiento muestra precisamente por qué no puede deducirse que el límite fuerte de normales sea normal.
\end{solucion}

\begin{ejercicio}\label{ej:v3-l31-7}
Sean $T_\alpha,T$ normales, $T_\alpha\to T$ en SOT, y $\norm{T_\alpha},\norm T\le r$. Si $f:\C\to\C$ es continua, pruebe $f(T_\alpha)\to f(T)$ en SOT.
\end{ejercicio}
\begin{solucion}
Por el ejercicio anterior, los adjuntos también convergen fuertemente. La continuidad conjunta del producto para familias uniformemente acotadas prueba por inducción la convergencia fuerte de cada polinomio $p(T_\alpha,T_\alpha^*)$ hacia $p(T,T^*)$. Stone--Weierstrass aproxima $f$ uniformemente sobre el disco compacto $|z|\le r$ por un polinomio en $z,\overline z$. Para un vector $\xi$,
\[
 \norm{(f(T_\alpha)-f(T))\xi}
 \le2\norm{f-p}_{\infty,\,|z|\le r}\norm\xi
 +\norm{(p(T_\alpha,T_\alpha^*)-p(T,T^*))\xi}.
\]
Primero haga pequeño el primer término eligiendo $p$ y después el segundo haciendo avanzar la red. La normalidad del límite es necesaria para que $f(T)$ tenga este significado y para la convergencia de adjuntos; no estaba incluida explícitamente en la lectura original.
\end{solucion}

\section{Redes monótonas de operadores}\label{lec:v3-32}
\fuentelectura{32}{El teorema del doble conmutador}

Recordemos que $T\ge0$ significa $\ip{T\xi}{\xi}\ge0$ para todo $\xi$, con signo \emph{no negativo}. Una red autoadjunta es creciente si $\alpha\le\beta$ implica $T_\alpha\le T_\beta$, y está acotada superiormente si existe un autoadjunto $S$ con $T_\alpha\le S$ para todo $\alpha$.

\begin{teorema}[Vigier]\label{v3:vigier}
Una red creciente de operadores autoadjuntos converge en SOT si y sólo si está acotada superiormente. Su límite es su supremo en el orden de los autoadjuntos.
\end{teorema}
\begin{proof}
Si $T_\alpha\to T$ fuertemente, para cada $\alpha$ y $\beta\ge\alpha$ se tiene $\ip{(T_\beta-T_\alpha)\xi}{\xi}\ge0$. Pasar al límite en $\beta$ da $T_\alpha\le T$. El límite es autoadjunto porque los coeficientes conservan la simetría hermitiana.

Recíprocamente, fije $\alpha_0$ y trabaje en la cola $\alpha\ge\alpha_0$. Al restar $T_{\alpha_0}$ podemos suponer $0\le T_\alpha\le S$ con $S\ge0$. Entonces $\norm{T_\alpha}\le M=\norm S$. Para cada $\xi$, los números $\ip{T_\alpha\xi}{\xi}$ crecen y están acotados, por lo que convergen. La identidad de polarización, con nuestra convención lineal en la primera variable,
\[
 \ip{T_\alpha\xi}{\eta}
 =\frac14\sum_{k=0}^3 i^k\ip{T_\alpha(\xi+i^k\eta)}{\xi+i^k\eta},
\]
proporciona un límite $b(\xi,\eta)$. Las identidades sesquilineales y la cota $|b(\xi,\eta)|\le M\norm\xi\norm\eta$ pasan al límite, así Riesz define un operador $T$ autoadjunto con esos coeficientes y $\norm T\le M$. Además $0\le T-T_\alpha\le M1$. Por cálculo continuo, para $D=T-T_\alpha$ se tiene $D^2\le MD$; luego
\[
 \norm{(T-T_\alpha)\xi}^2\le M\ip{(T-T_\alpha)\xi}{\xi}\longrightarrow0.
\]
Así la convergencia es fuerte. Cualquier otra cota superior $R$ satisface $\ip{(R-T)\xi}{\xi}=\lim_\alpha\ip{(R-T_\alpha)\xi}{\xi}\ge0$, por lo que $T\le R$. Volver a sumar $T_{\alpha_0}$ da la conclusión para la red original. El signo del último producto interno debe ser $T-T_\alpha$, no su opuesto.
\end{proof}

\begin{proposicion}\label{v3:limites-proyecciones}
Un límite fuerte de proyecciones ortogonales es una proyección. Si $(P_\alpha)$ es creciente, converge fuertemente a la proyección sobre $\overline{\bigcup_\alpha P_\alpha H}$; si es decreciente, converge a la proyección sobre $\bigcap_\alpha P_\alpha H$.
\end{proposicion}
\begin{proof}
Las proyecciones tienen norma a lo sumo uno. Si $P_\alpha\to P$ fuertemente, los coeficientes dan $P=P^*$ y la continuidad fuerte del producto bajo esa cota da $P^2=P$. Para proyecciones, $P\le Q$ equivale a $PH\subset QH$: la inclusión da $QP=PQ=P$; recíprocamente, si $P\le Q$ y $P\xi=\xi$, entonces $\norm\xi^2\le\ip{Q\xi}{\xi}\le\norm\xi^2$, de donde $(I-Q)\xi=0$. Una red creciente está acotada por $I$, así Vigier da una proyección límite. Su imagen contiene cada $P_\alpha H$ y, porque $P_\alpha\xi\to P\xi$, está contenida en la clausura de su unión. Para una red decreciente aplique el resultado a $I-P_\alpha$ y tome complementos ortogonales.
\end{proof}

\section{El doble conmutante}

\begin{definicion}
Para $\mathcal S\subset\Lcal(H)$, su \emph{conmutante} es
\[
 \mathcal S'=\{T:TS=ST\text{ para todo }S\in\mathcal S\}.
\]
Su doble conmutante es $\mathcal S''=(\mathcal S')'$. Una \emph{álgebra de von Neumann unital sobre $H$} es una subálgebra autoadjunta de $\Lcal(H)$ que contiene $I_H$ y es WOT-cerrada.
\end{definicion}
\index{conmutante}\index{algebra de von Neumann@álgebra de von Neumann}

El conmutante siempre es una subálgebra que contiene $I$, y $\mathcal S\subset\mathcal S''$. Si $\mathcal S$ es estable bajo adjuntos, también lo es $\mathcal S'$: tome adjuntos de $TS^*=S^*T$. Los conmutantes son SOT y WOT cerrados, porque para $S$ fijo el mapa $T\mapsto TS-ST$ es continuo en ambas topologías.

\begin{teorema}[Doble conmutante de von Neumann]\label{v3:doble-conmutante}
Si $\mathcal A\subset\Lcal(H)$ es una subálgebra autoadjunta que contiene $I_H$, entonces
\[
 \overline{\mathcal A}^{\mathrm{SOT}}
 =\overline{\mathcal A}^{\mathrm{WOT}}=\mathcal A''.
\]
En particular, para tal subálgebra son equivalentes ser SOT-cerrada, ser WOT-cerrada y ser igual a su doble conmutante.
\end{teorema}
\begin{proof}
Como $\mathcal A''$ es cerrado y contiene a $\mathcal A$, basta aproximar fuertemente cada $T\in\mathcal A''$ por elementos de $\mathcal A$. Fije primero un vector $\xi$ y sea $K=\overline{\mathcal A\xi}$. Es un subespacio reductor para $\mathcal A$, pues ésta es autoadjunta. La proyección $P_K$ conmuta con todos sus elementos, ya que éstos preservan tanto $K$ como $K^\perp$. Así $P_K\in\mathcal A'$. Como $I\in\mathcal A$, $\xi\in K$ y $P_K\xi=\xi$. Entonces $P_KT\xi=TP_K\xi=T\xi$, de modo que $T\xi\in K$. Podemos aproximar $T\xi$ por $S\xi$ con $S\in\mathcal A$.

Falta aproximar simultáneamente un número finito de vectores con \emph{el mismo} operador. Para $\xi_1,\ldots,\xi_n$, trabaje en $H^n$ con la amplificación diagonal $S^{(n)}=\operatorname{diag}(S,\ldots,S)$. Si una matriz de operadores $R=(R_{ij})$ conmuta con todos los $S^{(n)}$, cada entrada $R_{ij}$ conmuta con todos los $S\in\mathcal A$, así $R_{ij}\in\mathcal A'$. Por tanto $T$ conmuta con cada entrada y $T^{(n)}\in(\mathcal A^{(n)})''$. El argumento de un vector aplicado a $\xi=(\xi_1,\ldots,\xi_n)$ produce $S\in\mathcal A$ con
\[
 \sum_{k=1}^n\norm{(T-S)\xi_k}^2<\eps^2.
\]
En particular, $\norm{(T-S)\xi_k}<\eps$ para cada $k$. Éstas son exactamente las vecindades básicas de SOT, así $T$ pertenece a la clausura fuerte. La igualdad entre las dos clausuras también resulta del ejercicio~\ref{ej:v3-l31-5}, porque un álgebra es convexa.
\end{proof}

La condición $I_H\in\mathcal A$ no puede omitirse sin reemplazarla por no degeneración y ajustar el argumento: para $\mathcal A=\{0\}$, su clausura es cero, pero $\mathcal A''=\C I$. Algunas convenciones llaman también álgebras de von Neumann a las subálgebras autoadjuntas WOT-cerradas con identidad propia $P\ne I_H$. En el capítulo siguiente veremos que esa identidad siempre es una proyección y que el teorema anterior se aplica al restringir la acción a $PH$. Indicaremos cuándo se usa esta convención más amplia.

\section{Ejercicios sobre conmutantes y clausuras}

\begin{ejercicio}\label{ej:v3-l32-1}
Si $\mathcal A$ es una subálgebra $C^*$ de $\Lcal(H)$, pruebe que $\mathcal A'$ es un álgebra de von Neumann unital.
\end{ejercicio}
\begin{solucion}
Contiene $I$, es una subálgebra y es autoadjunta porque $\mathcal A$ lo es. Para cada $S\in\mathcal A$, la ecuación $TS=ST$ define un cerrado WOT. Su intersección sobre todos los $S$ es $\mathcal A'$, así es WOT-cerrada. No se necesita que $\mathcal A$ contenga la identidad para esta afirmación.
\end{solucion}

\begin{ejercicio}\label{ej:v3-l32-2}
Si $\mathcal A_j\subset\Lcal(H_j)$ son álgebras de von Neumann unitarias, pruebe que el álgebra de operadores diagonales
\[
 \mathcal M=\{\bigoplus_jT_j:T_j\in\mathcal A_j,\ \sup_j\norm{T_j}<\infty\}
\]
es un álgebra de von Neumann sobre $\bigoplus_jH_j$. Aclare qué suma de álgebras se utiliza.
\end{ejercicio}
\begin{solucion}
El producto acotado coordenado es un álgebra $C^*$ y actúa diagonalmente con norma $\sup_j\norm{T_j}$. Contiene la identidad. Si una red de sus operadores converge en WOT a $T$, todos los coeficientes entre componentes diferentes siguen siendo cero, así $T$ es diagonal. Sus compresiones $T_j$ son límites WOT de elementos de $\mathcal A_j$, luego pertenecen a ellas; además $\norm{T_j}\le\norm T$. Por tanto $T\in\mathcal M$ y el álgebra es WOT-cerrada.

La notación de suma directa aquí significa \emph{producto acotado diagonal}. La suma algebraica, ni en general la suma $c_0$, es WOT-cerrada. Por ejemplo, las proyecciones de soporte finito en $\bigoplus_{n\ge1}\C=\ell^2$ pertenecen al álgebra diagonal $c_0$ y convergen fuertemente a $I$, que no pertenece a $c_0$.
\end{solucion}

\begin{ejercicio}\label{ej:v3-l32-3}
Pruebe que una álgebra de von Neumann es cerrada en norma y, por tanto, es una álgebra $C^*$.
\end{ejercicio}
\begin{solucion}
Una red convergente en norma converge en WOT. Si sus elementos pertenecen al álgebra WOT-cerrada, el límite también pertenece a ella. Así es normocerrada en el álgebra de Banach $\Lcal(H)$, por lo que es completa; hereda la involución y la identidad $C^*$ de los operadores.
\end{solucion}

\begin{ejercicio}\label{ej:v3-l32-4}
Pruebe que la bola unitaria cerrada de una álgebra de von Neumann es WOT-compacta y SOT-cerrada. Determine si siempre es SOT-compacta.
\end{ejercicio}
\begin{solucion}
La bola del álgebra es la intersección de un cerrado WOT con la bola WOT-compacta de $\Lcal(H)$, y por tanto es WOT-compacta. Es WOT-cerrada porque los compactos en un Hausdorff son cerrados; luego también es SOT-cerrada. No siempre es SOT-compacta. En $\Lcal(\ell^2)$ considere $T_n\xi=\ip\xi{e_1}e_n$. Todos tienen norma uno, pero $T_ne_1=e_n$. Si la bola fuera SOT-compacta, su imagen por la evaluación continua $T\mapsto Te_1$ sería compacta en norma y contendría la sucesión ortonormal, imposible porque sus términos están a distancia $\sqrt2$ unos de otros. En dimensión finita, en cambio, las tres topologías coinciden y la bola sí es compacta.
\end{solucion}

El doble conmutante identifica el cierre de un álgebra mediante sus simetrías de conmutación. Todavía falta controlar las normas de los operadores aproximantes. El teorema de Kaplansky hará posible esa mejora y preparará la extensión del cálculo continuo a funciones borelianas. Referencias complementarias: \cite{conway1990,murphy1990,zhu1993}.

\chapter{Densidad de Kaplansky y cálculo funcional boreliano}\label{cap:v3-17}

El cierre fuerte de un álgebra de operadores contiene límites que el cierre en norma no detecta. Las proyecciones sobre rangos cerrados y las funciones indicadoras de conjuntos espectrales son ejemplos centrales. Este capítulo demuestra cómo aparecen y cómo aproximar operadores del cierre sin perder el control de sus normas.

\section{Proyecciones de soporte e identidad propia}\label{lec:v3-33}
\fuentelectura{33}{Propiedades de las álgebras de von Neumann}

Cuando no se exija contener $I_H$, llamaremos álgebra de von Neumann \emph{en sentido amplio} a una subálgebra autoadjunta WOT-cerrada de $\Lcal(H)$. Tendrá una identidad propia que es una proyección; restringida al rango de esa proyección, será un álgebra de von Neumann unital en el sentido del capítulo anterior.

\begin{teorema}[Proyección de soporte]\label{v3:proyeccion-soporte}
Si $0\le T\le I$, la sucesión creciente $T^{1/n}$ converge fuertemente a la proyección sobre $\overline{TH}$. En consecuencia, una álgebra de von Neumann que contiene a $S$ contiene la proyección sobre $\overline{SH}$.
\end{teorema}
\begin{proof}
Para $0\le t\le1$, $t^{1/n}$ crece con $n$ y está entre cero y uno. El cálculo continuo traslada estas desigualdades a operadores, y Vigier da un límite fuerte autoadjunto $P$. Por la cota uniforme, $T^{2/n}=(T^{1/n})^2\to P^2$ fuertemente. La subsucesión de índices pares tiene límite $P$, porque $T^{2/(2n)}=T^{1/n}\to P$. La unicidad de límite implica $P=P^2$.

Si $T\xi=0$, todo polinomio en $T$ sin término constante anula $\xi$. Aproximar $t^{1/n}$ uniformemente por tales polinomios muestra $T^{1/n}\xi=0$, y luego $P\xi=0$. Recíprocamente, $0\le T\le P$ porque la sucesión creciente comienza en $T$; si $P\xi=0$, entonces $0\le\ip{T\xi}{\xi}\le\ip{P\xi}{\xi}=0$, y $\norm{T^{1/2}\xi}^2=0$, de donde $T\xi=0$. Así $\ker P=\ker T$, y $PH=(\ker T)^\perp=\overline{TH}$.

Para $S\ne0$, aplique lo anterior a $T=SS^*/\norm{SS^*}$. Sus raíces pertenecen al álgebra, incluso si ésta carece de $I_H$, porque las funciones de raíces se anulan en cero. El límite pertenece por cerradura fuerte. Finalmente $\ker(SS^*)=\ker S^*$, pues $\ip{SS^*\xi}{\xi}=\norm{S^*\xi}^2$; por tanto $\overline{SS^*H}=\overline{SH}$. Para $S=0$, la proyección es cero. Se necesitan estas clausuras de rangos y el núcleo de $S^*$, no el de $S$.
\end{proof}

\begin{teorema}\label{v3:identidad-propia}
Toda subálgebra autoadjunta WOT-cerrada $\mathcal M\subset\Lcal(H)$ tiene una identidad propia $P$, que es la proyección sobre $\overline{\mathcal M H}$.
\end{teorema}
\begin{proof}
Como es convexa, $\mathcal M$ también es SOT-cerrada. Considere los subconjuntos finitos $F$ de sus proyecciones y, para cada uno, la proyección $P_F$ sobre el rango cerrado de $\sum_{Q\in F}Q$. El teorema anterior asegura $P_F\in\mathcal M$. Su rango es la clausura de la suma de los rangos de los $Q\in F$: los complementos ortogonales coinciden porque
$\ker(\sum Q)=\bigcap\ker Q$, como se ve de $\ip{(\sum Q)\xi}{\xi}=\sum\norm{Q\xi}^2$.

La red $P_F$, ordenada por inclusión de los conjuntos finitos, es creciente y está acotada por $I$. Converge fuertemente a una proyección $P\in\mathcal M$ que domina todas las proyecciones de $\mathcal M$. Para $T\in\mathcal M$, la proyección sobre $\overline{TH}$ pertenece a $\mathcal M$, así $PT=T$. Aplicando esto a $T^*$ y tomando adjuntos se obtiene $TP=T$. De aquí $\mathcal M H\subset PH$, mientras $P\in\mathcal M$ da $PH\subset\mathcal M H$; en particular $PH=\overline{\mathcal M H}$. El mismo argumento incluye $\mathcal M=0$, con $P=0$.
\end{proof}

\section{Funciones acotadas y convergencia fuerte}

\begin{teorema}\label{v3:funcion-acotada-sot}
Si $f:\R\to\C$ es continua y acotada y $T_\alpha,T$ son autoadjuntos con $T_\alpha\to T$ fuertemente, entonces $f(T_\alpha)\to f(T)$ fuertemente. No se requiere una cota uniforme previa para $\norm{T_\alpha}$.
\end{teorema}
\begin{proof}
Para $z=i$ o $z=-i$, los resolventes de un autoadjunto satisfacen $\norm{(zI-T_\alpha)^{-1}}\le1$, por el cálculo continuo sobre un espectro real. La identidad resolvente da
\[
 (zI-T_\alpha)^{-1}-(zI-T)^{-1}
 =(zI-T_\alpha)^{-1}(T_\alpha-T)(zI-T)^{-1}.
\]
Al aplicarla a un vector fijo, el miembro derecho tiende a cero en norma. Así las funciones $r_\pm(t)=(\pm i-t)^{-1}$ tienen la propiedad requerida.

Dentro de $C_0(\R)$, las funciones con esa propiedad forman una subálgebra estable bajo conjugación y cerrada en norma uniforme. La estabilidad de productos se sigue de la continuidad fuerte conjunta para factores uniformemente acotados, pues $\norm{f(T_\alpha)}\le\norm f_\infty$. La estabilidad bajo conjugación se obtiene porque los operadores $f(T_\alpha)$ y $f(T)$ son normales y se aplica el ejercicio~\ref{ej:v3-l31-6}; también puede comprobarse primero en los generadores y luego por clausura. Para la cerradura uniforme se usa el argumento de tres términos con una aproximación uniforme y un vector fijo. Las funciones $r_+,r_-$ separan puntos, no se anulan simultáneamente y son conjugadas entre sí. Stone--Weierstrass en la compactificación por un punto de $\R$ implica que generan densamente $C_0(\R)$. Por tanto todas las funciones de $C_0(\R)$ tienen la propiedad.

Finalmente tome $f\in C_b(\R)$ y elija $g\in C_c(\R)$ igual a uno en un intervalo que contenga $\sigma(T)$. Entonces $g(T)=I$, $g(T_\alpha)\to I$ fuertemente y $(fg)(T_\alpha)\to(fg)(T)=f(T)$ fuertemente. Para cada $\xi$,
\[
 \norm{(f(T_\alpha)-f(T))\xi}
 \le\norm f_\infty\norm{(I-g(T_\alpha))\xi}
 +\norm{((fg)(T_\alpha)-f(T))\xi}\to0.
\]
Esta localización controla las posibles normas grandes de los $T_\alpha$ sin suponer que la red completa sea acotada.
\end{proof}

\begin{teorema}[Densidad de Kaplansky]\label{v3:kaplansky}
Sea $\mathcal A$ una subálgebra $C^*$ de $\Lcal(H)$ y $\mathcal M=\overline{\mathcal A}^{\mathrm{SOT}}$. Entonces:
\begin{enumerate}[label=(\roman*)]
\item $\mathcal A_{\mathrm{sa}}$ es SOT-densa en $\mathcal M_{\mathrm{sa}}$;
\item su bola unitaria autoadjunta es SOT-densa en la bola unitaria autoadjunta de $\mathcal M$;
\item la bola unitaria de $\mathcal A$ es SOT-densa en la bola unitaria de $\mathcal M$.
\end{enumerate}
\end{teorema}
\begin{proof}
Primero, $\mathcal M$ es también la clausura WOT porque $\mathcal A$ es convexa. La continuidad WOT del adjunto y la continuidad por separado del producto muestran que esa clausura es una subálgebra autoadjunta: para el producto se aproxima primero un factor de $\mathcal M$ manteniendo el otro en $\mathcal A$, y después se aproxima el segundo. Por tanto $\mathcal M$ es una álgebra de von Neumann en sentido amplio.

(i) Si $T=T^*\in\mathcal M$, una red $S_\alpha\in\mathcal A$ converge a $T$ en WOT. Entonces $(S_\alpha+S_\alpha^*)/2\to T$ en WOT. El conjunto real convexo $\mathcal A_{\mathrm{sa}}$ tiene iguales clausuras fuerte y débil, por lo que $T$ pertenece a su clausura fuerte.

(ii) Sea $T=T^*$, $\norm T\le1$, y elija una red autoadjunta $T_\alpha\in\mathcal A$ que converge fuertemente a $T$ por (i). La función continua
\[
 f(t)=\begin{cases}t,&|t|\le1,\\1/t,&|t|>1\end{cases}
\]
pertenece a $C_0(\R)$, satisface $f(0)=0$, es real y tiene norma uno. Por el teorema anterior $f(T_\alpha)\to f(T)=T$ fuertemente. Los aproximantes son autoadjuntos, tienen norma a lo sumo uno y pertenecen a $\mathcal A$; la condición $f(0)=0$ asegura pertenencia aun cuando $\mathcal A$ no contenga $I_H$.

(iii) El álgebra $M_2(\mathcal A)$ es una subálgebra $C^*$ de $\Lcal(H\oplus H)$ y es SOT-densa en $M_2(\mathcal M)$. Para justificar esta densidad, una vecindad fuerte impone un número finito de pruebas sobre vectores de $H\oplus H$; basta aproximar las cuatro entradas en los correspondientes conjuntos finitos de componentes y sumar las estimaciones. Para una contracción $T\in\mathcal M$, el operador
\[
 S=\begin{pmatrix}0&T\\T^*&0\end{pmatrix}
\]
es autoadjunto y tiene norma $\norm T\le1$, porque $S^2=\operatorname{diag}(TT^*,T^*T)$. Por (ii) en el álgebra de matrices, existen contracciones autoadjuntas $S_\alpha\in M_2(\mathcal A)$ que convergen fuertemente a $S$. Sus entradas $(1,2)$ pertenecen a la bola de $\mathcal A$ y convergen fuertemente a $T$, pues se obtienen aplicando a vectores $(0,\xi)$ y proyectando a la primera componente.
\end{proof}

\section{Medidas espectrales e integración}\label{lec:v3-34}
\fuentelectura{34}{El cálculo funcional boreliano}

\begin{definicion}
Sea $X$ compacto Hausdorff. Una \emph{medida espectral} $E$ en $X$ con valores en $\Lcal(H)$ asigna una proyección ortogonal $E(B)$ a cada Borel $B$, con $E(\varnothing)=0$, $E(X)=I$, $E(B\cap C)=E(B)E(C)$, y tal que
$E_{\xi,\eta}(B)=\ip{E(B)\xi}{\eta}$ es una medida compleja regular para cada $\xi,\eta\in H$.
\end{definicion}
\index{medida espectral}\index{cálculo funcional!boreliano}

Si $(B_n)$ son disjuntos, las proyecciones $E(B_n)$ son ortogonales entre sí. La aditividad escalar y el producto interno dan
\[
 E\Bigl(\bigcup_nB_n\Bigr)\xi=\sum_nE(B_n)\xi
\]
con convergencia en norma: la norma cuadrada de la cola es $E_{\xi,\xi}$ de la unión de la cola y tiende a cero. Así la aditividad es fuerte, no en general en norma de operadores. El ejemplo modelo es $E(B)g=\mathbf1_Bg$ en $L^2(X,\mu)$ para una medida positiva finita regular.

El espacio $B^\infty(X)$ de funciones borelianas acotadas, con norma del supremo puntual y conjugación, es una álgebra $C^*$. Una sucesión uniforme de Cauchy converge uniformemente a una función acotada, que es boreliana como límite puntual; esto prueba completitud, y la identidad $C^*$ es puntual. Aquí no se identifican funciones casi iguales: no hay una medida escalar única fijada de antemano.

\begin{proposicion}\label{v3:integral-espectral}
Para una medida espectral $E$, existe un $*$-homomorfismo unital contractivo
\[
 \Phi_E:B^\infty(X)\to\Lcal(H),\qquad f\mapsto\int_Xf\,dE,
\]
caracterizado por
$\ip{\Phi_E(f)\xi}{\eta}=\int_X f\,dE_{\xi,\eta}$.
\end{proposicion}
\begin{proof}
Si $f=\sum_{j=1}^mc_j\mathbf1_{B_j}$ con $B_j$ disjuntos, defina $\Phi_E(f)=\sum_jc_jE(B_j)$. La ortogonalidad de los rangos da
\[
 \norm{\Phi_E(f)\xi}^2=\sum_j|c_j|^2\norm{E(B_j)\xi}^2
 \le\norm f_\infty^2\norm\xi^2.
\]
Refinar dos particiones a sus intersecciones demuestra que la definición no depende de la representación simple y que preserva productos, sumas y adjuntos. Toda función boreliana acotada se aproxima uniformemente por funciones simples, particionando su rango acotado en cuadrados pequeños. La cota anterior permite definir $\Phi_E(f)$ como límite en norma; las propiedades algebraicas pasan al límite.

Para identificar los coeficientes, note que la variación total de $E_{\xi,\eta}$ es a lo sumo $\norm\xi\norm\eta$. En una partición finita,
\begin{align*}
 \sum_j|\ip{E(B_j)\xi}{\eta}|
 &=\sum_j|\ip{E(B_j)\xi}{E(B_j)\eta}|\\
 &\le\Bigl(\sum_j\norm{E(B_j)\xi}^2\Bigr)^{1/2}
 \Bigl(\sum_j\norm{E(B_j)\eta}^2\Bigr)^{1/2}
 \le\norm\xi\norm\eta.
\end{align*}
La fórmula de coeficientes vale para funciones simples y pasa al límite uniforme usando esta cota. En particular la forma sesquilineal $\int f\,dE_{\xi,\eta}$ es continua y determina un único operador, como exige la definición por formas.
\end{proof}

\begin{teorema}[Resolución espectral de un operador normal]\label{v3:espectral-borel}
Si $T\in\Lcal(H)$ es normal, existe una única medida espectral regular $E$ sobre $K=\sigma(T)$ tal que
\[
 T=\int_K z\,dE(z).
\]
La aplicación $f\mapsto f(T)=\int f\,dE$ extiende el cálculo continuo a $B^\infty(K)$. Si una álgebra de von Neumann unital contiene a $T$, contiene a cada $f(T)$.
\end{teorema}
\begin{proof}
Para $\xi,\eta\in H$, el funcional $f\mapsto\ip{f(T)\xi}{\eta}$ de $C(K)$ tiene norma a lo sumo $\norm\xi\norm\eta$. Riesz--Markov produce una única medida compleja regular $\mu_{\xi,\eta}$ que lo representa, con variación total acotada por ese mismo producto. La unicidad implica que estas medidas dependen sesquilinealmente de $\xi,\eta$, y $\mu_{\eta,\xi}=\overline{\mu_{\xi,\eta}}$. Además $\mu_{\xi,\xi}$ es positiva, porque $f(T)\ge0$ para $f\ge0$ continua.

Para $f$ boreliana acotada, la fórmula $b_f(\xi,\eta)=\int f\,d\mu_{\xi,\eta}$ define una forma sesquilineal acotada por $\norm f_\infty\norm\xi\norm\eta$. Riesz produce un operador $\Psi(f)$ con esos coeficientes. Es lineal y contractivo; extiende el cálculo continuo por unicidad y conserva adjuntos porque $\mu_{\eta,\xi}=\overline{\mu_{\xi,\eta}}$.

Detallamos la multiplicatividad, donde no basta extender formalmente una identidad de funciones continuas. Para $g\in C(K)$, la unicidad de medidas, probando contra funciones continuas, da
\[
 \mu_{g(T)\xi,\eta}=g\mu_{\xi,\eta}
 =\mu_{\xi,g(T)^*\eta}.
\]
Por tanto, para cualquier $f\in B^\infty(K)$,
\[
 \Psi(f)g(T)=\Psi(fg)=g(T)\Psi(f).
\]
Ahora, al probar la medida $\mu_{\Psi(f)\xi,\eta}$ contra una función continua $g$, se obtiene
\[
 \int g\,d\mu_{\Psi(f)\xi,\eta}
 =\ip{g(T)\Psi(f)\xi}{\eta}
 =\int gf\,d\mu_{\xi,\eta}.
\]
La medida $f\mu_{\xi,\eta}$ es regular y finita; en el compacto métrico $K\subset\C$, toda medida de Borel finita es regular. Riesz y unicidad dan $\mu_{\Psi(f)\xi,\eta}=f\mu_{\xi,\eta}$. Para $h$ boreliana acotada,
\[
 \ip{\Psi(h)\Psi(f)\xi}{\eta}
 =\int h\,d\mu_{\Psi(f)\xi,\eta}
 =\int hf\,d\mu_{\xi,\eta}
 =\ip{\Psi(hf)\xi}{\eta}.
\]
Así $\Psi$ es multiplicativa en todas las funciones borelianas acotadas.

Defina $E(B)=\Psi(\mathbf1_B)$. Las identidades de indicadores y el $*$-homomorfismo prueban que son proyecciones, que se multiplican por intersección y que $E(K)=I$. Sus medidas de coeficientes son precisamente $\mu_{\xi,\eta}$, por definición, así satisfacen regularidad y aditividad. La integral espectral de cualquier $f$ tiene los mismos coeficientes que $\Psi(f)$ y coincide con ella; en particular $T=\Psi(z)=\int z\,dE$.

Para la pertenencia a una álgebra de von Neumann $\mathcal M$ que contiene a $T$, sea $S\in\mathcal M'$. Como $S$ conmuta con $T,T^*$, conmuta con el cálculo continuo. La igualdad de coeficientes para funciones continuas implica $\mu_{S\xi,\eta}=\mu_{\xi,S^*\eta}$; integrar cualquier función boreliana acotada da $\Psi(f)S=S\Psi(f)$. Por el doble conmutante, $\Psi(f)\in\mathcal M''=\mathcal M$.

Finalmente, si $F$ es otra medida espectral en $K$ con integral coordenada $T$, la proposición~\ref{v3:integral-espectral} muestra que su integral coincide con el cálculo continuo en los polinomios en $z,\overline z$, y por Stone--Weierstrass en todo $C(K)$. La unicidad de Riesz implica que sus medidas de coeficientes son $\mu_{\xi,\eta}$; evaluar en cada Borel y en todos los vectores da $F(B)=E(B)$.
\end{proof}

\begin{corolario}\label{v3:proyecciones-generan}
El espacio lineal generado por las proyecciones de una álgebra de von Neumann es denso en norma en el álgebra. Todo unitario de una álgebra de von Neumann unital tiene la forma $e^{iS}$ con $S=S^*$ en el álgebra y $\norm S\le\pi$.
\end{corolario}
\begin{proof}
Si $T$ es normal, aproxime uniformemente la función coordenada sobre su espectro por funciones simples borelianas. Sus integrales son combinaciones finitas de proyecciones del álgebra y convergen en norma a $T$. Un elemento general es $H+iK$ con $H,K$ autoadjuntos, y se aproxima cada parte. Para el segundo enunciado, elija una determinación boreliana $\theta:\mathbb T\to[-\pi,\pi]$ del argumento, y ponga $S=\theta(U)$ mediante el cálculo boreliano. Es autoadjunto y su norma es a lo sumo $\pi$. La identidad $e^{iS}=U$ se verifica aproximando la exponencial uniformemente por polinomios sobre $[-\pi,\pi]$ y usando la multiplicatividad del cálculo: el resultado es $(e^{i\theta})(U)=z(U)=U$.
\end{proof}

\section{Ejercicios sobre álgebras de von Neumann y densidad}

\begin{ejercicio}\label{ej:v3-l33-1}
Sea $\mathcal M$ una álgebra de von Neumann. (i) Si $\mathcal B\subset\mathcal M$ es una subálgebra $C^*$ SOT-cerrada y hereditaria, pruebe que existe una única proyección $P\in\mathcal B$ con $\mathcal B=P\mathcal M P$. Hereditaria significa que $0\le a\le b$, $a\in\mathcal M$, $b\in\mathcal B$, implica $a\in\mathcal B$. (ii) Si $I$ es un ideal bilateral SOT-cerrado, pruebe $I=\mathcal M P$ para una única proyección $P\in I\cap\mathcal M'$.
\end{ejercicio}


\begin{ejercicio}\label{ej:v3-l33-2}
Si $I_H\in\mathcal A$ y $\mathcal M=\overline{\mathcal A}^{\mathrm{SOT}}$, pruebe que los unitarios de $\mathcal A$ son SOT-densos en los unitarios de $\mathcal M$.
\end{ejercicio}


\begin{ejercicio}\label{ej:v3-l33-3}
Para una subálgebra $C^*$ $\mathcal A\subset\Lcal(H)$, pruebe la equivalencia entre: (i) ser WOT-cerrada; (ii) tener bola unitaria SOT-cerrada; (iii) tener bola unitaria WOT-cerrada. Si se usa la definición unital de álgebra de von Neumann, suponga también $I_H\in\mathcal A$.
\end{ejercicio}


\begin{ejercicio}\label{ej:v3-l33-4}
Pruebe las versiones de Kaplansky (i) y (ii) para elementos positivos y para contracciones positivas.
\end{ejercicio}


\begin{ejercicio}\label{ej:v3-l33-5}
Sea $\Phi:\Lcal(H_1)\to\Lcal(H_2)$ un $*$-monomorfismo WOT-continuo. Pruebe que lleva álgebras de von Neumann a álgebras de von Neumann en sentido amplio; si $\Phi$ es unital, la imagen es unital.
\end{ejercicio}


\section{Ejercicios sobre integrales espectrales y compacidad}

\begin{ejercicio}\label{ej:v3-l34-1}
Para una medida espectral $E$ en $X$, pruebe que $f\mapsto\int f\,dE$ es un $*$-homomorfismo de $B^\infty(X)$ en $\Lcal(H)$.
\end{ejercicio}


\begin{ejercicio}\label{ej:v3-l34-2}
Pruebe que la resolución de la identidad de un operador normal es única.
\end{ejercicio}


\begin{ejercicio}\label{ej:v3-l34-3}
Si $0\le T\le I$ pertenece a una álgebra de von Neumann $\mathcal M$, pruebe que existen proyecciones $P_n\in\mathcal M$ con $T=\sum_{n\ge1}2^{-n}P_n$ en norma.
\end{ejercicio}


\begin{ejercicio}\label{ej:v3-l34-4}
Si $\mathcal M$ es una álgebra de von Neumann unital, pruebe que $T\in\mathcal M$ si y sólo si conmuta con todas las proyecciones de $\mathcal M'$.
\end{ejercicio}


\begin{ejercicio}\label{ej:v3-l34-5}
Sea $0\le T\le I$. Si todo subespacio cerrado de $H$ contenido en $TH$ tiene dimensión finita, pruebe que $T$ es compacto. Como paso auxiliar, pruebe que $UU^*\le VV^*$ implica una factorización $U=VW$ con $\norm W\le1$.
\end{ejercicio}


El cálculo boreliano añade al álgebra las proyecciones que recortan conjuntos espectrales, y esas proyecciones bastan para aproximar en norma todos sus elementos. En el caso conmutativo este fenómeno tendrá una descripción familiar: multiplicación por funciones esencialmente acotadas. Referencias para ampliar: \cite{conway1990,murphy1990,zhu1993}.

\chapter{Álgebras de multiplicación y el caso conmutativo}\label{cap:v3-18}

Las funciones esencialmente acotadas actúan sobre $L^2$ por multiplicación. Esta acción conserva normas, adjuntos y productos, y convierte la topología débil estrella de $L^\infty$ en la topología débil del operador. Al final veremos que toda álgebra de von Neumann conmutativa sobre un Hilbert separable admite una descripción de este tipo, aunque su representación inicial pueda tener multiplicidades.

\section{Multiplicación en un espacio de medida}\label{lec:v3-35}
\fuentelectura{35}{L∞(X) como álgebra de von Neumann}

Fijamos un compacto Hausdorff no vacío $X$ y una medida de Borel positiva, \emph{finita y regular} $\mu$, con soporte todo $X$. La última condición equivale a que cada abierto no vacío tenga medida positiva: el complemento de un abierto nulo sería un cerrado propio de medida total, y recíprocamente. La finitud no se deduce de la sola compacidad para una medida de Borel arbitraria; por ejemplo, la medida de conteo en $[0,1]$ tiene masa infinita. Es automática si se parte de una medida de Radon localmente finita sobre un compacto, y aquí la incluimos explícitamente.

Los elementos de $L^\infty(X,\mu)$ son clases de funciones iguales casi en todas partes. Un representante esencialmente acotado puede modificarse en un conjunto nulo para ser acotado en todos los puntos. La dualidad $L^\infty=(L^1)^*$ vale en este espacio finito de medida, y su topología débil estrella es la convergencia de las integrales contra todos los elementos de $L^1$.

\begin{proposicion}\label{v3:multiplicacion-isometrica}
La aplicación $f\mapsto M_f$, $M_fg=fg$, es un $*$-monomorfismo isométrico de $L^\infty(\mu)$ en $\Lcal(L^2(\mu))$. La inclusión de $C(X)$ en $L^\infty(\mu)$ también es isométrica.
\end{proposicion}
\begin{proof}
La estimación $\norm{fg}_2\le\norm f_\infty\norm g_2$ prueba $\norm{M_f}\le\norm f_\infty$. Si $c<\norm f_\infty$, el conjunto $E=\{|f|>c\}$ tiene medida positiva y finita. Para $g=\mathbf1_E/\sqrt{\mu(E)}$, $\norm g_2=1$ y $\norm{M_fg}_2\ge c$. Dejar $c$ acercarse al supremo esencial da la igualdad. Productos, sumas e identidad se conservan por las operaciones puntuales, y $M_f^*=M_{\overline f}$ por la identidad de productos internos. La isometría implica inyectividad en clases casi iguales.

Si $f\in C(X)$, su máximo absoluto se alcanza. Para cualquier $c<\norm f_{\infty,X}$, el conjunto abierto $\{|f|>c\}$ es no vacío y, por soporte total, tiene medida positiva. Así el supremo esencial coincide con el máximo puntual. En particular, una función continua nula casi en todas partes es nula en todo $X$.
\end{proof}

Denotaremos $\mathcal C=\{M_f:f\in C(X)\}$ y $\mathcal M_\mu=\{M_f:f\in L^\infty(\mu)\}$.

\begin{teorema}\label{v3:weakstar-wot-multiplicadores}
La isometría $f\mapsto M_f$ es un homeomorfismo de $L^\infty(\mu)$ con su topología débil estrella sobre $\mathcal M_\mu$ con WOT.
\end{teorema}
\begin{proof}
Cada coeficiente es $\ip{M_fg}{h}=\int f g\overline h\,d\mu$, con $g\overline h\in L^1$ por Cauchy--Schwarz. Recíprocamente, para $k\in L^1$ escriba $k=g\overline h$, donde $h=\sqrt{|k|}$ y $g=\theta\sqrt{|k|}$, con $\theta=k/|k|$ donde $k\ne0$ y cero donde se anula. Ambos están en $L^2$ y el producto es exactamente $k$. Así las pruebas que definen ambas topologías coinciden, para redes arbitrarias y sin restringirlas a una bola.
\end{proof}

\begin{teorema}[Densidad de funciones continuas]\label{v3:continua-densa-linfty}
La bola unitaria de $C(X)$ es débil estrella densa en la bola unitaria de $L^\infty(\mu)$. Por tanto $\mathcal C$ es WOT-densa y SOT-densa en $\mathcal M_\mu$.
\end{teorema}
\begin{proof}
Primero sea $s=\sum_{i=1}^nc_i\mathbf1_{E_i}$ una función simple de módulo a lo sumo uno, con $(E_i)$ una partición boreliana de $X$. Fije $g_1,\ldots,g_m\in L^1$ y $\eps>0$. Por regularidad y continuidad absoluta de las integrales, elija compactos disjuntos $K_i\subset E_i$ con
\[
 \int_{E_i\setminus K_i}|g_j|\,d\mu<\frac{\eps}{2n}
 \quad(1\le i\le n,\ 1\le j\le m).
\]
La continuidad absoluta significa que una medida suficientemente pequeña hace pequeñas todas esas integrales; la regularidad interior produce los compactos con ese control.

La normalidad de un compacto Hausdorff permite elegir vecindades abiertas disjuntas de los compactos $K_i$ y funciones continuas $0\le u_i\le1$, iguales a uno en $K_i$ y con soportes dentro de esas vecindades. Entonces $h=\sum_i c_i u_i$ satisface $\norm h_\infty\le1$, porque en cada punto hay a lo sumo un sumando no nulo, y $h=s$ en $\bigcup_iK_i$. Por tanto
\[
 \left|\int(s-h)g_j\,d\mu\right|
 \le2\sum_i\int_{E_i\setminus K_i}|g_j|\,d\mu<\eps.
\]
Este argumento construye una extensión compleja con control del módulo; extender por separado partes real e imaginaria no daría automáticamente la misma cota uno.

Para una clase $f$ con $\norm f_\infty\le1$, elija un representante de módulo a lo sumo uno y aproxímelo uniformemente por simples con valores en el disco cerrado unitario. El error contra cada $g_j$ está acotado por el error uniforme multiplicado por $\norm{g_j}_1$. Combinando esa aproximación con la construcción anterior se obtiene un elemento de $C(X)$ en cualquier vecindad débil estrella de $f$. La equivalencia de topologías del teorema anterior da densidad WOT de los multiplicadores. Como las bolas son convexas, sus clausuras WOT y SOT coinciden, y se obtiene también la afirmación fuerte.
\end{proof}

\begin{teorema}[Álgebra maximal conmutativa de multiplicadores]\label{v3:masa-multiplicadores}
Se tiene $\mathcal M_\mu'=\mathcal M_\mu$. En particular $\mathcal M_\mu$ es una álgebra de von Neumann conmutativa maximal.
\end{teorema}
\begin{proof}
Los multiplicadores conmutan entre sí, así $\mathcal M_\mu\subset\mathcal M_\mu'$. Para la otra inclusión, sea $T\in\mathcal M_\mu'$ y ponga $g=T1\in L^2$; aquí se usa que $\mu$ es finita. Para $f\in L^\infty$,
\[
 Tf=T(M_f1)=M_fT1=fg.
\]
Afirmamos que $|g|\le\norm T$ casi en todas partes. Si $E=\{|g|>\norm T+\delta\}$ tuviera medida positiva, entonces
\[
 \norm{T\mathbf1_E}_2^2=\int_E|g|^2\,d\mu
 >(\norm T)^2\mu(E),
\]
contradiciendo la cota de operador. Así $g\in L^\infty$, y $T=M_g$ sobre el subespacio denso $L^\infty\subset L^2$. La densidad se obtiene truncando cualquier función de $L^2$; por continuidad, $T=M_g$ en todo $L^2$.

El conmutante es WOT-cerrado, por lo que la igualdad prueba que $\mathcal M_\mu$ es de von Neumann. Si otra subálgebra conmutativa la contuviera, cada uno de sus elementos conmutaría con $\mathcal M_\mu$ y por tanto pertenecería a $\mathcal M_\mu'=\mathcal M_\mu$. Esto prueba maximalidad.
\end{proof}

\section{Cuándo dos medidas dan la misma álgebra sobre \texorpdfstring{$C(X)$}{C(X)}}

\begin{teorema}\label{v3:unicidad-clase-medida}
Suponga $X$ compacto metrizable y $\mu,\nu$ medidas positivas finitas regulares con soporte $X$. Si un $*$-isomorfismo
$\Phi:L^\infty(X,\mu)\to L^\infty(X,\nu)$ fija las funciones de $C(X)$, entonces $\mu$ y $\nu$ tienen los mismos conjuntos nulos. Las clases de $L^\infty$ se identifican canónicamente para ambas medidas y, bajo esa identificación, $\Phi$ es la identidad.
\end{teorema}
\begin{proof}
Un $*$-isomorfismo y su inversa preservan el orden. Por tanto preservan todos los supremos e ínfimos que existan: si $a$ es la menor cota superior de un conjunto, su imagen es una cota superior, y cualquier otra cota superior se devuelve por la inversa a una cota superior del conjunto original, que domina a $a$. El argumento para ínfimos es análogo.

Si $K\subset X$ es cerrado no vacío y $d$ induce la topología, las funciones continuas $f_n(x)=\max(0,1-nd(x,K))$ decrecen puntualmente a $\mathbf1_K$. En cada espacio $L^\infty$ su ínfimo de orden es la clase de $\mathbf1_K$: cualquier cota inferior queda debajo de todos los $f_n$ fuera de una unión numerable de nulos y por tanto debajo del límite. Como $\Phi$ fija los $f_n$ y conserva ínfimos,
\[
 \Phi([\mathbf1_K]_\mu)=[\mathbf1_K]_\nu.
\]
El caso $K=\varnothing$ es inmediato. Esta aproximación no vale para un Borel arbitrario $E$: $d(x,E)=d(x,\overline E)$, de modo que el límite sería $\mathbf1_{\overline E}$, precisión que corrige el argumento original.

Sea $\mathcal D$ la familia de Borel $E$ para los que la igualdad de indicadores vale. Contiene los cerrados, es estable bajo complementos porque $\Phi(1)=1$, y bajo uniones finitas porque $\mathbf1_{E\cup F}=\mathbf1_E+\mathbf1_F-\mathbf1_E\mathbf1_F$. Es estable bajo uniones numerables al escribirlas como supremos de sus uniones finitas y usar la preservación de supremos. Por tanto contiene la $\sigma$-álgebra de Borel generada por los cerrados.

Si $\mu(E)=0$, $[\mathbf1_E]_\mu=0$ y la identidad da $\nu(E)=0$. La inversa de $\Phi$ proporciona la implicación recíproca. Así las nociones de igualdad casi en todas partes y de acotación esencial coinciden. El mapa fija todas las funciones simples y, por isometría y aproximación uniforme por simples, fija cada clase de $L^\infty$.
\end{proof}

Este resultado no afirma que todas las medidas de soporte total sean equivalentes. Por ejemplo, en $[0,1]$ tanto la medida de Lebesgue como una probabilidad concentrada en un subconjunto numerable denso tienen soporte total, pero son mutuamente singulares. La hipótesis decisiva es la existencia del isomorfismo que fija $C(X)$.

\section{Ejercicios sobre multiplicadores y orden}

\begin{ejercicio}\label{ej:v3-l35-1}
Pruebe que una subálgebra $C^*$ $\mathcal A\subset\Lcal(H)$ es una álgebra de von Neumann conmutativa maximal si y sólo si $\mathcal A=\mathcal A'$.
\end{ejercicio}


\begin{ejercicio}\label{ej:v3-l35-2}
Si $\mathcal A$ es una subálgebra $C^*$ conmutativa, pruebe que su clausura WOT es una álgebra de von Neumann conmutativa en sentido amplio. Deduzca que el álgebra de von Neumann unital generada por un operador normal es conmutativa.
\end{ejercicio}


\begin{ejercicio}\label{ej:v3-l35-3}
Con la notación de esta lectura, pruebe $\mathcal C'=\mathcal M_\mu$. Deduzca que $L^\infty(\mu)$ es la clausura débil estrella de $C(X)$.
\end{ejercicio}


\begin{ejercicio}\label{ej:v3-l35-4}
Pruebe que todo $*$-isomorfismo preserva los supremos existentes en el orden de los autoadjuntos.
\end{ejercicio}


\begin{ejercicio}\label{ej:v3-l35-5}
Si $\mu,\nu$ son medidas positivas regulares mutuamente absolutamente continuas sobre $X$, pruebe la identificación canónica $L^\infty(X,\mu)=L^\infty(X,\nu)$, sin exigir soporte total.
\end{ejercicio}


\section{Vectores cíclicos y separantes}\label{lec:v3-36}
\fuentelectura{36}{Caracterización de las álgebras de von Neumann conmutativas}

En el resto del capítulo $H$ es separable. Un vector $\xi$ es \emph{cíclico} para $\mathcal A$ si $\overline{\mathcal A\xi}=H$, y es \emph{separante} si $T\xi=0$, $T\in\mathcal A$, implica $T=0$. La primera propiedad dice que un solo vector genera el espacio; la segunda, que la acción sobre ese vector permite distinguir operadores del álgebra.
\index{vector!cíclico}\index{vector!separante}

\begin{proposicion}\label{v3:ciclico-separante}
Si $\mathcal A$ es una subálgebra autoadjunta que contiene $I$, entonces $\xi$ es cíclico para $\mathcal A$ si y sólo si es separante para $\mathcal A'$.
\end{proposicion}
\begin{proof}
Si $\mathcal A\xi$ es denso y $T\in\mathcal A'$ anula $\xi$, entonces $TS\xi=ST\xi=0$ para cada $S\in\mathcal A$. Por continuidad, $T$ anula todo $H$ y es cero. Recíprocamente, sea $K=\overline{\mathcal A\xi}$. Es reductor porque $\mathcal A$ es autoadjunta, así su proyección $P_K$ pertenece a $\mathcal A'$. Como $I\in\mathcal A$, $\xi\in K$, y $(I-P_K)\xi=0$. Si $\xi$ es separante para $\mathcal A'$, se obtiene $I-P_K=0$, así $K=H$.
\end{proof}

\begin{teorema}\label{v3:separante-abeliana}
Toda álgebra de von Neumann conmutativa unital sobre un Hilbert separable no nulo tiene un vector separante. Si es maximal conmutativa, tiene un vector cíclico.
\end{teorema}
\begin{proof}
La descomposición en espacios cíclicos del teorema~\ref{v3:descomposicion-ciclica} da vectores unitarios $(\xi_j)$ cuyas clausuras $H_j=\overline{\mathcal A\xi_j}$ son ortogonales y generan $H$. Como los propios $\xi_j$ forman un conjunto ortonormal en un Hilbert separable, la familia es a lo sumo numerable. Para justificarlo, bolas de radio menor que $\sqrt2/2$ alrededor de vectores ortonormales son disjuntas y cada una contiene un punto de un conjunto numerable denso. Enumere la familia y ponga $\xi=\sum_j2^{-j}\xi_j$; la serie converge en norma.

Si $T\in\mathcal A$ y $T\xi=0$, los vectores $T\xi_j$ pertenecen a subespacios mutuamente ortogonales. Por tanto
$0=\norm{T\xi}^2=\sum_j4^{-j}\norm{T\xi_j}^2$, y cada $T\xi_j=0$. La conmutatividad implica $TS\xi_j=ST\xi_j=0$ para todo $S\in\mathcal A$, así $T$ anula cada $H_j$ y, por densidad de su suma, todo $H$. Entonces $\xi$ es separante. Si $\mathcal A$ es maximal conmutativa, $\mathcal A=\mathcal A'$; la proposición anterior convierte ese vector separante en cíclico.
\end{proof}

\section{Representación de un álgebra conmutativa cíclica}

\begin{teorema}\label{v3:representacion-abeliana-ciclica}
Sea $\mathcal A$ una álgebra de von Neumann conmutativa unital sobre un Hilbert $H\ne\{0\}$ y suponga que tiene un vector cíclico. Existen un compacto Hausdorff $X$, una probabilidad regular $\mu$ con soporte $X$ y un unitario $U:H\to L^2(X,\mu)$ tales que
\[
 U\mathcal A U^*=\{M_f:f\in L^\infty(X,\mu)\}.
\]
Si $H$ es separable, puede elegirse $X$ compacto metrizable.
\end{teorema}
\begin{proof}
Primero construimos una representación usando $X=\Delta(\mathcal A)$, sin afirmar aún metrizabilidad. Sea $\mathcal G:\mathcal A\to C(X)$ el isomorfismo de Gelfand y normalice el vector cíclico $\xi$ para que $\norm\xi=1$. El funcional positivo
\[
 f\longmapsto\ip{\mathcal G^{-1}(f)\xi}{\xi}
\]
vale uno en la identidad. Riesz--Markov lo representa por una probabilidad regular $\mu$ en $X$.

El vector $\xi$ también es separante para $\mathcal A$, porque es separante para $\mathcal A'$ y $\mathcal A\subset\mathcal A'$. Si un abierto no vacío $V$ tuviera medida cero, elija una función continua no nula $0\le f\le1$ con soporte en $V$; se obtiene separando un punto de $V$ del cerrado $X\setminus V$ por Urysohn. Entonces
$\norm{\mathcal G^{-1}(\sqrt f)\xi}^2=\int f\,d\mu=0$. La separación obligaría $\mathcal G^{-1}(\sqrt f)=0$ y, por inyectividad, $f=0$, contradicción. Así $\mu$ tiene soporte total.

Defina en el subespacio denso $\mathcal A\xi$ la aplicación $U_0(T\xi)=\mathcal G(T)$. Es isométrica porque
\[
 \int_X|\mathcal G(T)|^2\,d\mu
 =\ip{T^*T\xi}{\xi}=\norm{T\xi}^2.
\]
La misma igualdad prueba que es independiente de la elección de $T$. Su imagen es $C(X)$, densa en $L^2(\mu)$ por regularidad; así se extiende a un unitario $U:H\to L^2(\mu)$. Para $S,T\in\mathcal A$,
\[
 UTU^*(\mathcal G(S))=UTS\xi=\mathcal G(TS)
 =\mathcal G(T)\mathcal G(S).
\]
Por densidad, $UTU^*=M_{\mathcal G(T)}$.

La conjugación por un unitario es un $*$-isomorfismo y un homeomorfismo para SOT y WOT: las normas de vectores se conservan, y los coeficientes se transforman como
$\ip{UTU^*g}{h}=\ip{TU^*g}{U^*h}$. Por tanto $U\mathcal A U^*$ es WOT-cerrada. Acabamos de identificarla con los multiplicadores continuos, y éstos son WOT-densos en todos los multiplicadores de $L^\infty$. De aquí la igualdad afirmada.

Para la elección métrica cuando $H$ es separable, se necesita una construcción adicional. La bola unitaria de $\mathcal A$ es WOT-compacta y metrizable: una familia numerable densa de vectores de $H$ proporciona un número numerable de coeficientes que, en una bola acotada, determinan todos los demás por aproximación. Una suma con pesos geométricos de las diferencias truncadas de esos coeficientes define una métrica. Elija una sucesión $(T_n)$ WOT-densa en esa bola y sea $\mathcal B=C^*(I,T_1,T_2,\ldots)$. Es normoseparable, pues los $*$-polinomios en un número finito de generadores con coeficientes racionales complejos forman un conjunto numerable denso. Su clausura WOT, y por convexidad también SOT, es $\mathcal A$.

El mismo vector $\xi$ es cíclico para $\mathcal B$: si $T\in\mathcal A$, la aproximación fuerte por elementos de $\mathcal B$ aproxima $T\xi$. Además es separante para $\mathcal B$, por inclusión en $\mathcal A$. Repita la construcción anterior con $X=\Delta(\mathcal B)$, que ahora es compacto metrizable por separabilidad. Produce una probabilidad de soporte total y un unitario que lleva $\mathcal B$ a los multiplicadores continuos. Tomar clausuras WOT lleva $\mathcal A$ a todos los multiplicadores de $L^\infty(X,\mu)$, como se requería.
\end{proof}

\begin{observacion}[Dos espacios compactos diferentes]
El espacio $\Delta(\mathcal A)$ de toda el álgebra de von Neumann no tiene por qué ser segundo numerable, aunque $H$ sea separable. Por ejemplo, $L^\infty([0,1])$ actúa en el Hilbert separable $L^2([0,1])$, pero no es separable en norma: los indicadores $\mathbf1_{[0,t]}$, $0<t<1$, están a distancia uno para parámetros distintos. Si su espacio de Gelfand fuera compacto metrizable, $C(\Delta(\mathcal A))$ sería separable en norma, contradiciendo el isomorfismo de Gelfand. La elección métrica del teorema usa $\Delta(\mathcal B)$ para una subálgebra normoseparable fuertemente densa, no $\Delta(\mathcal A)$.
\end{observacion}

\section{Caracterización abstracta del caso conmutativo}

\begin{teorema}\label{v3:caracterizacion-abeliana}
Toda álgebra de von Neumann conmutativa no nula sobre un Hilbert separable es $*$-isomorfa isométricamente a $L^\infty(X,\mu)$ para algún compacto metrizable $X$ y alguna probabilidad regular de soporte $X$.
\end{teorema}
\begin{proof}
Si la identidad propia es $P\ne I_H$, restrinja primero la acción a $PH$. Como todo elemento satisface $T=PTP$, esa restricción es fiel e isométrica; su imagen es WOT-cerrada, pues un límite de operadores restringidos se extiende por cero a $P^\perp H$ y los coeficientes muestran convergencia WOT de las extensiones. Así podemos suponer identidad $I$.

El teorema~\ref{v3:separante-abeliana} da un vector separante $\xi$. Sea $H_0=\overline{\mathcal A\xi}$, subespacio reductor, y defina $\rho(T)=T|_{H_0}$. Es un $*$-homomorfismo: la invariancia de $H_0$ y $H_0^\perp$ permite restringir tanto operadores como adjuntos. Si $\rho(T)=0$, en particular $T\xi=0$, y la separación implica $T=0$. Así $\rho$ es inyectivo y por tanto isométrico.

Su imagen $\mathcal B=\rho(\mathcal A)$ es normocerrada y tiene a $\xi$ como vector cíclico. Es además WOT-cerrada. En efecto, la restricción es WOT-continua, porque sus coeficientes son coeficientes originales con vectores de $H_0$. La imagen de la bola WOT-compacta de $\mathcal A$ es, por isometría, exactamente la bola de $\mathcal B$, que por tanto es WOT-compacta y cerrada. La caracterización de bolas cerradas del ejercicio~\ref{ej:v3-l33-3} demuestra que $\mathcal B$ es una álgebra de von Neumann.

Aplique el teorema~\ref{v3:representacion-abeliana-ciclica} a $\mathcal B$ sobre el Hilbert separable $H_0$. Se obtiene una representación unitaria de $\mathcal B$ como los multiplicadores de $L^\infty(X,\mu)$, con $X$ compacto metrizable y soporte total. Componer esa identificación con $\rho$ produce el $*$-isomorfismo isométrico buscado para $\mathcal A$. Si se desea incluir el álgebra cero, puede representarse tomando el espacio compacto vacío y la medida cero, sin la normalización de probabilidad.
\end{proof}

La conclusión final es una identificación de álgebras, no necesariamente una equivalencia unitaria de la representación inicial sobre todo $H$ con la representación estándar de multiplicación. Por ejemplo, $f\mapsto M_f\oplus M_f$ actúa en $L^2(\mu)\oplus L^2(\mu)$ y tiene la misma álgebra abstracta $L^\infty(\mu)$, pero para una medida no nula esa representación no es cíclica, como demuestra el ejercicio~\ref{ej:v3-l30-1}. La restricción a un subespacio cíclico generado por un vector separante elimina esa multiplicidad sin perder información algebraica.

El recorrido del volumen termina así en una relación precisa entre tres objetos: álgebras abstractas con norma e involución, operadores sobre espacios de Hilbert y funciones medibles. El cálculo funcional conserva las operaciones; los estados producen representaciones; y las topologías de operadores determinan cuándo aparecen todas las funciones esencialmente acotadas. Para continuar el estudio de las álgebras de von Neumann y de sus representaciones, véanse \cite{conway1990,murphy1990,zhu1993,rudin1991}.

